% Author Thomas C. Hales
% Copyright Thomas C. Hales
% Format latex.
% 

%!TEX TS-program = latex    
%% This line is for TexShop. 

% Revision history. See svn.
% Document created Dec 6, 2002
% Revision started Jan 2007, from published DCG.

\documentclass[cup9a]{cupbook}
%% Cambridge University Press Macros from
%% https://authornet.cambridge.org/information/productionguide/laTex_files/

%

\usepackage{graphicx}
\usepackage{verbatim}
\usepackage{latexsym}
\usepackage{amsfonts}
%\usepackage{amsthm}
\usepackage{amsmath}
%\usepackage{mathidx}
\usepackage{makeidx}
\usepackage{multicol}
\usepackage{crop}
\usepackage{txfonts}
%\usepackage{pdfsync}  %for TexShop sync.
\usepackage[letterpaper,colorlinks=true,
            ps2pdf,hyperindex=true]{hyperref}
%\usepackage{mparhack} %http://www.tex.ac.uk/cgi-bin/texfaq2html?label=marginparside
\usepackage{multind}
\usepackage{url}






% This file contains local settings and system dependencies



% Auxiliary directories
\def\ps{../../../../Pictures/kep07ps}  % flypaper graphics
\def\pdf{/Users/thomashales/Pictures/collect_geom} % tarski graphics

\def\showgraphics{f}  
% t: display graphics (there are none to show yet)
% f (default): print a "no graphics logo" where graphics would normally go.


\def\displayallproof{t} 
% t (default): display all proofs.
% f: print documents without the proofs-- theorem statements only

\def\displayrating{f}
% t (default): display all ratings (verbose is also true)
% f : don't show them.

\def\verbose{f}
% f (default): do not display debugging information,
% t : display debug information and information about the formalization.

\def\tarskipagesep{f}
% f (default): no page breaks between lemmas in tarski collection.
% t :add page breaks

\def\tarskistrip{f}
% f (default): preserves control sequences in tarski collection
% t: control sequences in tarski collection are stripped and printed to an auxiliary file.

\def\tarskidump{f}
% t : dump tarski guid data to external write file.
% f (default) : don't use external write file.
 
%-%
% --Repository--
%-%
% generate revision number by
% svn propset svn:keywords "LastChangedRevision" kepmacros.tex
\def\svninfo{%
  TeXed on \today; \hfill\break
  Repository Root: https://flyspeck.googlecode.com/svn \hfill\break
  SVN $LastChangedRevision$
  }

%-%
% --Fonts--
%-%
\font\twrm=cmr8

%-%
% --Graphics--
%-%
%set \showgraphics option in flag_fly.tex
% flypaper graphics
\def\myincludegraphics#1{%
      \if\showgraphics t{\includegraphics{#1}}%
      \else{\includegraphics{noimage.eps}}\fi}
\def\szincludegraphics[#1]#2{%
      \if\showgraphics t{\includegraphics{#2}}%
      \else{\includegraphics{noimage.eps}}\fi}

% kepler graphics
\def\pdffigtemplatex#1#2#3{%
\begin{figure}[htb]%
  \centering
  \myincludegraphics{\pdfp/#1}
  \caption{#3}
  \label{fig:#2}%
\end{figure}%
}
\def\pdfg#1#2#3{\if\showgraphics t{\pdffigtemplatex{#1}{#2}{#3}}\else{}\fi}


% tarski graphics
\def\pdffigtemplate#1#2#3{%
\begin{figure}[htb]%
  \centering
  \myincludegraphics{\pdf/#1}
  \caption{#3}
  \label{tarski:fig:#2}%
\end{figure}%
}
\def\pdffig#1#2#3{\if\showgraphics t{\pdffigtemplate{#1}{#2}{#3}}\else{}\fi}

%-%
% --Footnotes--
%-%
% http://help-csli.stanford.edu/tex/latex-footnotes.shtml
\long\def\symbolfootnote[#1]#2{\begingroup%
\def\thefootnote{\ensuremath{\fnsymbol{footnote}}}\footnote[#1]{#2}\endgroup}

%-%
% --Special Formatting--
%-%
% http://en.wikibooks.org/wiki/LaTeX/Formatting#List_Structures
\renewcommand{\labelitemii}{$\star$}


%-%
% --Indexing, References, Citations--
%-%
\def\indy#1#2{\index{index/#1}{#2}}

%-%
% --Proof Display--
%-%
% set with \displayallproof in flag_fly. If f, then proofs are swallowed.
%% "proved" environment. toggle with \displayallproof
%
\def\hide#1{}
\def\swallowed{\relax}
\def\swallow#1\swallowed{}
\newenvironment{iproved}{}{}
\newenvironment{proved}{\resetproved\begin{iproved}}{\end{iproved}}
\def\hideproof{\renewenvironment{iproved}{%
   \centerline{\it -- Proof Proofed --}
  \renewenvironment{itemize}{}{}
  \renewenvironment{enumerate}{}{}
  \def\item{\relax}
  \catcode13=12
  \swallow
}{}}
\def\showproof{\renewenvironment{iproved}{\begin{proof}}{\end{proof}}}
\def\resetproved{\if\displayallproof t\showproof\else\hideproof\fi}



%-%
% --Debugging Information--
%-%
\def\uf#1{{\par\narrower\it #1}} % unfinished manuscript
%% verbose:
\def\rating#1{\if\displayrating t{{\textsc {[rating={\ensuremath {#1}}].\ }}}\else{}\fi}
\def\oldrating#1{\if\displayrating t{{\textsc {[former rating={\ensuremath {#1}}].\ }}}\else{}\fi}
\def\formal#1{\if\verbose t{{\tt [formal: #1].\ }}\else{}\fi}
\def\formalauthor#1{\if\verbose t{{\tt [formal proof by: #1].\ }}\else{}\fi}
\def\footformal#1{\if\verbose t{\footnote{\formal{#1}}}\else{}\fi}
\def\usage#1{\if\verbose t{\symbolfootnote[1]{#1}}\else{}\fi}
\def\move#1#2{\begin{#1} See \ref{#2}.\end{#1}}
\def\FIXX#1{\if\verbose t{\footnote{\tt [#1]}}\else{}\fi}
\def\page#1#2{#2 {\it [dcgp.#1]}}
% first entry Lemma or Sec of DCG, second entry page.
\def\dcg#1#2{{\if\verbose t{{\tt{[DCG-#1]}}\indy{References}{ZC{#2 #1}@{DCG-#1}|page{#2}}}\else{}\fi}}
\def\tlabel#1{\label{#1}\if\verbose t{{\tt [#1].\ }%
   \indy{References}{#1|itt}}\else{}\fi}
\def\ifverbose#1{\if\verbose t{{#1}}\else{}\fi}




%% Indexing
\def\refXX{0} % for unknown links.
\def\showref#1{\ref{#1}{\tt[#1]}\indy{References}{#1}}
\def\oldlabel#1{\label{x-#1}}

\def\tref#1{\ref{#1}\indy{References}{#1}}
\def\itt#1{{\bf #1}}
\def\guid#1{{\tt[#1].\ }\indy{References}{ZA{#1}@{#1}|itt}}
% textsc
\def\calc#1{{\textsc{calc-#1}}\indy{Interval}{{#1}@{#1}}}
\def\newcalc#1{{\textsc{newcalc-#1}}\indy{Interval}{{#1}@{#1 (new XX)}}}
\def\assembly#1{{\textsc{assembly-#1}}\indy{References}{ZD{#1}@{#1}}}
%\def\conseq#1{{\textsc{conseq-#1}}}

%% 


%% Tarski indexing. Write summary data to external file "\mywrite",
% which is generated in tarski.tex
% Need to extend cs character set.  Use this only within a block.
\def\setcat{
\catcode`\-=11
\catcode`\.=11
\catcode`\:=11
\catcode`\0=11
\catcode`\1=11
\catcode`\2=11
\catcode`\3=11
\catcode`\4=11
\catcode`\5=11
\catcode`\6=11
\catcode`\7=11
\catcode`\8=11
\catcode`\9=11
}
\def\tarfe#1#2{{\foote{{\tt[L.E.G:#1]}~\csname #1-sum\endcsname%
  \expandafter\ifx \csname #1-guid\endcsname \relax [XX-BAD-IDENTIFIER]\else{}\fi
  }
  {\csname #1-guid\endcsname{#2}}}}
\def\tarf#1{\tarfe{#1}{}}
\def\tarfE#1{\tarfe{EE}{-\csname #1\endcsname}}
% control sequences will contain [-.0-9:] \catcode`\-=11, etc.
\newtoks\mysummary
\newtoks\myguid
\newtoks\myname
\def\separator{\if\tarskipagesep t{\clearpage}\else%
  \smallskip\hrule height 0.5pt depth 0pt width 60pt\bigskip\fi}
\def\identity#1{#1}
\def\foote#1#2{\footnote{#1~{\it (#2)}}}
\def\marker#1#2{\par\hangafter1\hangindent=1em\noindent{\bf #1:}\quad #2}
\def\name#1{\marker{Name}{#1}\myname={#1}
  \write\myhtml{  "#1", //name}
}
\def\summary#1{\marker{Summary}#1\mysummary={#1}
  % cleanse control sequences for html
  \if\tarskistrip t{
  \def\sqr{sqrt}
  \def\sqrt{sqrt}
  \def\op{\identity}
  \def\epsilon{epsilon}
  \def\Delta{Delta}
  \def\ups{upsilon}
  \def\beta{beta}
  \def\rho{rho}
  \def\CalE{E}
  \def\chi{chi}
  \def\eta{eta}
  \write\myhtml{  "#1", // summary}
  }
  \else\relax\fi
}
\def\guid#1{\if\tarskidump f{{\tt [#1]}}\else{\marker{ID}#1
   \indy{guid}{#1}
   \immediate\write\mywrite{\global\expandafter\def\csname \the\myname-sum \endcsname{\the\mysummary}}
   \immediate\write\mywrite{\global\expandafter\def\csname \the\myname-guid \endcsname{#1}}
   \immediate\write\myhtml{  "#1" // guid}}%
   \fi
}
\def\tag#1{\marker{Tags}#1
   \immediate\write\myhtml{  "#1", // tag}
}
%\def\rating#1{\marker{Rating}#1
%   \write\myhtml{  <rating>#1</rating>}
%}
\newenvironment{tarski}{%
\write\myhtml{new Tarski(}
%\def\name   {\marker{Name}}
%\def\summary{\marker{Summary}}
%\def\tag    {\marker{Tags}}
\def\rating {\marker{Rating}}
\def\hol{\marker{HOL-Light}}
%\def\guid   {\marker{ID}}
\def\tlabel{\label}
}{\separator\immediate\write\myhtml{),}
}
\newenvironment{tarskidata}{\write\myhtml{var tarskis=[}}
{
\immediate\write\myhtml{];}
}
%% tarski indexing
%\def\tarf#1{\footnote{{\tt geom-#1~\ref{#1}}\indy{ZA{#1}@{#1}}}}




%-%
% --Symbols--
%-%
% norm
\def\|{{\hskip0.1em|\hskip-0.15em|\hskip0.1em}}
\def\mid{\ :\ }
\def\norm#1#2{\|#1 - #2\|}
\def\normo#1{{\|#1\|}}
\def\sland{\ \land\ }
% Sam's
\def\myscorept{\text{ \sl pt}}
\def\qrtet{{quasi-regular tetrahedron}}
\def\qrtets{{quasi-regular tetrahedrons}}
\def\tomcite{{}}
\DeclareMathOperator{\myscore}{\sigma}
 \DeclareMathOperator{\gma}{gma}
 \DeclareMathOperator{\score}{\sigma}
 \DeclareMathOperator{\sol}{sol}
 \DeclareMathOperator{\dih}{dih}
\newtheorem{calcf}{Calculation}[subsection]

% mathcal
\def\CalB{{\mathcal B}}
\def\CalC{{\mathcal C}}
\def\CalD{{\mathcal D}}
\def\CalE{{\mathcal E}}
\def\FF{{\mathcal F}}
\def\CalF{{\mathcal F}}
\def\CalQ{{\mathcal Q}}
%\def\CalW{{\mathcal W}}
\def\CalR{{\mathcal R}}
\def\CalS{{\mathcal S}}
\def\CalV{{\mathcal V}}
\def\T{{\mathcal T}}
\def\Q{{\mathcal Q}}


% brackets
\def\leftopen{]}
\def\leftclosed{[}
\def\rightopen{[}
\def\rightclosed{]}

%% obsolete
\def\leftb{[}
\def\rightb{]}

% squiqqly relations
\def\seq{\approx}
\def\sle{\preceq}
\def\sge{\succeq}
\def\slt{\prec}
\def\sgt{\succ}

% mathbb
\def\R{{\mathbb R}}
\def\N{{\mathbb N}}
\newcommand{\ring}[1]{\mathbb{#1}}
\def\A{{\mathbf A}}
\def\Rp{\ring{R}^{3\,\prime}}

% operatorname
\def\op#1{{\operatorname{#1}}}
\def\optt#1{{\operatorname{{\texttt{#1}}}}}

\def\opat{\op{@}}
\def\atn{\op{arctan\ensuremath{_2}}}
\def\azim{\op{azim}}
\def\sol{\operatorname{sol}}
\def\vol{\op{vol}}
\def\quo{\operatorname{quo}}
\def\anc{\operatorname{anc}}
\def\cro{\operatorname{crown}}
%\def\vor{\operatorname{vor}}
%\def\svor{\operatorname{svor}}
\def\sign{\operatorname{sign}}
\def\octavor{\operatorname{octavor}}
\def\dih{\operatorname{dih}}
\def\Adih{\operatorname{Adih}}
\def\arc{\operatorname{arc}}
\def\cosarc{\operatorname{cosarc}}
\def\rad{\operatorname{rad}}
\def\gap{\operatorname{gap}}
\def\sc{{\operatorname{sc}}}
\def\geom{{\operatorname{g}}}
\def\anal{{\operatorname{an}}}
\def\PM{\operatorname{PM}}
\def\bool{\operatorname{bool}}
\def\true{\op{true}}
\def\false{\op{false}}
\def\flat{\operatorname{flat}}
\def\tangle#1{\langle #1\rangle}
\def\ceil#1{\lceil #1\rceil}
\def\floor#1{\lfloor #1\rfloor}
\def\ups{\upsilonup} % Needs txfonts; else use \upsilon
\def\orgn{\varthetaup} % center of packing
%\def\comp#1{\llbracket #1 \rrbracket}
\def\comp#1{[#1]}
\def\Wdart{W_{\text{dart}}}
\def\Wedge{W_{\text{edge}}}
\def\cell{\operatorname{cell}}

\def\SA{A}
\def\SB{B}
\def\SC{C}
\def\SD{D}
\def\SE{E}
\def\del{\partial}
\def\doct{\delta_{oct}}
\def\dtet{\delta_{tet}}
\def\pt{\hbox{\it pt}}
\def\hm{{h_0}} % 1.26
%\def\Vol{\hbox{vol}}
\def\scoregoal{8\,\pt}
\def\maxpi{\pi_{\max}}
\def\tausc{{\tau\!\operatorname{sc}}}
\def\piF{{\pi_F}}
\def\xiG{\xi_\Gamma}
\def\piG{\pi_\Gamma}
\def\xiV{\xi_V}
\def\xik{\xi_\kappa}
\def\xikG{\xi_{\kappa,\Gamma}}
\def\piV{\pi_V}
\def\tauLP{{\tau_{\hbox{\twrm LP}}}}
\def\DLP{\operatorname{D}_{\hbox{\twrm LP}}}
\def\ZLP{\operatorname{Z}_{\hbox{\twrm LP}}}
\def\tlp{\tau_{\hbox{\twrm LP}}}  % 2 args (p,q) tri, quad
\def\sLP{\sigma_{\hbox{\twrm LP}}}  % 2 args (p,q) tri, quad
\def\squander{(4\pi\zeta-8)\,\pt}
\def\trgt{\operatorname{\it{target}}}
\def\tildeF{{\hbox{$\tilde F$}}}
\def\zloop{{z}_{loop}}
\def\dloop{{\delta}_{loop}}
\def\sqr{\sqrt}
\def\tildeV{V^{\CalS}}
%\def\BigD{\Delta}
%\def\bigd{\Delta^0}
\def\FCinner{\op{FC}_{\text{inner}}}
\def\FC{\op{FC}}
\def\FCR{\op{FCR}}
\def\rogFC{\op{rog}^0_{\text{FC}}}
\def\chino{\chi_0}

\def\x#1{\op{\bf c}{#1}}  % marks context \x(p,r).
\def\xs#1{\op{\bf sc}{#1}}  % marks special context \x(p,r,s).
\def\xtazimlt{\x{(3,1: \op{azim}(\op{gap}) <\pi)}}
\def\xtazimge{\x{(3,1: \op{azim}(\op{gap}) \ge \pi)}}
\def\xfgap{\x{(4,1:\op{gap})}}
\def\y#1{\op{\bf ep}{#1}}  % marks edge parameters \y(n,k).
\def\pqr#1{#1} % marks type (p,q,r).
%\def\pqr#1{\op{\bf typ}{#1}} % marks type (p,q,r).

%% HYPERMAP macros:
% avoid e for both hypermap edge and edge {v,w}
\def\e{\varepsilonup}
\def\ocirc{}



%% FORMULATION macros:
\def\lam{\lambda}
\def\Lam{\Lambda}
\def\bl{{\underline{\lam}}}
\def\bm{{\underline{\mu}}}
\def\angle#1#2#3{{\op{angle}(#1,{#2,#3})}}


\def\card{\op{card}}

% accent used in India citation in Overview.
\def\=#1{\accent"16 #1}
\def\i{I}
%\def\={\relax}



%% Enclosed CALE control sequences:
\expandafter\def\csname 2t0-doesnt-pass-through\endcsname{01}
\expandafter\def\csname E1453\endcsname{02}   
\expandafter\def\csname dcg-p89\endcsname{03}   
\expandafter\def\csname enclosed-v\endcsname{04}   
\expandafter\def\csname 277\endcsname{05}   
\expandafter\def\csname 245\endcsname{06}   
\expandafter\def\csname 245bis\endcsname{07}   
\expandafter\def\csname node\endcsname{08}   
\expandafter\def\csname E:part4:1\endcsname{09}   
\expandafter\def\csname E:part4:2\endcsname{10}   
\expandafter\def\csname E:part4:3\endcsname{11}   
\expandafter\def\csname E:part4:4\endcsname{12}   
\expandafter\def\csname E:part4:5\endcsname{13}   
\expandafter\def\csname E:part4:6\endcsname{14}   
\expandafter\def\csname E:part4:7\endcsname{15}   
\expandafter\def\csname E:part4:8\endcsname{16}   
\expandafter\def\csname E:part4:9\endcsname{17}   
\expandafter\def\csname E:part4:10\endcsname{18}   
\expandafter\def\csname dcg-p142\endcsname{19}   
\expandafter\def\csname convex-quad\endcsname{20}   
\expandafter\def\csname last:E\endcsname{21}
\expandafter\def\csname rem2.7\endcsname{22}     
\expandafter\def\csname pass-anchor\endcsname{23}

%% 
 
 
 
 
 
 
 


 \crop
%\makeindex

% CUP BOOK CLASS SPECIFIC
\newtheorem{lemma}{Lemma}[chapter]
\newtheorem{definition}{Definition}[chapter]
\newtheorem{remark}{Remark}[chapter]
\newtheorem{theorem}{Theorem}[chapter]
\newtheorem{corollary}{Corollary}[chapter]
\newtheorem{example}{Example}[chapter]
\newtheorem{claim}{Claim}[chapter]
\newtheorem{notation}{Notation}[chapter]
\newtheorem{assumption}{Assumption}[chapter]
\newtheorem{interpretation}{Interpretation}[chapter]
\newtheorem{conjecture}{Conjecture}[chapter]
%%%%%%%%%%%%%%%%%%%%%%%%%%%%%%%%%%%%%%%%%%%%%%%%%%%%

\begin{document}
\raggedbottom  % for now.
%\raggedright  % don't worry for now.

    \title{Flyspeck :
      %\\ \phantom{Hello}
      \\  A Blueprint for the Formal Proof of the Kepler Conjecture
      \\ (Marchal Version)}
    \author{Thomas C. Hales}
    \maketitle
    \frontmatter
    \tableofcontents
    \thanks

\mainmatter

\noindent

\bigskip\noindent
This material is based upon work supported by the National Science
Foundation under
Grants 0503447 and 0804189.

\bigskip\noindent\svninfo 

%\bigskip\noindent
%A pdf of an earlier version of this work appears at\hfill\break 
%\url{http://www.math.pitt.edu/~thales/papers/Flyspeck_Blueprint_2008.pdf}.

\bigskip\noindent
%% XX Thanks to Tran Nam Trung on the hypermap chapter.

%% INDEX OF FORMAL TERMS.

\def\idx#1#2{\indy{Spec}{{\tt #1}; #2}}
\def\idz#1#2#3{\indy{Spec}{{\tt #1}; #2; {\tt \textbackslash \string#3}}}
\idx{pt}{$\pt$}
\idz{atn2~x~~y}{$\atn(x,y)$}{atn}
\idx{sqrt8}{$\sqrt8$}
\idx{sqrt2}{$\sqrt2$}
\idx{t0}{$t_0$}
\idx{two\_t0}{$2t_0$}
\idx{square\_2t0}{$(2t_0)^2$}
\idx{square\_4t0}{$(4t_0)^2$}
\idx{zeta}{$\zeta$}
\idz{doct}{$\doct$}{doct}
\idz{dtet}{$\dtet$}{dtet}
\idx{pi\_rt18}{$\pi/\sqrt{18}$}
\idx{rogers}{$\sqrt2/\zeta$}
\idx{compression\_cut}{$1.41$}
\idz{squander\_target}{$\squander$}{squander}
\idz{xi\_gamma}{$\xiG$}{xiG}
\idz{xiV}{$\xiV$}{xiV}
\idx{xi'\_gamma}{$\xiG'$}
\idz{xi\_kappa}{$\xik$}{xik}
\idz{xi\_kappa\_gamma}{$\xikG$}{xikG}
\idz{pi\_max}{$\maxpi$}{maxpi}
\idx{t4}{$t_4$}
\idx{t5}{$t_5$}
\idx{t6}{$t_6$}
\idx{t7}{$t_7$}
\idx{t8}{$t_8$}
\idx{t9}{$t_9$}
\idx{t10}{$t_{10}$}
\idx{s5}{$s_5$}
\idx{s6}{$s_6$}
\idx{s7}{$s_7$}
\idx{s8}{$s_8$}
\idx{s9}{$s_9$}
\idx{s10}{$s_{10}$}
\idx{eps0}{$\epsilon_0$}
\idx{Z31}{$Z_{3,1}$}
\idx{D31}{$D_{3,1}$}
\idx{Z31}{$Z(\y{(3,1)})$}
\idx{Z32}{$Z(\y{(3,2)})$}
\idx{Z33}{$Z(\y{(3,3)})$}
\idx{Z41}{$Z(\y{(4,1)})$}
\idx{Z42}{$Z(\y{(4,2)})$}
\idx{D31}{$D(\y{(3,1)})$}
\idx{D32}{$D(\y{(3,2)})$}
\idx{D33}{$D(\y{(3,3)})$}
\idx{D41}{$D(\y{(4,1)})$}
\idx{D42}{$D(\y{(4,2)})$}
\idx{D51}{$D(\y{(5,1)})$}
\idz{delta\_x x1 ... }{$\Delta(x_1,\ldots)$}{Delta}
\idx{delta\_x4 x1 ... }{$\Delta_4(x_1,\ldots)$}
\idx{delta\_x6 x1 ... }{$\Delta_6(x_1,\ldots)$}
\idz{ups\_x x1 ...}{$\ups(x_1,\ldots)$}{ups}
\idz{eta\_x x1 x2 x3}{$\eta(\sqrt{x_1},\sqrt{x_2},\sqrt{x_3})$}{eta}
\idz{eta\_y y1 y2 y3}{$\eta(y_1,y_2,y_3)$}{eta}
\idz{rho\_x x1 ... }{$\rho(x_1,\ldots)$}{rho}
\idx{rad2\_y y1 ... }{$(\rad_V\{v_0,v_1,v_2,v_3\})^2,\ y_1 = |v_0-v_1|,\ldots$}
\idz{chi\_x x1 ... }{$\chi(x_1\ldots)$}{chi}
\idz{dih\_x x1 ... }{$\dih(\sqrt{x_1},\ldots)$}{dih}
\idz{dih\_y y1 ... }{$\dih(y_1,\ldots)$}{dih}
\idx{dih2\_x x1 x2 ... }{$\dih(\sqrt{x_2},\ldots)$}
\idx{dih2\_y x1 x2 ... }{$\dih(y_2,\ldots)$}
\idx{dih3\_x x1 x2 ... }{$\dih(\sqrt{x_3},\ldots)$}
\idx{dih3\_y x1 x2 ... }{$\dih(y_3,\ldots)$}
\idz{sol\_x x1 ... }{$\sol(v_0,\ldots),\ x_1=|v_0-v_1|^2,\ldots$}{sol}
\idz{sol\_y y1 ... }{$\sol(v_0,\ldots),\ y_1=|v_0-v_1|,\ldots$}{sol}
\idx{vol\_x x1 ... }{$\sqrt{\Delta(x_1,\ldots)}/12$}
\idz{beta psi theta}{$\beta_\psi\theta$}{beta}
\idz{arclength a b c}{$\arc(a,b,c)$}{arc}
\idx{volR a b c}{$\op{volR}(a,b,c)$}
\idx{solR a b c}{$\op{solR}(a,b,c)$}
\idx{dihR a b c}{$\op{dihR}(a,b,c)$}
\idx{vorR a b c}{$\op{sovoR}(a,b,c,\lambda_{oct})$}
\idx{denR a b c}{$\delta(a,b,c)$}
\idx{tauR a b c}{$\op{sovoR}(a,b,c,\lambda_{sq})$}
\idx{qy y1 y2 y3 t}{$\op{quovol}(a,b,c)$}
\idx{quo\_x x y z}{}
\idx{qn y1 y2 z t}{}
\idx{phi h t}{}
\idx{phi0}{}
\idx{eta0 h}{}
\idx{crown h}{}
\idx{anc y1 y2 y6}{}
\idx{K0 y1 y2 y6}{}
\idx{AH h t}{}
\idx{BHY y}{}
\idx{KY y1 y2 y3 y4 y5 y6}{}
\idx{KX x1 x2 x3 x4 x5 x6}{}
\idx{vor\_analytic\_x x1 x2 x3 x4 x5 x6}{}
\idx{vor\_analytic\_x\_flipped x1 x2 x3 x4 x5 x6}{}
\idx{octavor\_analytic\_x x1 x2 x3 x4 x5 x6}{}
\idx{tau\_analytic\_x x1 x2 x3 x4 x5 x6}{}
\idx{kappa y1 y2 y3 y4 y5 y6}{}
\idx{kappa\_dih\_y y1 y2 y3 y5 y6 d}{}
\idx{level\_at h t}{}
\idx{vorstar h1 h2 h3 a1 a2 a3 t}{}
\idx{vort\_y y1 y2 y3 y4 y5 y6 t}{}
\idx{vor\_0\_y y1 y2 y3 y4 y5 y6}{}
\idx{tau\_0\_y y1 y2 y3 y4 y5 y6}{}
\idx{vor\_0\_x x1 x2 x3 x4 x5 x6}{}
\idx{tau\_0\_x x1 x2 x3 x4 x5 x6}{}
\idx{vort\_x x1 x2 x3 x4 x5 x6 t}{}
\idx{tauVt\_x x1 x2 x3 x4 x5 x6 t}{}
\idx{vorA\_x x1 x2 x3 x4 x5 x6}{}
\idx{tauA\_x x1 x2 x3 x4 x5 x6}{}
\idx{vorC0\_x x1 x2 x3 x4 x5 x6}{}
\idx{tauC0\_x x1 x2 x3 x4 x5 x6}{}
\idx{vorC\_x x1 x2 x3 x4 x5 x6}{}
\idx{tauC\_x x1 x2 x3 x4 x5 x6}{}
\idx{v0x x1 x2 x3 x4 x5 x6}{}
\idx{v1x x1 x2 x3 x4 x5 x6}{}
\idx{gamma\_x x1 x2 x3 x4 x5 x6}{}
\idx{tau\_gamma\_x x1 x2 x3 x4 x5 x6}{}
\idx{rad2\_x x1 x2 x3 x4 x5 x6}{}
\idx{sigma\_qrtet\_x x1 x2 x3 x4 x5 x6}{}
\idx{sigma1\_qrtet\_x x1 x2 x3 x4 x5 x6}{}
\idx{tau\_sigma\_x x1 x2 x3 x4 x5 x6}{}
\idx{sigma32\_qrtet\_x x1 x2 x3 x4 x5 x6}{}
\idx{mu\_flat\_x x1 x2 x3 x4 x5 x6}{}
\idx{taumu\_flat\_x x1 x2 x3 x4 x5 x6}{}
\idx{mu\_upright\_x x1 x2 x3 x4 x5 x6}{}
\idx{mu\_flipped\_x x1 x2 x3 x4 x5 x6}{}
\idx{vor\_0\_x\_flipped x1 x2 x3 x4 x5 x6}{}
\idx{octavor0\_x x1 x2 x3 x4 x5 x6}{}
\idx{nu\_x x1 x2 x3 x4 x5 x6}{}
\idx{nu\_gamma\_x x1 x2 x3 x4 x5 x6}{}
\idx{taunu\_x x1 x2 x3 x4 x5 x6}{}
\idx{octa\_x x1 x2 x3 x4 x5 x6}{}
\idx{sigmahat\_x x1 x2 x3 x4 x5 x6}{}
\idx{sigmahat\_x' x1 x2 x3 x4 x5 x6}{}
\idx{sigmahatpi\_x x1 x2 x3 x4 x5 x6}{}
\idx{tauhat\_x x1 x2 x3 x4 x5 x6}{}
\idx{tauhatpi\_x x1 x2 x3 x4 x5 x6}{}
\idx{pi\_prime\_tau k0 k1 k2}{}
\idx{pi\_prime\_sigma k0 k1 k2}{}



  % Formal Spec Index.

\smallskip
\newpage



%%%%%%%%%%%%%%%%%%%%%%%%%%%%%%%%%%%%%%%%%%%%%%%%%%%%


%%  \label{part:intro}
   \newpage

   \setcounter{chapter}{-1}
   \chapter{Preface}
   \label{XX}  % for missing references. XX.
   %%------------------------------------------------------------
% Author: Thomas C. Hales
% Format: LaTeX
% Book Chapter: Dense Sphere Packings
%------------------------------------------------------------

\chapter*{Preface}


%{
%
%\narrower
%
%{\it ``A personality which has been veiled by a formal method
%  throughout many chapters is suddenly seen face to face in the
%  Preface.'' }
%% Introductory note, p3, Famous Prefaces, Harvard Classics, Vol
%% 39. P.F. Collier & Son, 1910.
%
%}

%\bigskip

%\centerline{\it ``Those who justify themselves do not convince.''
%  --Lao-Tzu}
%% quoted in Watts's essay in "Modern Buddhism" ed. Donald Lopez,
%% page 160


{

\narrower

{\it

  ``I think there's a revolution in mathematics around the corner. I
  think that~\dots{}   %in later times
  people will look back on the fin-de-siecle of the twentieth century
  and say `Then is when it happened' (just like we look back at the
  Greeks for inventing the concept of proof and at the nineteenth
  century for making analysis rigorous). I really believe that. And it
  amazes me that no one seems to notice\dots.

  ``Never before have the Platonic mathematical world and the physical
  world been this similar, this close. Is it strange that I expect
  leakage between these two worlds? That I think the proof strings
  will find their way to the computer memories?\dots 

  ``What I expect is that some kind of computer system will be
  created, a proof checker, that all mathematicians will start using
  to check their work, their proofs, their mathematics. I have no idea
  what shape such a system will~\dots{}  take. But I expect some
  system to come into being that is past some threshold so that it is
  practical enough for real work, and then quite suddenly some kind of
  `phase transition' will occur and everyone will be using that
  system.''

{\hfill\rm--Freek Wiedijk~\cite{FWR}} % http://www.cs.ru.nl/~freek/jordan/index.html

}

}

\newpage

{

\narrower\parindent=0pt
\parskip=0.4\baselineskip

{\it

Alecos: Christos has a problem with the `foundational quest'!

Christos:  Wrong!  I have two problems with your  {\rm {version}} of it!  One, it
didn't fail and, two, it wasn't a tragedy!  Granted, there are some tragic
parts!  But the ending is happy, as in the `Oresteia'!  

Apostolos:  Happy for whom?  Cantor, going insane?  G\"odel starving himself to
death out of paranoia? Hilbert or Russell and their psychotic sons? Or Frege with--

Christos: `The meaning
is in the ending!' you said so yourself!  So, follow the quest for ten more years and
you get a brand-new triumphant finale with the creation of the computer, which is the
quest's real hero!   Your problem is, simply, that you see it as a story of people!

Apostolos: Well, stories do tend to be about people!

Christos:  So, choose the right people!  And show what they really did!  All we we learn
of the great von Neumann is he said `It's over'  when he heard G\"odel!

Alecos: But it was over in a sense, wasn't it?  Pop went Hilbert's `no ignorabimus'!

Christos:  But then came the quest's jeune premier, its parsifal~\dots{}  Alan Turing!
He said `Ok, we can't prove everything! So, let's see what we can prove!' and to define
proof, he invented, in 1936, a theoretical machine which contains all the ideas of the
computer!\dots{}  which, after the war, he and von Neumann, the quest's proudest sons,
brought to full life!

{\hfill\rm--Doxiadis and Papadimitriou, ~~Logicomix~\cite{Logicomix}} % page 303.

}

}




\newpage

{

\narrower\parindent=0pt
\parskip=0.4\baselineskip

{\it

``Ever hear of the Kepler Conjecture?''

``Nope.''

I laid the notebook on the table and flipped through the pages. ``It was first stated
in 1611 by Johannes Kepler,'' I said.  ``Kepler became interested in the problem while
he was corresponding with an Englishman named Thomas Harriot, who was trying
to help his friend Sir Walter Raleigh figure out the best way to stack cannonballs on ship
decks.  The goal was to find the densest possible spherical arrangements~\dots{}  basically,
the way grocers stack oranges''

``Okay,'' he said, nodding.

``Kepler's conjecture {\rm{seems}} perfectly sound,'' I said.

``That is does,'' Ben said.

``But here's the thing\dots.   I looked it up and discovered that, in 1998, a proof
had finally been put forward by an American mathematician named Thomas Hales.  In 2003,
a committee that had been assigned to verify Hales's work confirmed that they were ninety-nine
percent certain of the proof's correctness.  But that one percent was key.  The mathematical
world is still waiting for the publication of the data that will prove the Kepler Conjecture definitively.''

``Sucks for Thomas Hales,'' Ben said.

``I agree.  But it makes sense that they have to be certain, doesn't it?''


{\hfill\rm--Michelle Richmond, No One You Know~\cite{Richmond}} % page 177.

}

}

%%DD Figure of cannonballs.

\bigskip

{

\narrower\parindent=0pt
\parskip=0.4\baselineskip

{\it

``Sometimes fixing a 1 percent defect takes 500 percent effort.''

{\hfill\rm-- Joel Spolsky, Joel on Software~\cite{Spolsky}} % page 122

}

}

\bigskip

{

\narrower\parindent=0pt
\parskip=0.4\baselineskip

{\it

``Every one fully persuaded is a fool.''

{\hfill\rm-- Barthasar Graci\'an, the Art of Worldly Wisdom~\cite{Gracian}} % p110

}

{

\it

%Believe not everything, but only what is proven: the former is foolish, 
%the latter the act of a sensible man. 

%{\hfill -- Democritus}

}

}



\newpage

\section*{The Kepler Conjecture}

In 1611, Johannes Kepler wrote a booklet in which he asserted that the familiar
cannonball arrangement of congruent balls in space achieves the
highest possible density.  This assertion has become known as the
Kepler conjecture.  This book presents a proof.

As early as 1831, Gauss established a special case of the conjecture,
by proving that the cannonball arrangement is optimal among all
\newterm{lattices}~\cite{Gau31}.  Later in the nineteenth century,
Thue solved the corresponding problem in two dimensions, showing that
the hexagonal arrangement of disks in the plane achieves optimal
density~\cite{Thu92}, \cite{Thu10}.  Hilbert, in his famous list of
mathematical problems, made the Kepler conjecture part of his eighteenth
problem.  In 1953, Fejes T\'oth formulated a general strategy to
confirm the Kepler conjecture, but lacked the computational resources
to carry it out~\cite{Fej53}.  The conjecture was finally resolved in
1998, even though the full proof was not published until
2006~\cite{Hales:2006:DCG}.  Section~\ref{sec:history} gives
additional historical background.

The Kepler conjecture has become a test of the capability of computers to
deliver a reliable mathematical proof.  The original proof 
involved many long computer calculations that
exhausted the efforts of a team of referees.  This book has
 redesigned the proof in a way that makes the correctness of
the computer proof as transparent as possible.


\section*{Formal Proofs}

After all is said and done, a proof is only as reliable as the
processes that are used to verify its correctness.  The ultimate
standard of proof is a formal proof, which is nothing other
than an unbroken chain of logical inferences from an explicit set of
axioms.  While this may be the mathematical ideal of proof, actual
mathematical practice generally deviates significantly from the ideal.



%More than ten years have passed since a proof was first
%obtained. Why give a new presentation of the proof?
%
%The original proof was not widely understood.  The complexity was not
%because of conceptual challenges.  In fact, the proof makes only
%modest demands on the theoretical training of the reader.  It is
%possible to read and understand the proof with a knowledge of a
%limited body of mathematics, such as basic calculus and elementary
%Euclidean geometry.
%
%Nevertheless, the proof involves many long calculations. Even worse,
%it it relies on computer calculations.  An error in any calculation or
%a bug in the computer code has the potential to topple the entire
%proof.
%
%The referees were conscientious and checked many of the calculations.
%However, for the most part, the computer code lay beyond the scope of
%referee review, and even careful quality control can let a
%occasional bug slip through undetected.
%
%After all is said and done, no proof is more reliable than the
%reliability of the processes that are used to verify its
%correctness.  These processes include the checking that the author
%makes before releasing the proof for public scrutiny, the checking
%of the referees, and the checking done by readers after publication.

In recent years, as part of this project, I have been increasingly preoccupied by the
processes that mathematicians rely on to ensure the correctness of complex
proofs. A century ago, Russell's paradox and other antinomies threatened
 set theory  with fires of destruction.
Researchers from Frege to G\"odel solved the problem of
rigor in mathematics and found a theoretical solution but did not
extinguish the fire at the foundations of mathematics
because they omitted the practical implementation. Some, such as
Bourbaki, have even gone so far as to claim that ``formalized
mathematics cannot in practice be written down in full'' and call
such a project
``absolutely unrealizable'' \cite[pp.~10--11]{Bour:68:Sets}. % Theory of
                                                          % Sets, page
                                                          % 10,11.

While it is true that formal proofs may be too long to print,
computers -- which do not have the same limitations as paper -- have
become the natural host of formal mathematics. In recent decades,
logicians and computer scientists have reworked the foundations of
mathematics, putting them in an efficient form designed for real use
on real computers.

For the first time in history, it is possible to generate and verify
every single logical inference of major mathematical theorems.  This
has now been done for many theorems, including the four-color theorem, the prime number
theorem, the Jordan curve theorem, the Brouwer fixed point theorem,
and the fundamental theorem of calculus.  Freek Wiedijk
reports that 83\% of a list of one hundred famous theorems have now been
checked formally \cite{wiedijk:100}.  The list of seventeen remaining
theorems contains two particular challenges: the independence of the
Continuum Hypothesis and Fermat's Last theorem.
% list as of 5/5/2010.

Some mathematicians remain skeptical of the process because computers
have been used to generate and verify the logical inferences.
Computers are notoriously imperfect, with flaws ranging from software
bugs to defective chips.  Even if a computer verifies the inferences,
who  verifies the verifier, or then verifies the verifier of the
verifier?  Indeed, it would be unscientific of us to place an
unmerited trust in computers.

The choice comes down to two competing verification processes.  The
first is the traditional process of referees, which depends largely on
the luck of the draw -- some referees are meticulous, while others are
careless.  The second process is formal computer verification, which
is less dependent on the whims of a particular referee.  In my view,
the choice between the conventional  process by human referee and computer
verification is as evident as the choice between a sundial and an atomic
clock in contemporary science.

The standard of proof I have adopted is the highest scientific
standard available by current technology.  
%That standard, which
%continues to evolve with the advancement of technology, is formal
%verification by computer.
%The boundary that separates an easy proof from a difficult
%one shifts with current technology.  
The introduction of steel in
architecture is not a mere reinforcement of wood and stone; it changes
the world of structural possibilities.  There is no longer any
reason to limit proofs to ten-thousand-pages when our
technology supports a million pages.

The style of formal proofs is different from that of conventional
ones.  It is easier to formalize several short snappy proofs
than a few intricate ones.  Humans enjoy surprising new perspectives,
but computers benefit from repetition and standardization.  Despite
these differences, I have sought  proofs that might bring
pleasure to the human reader while providing precise instructions for
the implementation in silicon.

\section*{Conventions}

To make formalization proceed more smoothly, long proofs have been
broken into a series of smaller claims.  Each claim starts a new paragraph
and is set in italics.  The second sentence of the paragraph begins with
the word \newterm{indeed} when the proof of the claim is direct and with
the word \newterm{otherwise} when the proof is indirect by contradiction.

The pronoun \newterm{we} is used inclusively for the author and reader as
we work our way through the proofs in this book.  The pronoun \newterm{I} 
refers to the author alone, as in the first sentence of the following section.

The asterisk $\wild$ is used as a wildcard symbol in patterns.  It
replaces a term in contexts where the name of the term is not
relevant.  It can also denote a bound variable.  For example, the
function $f(\wild,y)$ of a single variable is obtained from $f$ by
evaluating the second argument at a fixed value $y$.

The union of the family $X$ of sets is written as $\bigcup\ X$ or as
$\bigcup_{x\in X} x$ without any difference in meaning.  The first form
is preferred because of its economy.  We also use both expressions
$\bigcap\ X$ and $\bigcap_{x\in X} x$ for the intersection of a family of sets.

The documentation of the computer calculations for the Kepler
conjecture has evolved over time.  The 1998 preprint version of the
proof of the Kepler conjecture contains long appendices of doubtful
value that list hundreds of calculations that enter into the proof.
These appendices were cut from the published version of the proof
because it is more useful to store the computer part of the proof at a
computer code repository that is permanent, versioned, and freely
available.  The computer code and documentation are housed at
\newterm{Google Code project hosting}.  A separate guide, which is
available at the project site, describes the computer calculations
that appear in this book.  When this book uses an external
calculation, it is marked in italic font as a \cc{This explains notation}.

\section*{A Blueprint}

I dream of a fully formally verified solution to the packing problem.
This project is still unfinished, but significant progress is being
made.  The book is a blueprint for formal proofs because it gives the
design of the formal proof to be constructed.  Decisions about what to
include in this book have been shaped by the list of theorems already
available in the library of the proof assistant \newterm{HOL Light}.
For example, this book accepts basic point-set topology and measure
theory because they have been formalized by Harrison~\cite{HOLL}.


The book is divided into four parts, the first of which
 describes the major ideas, methods, and
organization of the proof.  

%There is an essay on each major computer
%component of the proof. The purpose is to provide a panoramic view of
%proof, to provide intuition about proof strategies.  After reading
%this part of the book, the reading should understand what the proof is
%all about, without yet dipping into technical details.

The part on foundations provides background material about
constructions in discrete geometry.  The first of these chapters
covers trigonometric identities and basic vector geometry.  The second
treats volume from an elementary point of view.  The third chapter
covers planar graph theory from a purely combinatorial perspective.
The fourth chapter continues with planar graphs, now from a geometric
perspective.

The next part of the book gives the solution to the packing problem.
The first chapter  gives a top-level overview of the major
steps of the proof,   describing how the problem can be reduced from
a problem with infinitely many variables to one in finitely many
variables.  The remaining chapters in this part flesh out the proof.

The final part of the book views dense sphere packings from a larger perspective.
One chapter resolves some other longstanding conjectures in
discrete geometry: Bezdek's strong dodecahedral conjecture and Fejes
T\'oth's full contact conjecture.  Another chapter describes  formal proof, 
the HOL Light proof assistant, and the Flyspeck formalization project.
An appendix  gives an introduction  to the methods used in the computer calculations.

\section*{Simplifications}

Many simplifications of the original proof have been found over the
past several years.  These simplifications are published here for the
first time.  Gonthier has reworked the proof of the four-color
theorem to avoid the use of the Jordan curve theorem, using instead
the much simpler notion of M\"obius contour from the theory of
hypermaps.  I have followed Gonthier's lead.

The optimality of the face-centered cubic packing is an assertion
about infinite space-filling packings.  For computational purposes, it is
useful to reduce the sphere packing problem to finite packings.  A
\newterm{correction term} is associated with each different reduction from
infinite packings to finite packings.  Ferguson and I worked together to
find the original proof of the Kepler conjecture.  The two of us considered a
large number of different correction terms, seeking one that
would simplify the computations as much as possible.  In a discussion
of the solution of the packing problem, I wrote that ``correction
terms are extremely flexible and easy to construct, and soon Samuel
Ferguson and I realized that every time we encountered difficulties in
solving the minimization problem, we could adjust $f$ [the correction
term] to skirt the difficulty\dots.  If I were to revise the proof
to produce a simpler one, the first thing I would do would be to
change the correction term once again.  It is the key to a simpler
proof''~\cite{Hales:2000:cannonballs}.  Marchal has recently found a remarkably simple 
correction term, giving a new way  to  reduce from infinite packings
to finite packings~\cite{marchal:2009}.  This book implements his reduction step.

There are many other improvements of the proof that are not visible in
the book because they are implemented in computer code, including
a reduction of  the number of lines of computer code from over 187,000
to well under 10,000.  Needless to say, the quickest way to verify that a block
 of computer code is free of bugs is to delete the code altogether.




\bigskip
\hbox{}



\bigskip
\hbox{}

{
\parindent=0pt
\obeylines

Thomas C. Hales
Pittsburgh, PA
May 2010

}







 %proofed
   
   

%%%%%%%%%%%%%%%%%%%%%%%%%%%%%%%%%%%%%%%%%%%%%%%%%%%%
  
    
    



In this \chap, trigonometry is developed analytically.  The basic
trigonometric functions are defined by their power series
representations, and calculus of a single real variable is used to
develop the basic properties of these functions.  We give a brief
review of the definitions and properties.


\subsection{sine and cosine}

The cosine and sine functions are defined by their infinite series:
    \begin{equation}\label{eqn:cos-def}\cos(x) = 1 - x^2/2! + x^4/4! \cdots,\qquad
  \sin(x) = x - x^3/3! + x^5/5! \cdots.
    \end{equation}
Convergence is absolute for every real number $x$.
The series can be evaluated at $0$:
    \begin{equation}\label{eqn:cos0}
    \cos(0) = 1,\qquad \sin(0) = 0.
    \end{equation}

The series may be differentiated term-by-term to establish the identities:
    \begin{equation}\label{eqn:cos'}
    \frac{d \cos(x)}{dx} = -\sin(x),\qquad \frac{ \sin(x)}{dx} = \cos(x).
    \end{equation}

The powers $(\cos(x))^n$ and $(\sin(x))^n$ are conventionally written
$\cos^n(x)$ and $\sin^n(x)$.

Trigonometric identities follow easily from these definitions.  We
give a couple of examples.

\begin{lemma}\label{lemma:circle}
   $$\cos^2(x) + \sin^2(x) = 1.$$
\end{lemma}

\begin{proof}
By basic properties of derivatives and Equation~\ref{eqn:cos'},
the derivative of the function $f(x) = \cos^2(x) +\sin^2(x)$ is
identically zero.   So the function is constant.  From
Equation~\ref{eqn:cos0}, it follows that $f(x)=f(0)=1$.
\end{proof}

\begin{lemma}\label{lemma:sin_add}
  $$\begin{array}{lll}
  \sin(x+y) &= \sin(x)\cos(y) + \cos(x)\sin(y)\\
  \cos(x+y)  &= \cos(x)\cos(y) - \sin(x)\sin(y).
  \end{array}$$
\end{lemma}

\begin{proof}
Fix $y$.  Let
    $$\begin{array}{lll}
    f(x) = &(\cos(x+y) - \cos(x)\cos(y) +
    \sin(x)\sin(y))^2 +\\ & (\sin(x+y) -\sin(x)\cos(y) -
    \cos(x)\sin(y))^2.
    \end{array}$$
The derivative of $f$ is identically zero.  The function is
therefore constant. Also, $f(0)=0$.  Thus, $f$ is identically zero.
If a sum of real squares is zero, the individual terms are zero. The
identities follow.
\end{proof}

\begin{lemma}\label{lemma:cos_even}
The cosine is an even function.  The sine is an odd
function.
    $$\cos(-x) = \cos(x),\quad\sin(-x) = -\sin(x).$$
\end{lemma}

\begin{proof}
The derivative of
    $$(\cos(-x) - \cos(x))^2 + (\sin(-x) +\sin(x))^2$$
is identically zero.  Argue as in the previous lemmas.
\end{proof}

\subsection{$\pi$}
\label{sec:pi}

It is known that the cosine function has a unique root between $0$
and $2$. The constant $\pi$ is defined as twice that root.  Thus, by
definition $\cos(\pi/2) = 0$, and the cosine has no root greater
than $0$ and less than $\pi/2$. The $\cos$ function is in fact
positive on the interval $\leftclosed 0,\pi/2\rightopen$.


\begin{lemma}\label{lemma:sin_pi2}
    $\sin (\pi/2) = 1.$
\end{lemma}

\begin{proof}
The derivative of $\sin$ is non-negative between $0$ and $\pi/2$.
Its value at $0$ is $0$.  Thus, $\sin$ is non-negative on
$[0,\pi/2]$.  We can check the identity of the lemma by checking
that the squares of the two sides are equal. Then $\sin^2(\pi/2) =
{1-\cos^2(\pi/2)} = 1$.
\end{proof}

\begin{lemma}\label{lemma:sin_cos}
    $$\sin(\pi/2 - x) = \cos(x).$$
\end{lemma}

\begin{proof}
    Apply the addition law for the sine function,
    $$\sin(\pi/2 - x) = \sin(\pi/2)\cos(-x) + \cos(\pi/2)\sin(-x)$$
    and use
    $\sin(\pi/2) = 1$, $\cos(\pi/2) = 0$.  Then use that $\cos$ is
    an even function.
\end{proof}

Similarly, $\cos(\pi/2 - x) = \sin(x)$ or $\cos(\pi/2 + x) =
-\sin(x)$, $\sin(\pi/2 + x) = \cos(x)$.  Further,
    $$\begin{array}{lll}
      \sin(\pi + x) &= \phantom{-}\cos(\pi/2 + x) &= -\sin(x),\\
      \cos(\pi + x) &= -\sin(\pi/2 + x) &= -\cos(x),\\
      \sin(2\pi + x) &= -\sin(\pi + x) &= \phantom{-}\sin(x),\\
      \cos(2\pi + x) &= -\cos(\pi + x) &= \phantom{-}\cos(x).
      \end{array}$$


\subsection{tangent}
\label{sec:tangent}

\begin{definition}\label{def:tan}
Let $\tan(x) = \sin(x)/\cos(x)$, defined when $\cos(x)\ne0$.
\end{definition}


\begin{lemma}\label{lemma:tan_add}
    If $\cos(x)\ne 0$, $\cos(y)\ne 0$, and $\cos(x+y)\ne0$ then
    $$\tan(x+y) = \frac{\tan(x) + \tan(y) }{ 1 - \tan(x)\tan(y)}$$
\end{lemma}

\begin{proof}
  Divide the first line of Lemma~\ref{lemma:sin_add} by the second
  line of the same lemma.
\end{proof}

\begin{lemma}\label{lemma:tan_pi4}
    $$\tan(\pi/4) = 1.$$
\end{lemma}

\begin{proof}
    $$\tan(\pi/4) = \sin(\pi/2-\pi/4)/\cos(\pi/4) =
    \cos(\pi/4)/\cos(\pi/4) = 1.$$
\end{proof}


\subsection{arctangent}

We review the properties of the arctangent function.

\begin{definition}\label{definition:arctan}
There is a unique function $\arctan:\ring{R}\to\ring{R}$ with
image $(-\pi/2,\pi/2)$ such that
    $$\tan(\arctan x) =x.$$
\end{definition}

If $-\pi/2 < x < \pi/2$, then also $\arctan(\tan(x)) = x$. In
particular,
    $$\arctan(1) = \arctan(\tan(\pi/4)) = \pi/4.$$


The function $\arctan$ is differentiable with derivative
    $$\frac{d \arctan(x)}{dx} = \frac{1}{1 + x^2}.$$
The derivative is everywhere positive, and the function $\arctan$ is
strictly increasing.


\subsection{other inverse trig functions}

Other inverse trigonometric functions will generally be reduced to
the arctangent.

\begin{definition}\label{definition:arccos}
$\arccos y$ is the function on the interval $[-1,1]$
taking values in $[0,\pi]$ that is the inverse function of $\cos$:
    $$\begin{array}{lll}
        y\in [-1,1] &\Rightarrow \cos(\arccos y) = y\\
        x\in[0,\pi] &\Rightarrow \arccos(\cos x) = x
    \end{array}$$
\end{definition}

\begin{lemma}\label{lemma:sin_arccos} If $y\in[-1,1]$, then
    $$\sin(\arccos(y)) = \sqrt{1-y^2}.$$
\end{lemma}

\begin{proof}
    The range of $\arccos(y)$ is $[0,\pi]$.  On this interval, $\sin$
    is nonnegative.  Thus, we can check the identity by squaring
    both sides.  It then follows from the circle identity
    (Lemma~\ref{lemma:circle}).
\end{proof}


Our preference is to remove the $\arccos$ function whenever
possible, by replacing it with the $\arctan$ function through the
following identity.  We will occasionally be forced to revert to
$\arccos$ when we are near the endpoints $-1$ and $1$ of its domain.

\begin{lemma}\label{lemma:arccos_arctan}  If $y\in (-1,1)$, then
    $$\arccos(y) = \pi/2 - \arctan(\frac{y}{ \sqrt{1-y^2}}).$$
\end{lemma}

\begin{proof} If $y\in (-1,1)$,  $x = \arccos(y)$, 
$z = \arctan(y/\sqrt{1-y^2})$, then
    $-\pi/2 < \pi/2 - x < \pi/2$, and $-\pi/2 < z < \pi/2$.  It is
    therefore enough to check that
        $\tan(\pi/2 - x) = \tan(z)$.
    But
        $$\tan(\pi/2-x) = \frac{\cos(x)}{\sin(x)} = \frac{y}{
        \sin(\arccos(y))} = \frac{y}{ \sqrt{1-y^2}} = \tan(z).$$
\end{proof}

\begin{definition}\label{def:arcsin}
$\arcsin y$ is the function on the interval $\leftclosed
-1,1\rightclosed$ taking values in $[-\pi/2,\pi/2]$ that is the
inverse function of $\sin$:
    $$\begin{array}{lll}
        y\in [-1,1] &\Rightarrow \sin(\arcsin y) = y\\
        x\in[-\pi/2,\pi/2] &\Rightarrow \arcsin(\sin x) = x
    \end{array}$$
\end{definition}

As we do not have need for the arcsine function, we do not develop
any of its properties.





\section{Triangle Inequality}

The set $\ring{R}^n$ is defined as the set of functions
$f:\ring{N}\to\ring{R}$ such that $f(i) = 0$ when $n\le i$. Defined
this way, there are natural inclusions
    $$\ring{R}^0 \subset \ring{R}^1 \subset \ring{R}^2 \cdots$$

We let $\ring{R}^\infty$ be the union of all $\ring{R}^n$.  It is
the set of function $f:\ring{N}\to\ring{R}$ that have finite
support.  We write $v_i$ for the value of the function $v$ at $i$.
We call elements of $\ring{R}^\infty$ {\it vectors}.  There is a
zero vector $0$, defined as the function that is identically zero.

There is an addition and scalar multiplication on $\ring{R}^\infty$
defined by addition and multiplication of functions.
    $$\begin{array}{lll}
    (u + v)_i &= u_i + v_i.\\
    (s u)_i & s u_i,\quad s\in\ring{R}.\\
    \end{array}
    $$
The addition is commutative and associative, and gives
$\ring{R}^\infty$ the usual properties of a vector space. If $u$ and
$v$ both lie in some $\ring{R}^n$, then so do $u+v$ and $s u$.

We define the difference of two vectors as $u - v = u + (-1) v$.

%\begin{definition}
%A set of vectors $V$ is collinear, if there is vector $u$ such that
%for every $v,v'\in V$ there exists $t$ such that $v - v' = t u$.
%\end{definition}

There is a dot product
$$(\cdot):\ring{R}^\infty\to\ring{R}^\infty\to\ring{R}$$ defined by
    $$u\cdot v = \sum_{i=0}^\infty u_i v_i.$$
Only finitely many terms in the sum are nonzero, so it is actually a
finite sum.  The dot product satisfies the following
properties:
    $$\begin{array}{lll}
        u \cdot (v + w) &= u \cdot v + u \cdot w\\
        (u + v)\cdot w &= u \cdot w + v \cdot w\\
        (s u)\cdot w &= s(u \cdot w) = u \cdot (s w)\\
        0 &\le u\cdot u\\
    \end{array}$$


There is a norm
$$|u| = \sqrt{u\cdot u}.$$
It satisfies $| s u | = |s| \, |u|$.  We have $|u|=0$  if and
only if $u=0$, the zero function.

\begin{lemma}[Cauchy-Schwartz inequality]
    $$|u \cdot v| \le |u|\,|v|.$$
Equality holds exactly when $|v|u = \pm |u|v$.  Furthermore,
$u\cdot v = |u|\,|v|$ holds exactly when $|v| u = |u| v$.
\end{lemma}

\begin{proof}
   Let $w = |v| u \pm |u| v$.  Then expanding $w\cdot w$ we get
    $$0\le w\cdot w = 2|u|^2|v|^2 \pm 2|u|\, |v| (u\cdot v) =
    2|u|\,
    |v| (|u|\, |v| \pm (u \cdot v)).$$
    If $2|u| \,|v| = 0$, then $u$ or $v$ is zero, and the result
    easily follows.  Otherwise divide both sides of the inequality
    by the positive quantity $2 |u| \,|v|$ to get the result.
\end{proof}

\begin{lemma}[triangle inequality]
  $$
  |u + v| \le |u| + |v |.
  $$
Equality holds exactly when $|v|u = |u|v$.
\end{lemma}

\begin{proof}
Both sides are nonnegative; it is enough to compare the squares of
both sides.  By the Cauchy-Schwartz inequality,
    $$|u + v|^2 = u\cdot u + 2 u\cdot v + v\cdot v \le
      u\cdot u + 2 | u|\,|v| + v\cdot v = (|u|+|v|)^2.
    $$
The case of equality follows from the case of equality in the
Cauchy-Schwartz inequality.
\end{proof}

There is a distance function $d(u,v) = | u - v|$ that makes
$\ring{R}^\infty$ into a metric space.  The subsets $\ring{R}^n$ are
metric subspaces.






\subsection{coordinates}

Let $\{v_1,v_2,v_3,v_4\}$ be a set of vectors of cardinality four.
There are $6$ (four choose two) different pairs of vectors. We let
$y_{ij} = |v_i-v_j|$ and $x_{ij}=y_{ij}^2$.


%% WW repeated DEF.
\begin{definition}\label{def:delta}  Let 
$$
\begin{array}{lll}
\Delta(x_1,\ldots,x_6) &= x_1 x_4 (- x_1+x_2+x_3- x_4+x_5+x_6)+\\&
            x_2 x_5 (x_1- x_2+x_3+x_4- x_5+x_6)
            +x_3 x_6 (x_1+x_2- x_3+x_4+x_5- x_6)
            - \\&x_2 x_3 x_4- x_1 x_3 x_5- x_1 x_2 x_6- x_4 x_5 x_6
\end{array}
$$
\end{definition}

We write $\Delta_j$ for the $j$th partial derivative of $\Delta$. 
Let $D = \det(v_1-v_0,v_2-v_0,v_3-v_0)$.
Lemma~\ref{tarski:cm4} asserts that 
  $$
  4 D^2 = \Delta(x_{12},x_{13},x_{14},x_{34},x_{24},x_{23}).
  $$
In particular, the polynomial $\Delta$ is positive on values $x_{ij}$
that come from a simplex of positive volume.





\subsection{affine geometry}

%% WW This section is repeated verbatim in tarski.tex

\begin{definition}\label{def:aff} 
 If $S = \{v_1,v_2,\ldots,v_k\}$ 
and $S'=\{v_{k+1},\ldots,v_n\}$ are  finite sets, then
set
	$$\begin{array}{lll}
      \op{aff}\, S &= \{t_1 v_1 +\cdots t_k v_k \mid
	t_1 +\cdots+t_k = 1\}.\\
        \op{aff}_{\pm} (S,S') &= \{t_1 v_1 +\cdots t_n v_n \mid
	t_1 +\cdots+t_n = 1, \pm t_j \ge 0, \text{ for } j>k.\}.\\
        \op{aff}^0_{\pm} (S,S') &= \{t_1 v_1 +\cdots t_n v_n \mid
	t_1 +\cdots+t_n = 1, \pm t_j > 0, \text{ for } j>k.\}.\\
		\end{array}
        $$
\end{definition}

\begin{definition}  If $S = \{v_1,v_2,\ldots,v_n\}$ is a finite set
of points in $\ring{R}^3$, then
set
	$$
        \begin{array}{lll}
          \op{conv}\, S &= \op{aff}_+\, (\emptyset,S)\\
	   \op{conv}^0 S &= \op{aff}^0_+\, (\emptyset,S).\\
           \end{array}
        $$
\end{definition}


\begin{definition}
Let $S=\{v_1,\ldots,v_n\}$ be a finite set of points in 
$\ring{R}^3$.  Let $v\in\ring{R}^3$. Set
  $$\begin{array}{lll}
  \op{cone}(v,S) &= \op{aff}_+(\{v\},S)\\
  \op{cone}^0(v,S) &= \op{aff}^0_+(\{v\},S)\\
  \end{array}
  $$
\end{definition}

\begin{definition} 
Let $S$ be a finite set of points in 
$\ring{R}^3$.  Let $v\in\ring{R}^3$. Set
  $$
  \Omega(v,S) = \{x \mid |v-x| < |w-x|, \forall w\in S\setminus\{v\}\}
  $$
\end{definition}
	
\begin{definition}	
Any set of the form $\op{aff}\{v,w\}$ for some $v\ne w$ is called a 
 {\it line.}
\end{definition}

\begin{definition}  A set $S$ is collinear if there exists
a line that contains every point of $S$.
\end{definition}

\begin{definition}	
If $A=\op{aff}\{u,v,w\}$ for some set $\{u,v,w\}$ that is not collinear,
then $A$ is a {\it plane.}  Sets $\op{aff}(\{u,v\},\{w\})$, with
$\{u,v,w\}$ not collinear, are called half-planes.
\end{definition}

\begin{definition} A set $S$ is  coplanar if there exists
a plane that contains every point of $S$.
\end{definition}

\begin{definition} A set $\op{aff}_{\pm}(\{u,v,w\},\{v'\})$,
with $\{u,v,w,v'\}$ not coplanar, is called a half-space.  If
we replace $\op{aff}_{\pm}$ with $\op{aff}_{\pm}^0$ we have an
{\it open half-space}.
\end{definition}

\section{Angle}



If $u,v$ are given with $u$ and $v$ nonzero, then by the
Cauchy-Schwartz inequality,
    $$-1 \le \frac{u\cdot v}{|u|\,|v|} \le 1.$$

\begin{definition}\label{def:angle}
There is a unique $\theta$ with $0\le\theta\le\pi$ such that
    $$\cos\theta = \frac{u\cdot v}{|u|\,|v|}.$$
It is the {\it angle} formed by $u$ and $v$ at the origin. More
generally, the angle formed by $u$, and $v$ at $w$ is defined to be
the angle formed by $u-w$ and $v-w$ at the origin.
\end{definition}

By the relation between $\arccos$ and $\arctan$
(Lemma~\ref{lemma:arccos_arctan}), if $|u\cdot v|\ne |u|\,|v|$,
then the angle $\theta$ formed by $u$ and $v$ at the origin is
    \begin{equation}\label{eqn:angle}
    \theta = \frac{\pi}2 - \arctan\left (\frac {u\cdot v}{\sqrt{(|u|^2|v|^2 -
    (u\cdot v)^2)}}\right).
    \end{equation}

Let $\ups$ (the symbol is a greek upsilon, which is written with a
wider stroke than a roman vee) be the polynomial
    $$\ups(x,y,z) = -x^2 - y^2 - z^2 + 2 x y + 2 y z + 2 z x.$$
%% WW Repeated def (tarski.tex)
This polynomial factors
    $$\ups(a^2,b^2,c^2) = (-a + b + c) (a - b + c) (a + b - c) (a + b +
    c).$$
If there are vectors $v_a$, $v_b$, and $v_c$ such that $a = |v_b
- v_c|$, $b = |v_a - v_c|$, $c = |v_a - v_b|$, then $0\le
\ups(a^2,b^2,c^2)$ follows from the triangle inequality applied to
each factor.  By the case of equality in Cauchy-Schwartz
inequality, the polynomial is strictly positive unless $v_a$,
$v_b$, and $v_c$ are collinear.



\begin{lemma}[law of cosines]
    Let $\gamma$ be the angle formed by vectors $v_a$ and $v_b$ at $v_c$.  Let $a
    = |v_b - v_c|$, $b = |v_a - v_c|$, and $c = |v_a - v_b|$.  Then
        $$c^2 = a^2 + b^2 - 2 a b \cos\gamma.$$
Or in terms of $\gamma$,
    $$\gamma = \arccos(\frac{a^2 + b^2 - c^2}{2 a b}).$$
Moreover, if $v_a$, $v_b$, and $v_c$ are not collinear then
    $$\gamma=
    \pi/2 - \arctan(\frac{ a^2 + b^2 - c^2}{\sqrt{\ups(a^2,b^2,c^2)}}).
    $$
\end{lemma}\index{arc}\index{law of cosines} 

\begin{definition}
We write $\arc(a,b,c)$ for the angle $\gamma$, expressed as a function
of the lengths $a,b,c$.  Set 
  $$\arc_V(u,v,w) = \arc(|u-v|,|u-w|,|v-w|) = 
    \arccos\left(\frac{(v-u)\cdot (w-u)}{|v-u| |w-u|}\right).$$
\index{arc}
\end{definition}

\begin{proof}
    $$\begin{array}{lll}c^2 &= |v_b - v_a|^2 = ((v_b - v_c) - (v_a - v_c))\cdot ((v_b - v_c) - (v_a -
    v_c)) \\ &= a^2 + b^2 - 2 (v_b - v_c)\cdot (v_a - v_c) = a^2 +b ^2 - 2 a b
    \cos\gamma.
    \end{array}$$
    Solve this equation for $\gamma$.  The final formula
    comes by the $\arccos$-$\arctan$ relation Lemma~\ref{lemma:arccos_arctan}.  We
    have to exclude collinear vectors in the final formula to avoid
    a zero in the denominator.
\end{proof}

\begin{lemma}[law of sines]
   Let $\gamma$ be the angle formed by $v_a$ and $v_b$ at $v_c$.  Let $a
    = |v_b - v_c|$, $b = |v_a - v_c|$, and $c = |v_a - v_b|$.  Then
        $$2 a b \sin\gamma = \sqrt{\ups(a^2,b^2,c^2)}.$$
\end{lemma}

\begin{proof}  Both sides are non-negative.  It is enough to check
that their squares are equal.  By the law of cosines:
      $$4 a^2 b^2 \sin^2\gamma = 4 a^2 b^2 (1-\cos^2\gamma) = (4 a^2 b^2 - (a^2 + b^2 -
      c^2)^2) = \ups(a^2,b^2,c^2).$$
\end{proof}

\begin{remark}  Although we have not yet discussed area, we remark
that $(a b \sin\gamma)/2$ can be interpreted as the area of the
triangle with vertices $v_a$, $v_b$, and $v_c$.  The lemma then
computes the area of a triangle.
\end{remark}





\subsection{cross product}% uses angle.

\begin{definition}   Let $u =(x,y,z)$ and $u' = (x',y',z')$.  
Let the cross product be defined
by
    $$
    u \times u' = (y z' - y' z, z x' - x z', x y' - y x').
    $$
\end{definition}

\begin{lemma}  
Suppose that $u$ and $v$ form angle $\theta$ at the origin, then
    $$|u \times v| = |u|\,|v|\sin\theta.$$
\end{lemma}

\begin{proof}
   Both sides are non-negative, so it is enough to compare the
   squares of both sides.  The square of the left-hand side is
   $$
   \begin{array}{lll}
   (y z'- y'z)^2 + (z x' - x z')^2 + (x y' - y x')^2 &=
   (x^2 + y^2 + z^2)(x'^2 + y'^2 + z'^2) - (x x' + y y' + z z')^2
   \\&= |u|^2|v|^2 - (u\cdot v)^2 = |u|^2|v|^2 \sin^2\theta.
   \end{array}
   $$
\end{proof}


\begin{lemma} 
    $$
    u\times v = -v\times u,\quad
    (u\times v)\cdot w = (v\times w)\cdot u.
    $$
\end{lemma}

\begin{proof} These are elementary calculations.
\end{proof}



\subsection{dihedral angle}


\begin{definition} Assume that $w_0\ne w_1$.
We write $\dih_V(\{w_0,w_1\},\{w_2,w_3\})$ for the angle $\gamma\in[0,\pi]$
formed
by 
    $$
    v'_a = (v_c\cdot v_c) v_a - (v_c\cdot v_a) v_c\quad\text{and }v'_b =
            (v_c\cdot v_c) v_b - (v_c\cdot v_b) v_c,
    $$
where $v_a = w_2-w_0$, $v_b=w_3-w_0$ and $v_c=w_1-w_0$.  We call it
the dihedral angle formed by $w_2$ and $w_3$ along $\{w_0,w_1\}$.
\end{definition}

Up to positive scalars, $v'_a$ and $v'_b$ are the projections of
$v_a$ and $v_b$ to a hyperplane perpendicular to $v_c$.  The
dihedral angle is the angle between the projections.

The dihedral angle is unchanged if $v_c$ is replaced by $s v_c$ with
$s\ne0$. The dihedral angle is unchanged if $v_a$ is replaced with
$s v_a + t v_c$ with $0 < s$ and $t$ arbitrary.  It is unchanged if
$v_b$ is replaced with $s v_b + t v_c$ with $0 < s$ and $t$
arbitrary.  In particular, the dihedral angle formed by $v_a$ and
$v_b$ along $v_c$ is the same as that formed by $v_a/|v_a|$ and
$v_b/|v_b|$ along $v_c/|v_c|$.

The dihedral angle is degenerate and will not be used when $v_c =
0$, $v'_a = 0$, or $v'_b = 0$.



\begin{lemma}[spherical law of cosines]
Let $\gamma$ be the dihedral angle,  let $a$, $b$, and $c$ be the
angle between $v_b$ and $v_c$, $v_a$ and $v_c$, and $v_a$ and
$v_b$, respectively. Assume $v_c\ne 0$.  Assume that $v_a$ is not
a multiple of $v_c$. Assume that $v_b$ is not a multiple of $v_c$.
Then
    $$\cos\gamma = \frac{\cos c - \cos a \cos b}{\sin a\sin b}.$$
\end{lemma}

\begin{remark*}  Although we have not yet discussed spherical
triangles, we remark that $a$, $b$, and $c$ may be interpreted as
the geodesic lengths of the sides of a spherical triangle with
vertices $v_a/|v_a|$, $v_b/|v_b|$, and $v_c/|v_c|$. Also,
$\gamma$ measures the angle of the spherical triangle opposite the
side $c$.
\end{remark*}

\begin{proof}
By examining the formula giving the dihedral angle, we see that the
dihedral angle is unchanged if $v_a$, $v_b$, and $v_c$ are replaced
by $v_a/|v_a|$, $v_b/|v_b|$, $v_c/|v_c|$, respectively.  Hence
may now assume that $|v_a|=|v_b|=|v_c|=1$.

By the law of cosines, we have
        $$\cos\gamma = \frac{v'_a\cdot v'_b}{|v'_a|\,|v'_b|}$$
We compute the length of $v'_a$ (keeping in mind that $v_a$, $v_b$,
and $v_c$ all have unit norm):
        $$
        |v'_a|^2 = v'_a\cdot v'_a =
        (v_a - (v_c\cdot v_a)v_c)\cdot (v_a - (v_c\cdot v_a) v_c) =
        1 - (v_c\cdot v_a)^2 = \sin^2 b.
        $$
So $|v'_a| =\sin b$. Similarly, $|v'_b| = \sin a$.  This gives
the denominator in the spherical law of cosines. To compute the
numerator, we expand the dot product:
    $$
    v'_a\cdot v'_b = (v_a - (v_c\cdot v_a) v_c)\cdot (v_b - (v_c\cdot v_b) v_c)
    = (v_a\cdot v_b) - (v_c\cdot v_a) (v_c\cdot v_b) = \cos c - \cos
    a \cos b.
    $$
The identity follows.
\end{proof}

\begin{lemma}[spherical law of cosines - second form]
Let $v_a,v_b,v_c$ be three points in $\ring{R}^3$.
Let $\alpha,\beta,\gamma$ be the dihedral angles: 
   $$
   \begin{array}{lll}
     \alpha &= \dih_V(\{0,v_a\},\{v_b,v_c\})\\
     \beta &= \dih_V(\{0,v_b\},\{v_a,v_b\})\\
     \gamma&= \dih_V(\{0,v_c\},\{v_b,v_a\})\\
     \end{array}
   $$
Let $c$ be the
angle between $v_a$ and $v_b$. 
Assume that $\{0,v_a,v_b,v_c\}$ is not coplanar.
Then
    $$
    \cos c = \frac{\cos \gamma + \cos \alpha \cos \beta}
     {\sin \alpha\sin \beta}.
    $$
\end{lemma}

\begin{proof}
Let $A=\cos a$, $B=\cos b$, $C=\cos c$,
$A'=\sin a$, $B'=\sin b$, $C'=\sin c$.  By the
spherical law of cosines, we have
   $$\sin^2\beta = 1-\left(\frac{B-A C}{A' C'}\right)^2
     = \frac{p}{A'^2 C'^2},$$
where $p=1-A^2 - B^2 - C^2 + 2 A B C$.
In particular, $p\ge 0$.
Computing $\sin^2\alpha$ and the remaining terms in the same way, we get
   $$
   \begin{array}{lll}
     \sin\alpha\sin\beta &= \frac{p}{A' B' C'^2}\\ 
      \\
     \cos\gamma + \cos\alpha \cos\beta &=
         \frac{C - A B}{A' B'} + \frac{A - B C}{B' C'} \frac{B - A C}{A' C'}
         = \frac{p C}{A' B' C'^2}.\\
    % C &= \frac{\cos\gamma + \cos\alpha \cos\beta}{\sin\alpha \sin\beta}\\
   \end{array}
   $$
The result follows.
\end{proof}




\begin{lemma} \label{lemma:dihform} Let $v_0,v_1,v_2,v_3$ 
be vectors with $v_3\ne0$, $v'_1\ne 0$, and $v'_2 \ne 0$.  
Let $\gamma$ be the dihedral angle formed
by $v_2$ and $v_3$ along $\{v_0,v_1\}$. Let
    $$(x_1,\ldots,x_6) = 
    (x_{01},x_{02},x_{03},x_{23},x_{13},x_{12}),
    \text{ where } x_{ij}=|v_i-v_j|^2.$$
The dihedral angle $\gamma=\dih_V(\{v_0,v_1\},\{v_2,v_3\}$
is given by
    $$
    \gamma=\arccos(\frac{\Delta_4(x_1,\ldots,x_6)}{\sqrt{
    \ups(x_1,x_2,x_6)\ups(x_1,x_3,x_5)}}).
    $$
Assuming that $\gamma\ne 0,\pi$, it is also given by
    $$
    \gamma=\frac{\pi}{2} - \arctan
     (\frac{\Delta_4(x_1,\ldots,x_6)}{\sqrt{4 x_1 \Delta(x_1,\ldots,x_6)}}).
    $$
\end{lemma}
%% pi/ 2. -  atan(  deltax4/ (sqrt (4. * x1 * delta)))

\begin{proof}
Let $\beta$ be the angle between $v_1$ and $v_2$ at $0$.
    By expanding definitions and dot products
    $$
    v'_2\cdot v'_2 = (v_1\cdot v_1) ((v_1\cdot v_1)(v_2\cdot v_2) -
    (v_1\cdot v_2)^2) =  x_1^2 x_2 \sin^2 \beta = \frac{1}{4}
    x_1
    \ups(x_1,x_2,x_6).
    $$
    Similarly,
    $$v'_3 \cdot v'_3 = \frac{1}{4} x_1 \ups(x_1,x_3,x_5).$$

Let $y_i = \sqrt{x_i}$. Then
    $$\begin{array}{lll}
    v'_2\cdot v'_3 &= (v_1\cdot v_1)((v_1\cdot v_1)(v_2\cdot v_3) -
    (v_1\cdot v_2)(v_1\cdot v_3) ) \vspace{6pt} \\  &
    = x_1 \left(\frac{ x_1 (x_2 + x_3 -
    x_1)}{2} - \frac{(x_1 + x_2 - x_6)(x_1 + x_3 -
    x_5)}{4} \right)\vspace{6pt}\\&
    = \frac{x_1}{4} (2 x_1 (x_2+x_3-x_4) -
    (x_1+x_2-x_6)(x_1+x_3-x_5)) \vspace{6pt}\\&
    = {x_1\Delta_4(x_1,\ldots,x_6)}/{4}.
    \end{array}
    $$
The result follows in terms of $\arccos$.

To translate to the $\arctan$, 
we use the $\arccos$-$\arctan$ identity
and the following polynomial identity
    $$
    \frac{16}{x_1^2}(|v'_2|^2 |v'_3|^2 - (v'_2\cdot v'_3)^2) =
    \ups(x_1,x_2,x_6)\ups(x_1,x_3,x_5) - \Delta_4(x_1,\ldots,x_6)^2
    = 4 x_1 \Delta(x_1,\ldots,x_6).
    $$
\end{proof}

\subsection{beta cones }


\begin{lemma}\label{lemma:beta-cone}
Let $v_0$, $v_1$, $v_2$, $v_3$ be points in $\ring{R}^3$.   
Assume that $v_0\ne v_1,v_2,v_3$.  Assume
that $\dih_V(\{v_0,v_3\},\{v_1,v_2\})=\pi/2$.  Let
$\psi = \arc_V(v_0,v_2,v_3)$ and $\theta=\arc_V(v_0,v_1,v_2)$. 
Then $\dih_V(\{v_0,v_1\},\{v_2,v_3\}) = \beta_\psi(\theta)\in[0,\pi/2]$,
where
  $$
      \cos^2\beta_\psi = (\cos^2\psi-\cos^2\theta)/(1-\cos^2\theta).
  $$
\end{lemma}

\begin{proof}  On the unit sphere centered at $v_0$, we see that
$\dih$ is the angle on a spherical right triangle,
with opposite side of arclength $\psi$ and adjacent sides $\alpha$
and $\theta$, for some $\alpha$. 
 Since this is a spherical right triangle with hypotenuse
$\theta$, by the spherical law of cosines, we have
   \begin{equation}\label{eqn:rt}
     \cos\theta = \cos\alpha \cos\psi.
     \end{equation}
By the spherical law of cosines again, we have
    $$
    \cos\beta = \frac{\cos\psi - \cos\alpha\cos\theta}{\sin\alpha\sin\theta}.
    $$
Square this equation and eliminate $\alpha$ with Equation~\ref{eqn:rt}.
The result follows.
\end{proof}




\subsection{Euler angle sum}
\begin{lemma}[Euler triangle]\label{lemma:euler}
Let $v_0,v_1,v_2,v_3$ be points in $\ring{R}^3$. 
Let 
  $$(y_1,\ldots,y_6) =(y_{01},y_{02},y_{03},y_{23},y_{13},y_{12}),
   \text{ where } y_{ij}=|v_i-v_j|.$$
Set
$x_i = y_i^2$.   
and
    $$
    p = y_1 y_2 y_3 + y_1 (w_2\cdot w_3) + y_2 (w_1\cdot w_3) + y_3
    (w_1\cdot w_2).
    $$
where $w_i = v_i- v_0$.
Let $$\alpha_i =\dih_V(\{v_0,v_i\},\{v_j,v_k\})$$
where $\{i,j,k\}=\{1,2,3\}$.
Assume that $\Delta(x_1,\ldots,x_6)>0$. 
Then
    $$
    \alpha_1+\alpha_2+\alpha_3 - \pi
     = {\pi} - 2\arctan(\frac{2 p}{\Delta(x_1,\ldots,x_6)^{1/2}}),
    $$
for three given dihedral angles $\alpha_i$.
\end{lemma}

\begin{proof}
%% I checked all the details of this proof in 
%% Math'ca on May12,2007
The angles are unchanged if the vectors $w_i$ are rescaled so that
$|w_i|=1$.  The given formula is also unchanged under rescalings:
the factor $a$ is homogeneous of degree $3$ under a change $w_i
\mapsto t w_i$ for $t>0$, and so is $\sqrt{\Delta}$ by the
formula for $\Delta$.  Thus, we may assume that $|w_i|=1$, for
$i=1,2,3$. We have $y_1=y_2=y_3=1$.  It is convenient to use
different notation $a=x_4$, $b=x_5$, $c=x_6$ for the other
variables. Expanding the dot products in $p$ by the law of cosines
we have
    $$2 p = 8 - (a+b+c).$$

We have $$\Delta(x_1,\ldots,x_6) = \Delta(1,1,1,a,b,c) =
    \ups(a,b,c) - a b c.$$
Since $0 <\Delta$, the arctangent formula
in Lemma~\ref{lemma:dihform} 
applies for the dihedral angles $\alpha_i$. Making
this substitution and clearing the $3$ from the denominator, the
desired identity now takes the form $f(a,b,c)=0$, where
    $$
    f(a,b,c)= -\pi/2 - \sum_{i=1}^3\arctan(u_i/\sqrt{\Delta}) +
    2\arctan(2 p/\sqrt{\Delta}),
    $$
for some rational functions $u_i$ of $a,b,c$.  We prove that this
trig identity holds whenever $\Delta>0$.

Fix $b,c$ and differentiate $f$ with respect
to $a$.  The partial derivative $\partial f/\partial a$ has the form
$g(a,b,c)/\sqrt{\Delta}$ for some rational function of $a,b,c$ (with
nonzero denominator).  By brute force, we see that $g(a,b,c)$ is
identically $0$.  (Euler himself did not shun brute force.  See
\cite{Euler}.)

By symmetry, the partial derivatives with respect to $b$ and $c$ are
also identically zero.  The function $f(a,b,c)$  is continuously
differentiable whenever $\Delta>0$.  Thus, its directional
derivatives are zero.  Thus, $f(a,b,c)$ is constant on connected
sets of the domain of $\Delta>0$.  Without loss of generality,
assume that $a\le b\le c$.  We have, for $a\le t\le b$ that
$d\Delta(1,1,1,t,b,c)/dt\ge 0$. Hence $f(a,b,c)=f(b,b,c)$. 
Similarly,
 $f(b,b,c)=f(c,c,c)$. 
To complete the lemma, we evaluate the constant $f(c,c,c)$
by taking $c$ small.
$\Delta=3c^2-c^3$, $2p= 8-3c$,  $u_1=u_2=u_3 = c -c^2/2$. With this,
we compute
    $$f(a,b,c)= f(c,c,c) = \lim_{c\mapsto0} f(c,c,c) = 
    3 \arctan(1/\sqrt3)-\pi/2 =0.$$
\end{proof}






\subsection{polar coordinates}
\label{sec:polar}

For every pair of real numbers $x$ and $y$,  there are real numbers
$r$ and $\theta$ such that
    \begin{equation}\label{eqn:polar}
    x = r\cos\theta,\quad y = r\sin\theta.
    \end{equation}
If $x$ and $y$ are both zero, then we take $r=0$, and the
Equations~\ref{eqn:polar} hold for all choices of $\theta$. If $x$
and $y$ are not both zero, then we can take $0<r$, and $\theta$ is
uniquely determined (up to multiples of $2\pi$).  We can choose
$0\le\theta < 2\pi$.

Let $W=\{p_1,\ldots,p_k\}$ be a finite set of
nonzero points in the plane, with
polar coordinates $p_i = (r_i\cos\theta_i,r_i\sin\theta_i)$.
We can order the points with the lexicographic order on their
polar coordinates.  Write $p_i \prec p_j$ if
$\theta_i < \theta_j$ or ($\theta_i=\theta_j$ and $r_i<r_j$).
This is a total order on the points.
We have a cyclic permutation $\sigma:W\to W$ which sends
$p\in W$ to the next larger element with respect to this order,
or back to the first element if $p$ is the largest.
We will call $\sigma$ the {\it polar cycle}\index{polar cycle} 
of the set $W$.



For $\psi\in\ring{R}$, let $T:\ring{R}^2\to\ring{R}^2$ be the
linear transformation
   $$
   (x,y) \mapsto  (x\cos\psi + y\sin\psi,-x\sin\psi+y\cos\psi).
   $$
Let $\sigma'$ be the polar cycle for $T(W)$.  Then it is easily
checked that
$$
   \sigma'(T p) = T (\sigma p),\quad \text{ for } p\in W. 
$$

\begin{lemma}\label{lemma:polar2} 
Let $\{p_1,p_2\}\subset \ring{R}^2$.  Suppose
that both points are nonzero.  Let 
$p_i =(r_i\cos\theta_i,r_i\sin\theta_i)$, for $i=1,2$.
Let $$\theta_{ji} = \theta_i - \theta_j + 2\pi k_{ij},$$
where we pick $k_{ij}$ so that $0\le \theta_{ji}< 2\pi$.
Then 
$$
  \theta_{12} + \theta_{21} = \begin{cases}
    2\pi, & \text{ if }\theta_i\ne\theta_j\\
    0,    & \text{ if }\theta_i=\theta_j.
    \end{cases}
$$
\end{lemma}

\begin{proof} This is elementary.
\end{proof}


\begin{lemma}\label{lemma:polar-sum} 
Let $W\subset\ring{R}^2$ be a finite set,
of cardinality $n$. Suppose each point in $W$ is nonzero.
Let $\sigma$ be the polar cycle on $W$.  
Let $p=(r(p)\cos\theta(p),r(p)\sin\theta(p))$, for $p\in W$, with
$0\le\theta(p)<2\pi$.
Write
   $$
   \theta({p,q}) = \theta(q) - \theta(p) + 2\pi k_{pq},
   $$
where we choose $k_{pq}$ so that $0\le \theta({p,q}) < 2\pi$.
Let $n=\card(W)$.  Then have for all $p\in W$,
and all $0\le i \le j < n$,
   $$
   \theta(p,\sigma^i(p)) +\theta(\sigma^i(p),\sigma^j(p)) =
   \theta(p,\sigma^j(p)).
   $$
Moreover, if there exists $p,q\in W$ such that $\theta(p)\ne\theta(q)$,
we have 
  $$
  \sum_{i=0}^{n-1} \theta(\sigma^{i}p,\sigma^{i+1} p) = 2\pi.
  $$
(If $\theta(p)=\theta(q)$ for all $p,q\in W$, then all the
summands are zero.)
\end{lemma}

\begin{proof}
Fix $p$.  For $0\le i<n$, define $\theta'_i$ by
   $\theta'_0=\theta_0$ and 
   $$\theta'_i = \theta(\sigma^i(p)) + 2\pi \ell_i,$$
where we choose $\ell_i$ so that $\theta'_0\le \theta'_i < \theta'_0+2\pi$.
It follows from the definition of the polar cycle that
$\theta'_i \le \theta'_j$ for $0\le i\le j < n$. We find that
$\theta(\sigma^i p ,\sigma^j p) = \theta'_j - \theta'_i$.
The first identity becomes
  $$
  (\theta'_i-\theta'_0) + (\theta'_j-\theta'_i) = (\theta'_j-\theta'_0),
  $$
which is certainly true.
The second identity becomes
  $$
  \sum_{i=0}^{n-2} (\theta'_{i+1}-\theta'_i) + \theta(\sigma^{n-1}p,p)
  = \theta(p,\sigma^{n-1}p) + \theta(\sigma^{n-1}p,p).
  $$
This is $0$ or $2\pi$, by the previous lemma.
\end{proof}

\subsection{spherical coordinates}
\label{sec:spherical}


\begin{definition}
Let $x,y,z$ be any real numbers.  A
triple $(r,\theta,\phi)$ such that
    \begin{equation}
    \label{eqn:spherical}
    x = r\cos\theta\sin\phi,\quad y = r\sin\theta\sin\phi,\quad
    z = r\cos\phi
    \end{equation}
with $0\le r$, $0\le\theta<2\pi$, and $0\le\phi\le\pi$ are called
polar coordinates of $(x,y,z)$.  (We follow the variable
naming conventions of American calculus textbooks, which differ
from the international scientific notation.)
\end{definition}

Polar coordinates of any $(x,y,z)$ exist.
We have $r = \sqrt{x^2+y^2+z^2}$.  In the degenerate case $r=0$,
the Equations~\ref{eqn:spherical} become independent of $\theta$
and $\phi$. In the degenerate case when $\phi = 0$ or $\phi =
\pi$, the equations become independent of $\theta$. If $0<r =
\sqrt{x^2+y^2+z^2}$, and $\phi\ne 0,\pi$,  then $\phi$ is uniquely
determined by $x,y,z$. Also, $\theta$ is uniquely determined.


\begin{remark}\label{def:azimuth}
We call $\theta$ the {\it azimuth angle\/} and $\phi$ the {\it
zenith angle\/} of $(x,y,z)$.  The azimuth angle is also known as
the longitude.  The zenith angle is also known as the latitude. The
azimuth angle is a polar coordinate of $(x,y)$:
    $$
    (x,y) = (r'\cos\theta,r'\sin\theta), \quad r' = r\sin\phi.
    $$
\end{remark}

\subsection{azimuth angles}

We generalize azimuth angle.

\begin{definition} Let $(v,w)$ be an ordered pair of points in
$\ring{R}^3$, with $v\ne w$.  Let $e_1$ and $e_2$ be orthonormal
vectors
such that $e_i\cdot (v-w)=0$ and $\det(e_1,e_2,w-v)>0$.  Every
$u\in\ring{R}^3$ that is not in the line $\op{aff}(v,w)$
can then be uniquely expressed as
   $$
   u = r\cos\theta\, e_1 + r\sin\theta\, e_2 + h (w-v),
   $$
for some $0< r$, $0\le \theta < 2\pi$, $h\in\ring{R}$
(cylindrical coordinates).  Assume that $w_1$ and $w_2$ do
not lie in the line $\op{aff}(v,w)$.
We define $\op{azim}(v,w,w_1,w_2)$ to be the unique $\theta$ such
that $0\le\theta < 2\pi$ and
  $$
  \begin{array}{lll}
    w_1 &= r_1\cos\psi\, e_1 + r_1\sin\psi\, e_2 + h_1(w-v),\\
    w_2 &=  r_2\cos(\psi+\theta)\, e_1 + r_2\sin(\psi+\theta)\, e_2 
     + h_2(w-v),\\
\end{array}
  $$
\end{definition}

\begin{definition} Let $(v,w)$ be an ordered pair of points in
$\ring{R}^3$, with $v\ne w$.
Let $W$ be a finite set of points in $\ring{R}^3$.
We say that $W$ is nondegenerate with respect to $(v,w)$ if
the following two conditions hold.
First, $p = q + h (w-v)$, with $p,q\in W$ and $h\in \ring{R}$
implies that $p=q$.  Second, 
  $p\ne h(w-v)$ for $p\in W$, $h\in\ring{R}$.
\end{definition}

Nondegeneracy is precisely the condition we need for the
set $W$ to map injectively to the nonzero points of the 
plane under the projection
onto the first two cylindrical coordinates.

\begin{definition}
Let $(v,w)$ be an ordered pair of points in
$\ring{R}^3$, with $v\ne w$.
Let $W$ be a finite set of points in $\ring{R}^3$ that is
nondegenerate  with respect to $(v,w)$.
Pick $e_1$ and $e_2$ as in the definition
of $\op{azim}$.  Map 
$W$ bijectively to $W'\subset\ring{R}^2$ by the following map $\phi$:
   $$p =  r\cos\theta\, e_1 + r\sin\theta\, e_2 + h (w-v) \mapsto^\phi
     (r\cos\theta,r\sin\theta).$$
These conditions give a polar cycle $\sigma':W'\to W'$.  Let
$\sigma:W\to W$ be given by $\phi\sigma(p) =\sigma'\phi(p)$.
Call $\sigma$ the azimuth cycle
on $W$ with respect to $(v,w)$.
\index{azimuth cycle} 
\end{definition}

\begin{lemma} The polar cycle $\sigma:W\to W$ does not depend
on the choice of $e_1$ (and $e_2$ is determined by $e_1$).
\end{lemma}

\begin{proof} This follows from independence of $\sigma'$ from
rotations of the coordinate axes $e_1$ and $e_2$.
\end{proof}


\begin{lemma} Let $(v,w)$ be an ordered pair of points in $\ring{R}^3$,
with $v\ne w$.  Assume that $\{w_1,w_2\}$ is nondegenerate
with respect to $(v,w)$.  Then
  $$
  \op{azim}(v,w,w_1,w_2) + \op{azim}(v,w,w_1,w_2) 
  = \begin{cases} 2\pi, & \text{if }\op{azim}(v,w,w_1,w_2)\ne 0,\\
    0, & \text{if }\op{azim}(v,w,w_1,w_2)=0.
    \end{cases}
    $$
\end{lemma}

\begin{proof} This follows immediately from Lemma~\ref{lemma:polar2}.
\end{proof}

\begin{lemma} Let $(v,w)$ be an ordered pair of points in $\ring{R}^3$,
with $v\ne w$.  Let $W$ be a finite set in $\ring{R}^3$ that
is nondegenerate with respect to $(v,w)$,
with azimuth cycle $\sigma$.
Let $n=\card(W)$.  Then have for all $p\in W$,
and all $0\le i \le j < n$,
   $$
   \op{azim}(v,w,p,\sigma^i(p)) +
    \op{azim}(v,w,\sigma^i(p),\sigma^j(p)) =
   \op{azim}(v,w,p,\sigma^j(p)).
   $$
Moreover, if there exists $q\in W$ such that 
$\op{azim}(u,v,p,q)\ne0$,
we have 
  $$
  \sum_{i=0}^{n-1} \op{azim}(\sigma^ip,\sigma^{i+1}p) = 2\pi.
  $$
(If $\op{azim}(u,v,p,q)$ for all $q\in W$, then all the
summands are zero.)
\end{lemma}

\begin{proof} This follows immediately from 
Lemma~\ref{lemma:polar-sum}.
\end{proof}


The azimuth and dihedral angles are closely related.  For simplicity,
we will take our base point $v=0$.

\begin{lemma}  Let $w\ne 0$ be a nonzero vector in $\ring{R}^3$.
  Assume that $v_1$ and $v_2$ do not lie in the line $\op{aff}(0,w)$.
Let
  $$\gamma = \dih_V(\{0,w\},\{v_1,v_2\}).$$
  Then
    $$
    \cos(\op{azim}(0,w,v_1,v_2)) = \cos\gamma.
    $$
\end{lemma}

\begin{proof}
Let $v_i' = (w\cdot w) v_i - (w\cdot v_i) w$.  
We have $v_1'\ne 0$.  Set $e_1 = v_1'/|v'_1|$.  Choose a unit vector
$e_2$ so that $\det(e_1,e_2,w)>0$ and $e_1\cdot e_2 = w\cdot e_2=0$.
Write $v_i$ in cylindrical coordinates as 
   $$
   \begin{array}{lllll}
     v_1 &= r_1 e_1 &    &+h_1 w\\
     v_2 &= r_2 \cos\theta\, e_1 &+ r_2 \sin\theta\, e_2 &+ h_2 w.
    \end{array}
   $$
By the definition of $\op{azim}$, we have $\op{azim}(0,w,v_1,v_2)=\theta$.  By definition, $\cos\gamma$ is the angle between $v_1'$ and $v_2'$.
We compute
   $$
   \begin{array}{lll}
     v_1' &= |v'_1| e_1 \\
     v_2' &= (w\cdot w) r_2 \cos\theta\, e_1 
       &+ (w\cdot w) r_2 \sin\theta\, e_2 \\
     \end{array}
   $$
The result is now a result of the definition of angle 
(Definition~\ref{def:angle}).
\end{proof}







  
    %% Quadratic Volumes
%% File Created 3/22/07.

\chapter{Volume}

\section{Measure}
\indy{Index}{measure}

Nowhere does this book need a notion
of integration.  Measure alone suffices.  (However, there are a few
volumes described below that I do not see how to calculate without
first writing them as an integral.)

This book needs a concepts of null set, measurable set, and volume in
three dimensions.  For the purposes of this book, one can take the
the three dimensional Lebesgue measure.   
The null sets can be defined
to be the sets of zero Lebesgue measure. The measurable sets can
be defined as the bounded Lebesgue measurable sets.  The volume of
a measurable set can be defined as its Lebesgue measure.
As it will become clear in a moment, one needs considerably less than Lebesgue measure.
Most of the volumes to be considered were known in antiquity.
\indy{Index}{measure!Lebesgue}


\subsection{null set}\label{sec:null}
\indy{Index}{null set}

Assume a notion of null set with the following
properties.

\begin{enumerate}%[Null Set]
\item A finite union of null sets is a null set.\\
 \item A plane is a null set.\\
 \item A sphere is a null set.\\
 \item A circular cone is a null set; that is, a union of all
  lines through a fixed point $P$ and forming fixed
 forming fixed angle with a line through $P$.
\tlabel{enum:null}
\indy{Index}{null set}
\indy{Index}{plane}
\indy{Index}{sphere}
\indy{Index}{circular cone}
\end{enumerate}
\indy{Notation}{P@$P$ (point)}


Write $X\equiv Y$ if sets $X$ and $Y$ are equal up to a null set.
That is, there exists a null set $Z$ such that
   $(X\setminus Y) \cup (Y\setminus X) \subset Z$.
\indy{Index}{null set}
\indy{Notation}{1@$\equiv$}

\subsection{measurability}\label{sec:measure}
\indy{Index}{measure}

Assume a notion of measurability that has the following properties.

\begin{enumerate}%[Measurable set]
 \item Null sets are measurable.
 \item The union of two measurable sets is measurable.
 \item The intersection of two measurable sets is measurable.
 \item The difference of two measurable sets is measurable.
 \item Primitive regions are measurable (Definition~\ref{def:primitive}).
\tlabel{enum:measure}
\indy{Index}{null set}
\indy{Index}{primitive regions}
\end{enumerate}

\subsection{volume}\label{sec:volume}

Certain primitive regions are given in Definition~\ref{def:primitive}.
Assume a notion of volume that has the following properties.
\indy{Index}{primitive}

\begin{enumerate}%[Volume]
 \item The volume is defined for every measurable set.  It is
    a non-negative real number.
 \item The volume of a null set is $0$.
 \item If $X$ and $Y$ are  measurable, and if
 the symmetric difference of
 $X$ and $Y$ is contained in a null set, then 
    $X$ and $Y$ have the same volume.
 \item If $X$ and $Y$ are measurable sets, and if $X\cap
 Y$ is contained in a null set, then
    $$
    \op{vol}(X\cup Y) = \op{vol}(X) + \op{vol}(Y).
    $$
\indy{Notation}{X@$X$ (set)}
\indy{Notation}{Y@$Y$ (set)}
  \item (linear stretch) For $X\subset \ring{R}^3$, $v\in\ring{R}^3$, 
    set 
      $$
      T_v(X) = \{ (v_1u_1,v_2u_2,v_3u_3) \mid u\in X\}.
      $$
\indy{Notation}{Tv@$T_v$ (linear stretch)}
    If $X$ and $T_v(X)$ are measurable, then
     $\op{vol}(T_v(X)) = |v_1v_2v_3|\op{vol}(X)$.
  \item (translation) If $X\subset \ring{R}^3$ and $v\in\ring{R}^3$, then let
    $X+v = \{p + v\mid p\in X\}$.  If $X$ is measurable, then $X+v$ is
    as well, and $\op{vol}(X) = \op{vol}(X+v)$.
  \item (primitives) If $X$ is a primitive, then $X$ is measurable
    and $\op{vol}(X)$ is given by the formulas of Lemma~\ref{lemma:prim-volume}.
\tlabel{enum:volume}
\end{enumerate}
\indy{Index}{volume}
\indy{Index}{linear!stretch}
\indy{Index}{translation}
Volume calculations of primitive regions are given in 
Lemmas~\ref{lemma:prim-volume} and~\ref{lemma:wedge-vol}.
%% NB Don't need lemma:wedge-sol because solid of a \op{FR} is a \op{SC}.
%% Don't need lemma:prim-sol, because solds are \op{SC} or \op{ST}.

\begin{lemma}\guid{ATOAPUN}\oldrating{300}
\formalauthor{John Harrison}  
Let $\op{vol}$ be the Lebesgue measure
on $\ring{R}^3$. Let $B$ be the collection of all bounded Lebesgue
measurable subsets of $\ring{R}^3$.  Let $N$ be the collection of all
Lebesgue measurable subsets of $\ring{R}^3$ of measure zero.
Then
\begin{itemize}
\item  The collection $N$ satisfies the null set properties of Section~\ref{sec:null}
\item The collection $B$ satisfies the measurability properties of Section~\ref{sec:measure}
\item  $\op{vol}$ on $B$ satisfies the volume properties of Section~\ref{sec:volume}.  
\indy{Index}{measure!Lebesgue}
\end{itemize}
\end{lemma}
\indy{Notation}{B@$B$ (Lebesgue measurable subsets)}
\indy{Notation}{vol@$\op{vol}$}

\subsection{radial set and solid angle}\label{sec:solid}

Surface integrals are not required in this book.  Although
the `solid angle' is traditionally defined as a surface integral,
it is simpler to give an alternative definition based on volume.


\begin{definition}[open ball]  The open ball $B(p,r)$ with center $p$ and
radius $r$ is the set
    $$
    \{ y\in\ring{R}^3 \mid \norm{p}{y} < r.\}
    $$
\indy{Index}{open ball}
\indy{Notation}{open ball@$B(p,r)$}
\end{definition}



\begin{definition}[radial]
    A set $C$ is $r$-radial at center $p$ if  $C\subset B(p,r)$
    and if
        $p + u \in C$ implies
        $p + t u \in C$ for all $t$ satisfying $0 < \normo{u} t < r$.
A set $C$ is eventually radial at center $p$ if $C\cap B(p,r)$ is
$r$-radial at center $p$, for some $r>0$.
\indy{Index}{radial}
\end{definition}
\indy{Notation}{C@$C$ (set)}

\begin{lemma}\guid{PUACSHX}\tlabel{lemma:r-r'}\oldrating{140}\formalauthor{Nguyen Tat Thang, Nov 2008}
Assume that $C$ is measurable and $r$-radial at $p$.  Let $0\le r'<r$,
then $C\cap B(p,r')$ is measurable and
$\op{vol}(C\cap B(p,r')) = \op{vol}(C) (r'/r)^3$.
\end{lemma}

\begin{proof}  The radial set $C$ transforms into $C\cap B(p,r')$ by
a series of translations and stretch transformations.
\end{proof}


\begin{definition}[solid angle]\tlabel{def:sol}
When $C$ is measurable and eventually radial at center $p$, 
define the solid angle of $C$ at $p$ to be
    $$
    \op{sol}(p,C) = 3 \op{vol}(C\cap B(p,r))/r^3,
    $$
where $r$ is as in the definition of eventually radial. 
By Lemma~\ref{lemma:r-r'}, this
definition is independent of any such $r$.  When the center $p$ is
clear from the context, write $\op{sol}(C)$ for
$\op{sol}(p,C)$.
\indy{Index}{solid angle}
\indy{Notation}{solid angle@$\op{sol}$}
\end{definition}



The following properties follow immediately from the definitions.
If $C$ is $r$-radial for some $r>0$ then it is eventually radial.
If $C$ is measurable and $r$-radial, then the volume of $C$
satisfies
    $$
    \op{vol}(C) = \op{sol}(C) r^3/3.
    $$
If $C$ is bounded away from $p$, then $C$ is eventually radial at
$p$, and $\op{sol}(C) = 0$.
\indy{Index}{radial}

\begin{lemma}\guid{KODOBRF}\rating{60}  If $C$ and $C'$ are  $r$-radial
at $p$, then $C\cap C'$ is also $r$-radial at
$p$.
\end{lemma}






\section{Primitive Volume}
\indy{Index}{volume!primitive}

This book accepts
certain elementary volume calculations as axioms.  
These regions will be called primitive volumes.  There are only
a few primitive volumes.
All further
volumes calculations will be obtained from these through the basic
properties of measure.   
This treatment of volume is hardly about measure at all.  The focus is rather on the geometry of the various regions and how to decompose them into primitives.

To fix a convention, this book prefers to take the volume of open sets whenever that can be arranged.  It begins with a description of some of the primitive regions.

\subsection{wedge}
\indy{Index}{wedge}


The set $\op{aff}^0_+(\{v_0,v_1\},\{v_2,v_3\})$ was defined
in Definition~\ref{def:aff}.  Call it a lune.  It is the intersection
of two open half-spaces
    $$
    \op{aff}^0_+(\{v_0,v_1\},\{v_2,v_3\})
    =\op{aff}^0_+(\{v_0,v_1,v_2\},v_3)\cap
    \op{aff}^0_+(\{v_0,v_1,v_3\},v_2)
    $$
\indy{Notation}{aff0@$\op{aff}^0$ (lune)}


A lune has a dihedral angle $\dih(\{v_0,v_1\},\{v_2,v_3\})$ between
$0$ and $\pi$.   For angles that are larger than $\pi$,   a wedge
$W(v_0,v_1,w_1,w_2)$ is used instead of the lune.  Assume that $v_0\ne v_1$ and that
$w_1$ and $w_2$ do not lie on
the line $\op{aff}\{v_0,v_1\}$.  Set
$$
W(v_0,v_1,w_1,w_2) = 
  \{p\not\in\op{aff}\{v_0,v_1\} \mid 
  0< \op{azim}(v_0,v_1,w_1,p) < \op{azim}(v_0,v_1,w_1,w_2)\}.
$$
When the angle is less than $\pi$, there is no difference between
a wedge and a lune:
\indy{Index}{lune}
\indy{Index}{half-space}
\indy{Index}{angle!dihedral}
\indy{Index}{wedge}
\indy{Notation}{dih@$\dih$ (dihedral angle)}
\indy{Notation}{W@$W$ (wedge)}

\begin{lemma}\guid{VICUATE} Let $\{v_0,v_1,v_2,v_3\}$ be a set of four points
in $\ring{R}^3$.  Assume that the set is not coplanar.
Assume that $\op{azim}(v_0,v_1,v_2,v_3)<\pi$.
Then,
   $$W(v_0,v_1,v_2,v_3) = \op{aff}^0_-(\{v_0,v_1\},\{v_2,v_3\}).$$
\indy{Index}{coplanar}
\end{lemma}


\subsection{solid triangle}
\indy{Index}{solid triangle}

\begin{definition}[solid triangle] The solid triangle $\op{ST}(v_0,S,r)$ is
determined by  $v_0\in\ring{R}^3$, a finite subset $S\subset\ring{R}^3$,
and a radius $r\ge0$. 
    $$
    \op{ST}(v_0,S,r) = 
    B(v_0,r)\cap \op{cone}(v_0,S).
    $$
(Usually, $S=\{v_1,v_2,v_3\}$, and this accounts for
the name.)
\indy{Index}{solid triangle}
\indy{Notation}{solid triangle@$\op{ST}$ (solid triangle)}
\indy{Notation}{r@$r$ (radius)}
\end{definition}



\subsection{conic cap}
\indy{Index}{conic cap}

% renamed from spherical cap.

\begin{definition}[conic cap]
The conic cap $\op{SC}(v_0,v_1,r,a)$ is specified by an apex
$v_0\in\ring{R}^3$,  radius $r\ge0$,  non-zero vector $v_1-v_0$ giving
direction, and constant $a$.  The conic cap is the intersection of
the ball $B(v_0,r)$ with a solid right-circular cone:
    $$
    \op{SC}(v_0,v_1,r,a)=\{y \in B(v_0,r) \mid (y-v_0)\cdot (v_1-v_0) > \norm{y}{v_0}\, \norm{v_1}{v_0}\, a\}.
    $$
\indy{Index}{conic cap}
\indy{Index}{apex}
\indy{Index}{cone}
\indy{Notation}{conic cap@$\op{SC}$ (conic cap)}
\end{definition}

\subsection{frustum}
\indy{Index}{frustum}

\begin{definition}[rcone]\tlabel{def:p:rcone}
%\tlabel{def:rcone} % from Tarski collection..
\indy{Index}{cone!right-circular}
Define the following collection of right-circular cones.
If $v$ and $w$ are points in $\ring{R}^3$, and
  $h\in\ring{R}$, then set
  $$\begin{array}{lll}
    \op{rcone}(v,w,h) &= \{p\mid (p-v)\cdot (w-v) \ge \norm{p}{v}\,\norm{w}{v} h\},\\
    \op{rcone}^0(v,w,h) &= \{p\mid (p-v)\cdot (w-v) > \norm{p}{v}\,\norm{w}{v} h\}.\\
    %\op{rcone}_-^0(v,w,h) &= \{p\mid (p-v)\cdot (w-v) < \norm{p}{v}\,\norm{w}{v} h\}.\\
    %\partial\op{rcone}(v,w,h) &= \{p\mid (p-v)\cdot (w-v) = \norm{p}{v}\,\norm{w}{v} h\}.\\
    \end{array}
    $$
\end{definition}
\indy{Notation}{rcone@$\op{rcone}$}
\indy{Notation}{rcone2@$\op{rcone}^0$}
%\indy{Notation}{rcone3@$\op{rcone}_-^0$}
%\indy{Notation}{rcone4@$\partial\op{rcone}$}
\indy{Notation}{v@$v$ (point)}
\indy{Notation}{w@$w$ (point)}


\begin{definition}[frustum, FR] The frustum
$\op{FR}(v_0,v_1,h',h,a)$ is specified by an apex $v_0\in\ring{R}^3$, heights
$0\le h'\le h$, a vector $v_1-v_0$ giving its direction, and
$a\in[0,1]$. The set $\op{FR}$ is given as
    $$
    \{ y \in\op{rcone}^0(v_0,v_1,a) \mid \ 
       h'\norm{v_1}{v_0} < (y-v_0) \cdot (v_1-v_0) < h\norm{v_1}{v_0} \}.
    $$
\indy{Index}{frustum}
\indy{Notation}{frustum@$\op{FR}$ (frustum)}
\indy{Notation}{apex@$v_0$ (apex)}
\indy{Notation}{height@$0\le h' le h$ (height)}
\indy{Notation}{v1@$v_1$ (vector)}
\indy{Notation}{v0@$v_0$ (vector)}
\end{definition}

That is, the frustum is the part of a right-circular cone between two
parallel planes that cut the axis of the cone at a right angle.
When $h'=0$, the frustum extends to the apex of the cone.
Write $\op{FR}(v_0,v_1,h,a)=\op{FR}(v_0,v_1,h',h,a)$.

\subsection{tetrahedron}
\indy{Index}{tetrahedron}

\begin{definition}[tetrahedron] A tetrahedron is a set of the form
$$\op{conv}^0\{v_1,v_2,v_3,v_4\}.$$
\indy{Index}{tetrahedron}
\end{definition}

By Tarski arithmetic, %tarski: {hedra-tope}, 
this set can also be described
as the intersection of four open half-spaces, with each bounding
plane defined by three of the four points.
By taking this into account, it follows that
the sets in this section have all been defined by linear and quadratic
constraints.
\indy{Index}{Tarski arithmetic}

\subsection{primitive}
\indy{Index}{primitive}

\begin{definition}[primitive]\label{def:primitive} 
A primitive region is any of the following.

\begin{enumerate}%%[Primitive Volumes]
 \item A solid triangle $\op{ST}$.
\indy{Notation}{ST@$\op{ST}$ (triangle)}
 \item A tetrahedron $conv^0\{v_0,v_1,v_2,v_3\}$.
 \item A wedge of a frustum;
that is, the intersection of a frustum with
 a wedge:
    $$
     \op{FR}(v_0,v_1,h_1,h_2,a) \cap W(v_0,v_1,v_2,v_3).
    $$
\item A wedge of a conic cap; that is, the intersection of a conic cap
with
    a wedge:
    $$
    \op{SC}(v_0,v_1,r,c) \cap W(v_0,v_1,v_2,v_3).
    $$
\item a rectangle
    $$
    \{u \mid a_i < u_i < b_i,\text { for } i=1,2,3.\}.
    $$
\item A wedge of a ball
  $$
  B(v_0,r) \cap W(v_0,v_1,v_2,v_3).
  $$
%\item An ellipsoid
 %  $$
  % \{v \mid t_1 x_1^2 + t_2 x_2^2 + t_3 x_3^2 < r\},
  % $$
%where $t_1,t_2,t_3,r>0$.
\indy{Index}{solid triangle}
\indy{Index}{tetrahedron}
\indy{Index}{wedge}
\indy{Index}{frustum}
\indy{Index}{conic cap}
\indy{Index}{rectangle}
\indy{Index}{ball}
%\indy{Index}{ellipsoid}
\tlabel{enum:volume-prim}
\end{enumerate}
One could find a smaller set of primitives.  In particular, the
rectangle and wedge of a ball can be obtained from the others.  But
this would only make the definition more difficult to apply.

\end{definition}

\subsection{primitive volume calculation}\label{sec:primitive}

\begin{lemma}\guid{PAZNHPZ}\tlabel{lemma:prim-volume} 
\begin{enumerate} 
 \item \rating{250} Let $v_1,v_2,v_3$ be unit vectors.
   A solid triangle $\op{ST}(v_0,\{v_1,v_2,v_3\},r)$ is measurable and has volume
   $$
   (\alpha_{123}+\alpha_{231}+\alpha_{312}-\pi)r^3/3,
   $$
   where $\alpha_{ijk} = \dih_V(\{v_0,v_i\},\{v_j,v_k\})$.
  \item \rating{250} The conic cap $\op{SC}(v_0,v_1,r,a)$ is measurable and has volume:
   $$
    2\pi(1-a) r^3/3,
   $$
 \item \rating{250} A frustum $\op{FR}(v_0,v_1,h,a)$ (with $h'=0$) 
is measurable and has volume:
   $$
   \pi (t^2-h^2) h/3,\quad h = t a, 
   $$
 \item\rating{250} A tetrahedron $\op{conv}^0\{v_1,v_2,v_3,v_4\}$ is measurable and has volume:
   $$
   \sqrt{\Delta(x_{12},x_{13},x_{14},x_{34},x_{24},x_{23})}/12,
   $$
   where $x_{ij} = \norm{v_i}{v_j}^2$.
\item\rating{100} A rectangle $\op{rect}(a,b)$ is measurable and
has volume $0$ unless $a_i<b_i$ for all $i$.  In this case, the
volume is
$$(b_3-a_3)(b_2-a_2)(b_1-a_1).$$
 \item\rating{250} The intersection of a ball $B(v_0,r)$ with a wedge
 $W(v_0,v_1,v_2,v_3)$ is measurable.  The volume of the intersection
is 
   $$
   2 \theta r^3/3,\text{ for } r \ge 0,
   $$
where $\theta = \op{azim}(v_0,v_1,v_2,v_3)$.

\end{enumerate}
\end{lemma}

Euler's formula (Lemma~\ref{lemma:euler}) gives an
equivalent expression for $(\alpha_{123}+\alpha_{231}+\alpha_{312}-\pi)$.
Euler's formula will often be used instead of this formula.
\indy{Index}{Euler's formula}

\begin{proof}
The formula for the volume of a solid triangle is $r^3/3$ times
its solid angle.  The formula 
   $$\alpha_{123}+\alpha_{231}+\alpha_{312}-\pi$$
for the area of a spherical triangle is classical.    
The conic cap volume is
$r^3/3$ times its solid angle.  
The volume of a right-circular cone is $1/3$ its base times height.
The volume of a tetrahedron is
   $$|\det(v_2-v_1,v_3-v_1,v_4-v_1)|/6.$$
By Tarski arithmetic, %tarski: {cayley-menger-pos} 
the square of this determinant is given by a polynomial
$\op{CM}_4(x_{ij})$ in the squares of the edge-lengths. This is
$\Delta/4$, with $\Delta\ge0$.  The result follows.
\end{proof}
\indy{Notation}{CM@$\op{CM}_4$ (square of edge lengths)}



\subsection{wedge}\label{sec:wedge}
\indy{Index}{wedge}

When the region is realized by revolution along an axis $\op{aff}\{v_0,v_1\}$, 
one can also give the volume of the intersection of the region
with a wedge $W=W(v_0,v_1,v_2,v_3)$.  These intersections are measurable
by the last statement of Lemma~\ref{lemma:prim-volume} (with $r$ sufficiently
large).
  In the following
let $\theta = \azim(v_0,v_1,v_2,v_3)$.
\indy{Notation}{aff@$\op(aff)$}
\indy{Notation}{ZZtheta@$\theta$ (azim)}

\begin{lemma}\guid{DFNVMFM}\tlabel{lemma:wedge-vol}  Let $C$ be $\op{SC}(v_0,v_1,r,a)$, $B(v_0,r)$, or
   $\op{FR}(v_0,v_1,h,a)$.  Let $m$ be the volume of $C$.  
   Then $C\cap W$ has volume $m\,\theta/(2\pi)$.   
\end{lemma}
\indy{Notation}{C@$C$ (wedge)}

\begin{lemma}\guid{FMSWMVO}\tlabel{lemma:wedge-sol}  Let $C$ be either $\op{SC}(v_0,v_1,r,a)$, $B(v_0,r)$, or
   $\op{FR}(v_0,v_1,h,a)$.  Then $C$ and $C\cap W$ are eventually 
radial at $v_0$. Furthermore,
    $C\cap W$ solid angle 
  $s\,\theta/(2\pi)$, where $s$ is the solid angle of $C$.
\end{lemma}


\begin{proof}
These are elementary integrals.
\end{proof}


\subsection{primitive solid angle}

All of the primitive sets are eventually radial at the natural
base point $v_0$, and have a
solid angle.  By Lemma~\ref{lemma:wedge-sol}, it is enough to compute
the solid angle before intersecting with a wedge.
\indy{Notation}{v0@$v_0$ (base point)}

\begin{lemma}\guid{FUPXNLC} \tlabel{lemma:prim-sol}
\begin{enumerate}
    \item  $\op{ST}(v_0,\{v_1,v_2,v_3\})$ is eventually radial at $v_0$
     with solid angle 
     $$
     (\alpha_{123}+\alpha_{231}+\alpha_{312}-\pi),\quad
     \alpha_{ijk}=\dih_V(\{v_0,v_i\},\{v_j,v_k\}).
     $$
    \item $\op{conv}^0\{v_0,v_1,v_2,v_3\}$ is eventually radial at $v_0$
      with solid angle
           $$
     (\alpha_{123}+\alpha_{231}+\alpha_{312}-\pi),\quad
     \alpha_{ijk}=\dih_V(\{v_0,v_i\},\{v_j,v_k\}).
     $$
    \item $\op{SC}(v_0,v_1,r,a)$ is eventually radial at $v_0$ with solid
      angle 
      $2\pi(1-a)$.
    \item $\op{FR}(v_0,v_1,h,a)$ is eventually radial at $v_0$ with solid
      angle
        $$
        2\pi (1-a).
        $$
\indy{Index}{eventually radial}
\end{enumerate}
\end{lemma}
\indy{Notation}{conv0@$\op{conv^0}$}
\indy{Notation}{SC@$\op{SC}$}
\indy{Notation}{FR@$\op{FR}$}
\indy{Notation}{ZZalpha@$\alpha$ (angle)}
\indy{Notation}{dih@$\dih$}

\begin{proof} In every case, the intersection of 
  the region with $B(v_0,r')$, for $r'>0$ sufficiently small, is
  a conic cap or a solid triangle.  These two volumes have
  already been calculated.  This gives the results as stated.
\end{proof}


\section{Scissor and Volume}
\tlabel{sec:measure-second}

There are many other volumes that can be computed from the
primitive ones enumerated in Definition~\ref{enum:volume-prim}.

\subsection{lune}  
\indy{Index}{lune}

Consider the simple example of a derived volume given by a
lune $A=\op{aff}_+^0(\{v_0,v_1\},\{v_2,v_3\})$.  It is eventually
radial at $v_0$ and has a solid angle.
\indy{Notation}{A@$A$ (lune)}

\begin{lemma}\guid{WFFVGVF}  $\sol(v_0,A) = 2\dih_V(\{v_0,v_1\},\{v_2,v_3\})$.
\end{lemma}

\begin{proof}
Let $B_{\pm} = \op{aff}_\pm(\{v_0,v_2,v_3\},v_1)$.  $B_- \cap B_+$
is a null set.  The intersections $A\cap B_{\pm}\cap B(v_0,r)$ 
are solid triangles.  This gives the solid angle of $A$ as
follows:
   $$\begin{array}{lll}
   \sol(v_0,A) &= \sol(v_0,A\cap B_+)+\sol(v_0,A\cap B_-) \\
   &= 
   \sol(\op{ST}(v_0,\{v_1,v_2,v_3\})) + \sol(\op{ST}(v_0,\{v_1',v_2,v_3\})) \\
   &=
   2\dih_V(\{v_0,v_1\},\{v_2,v_3\}).
   \end{array}
   $$
Here,  $v_1'= 2 v_0 - v_1$, the reflection of $v_1$
through $v_0$.  The usual calculation of the volume of a solid triangle
inverts this proof, 
and derives the volume from the solid angle of a lune.
\end{proof}



\begin{lemma}\guid{CZOQFNL}\tlabel{lemma:wedge:sol} 
Assume that the sets $\{v_0,v_1,w_1\}$ and
$\{v_0,v_1,w_2\}$ are not collinear. 
$$\sol(v_0,W(v_0,v_1,w_1,w_2)) = 2\op{azim}(v_0,v_1,w_1,w_2).$$
\indy{Index}{collinear}
\end{lemma}    

\begin{proof} Every wedge is a union of two lunes, up to a null set.
\end{proof}

\begin{lemma}\guid{OXQZBJG}  
   $$
   \begin{array}{lll}
    \op{sol}(v_0,B(v_0,t)) &= 4\pi\\
    \op{vol}(v_0,B(v_0,t)) &= 4\pi t^3/3\\
    \op{sovo}(v_0,B(v_0,t),\lambda) &= 4 \pi \phi(t,t,\lambda).
   \end{array}
   $$
\end{lemma}

\begin{proof}
The ball is $t$-radial at $v_0$, so the volume is given by
Definition~\ref{def:sol} in terms of solid angle.  It is enough
to check that a hemisphere has solid angle $2\pi$.  This follows
from Lemma~\ref{lemma:wedge:sol}.
\end{proof}  




\subsection{Rogers simplex}
\indy{Index}{Rogers simplex}

\begin{definition}[orthosimplex, orth] \tlabel{def:ortho}
An {\it orthosimplex} is a tetrahedron
    $$\op{conv}^0(p,p+v_1,p+v_1+v_2,p+v_1+v_2+v_3),$$
where $v_i\cdot v_j=0$, for $1\le i<j\le 3$.   Write
$\op{orth}^0(p,v_1,v_2,v_3)$ for this orthosimplex.
 \indy{Index}{orthosimplex}
\indy{Notation}{orthosimplex@$\op{orth}^0$}
\end{definition}

\begin{figure}[htb]
  \centering
  %\myincludegraphics{\ps/rogers.eps}
  \caption{The Rogers simplex is an orthosimplex.}
  \tlabel{fig:rogers}
\end{figure}


\begin{definition}[Rogers simplex, rog] \tlabel{def:rog}
Let $\{v_0,v_1,v_2,v_3\}$ be a set of four points in $\ring{R}^3$.
Assume that they are not coplanar.  Let $p$ be the circumcenter
of $\{v_0,v_1,v_2\}$ and $r$ its circumradius (see Definition~\ref{def:etaV}).  Let $c\ge r$.
By Tarski arithmetic, %tarski: {rog-exist}, 
there exists a unique point $p'$ in $A=\op{aff}_+(\{v_0,v_1,v_2\},v_3)$ at equal distance $c$
from $v_0,v_1,v_2$.
\indy{Notation}{p@$p$ (circumcenter)}
Let 
$$
    \op{rog}^0(v_0,v_1,v_2,v_3,c) = 
    \op{orth}^0(v_0,w_1,w_2,w_3),
$$
$$
     w_1=(v_0+v_1)/2,\quad w_1+w_2=p,\quad w_1+w_2+w_3=p'.
$$
(Define $\op{rog}(v_0,\ldots,v_3,c)$ similarly, with
$\op{conv}$ instead of $\op{conv}^0$.)
Take $\op{rog}^0$ to be the empty set, if $c< r$.
 \indy{Index}{Rogers simplex}
\indy{Notation}{rogers simplex@$\op{rog}^0$}
\indy{Notation}{rogers simplex1@$\op{rog}$}
\end{definition}

\begin{lemma}\guid{XQBOVQF} The vectors $w_1,w_2,w_3$ of Definition~\ref{def:ortho}
are indeed mutually orthogonal.
\indy{Index}{orthogonality!mutual}
\end{lemma}

\begin{proof} The orthogonality of $w_1$ and $w_2$ is 
Tarski arithmetic. %tarski: {eta-ortho}.  
The orthogonality of $w_3$ with the others is also Tarski arithmetic. %tarski: {rog-ortho}.
\end{proof}

\begin{definition}[circumradius, $\eta$]\label{def:etaV}
Let $\eta(x,y,z)$ be the circumradius of a triangle with sides
$x,y,z$, and let $\eta_V(v_0,v_1,v_2) = \eta(\norm{v_0}{v_1},\norm{v_0}{v_2},\norm{v_1}{v_2})$.
\indy{Index}{circumradius}
\indy{Notation}{ZZeta@$\eta$(circumradius)}
\indy{Notation}{ZZetav@$\eta_v$}
\end{definition}

\begin{definition}[$abc$ parameter]\tlabel{def:abc}
Associate with $\op{rog}^0(v_0,v_1,v_2,v_3,c)$ the constants
$a=\norm{v_1}{v_0}/2$, $b=\eta_V(v_0,v_1,v_2)$, and $c$.
Call these the $abc$ parameters of $\op{rog}^0$.
\indy{Index}{abc parameter} 
\indy{Notation}{a@$a$ (constant)}
\indy{Notation}{b@$b$ (constant)}
\indy{Notation}{c@$c$ (constant)}
\end{definition}

The Rogers simplex is a tetrahedron.  Hence it is one of our
primitive regions.  It is eventually radial at $v_0$, hence
it has a solid angle at $v_0$.  The dihedral
angle, mentioned without further context, is understood to refer to 
   $$
   \dih_V(\{v_0,v_1\},\{v_2,p'\})=\dih_V(\{v_0,(v_0+v_1)/2\},\{p,p'\}),
   $$
\indy{Notation}{dih@$\dih_V$}
where $p$ and $p'$ are the points 
constructed in Definition~\ref{def:rog}.

The squares of the edge lengths of the tetrahedron are
   $$
   (a^2,b^2,c^2,c^2-b^2,c^2-a^2,b^2-a^2).
   $$
Define the functions
   $$
   \begin{array}{lll}
     \op{volR}(a,b,c) &= \begin{cases}
       a\sqrt{(b^2-a^2)(c^2-b^2)}/6& 0 < a < b < c,\\
       0,&\text{otherwise}
       \end{cases}\\
     \op{solR}(a,b,c) &= \begin{cases}
      2 \atn\left(\sqrt{{(a+b)(b+c)}},
         \sqrt{{(b-a) (c-b)}}\right).& 0 < a < b < c,\\
      0,&\text{otherwise}
     \end{cases}\\
     \op{dihR}(a,b,c) &= \begin{cases}
      \atn\left(\sqrt{b^2 - a^2},\sqrt{c^2 - b^2}\right)
      & 0 < a < b < c,\\
      0,&\text{otherwise}
     \end{cases}
     \end{array}
   $$
\indy{Notation}{volR@$\op{volR}$ (function)}
\indy{Notation}{solR@$\op{solR}$ (function)}
\indy{Notation}{dihR@$\op{dihR}$ (function)}
Specialization of the formulas for dihedral angle, volume, and solid angle to this
setting gives following expressions for volume and solid angle.
(The calculation of the
 solid angle formula is based on Euler's formula in 
Lemma~\ref{lemma:euler}.)

\begin{lemma}\guid{JJFHOIW}\label{lemma:rog:abc} 
Let $R=\op{rog}^0(v_0,v_1,v_2,v_3,c)$ and let $a$, $b$,
$c$ be the $abc$-parameters of $R$.  Let $\op{dih}(R)$ be the dihedral
angle of $R$ along the edge extending along $\op{aff}\{v_0,v_1\}$.  Then
$$
\begin{array}{lll}
\op{vol}(R) &= \op{volR}(a,b,c)\\
\op{sol}(v_0,R) &= \op{solR}(a,b,c)\\
\op{dih}(R) &= \op{dihR}(a,b,c)\\
\end{array}
$$
\end{lemma}
\indy{Notation}{a@$a$ (constant)}
\indy{Notation}{b@$b$ (constant)}
\indy{Notation}{c@$c$ (constant)}
\indy{Notation}{R@$R$}

\section{Finiteness and Volume}

Previous sections have developed all of the volume calculations that will
be needed in this book.  This chapter concludes with some 
elementary estimates based on the volumes of  cubes and balls.

\begin{lemma}\guid{WQZISRI}
\oldrating{100}
\formalauthor{Nguyen Tat Thang}
\rating{0}
\tlabel{lemma:Zcount}
    For all $p\in\ring{R}^3$ and all $r\ge 0$, the set
    $\ring{Z}^3\cap B(p,r)$ is finite of cardinality at most
    $4\pi (r+\sqrt3)^3/3$.
\indy{Index}{finite cardinality}
\end{lemma}

\begin{proof}  If $v\in\ring{Z}^3\cap B(p,r)$, then the ith
coordinate $v_i$ of $v$ must lie in the finite range
    $$
    p_i - r \le v_i \le p_i + r.
    $$
Hence there are only finitely many possibilities for $v$.


Place an open unit cube at each point of $\ring{Z}^3\cap B(p,r)$.
The cubes are measurable, disjoint, and contained in
$B(p,r+\sqrt3)$.  Thus, the combined volume of the cubes, which is
$\card(\ring{Z}^3\cap B(p,r))$,  is no greater than the volume of the
containing ball.  The result follows.
\end{proof}

\begin{lemma}\guid{PWVIIOL}
\oldrating{100}
\formalauthor{Nguyen Tat Thang}
\rating{0}
\tlabel{lemma:Zlow-count}
  For all $p\in\ring{R}^3$ and all $r\ge\sqrt3$, the set
    $\ring{Z}^3\cap B(p,r)$ is finite of cardinality at least
    $4\pi (r-\sqrt3)^3/3$.
\end{lemma}

\begin{proof}  Lemma~\ref{lemma:Zcount} establishes finiteness.  Place a closed unit cube at each point
of $\ring{Z}^3\cap B(p,r)$.  The cubes are measurable and cover
$B(p,r-\sqrt3)$.  Thus, the combined volume of the cubes is at
least the volume of the covered ball.  The result follows.
\end{proof}

\begin{lemma}\guid{TXIWYHI}
\oldrating{50}
\rating{0}
\formalauthor{Nguyen Tat Thang}
\tlabel{lemma:Zr2}
For all $p\in\ring{R}^3$, and $r_0,r_1>0$, there exists a $C$ such
that for all $r\ge r_1$, 
    $$
    \ring{Z}^3 \cap (B(p,r+r_0) \setminus B(p,r-r_1)) \le C r^2.
    $$
\end{lemma}

\begin{proof}  When $r \ge r_1+\sqrt3$, the previous two lemmas show
that the cardinality is at most $4\pi/3$ times
    $$(r +r_0 + \sqrt3)^3 - (r - r_1 - \sqrt3)^3 \le C' r^2$$
for some $C'$.  Similarly, if $r_1\le r\le r_1+\sqrt3$, the
cardinality is at most some fixed constant $C''$.  The result
easily follows.
\end{proof}


    %\chapter{Hypermap}\label{chap:hypermap}
\indy{Index}{hypermap}%

\section{Background on Permutations}

\begin{definition}[permutation]\guid{IFPQAWD}
A \newterm{permutation} $f$ on a set
  $D$ is a bijection $f:D\to D$.
\end{definition}

For example, the identity map $I_D$ on a set $D$,
\begin{displaymath}
I_D(x)=x \text{ for all } x \in D,
\end{displaymath}
 is a permutation.
If $f:D\to D$ is a permutation then there is an inverse function $f^{-1}:D\to D$
that is also a permuation.  
It satisfies
\begin{displaymath}
f f^{-1} = f^{-1} f = I_D.
\end{displaymath}
(This chapter uses product notation $f g$ for the composition of maps
$f\circ g$.)
If $D$ is a finite set, and two maps
$f,g:D\to D$ satisfy $f g = I_D$ on $D$, then $f$ and $g$ are permutations and are
inverses of one another:
\begin{displaymath}
f g = g f = I_D.
\end{displaymath}

Natural number powers  $f^k$ of a permutation $f:D\to D$ are defined
recursively by
\begin{displaymath}
f^0 = I_D,\quad\text{ and } f^{k+1} = f f^k.
\end{displaymath}
Integer powers $f^m$ of a permutation are defined as
$$f^m = f^i (f^{-1})^j,$$ where $m = i -j$.  This is well-defined.
The usual rule of exponents holds:
\begin{displaymath}
f^{a+b} = f^a f^b.
\end{displaymath}

If $f:D\to D$ is a permutation on a finite set $D$, then there is a smallest
positive integer $k$ such that $f^k=I_D$.  The integer $k$ is the \newterm{order}
of the permutation $f$.  If $f^m=I_D$ for any some $m$, then $m = k i$ for some
integer $i$, where $k$ is the order of $f$. The inverse $f^{-1} = f^k f^{-1} = f^{k-1}$ can be written as a
non-negative power of $f$.

A permutation $f$ of a finite set $D$ is \newterm{cyclic}, if the order of $f$ is the cardinality
of $D$.  A permutation $f$ is cyclic if and only if for every $x,y\in D$, there exists an integer $i$
such that $f^i x = y$.

The set of all permutations of the set $\{0,1,2,\ldots,k-1\}$ is written $\op{Sym}(k)$.
The set $\op{Sym}(k)$ is finite and has cardinality $k!$.



\section{Definitions}



\begin{definition}[hypermap,~dart]\guid{ZIHYYRA}\label{def:hypermap}  
  A hypermap is a finite set $D$, together with three functions
  $e,n,f:D\to D$ whose composition is the identity:
  \begin{displaymath}
e\ocirc n\ocirc f = I_D.
\end{displaymath} The
elements of $D$ are called \newterm{darts}.  The functions $e,n$ and
$f$ are called the \newterm{edge map}, the \newterm{node map}, and
the \newterm{face map}, respectively.  \indy{Index}{hypermap}%
\indy{Index}{dart}%
\indy{Index}{edge!map}%
\indy{Index}{node!map}%
\indy{Index}{face!map}%
\indy{Notation}{edgemapz@$e$ (edge map)}%
\indy{Notation}{nodemap@$n$ (node map)}%
\indy{Notation}{facemap@$f$ (face map)}%
\indy{Notation}{D@$D$ (dart)}%
\end{definition}

%\pdf{dart.pdf}{dart}{The arrowhead represents a dart.}
\begin{figure}[htb]
\centering
\szincludegraphics[width=2mm]{\pdfp/dart.eps}
\caption{This symbol represents a dart.}
\label{fig:dart}
\end{figure}

\begin{remark}[planar graphs as hypermaps]\guid{IVPJYAG}\tlabel{rem:hypermap} A hypermap is an abstraction of
the concept of 
planar graph.  Place a dart at each angle of a planar graph.
One function, $f$, 
cycles counterclockwise around the angles of each face.  
Another function, $n$, 
rotates counterclockwise around the angles at each
node.  A third function, $e$, pairs angles at opposite ends of
each edge  (Figure~\ref{fig:hypermap_ex}).   The hypermap extracts
the data $(D,e,n,f)$ from the planar graph and discards the rest.
\indy{Index}{planar graph}%
\end{remark}

\begin{figure}[htb]
\centering
\szincludegraphics[width=80mm]{\pdfp/hypermap-ex.eps}
\caption{Darts mark the angles of a planar graph.  Darts may
be permuted about faces, nodes, and edges.}
\label{fig:hypermap_ex}
\end{figure}

By the background on permutations, $e,n,f$ are all permutations on $D$.
A hypermap satisfies 
\begin{equation}\tlabel{eqn:triality}
e \ocirc n\ocirc f = n\ocirc f\ocirc e = f\ocirc e\ocirc n = I_D.
\end{equation}
Inverted, this triality becomes
\begin{displaymath}
n^{-1} \ocirc e^{-1} \ocirc f^{-1} = (f \ocirc e \ocirc n)^{-1} = I_D.
\end{displaymath}
This inversion is the abstract form of the the duality between nodes
and faces in a planar graph.  Because of these symmetries in the
defining relation, there will be multiple versions of theorems about
hypermaps, all obtained from one proof by symmetry.


\begin{definition}[path,~list,~sublist,~visit,~dart~set]\guid{RRQWGAY} 
Let $D$ be a set (of darts), and let $S$ be a set of permutations of $D$.
A \newterm{path} with \newterm{steps} in $S$
from $x_0$ to $x_{k-1}$ is a \newterm{list}\xfootnote{The empty path $[]$ seems
to have an ancient origin: ``This is the path made known to me
when I had learned to remove all darts.'' --The Dhammapada} of
darts $[x_0;x_1;\ldots;x_{k-1}]$ such that for each $i$, $x_{i+1} = h_i x_i$,
for some $h_i \in S$.   A \newterm{sublist} of a list is a consecutive
subsequence  $[x_i;x_{i+1};x_{i+1};\ldots;x_j]$, with $0\le i\le j\le k-1$.
A \newterm{unit list} is a list of the form $[x]$.  A
path is \newterm{injective} if $x_i=x_j$ implies $i=j$. 
The \newterm{dart set} of $L$ is $\{x_0,x_1,\ldots,x_{k-1}\}$.  A path \newterm{visits}
a dart $x$, if $x$ is an element of the dart set of $L$.  A set of paths visits a
dart $x$, if some path in the set visits the dart.
\end{definition}

\begin{notation}[$\cooln$]
%Write $P[x_i:x_j]$ for $i<j$ for the sublist $[x_{i+1};\ldots;x_j]$ of
%$P=[x_0;\ldots;x_{k-1}]$.  (The notation is ambiguous when the path is
%not injective.)  
The infix operator $\cooln$ prepends an element $x$ to a list $[x_0;\ldots]$:
\begin{displaymath}
x\cooln[x_0;\ldots] = [x;x_0;\ldots].
\end{displaymath}
%The infix operator $\opat$  \newterm{concatenates} 
%lists:
%\begin{displaymath}
%[a;\ldots;b] \opat [c;\ldots;d]  = [a;\ldots;b;c;\ldots;d].
%\end{displaymath}
\end{notation}
\indy{Notation}{1@$\cooln$ (list operation)}%
\indy{Notation}{P@$P$ (dart path)}%
%\indy{Notation}{1@$[-:-]$ (dart sublist)}%
%\indy{Notation}{Z@$\opat$~(concatenation)}%


\begin{definition}[$\sim_S$]\guid{IENSLJP}
Let $D$ be a set, and let $S$ be a 
set of permutations on $D$.
Define a relation on the set of darts by $x\sim_S y$ when there is a
path from $x$ to $y$ with steps in $S$.
\end{definition}

\begin{lemma}[equal equivalences]\guid{YBGABWW}\rating{50}\label{lemma:er} %\guid{QLPBIKV}
% wording changed by thales Jan 7, 2010.
Let $(D,e,n,f)$ be a hypermap and let $S$ be a  set of permutations.
Then for each $h_1,h_2\in S$, 
the relation $\sim_S$ is the same as the relation $\sim_T$, where
\begin{displaymath}
T = S \cup \{h_1h_2\}.
\end{displaymath}
Moreover, for each $h\in S$, 
the relation $\sim_S$ is the same as the relation $\sim_T$, where
\begin{displaymath}
T = S \cup \{h^{-1}\}.
\end{displaymath}
Also,  the relation $\sim_S$ (that is, $\sim_T$) is an equivalence relation.  
\indy{Index}{equivalence relation}%
\end{lemma}

\begin{proof} If $x\sim_S y$ then clearly $x\sim_T y$.  Conversely,
if $x\sim_T y$, where $T = S\cup\{h_1,h_2\}$, pick a path $P$ from $x$ to $y$ with steps
in $T$ that contains the fewest $h_1h_2$-steps.  

\claim{$P$ does not contain any $h_1h_2$-steps}.  Otherwise, a sublist $[\ldots;u;h_1h_2u;\ldots]$
of $P$ can be expanded to a path $[\ldots;u;h_2u;h_1u;\ldots]$ that contradicts the minimal
property of $P$.

This proves the first conclusion of the lemma.  Fix $h$ in a set of permutations $R$.
By an induction that uses the first conclusion,  for all $i$, $\sim_R$ equals the relation $\sim_{R(h,i)}$,
where $R(h,i) = R \cup \{h,h^2,\ldots,h^i\}$.  If $h\in S$ is an element of order $k$, 
and $T = S\cup\{h^{-1}\}$, then
the second conclusion follows because the following sets give the same relation:
\begin{displaymath}
S,\quad S(h,k-1) = T(h,k-2),\quad T.
\end{displaymath}

By repeated action of the previous conclusion, $\sim_S=\sim_T$, where 
$T = S\cup S^{-1}\cup \{I_D\}$, and where $S^{-1} = \{h^{-1}\mid h\in S\}$.
The unit path $[x]$ yields reflexivity of $\sim_T$.  Also, $T^{-1} = T$ gives the symmetry.  Finally, concatenation of paths gives transitivity.  Thus, $\sim_T$ (i.e., $\sim_S$) is an equivalence relation.
\end{proof}

\begin{definition}[combinatorial~component,~connected]\guid{JVTRXQR}
A \newterm{combinatorial component} of a hypermap $(D,e,n,f)$ is an 
equivalence class of the relation $\sim_S$, where
$S=\{e,f,n\}$. 
(See Lemma~\ref{lemma:er} for other sets that define the same equivalence classes.)  
Write $\#c$ for the
number of combinatorial components.  The hypermap is \newterm{connected} if
$\#c=1$.  \indy{Index}{Dhammapada}%
\indy{Index}{path}%
\indy{Index}{connected}%
\indy{Index}{component!combinatorial}%
\indy{Notation}{1@$\#c$~ (number of components)}%
\end{definition}





\begin{definition}[orbit,~node,~face,~edge]\guid{JIOUCMV}
The \newterm{orbit} of $x\in D$ under a permutation $h$ on
a set $D$ is a set of the form $\{h^i x\mid i\in\ring{N}\}$.  A \newterm{node}
of a hypermap $(D,e,n,f)$ is the orbit of a dart $x\in D$ under $n$.  
A \newterm{face} is an orbit under $f$.  
An \newterm{edge} is an
orbit under $e$.  \indy{Index}{node}%
\indy{Index}{face}%
\indy{Index}{edge}%
\end{definition}

Write $\#h$ for the
number of orbits of a permutation $h$ on $D$.  
\indy{Notation}{h@$h$ (permutation)}%
\indy{Notation}{1@$\#h$~(number of orbits)}%


\begin{lemma}[orbit relation]\guid{PKRQTKP}
Let $D$ be a finite set.  The orbit of $x\in D$ of a permutation $h:D\to D$
is the equivalence class of $x$ under the relation $\sim_S$, when $S=\{h\}$.
\end{lemma}

\begin{definition}[plain]\guid{HFRNMIU}
A hypermap $(D,e,n,f)$ is \newterm{plain}
  (carefully note\xfootnote{This deliberate play on the homophonous
    {\it plane} privileges writing over speech.  Every plane (hyper)graph may
be planar, but not all plain hypermaps are planar.}
  the  spelling!) when $e$ is an involution on $D$ (that is, $e^2 = I_D$).
  \indy{Index}{planar}%
\end{definition}




\begin{definition}[degenerate]\guid{MKSZLRM}
 A dart in a hypermap $(D,e,n,f)$ 
is degenerate if it is a
fixed point of one of the maps $e,n,f$; otherwise it is nondegenerate.  
%%It is nondegenerate otherwise.
\indy{Index}{dart!degenerate}%
\indy{Index}{dart!nondegenerate}%
\end{definition}

\begin{definition}[simple]\guid{KMHUQNS} 
A hypermap is \newterm{simple} if the intersection of each face with
each node contains at most one dart.  \indy{Index}{simple}%
\end{definition}


% Moved from cup05_tame.tex section on tame plane graphs. 9/5/07:
\begin{lemma}[nodal fixed point]\guid{ZHQCZLX}\rating{50}\tlabel{lemma:nondegen} 
Let $(D,e,n,f)$ be a simple plain hypermap such that every face has
at least three darts.
Then $n$ has no fixed point.
\indy{Index}{fixed point}%
\end{lemma}

\begin{proof} For a contradiction, let $x$ be a fixed point of
$n$. 

\claim{The darts $e x$ and $f x$ lie in the same node and face, so are
equal in the simple hypermap.}  Indeed, they lie in the same node
because $n(f x) = e^{-1} x = e x$. They lie in the same face because
\begin{displaymath}f^2 (e x) = f (f e n x) = f x.\end{displaymath}
So $e x = f x$.

Thus, $f^2 (e x) = f x = e x$, and $e x$ lies on a
face with at most two darts.  This contradicts what is given.
\end{proof}




\section{Walkup}

To focus on a dart $x$ in a
hypermap, it can be useful to draw a hexagon around $x$ and place
the six darts $e x$,
$f x$, $e^{-1} x$, $n x$,  $f^{-1} x$, $e x$, $n^{-1} x$  at its corners
in Figure~\ref{fig:dart+}.  Some of these seven darts may be
equal to one another, even if the figure draws them apart.
Figure~\ref{fig:dart-fix} shows the layout of a degenerate dart.
\indy{Notation}{x@$x$ (dart)}%

\begin{figure}[htb]
\centering
\szincludegraphics[width=40mm]{\pdfp/dart+.eps}
\caption{A dart $x$ and its entourage}
\label{fig:dart+}
\end{figure}

\begin{figure}[htb]
\centering
\szincludegraphics[width=60mm]{\pdfp/dart-fix.eps}
\caption{A dart fixed under a face map.}
\label{fig:dart-fix}
\end{figure}

\subsection{single}

A \newterm{walkup} deletes
a dart from a hypermap and reforms the edge, node, and face
maps to produce a hypermap on the reduced set of darts.  Walkups
come in three flavors: edge walkups, face walkups,
and node walkups.

\begin{definition}[walkup,~degenerate]\guid{DAIZNHD}
The edge \newterm{walkup}
$W_e$ at  a dart $x\in D$ of a hypermap $(D,e,n,f)$ is the hypermap
$(D',e',n',f')$, where $D' = D\setminus\{x\}$ and the
the maps skip over $x$:
\begin{displaymath}
\begin{array}{lll}
f' y &= \text{ if } (f y =  x) \text{ then } f x \text{ else
} f y\\
n' y &= \text{ if } (n y = x) \text{ then } n x \text{ else
} n y\\
e' &= (n'\ocirc f')^{-1}
\end{array}
\end{displaymath}
A walkup at $x$ is said to be \newterm{degenerate} if the dart $x$ is
degenerate.  
\indy{Index}{walkup}%
\indy{Index}{edge!walkup}%
\indy{Index}{face!walkup}%
\indy{Index}{node!walkup}%
\indy{Notation}{Wh@$W_h$ (walkup)}%
\end{definition}

Figure~\ref{fig:walk} shows
the result of an edge walkup on the hexagon around a dart $x$.
The triality symmetry~\ref{eqn:triality}, applied to the definition
of edge walkups, yields the definition of
face walkup $W_f$ and node walkup $W_n$.  
% Figure~\ref{fig:walkfn} shows the result of the face and node
% walkups on the hexagon around a dart $x$.

At a degenerate dart $x$, all three walkups are equal:
$W=W_e=W_n=W_f$ (Figure~\ref{fig:walkdegen}).
\indy{Index}{walkup!degenerate}%
\indy{Notation}{x@$x$ (dart)}%

\begin{figure}[htb]
\centering
\szincludegraphics[width=80mm]{\pdfp/walk.eps}
\caption{The effect of a walkup at $x$}
\label{fig:walk}
\end{figure}


\begin{figure}[htb]
\centering
\szincludegraphics[width=80mm]{\pdfp/walkdegen.eps}
\caption{The effect of a walkup at a degenerate dart}
\label{fig:walkdegen}
\end{figure}


\begin{definition}[merge,~split]\guid{KJIOZBJ}\tlabel{def:merge-split} 
Let $(D,e,n,f)$ be a hypermap, and let $h=n,e$, or $f$.  Let $\op{orbit}(h,x)$
denote the orbit of $x\in D$ under $h$.  Let $(D',e',n',f')$ be the hypermap obtained
from $(D,e,n,f)$ by the walkup $W_h$ at $x\in D$.
Let $h'=e',n',f'$, respectively, according to the choice of $h$.
The walkup $W_h$ at $x$ \newterm{merges} when the walkup joins the
orbit of $h$ through $x$ with another orbit.  That is, there is an orbit $O$ of some
$y\in D'$ under $h':D'\to D'$ of the form
\begin{displaymath}
O \cup\{x\} = \op{orbit}(h,x) \cup \op{orbit}(h,y),
\end{displaymath}
where $y\not\in \op{orbit}(h,x)$.
It \newterm{splits}
when the walkup splits the orbit at $x$ into two orbits.  That is, there are 
distinct orbits $O_1,O_2$ under $h'$ in the hypermap $(D',e',n',f')$ such that
\begin{displaymath}
\{x\}\cup O_1\cup O_2 = \op{orbit}(h,x).
\end{displaymath}
\indy{Index}{split}%
\indy{Index}{merge}%
\indy{Index}{orbit}%
\end{definition}

\begin{lemma}[merge-split]\guid{ZMFKZNH}\rating{150}\tlabel{lemma:merge-split} 
  Let $(D,e,n,f)$ be a hypermap and let $W_h$ be a nondegenerate
  walkup at a dart $x$.  Then $W_h$ merges or splits. Moreover, it merges if
  and only if $x$ and $y$ lie in distinct $h$-orbits, where
  $(h,y)=(f,e x)$,  $(e,n x)$, or $(n,f x)$.
\end{lemma}

\begin{proof} The walkup $W_f$ splits if and only if $f x$ 
(or $x$)
and $e x$ lie in the same $f$-orbit before the split. 
Figure~\ref{fig:split} makes this clear.
The other cases $h=e,n$ hold by triality.
\end{proof}


\begin{figure}[htb]
\centering
\szincludegraphics[height=90mm]{\pdfp/split.eps}
\caption{The face walkup at $x$ mixes $f$-orbits.  If it mixes a
single orbit, the orbit splits. If it mixes two separate orbits, the
orbits merge. }
\label{fig:split}
\end{figure}

The following is a useful way to tell if a walkup merges.


\begin{lemma}[merge criterion]\guid{FKSNTKR}\rating{80}\tlabel{lemma:ng-merge}  
Suppose, in a simple plain hypermap $(D,e,n,f)$, that an edge $\{x,y\}$ consists
of two nondegenerate darts.  Then the walkup $W_f$ 
% (resp. $W_n$)  removed Jan 10, 2009.  Needed? Is it even true?
at $x$ merges.
\end{lemma}
\indy{Index}{merge}%

\begin{proof} 
The darts $f x$ and $e x$ lie in the same node: $n (f x) = e^{-1} x
= e x$. If they are also in the same face of a simple hypermap, then
$f x = e x = y$. So
\begin{displaymath}n y  = n f x = n f e y = y,\end{displaymath}
and $y$ is a fixed
point of $n$, hence degenerate, contrary to assumption.  
Thus, $f x$ and $e x$ are in different faces, and the walkup merges
by Lemma~\ref{lemma:merge-split}.  
\end{proof}


\subsection{double}
\indy{Index}{walkup!double}%

A double walkup is the composite of two walkups of the same type.  The
two darts for the two walkups are to be the members of an orbit of
size two (under $n$, $e$, or $f$).
%%XX?The first walkup is to be chosen so that it merges.  
By choosing the type of the walkups to be different from the type of
the orbit, the first walkup reduces the orbit to a singleton, forcing
the second walkup to be degenerate.

Here are some examples.
\begin{itemize}
\item A double $W_n$ along an edge deletes the edge and 
merges the two endpoints into
a single node (Figure~\ref{fig:doublenode}). 
\item A double $W_f$ along an edge 
deletes the edge and merges the two faces along the edge into
one (Figure~\ref{fig:doubleface}).
\item A double $W_e$ at a node of degree two
deletes the node and merges the two edges at the node into
one (Figure~\ref{fig:doubleedge}).
\end{itemize}


\begin{figure}[htb]
\centering
\szincludegraphics[width=90mm]{\pdfp/double-node-walkup.eps}
\caption{The double node walkup applied to an edge}
\label{fig:doublenode}
\end{figure}


\begin{figure}[htb]
\centering
\szincludegraphics[width=90mm]{\pdfp/double-face-walkup.eps}
\caption{The double face walkup applied to an edge}
\label{fig:doubleface}
\end{figure}


\begin{figure}[htb]
\centering
\szincludegraphics[width=80mm]{\pdfp/double-edge-walkup.eps}
\caption{The double edge walkup applied to a node}
\label{fig:doubleedge}
\end{figure}

\begin{figure}[htb]
\centering
\szincludegraphics[width=80mm]{\pdfp/double_edge.eps}
\caption{The double edge walkup preserves plainness.}
\label{fig:doubleplain}
\end{figure}


\begin{lemma}[plain walkup]\guid{HOZKXVW}\rating{150}\tlabel{lemma:dwalk-planar}  
The three preceding double walkups carry plain
hypermaps into plain hypermaps.
\end{lemma}
\indy{Index}{hypermap!plain}%

\begin{proof} The walkups $W_n$ and $W_f$ preserve the orbit structure
of edges, except for dropping one dart.  By dropping both darts from
the same edge, one edge is lost and all others edges remain
unchanged.

Figure~\ref{fig:doubleplain} illustrates the double $W_e$.  The two
edges $\{x,e x\}$, $\{y, e y\}$ meeting the node are fused by the
double walkup into $\{e x, e y\}$, which is still an edge of size
two.
\end{proof}

\begin{remark}[reverse double walkup]\guid{KPRURND}\label{rem:reverse-double-walkup}
Double walkup transformations can be run in reverse.
Let $H'=(D',e',n',f')$ be
a hypermap and let $D\supset D'$ be a set that contains two additional elements
$x,y$.  Fix distinct elements $x',y'\in D'$.  Define $n,e:D\to D$ as follows.
\begin{displaymath}
\begin{cases} 
e x = y, &\\
e y = x,&\\
e z = e' z,&\text{otherwise.}\\
\end{cases}
\qquad\qquad
\begin{cases} 
n x' = x, &\\
n x =  n' x',&\\
n y' = y,&\\
n y =  n' y',&\\
n z = n' z, &\text{otherwise.}\\
\end{cases}
\end{displaymath}
Define $f$ by forcing the hypermap identity $e n f = I_D$.  
The edge $\{x,y\}$ has been inserted by a reverse double walkup.  The insertion points of the edge
into the hypermap
depend on the data $\{(x',x),(y',y)\}$.   

Reverse double walkup transformations that
insert a node $\{x,y\}$ or a face $\{x,y\}$ into a hypermap are obtained similarly  by triality symmetry.  For example, to insert a node onto an edge $\{x',y'\}$ of cardinality $2$, use
\begin{displaymath}
\begin{cases} 
n x = y, &\\
n y = x,&\\
n z = n' z,&\text{otherwise.}\\
\end{cases}
\qquad\qquad
\begin{cases} 
f x' = x, &\\
f x =  f' x',&\\
\text{etc.}&\\%f
%f y' = y,&\\
%f y =  f' y',&\\
%f z = f' z, &\text{otherwise.}\\
\end{cases}
\end{displaymath}
\indy{Index}{reverse double walkup}%
\end{remark}

\begin{remark}[dart universe]\guid{SCYVYJW}\label{rem:dart-universe}
For reverse double walkups, we need a set from which to draw new darts $x,y$.
We will use a well-ordered set $\Omega$ from which we draw, as needed,
the minimal element of the complement in $\Omega$ of the set of darts
already in play.  We assume that darts can be supplied from $\Omega$,
without mentioning it explicitly.  

For example, if we insert an edge into $(D',e',n',f')$
using the
ordered pair $(x',y')$, we use the data $\{(x',x),(y',y)\}$, where $x$ is the
least element of $\Omega\setminus D'$, and $y$ is the least element of
$\Omega\setminus (D'\cup \{x\})$.  To insert a node into an edge of cardinality
two, it is enough to specify one dart $x'$ in the edge.  Then let $y'$ be the other
dart in the edge, and choose $x,y\in \Omega$ as above.
\indy{Notation}{zzZ@$\Omega$ (well-ordered universe of darts)}
\end{remark}

\begin{definition}[RDW]\guid{TPEZAAM}\label{def:R}  
  Let $H'=(D',e',n',f')$ be a hypermap and let $x'\in D'$.  
%  Let $r$ be the
%  cardinality of the face of $x$, and let $m,p,q$ be integers that
%  satisfy $0\le p$, $0\le m < q < r$, and $m+1 <
%  p+q$.
Let $m,q,p$ be natural numbers.  Assume that $y' =(f')^{m+1}x'$ is not
equal to $z'=(f')^{q+1}x'$.
Construct a hypermap $RDW(H',x',m,p,q)$  as follows.  
First
  add an edge into $H$ using the ordered pair $(y',z')$
(by the reverse double walkup of
  Remark~\ref{rem:reverse-double-walkup}).  Then insert $p$ new nodes
  (of degree $2$) again by reverse double
  walkup transformations, each time at the edge containing $y'$.  
This is $RDW(H,x',m,p,q)$.
\end{definition}
\indy{Notation}{RDW@$RDW$ (reverse double walkup)}


\section{Planarity}
\indy{Index}{walkup}%
\indy{Index}{planarity}%

\begin{definition}[planar]\guid{QVATKMJ}
A hypermap is \newterm{planar} (note the
spelling!) when the Euler relation holds:
\begin{displaymath}\# n + \# e + \# f = \# D + 2\, \#c.\end{displaymath}
\indy{Index}{planar}%
\end{definition}


\begin{remark}[Eulerian relation]\guid{YPVCMHI}\label{rem:Euler}   
The Euler relation for planar graphs can be translated into the
language of hypermaps.  Consider a connected planar graph that
satisfies the Euler relation for the alternating sum of Betti
numbers:
\begin{displaymath}b_0 - b_1 + b_2 = 2\end{displaymath} where $b_0$
is the number of vertices, $b_1$ the number of edges, and $b_2$ the
number of faces (including an unbounded face) of the planar
graph. The hypermap $(D,e,n,f)$, made from the planar graph in
Remark~\ref{rem:hypermap}, is plain, and the involution $e$ has no fixed points.  
Thus, $\# D = 2\#e$, according to the partition of $D$ into edges.  Moreover,
\begin{displaymath}\begin{array}{lll}
b_0 &= \# n\\
b_1 &= \# e\\
b_2 &= \# f\\
2b_1 &= \# D\\
1 &= \#c\\
b_0 - b_1 + b_2  &= \# n + (\#e - \#D) + \# f = 2\,\# c.
\end{array}
\end{displaymath}
Thus, the hypermap is also planar.
\indy{Index}{Euler relation} %
\end{remark}


\begin{lemma}[dart bound]\guid{TGJISOK}\rating{80}\label{lemma:dart-upper} 
Let $H$ be a connected plain planar hypermap such that every edge
has cardinality two.  Assume that there are at least three darts in
every node.  Then
\begin{displaymath}
\# D \le (6\, \#f - 12).
\end{displaymath}
\end{lemma}
\indy{Notation}{H@$H$ (hypermap)}%

\begin{proof}  In a plain planar hypermap, the Euler relation becomes
\begin{displaymath}6\, \#f - 12 = 3\,\#D - 6\,\#n,\end{displaymath}
so it is enough to show that
\begin{displaymath}
\# D \ge 3\,\#n.
\end{displaymath}
This follows directly by assumption: the set of darts can be
partitioned into nodes, with at least three darts per node.
\end{proof}


\begin{definition}[planar~index]\guid{ICAWSNK}
The planar index of a hypermap is
\begin{displaymath}\iota = \# f + \# e + \# n - \# D - 2\,\#
c.\end{displaymath}
(A hypermap with null index is planar.)
\indy{Index}{hypermap!planar index}%
\indy{Notation}{ZZiota@$\iota$ (planar index)}%
\end{definition}

\begin{lemma}[walkup index]\guid{IUCLZYI}\rating{400}\tlabel{lemma:index} 
Let $x$ be a nondegenerate dart of a hypermap $(D,e,n,f)$. Let
$(D',e',n',f')$ be the result of the face walkup $W$ at $x$.  The
walkup changes the size of some orbits.
\begin{displaymath}
\begin{array}{lll}
%\text{\bf Non-degenerate dart $x$: }&\\
\# f' &=\# f +\op{split}_f  \\  
\# e'&=\# e \\
\# n'&=\# n \\
\# D'&=\# D - 1 \\
\#c'&=\# c + \op{split}_c\\
\iota' &= \iota + 1+\op{split}_f - 2\op{split}_c,\\
\end{array}
\end{displaymath}
where
\begin{displaymath}
\op{split}_f = \begin{cases}
1,&\text{if $W$ splits }\\
-1,&\text{if $W$ merges}\\
\end{cases}
\end{displaymath}
and $\op{split}_c=1$ if $e x$ and $f^{-1} x$ belong to different
combinatorial components after the walkup $W$, and $\op{split}_c=0$
otherwise. Moreover, a walkup at a degenerate dart preserves the
planar index.  \indy{Notation}{splitc@$\op{split}_c$}%
\indy{Notation}{splitf@$\op{split}_f$}%
\indy{Notation}{W@$W$ (walkup)}%
\end{lemma}

\begin{proof} The figures make this clear.
\end{proof}

\begin{lemma}[index inequality]\guid{BISHKQW}\rating{100}\tlabel{lemma:planar-index2}
Let $\iota$ be the index of a hypermap $(D,e,n,f)$, and let $\iota'$
be the index after a walkup $W_h$ at a dart $x$.  Then $\iota \le
\iota'$.
\end{lemma} 


\begin{proof} Without loss of generality, by triality symmetry, the
walkup is a face walkup.  If $\op{split}_c=0$, the inequality is
immediate by Lemma~\ref{lemma:index}.  If $\op{split}_c=1$, 
then $e x$ and $f^{-1} x$ lie in
different components after the walkup, hence also in different
faces.  Thus, the walkup splits by Lemma~\ref{lemma:merge-split}.
Hence  $\op{split}_f = 1$.  The result
follows by Lemma~\ref{lemma:index}.
\end{proof}


\begin{lemma}[non-positive index]\guid{FOAGLPA}\rating{50}
\tlabel{lemma:planar-nonpos}  
The planar index
of a hypermap is never positive.
\end{lemma}

\begin{proof}  An face walkup never decreases the index.  A sequence
of face walkups leads to the empty hypermap, which has
index zero.
\end{proof}


\begin{lemma}[planar walkup]\guid{SGCOSXK}\rating{50}
\tlabel{lemma:walkup-planar}
Walkups take planar hypermaps to planar
hypermaps.
\end{lemma}

\begin{proof}  
A planar hypermap has maximum index.  The walkup
can only increase the index, but never beyond its maximum.  
Thus, the index remains at its maximum value.
\end{proof}





\section{Path}

\subsection{contour}

\begin{definition}[cyclic~list]\guid{MYJNYCZ}
A \newterm{cyclic list} $\lp{x_0;\ldots;x_{k-1}}$ is an equivalence class of lists under the transitive closure of the relation:
\begin{displaymath}
[x_0;x_1;x_2\ldots;x_{k-1}] \sim [x_1;x_2;\ldots;x_{k-1};x_0].
\end{displaymath}
A sublist of a cyclic list is a sublist of some representative of the equivalence class.
\end{definition}

\begin{definition}[contour~path,cyclic~list,~contour~loop]\guid{AUIDQRN}
 A \newterm{contour path} from
$x_0$ to $x_{k-1}$ is a path $[x_0;x_1;\ldots;x_{k-1}]$ such that
$x_{i+1} = n^{-1} x_i$ or $f x_i$ for each $i<k$.  (That is, each
step in the path is clockwise step around a node or a
counterclockwise step around a face.)  
A \newterm{contour loop} is an injective cyclic list
$\lp{x_0;x_1;\ldots;x_{k-1}}$ such that
for every $i$, there exists $h_i\in \{f,n^{-1}\}$ such that $x_{i+1} = h_i x_i$, 
where the subscripts are
read modulo $k$.
%$[x_1;\ldots;x_{k-1}]$ is injective and $x_0 = x_{k-1}$, then it is
%a \newterm{contour loop}.  
%A sublist of a contour loop $[x_0;\ldots;x_{k-1}]$ is a path
%$[y_0;
\indy{Index}{contour!path}%
\indy{Index}{contour!loop}%
\indy{Index}{loop}%
\end{definition}



\begin{remark}[contour path illustration]\guid{AWRGIPA}
 Figure~\ref{fig:hypermap_ex}
  constructs a hypermap from a planar graph by drawing darts next to
  each angle.  In this representation, the darts along a contour path
  lie to the left of the corresponding planar graph edges.  For that
  reason, a shaded region to the left of a curve depicts a contour
  path.
\end{remark}

\begin{figure}[htb]
\centering
\szincludegraphics[width=80mm]{\pdfp/shade_dart.eps}
\caption{A contour path as a sequence as dart is represented as a
shaded path.}
\label{fig:shade-dart}
\end{figure}

\begin{lemma}[injective path]\guid{QZTPGJV}\rating{50} 
An injective contour path from
  $x$ to $y$ can be constructed from an arbitrary contour path from
  $x$ to $y$ by dropping some darts from the path.
\end{lemma}

\begin{proof} Repeatedly replace $[\ldots;a;b;\ldots;b;c;\ldots]$ with
$[\ldots;a;b;c;\ldots]$.
\end{proof}





\begin{lemma}[contours-components]\guid{KDAEDEX}\rating{100}\tlabel{lemma:connect-contour}  
Let $H$ be a hypermap.
If $x$ and $y$ are darts in the same combinatorial component of $H$ if and only if
there exists a contour path from $x$ to $y$.
\end{lemma}

\begin{proof} 
Combinatorial components are defined by an equivalence relation $\sim_S$, where
$S = \{e,n,f\}$.  By Lemma~\ref{lemma:er}, this is the same equivalence relation as
$\sim_T$, where $T = \{n^{-1},f\}$.  By the definition of the equivalence relation $T$,
$x\sim_T y$ if and only if there is a contour path from $x$ to $y$.
\end{proof}
\indy{Index}{component!combinatorial}%

\begin{definition}[complement]\guid{GCACAFP} 
Let $(D,e,n,f)$ be a plain hypermap.
Let $P=\lp{x;y;\ldots}$ be a contour loop that does not visit any node
twice in a plain hypermap.   (That is, the dart set of $P$ intersected with a node
is the dart set of a maximal sublist $[z;n^{-1}z;\ldots;n^{-k}z]$ of $n^{-1}$ steps.)
 Replace each maximal sublist of
$n^{-1}$-steps
\begin{displaymath}
[z;n^{-1} z; \ldots; n^{-k} z]
\end{displaymath}
with the sublist
\begin{displaymath}
[n^{-(k+1)} z;n^{-(k+2)} z;\ldots; n z]
\end{displaymath}
Concatenate these new sublists in reverse order.  By the relation $n f = f^{-1} n^{-1}$,
the transitions between the new sublists are $f$-steps.
The resulting contour loop $P^c$
is the \newterm{complement}. 
\end{definition}
\indy{Notation}{1@$*^c$ (complement)}

\begin{figure}[htb]
\centering
\szincludegraphics[width=70mm]{\pdfp/complement.eps}
\caption{The complement contour traces the remaining darts
at the same nodes as the original contour loop. }
\label{fig:contour-comp}
\end{figure}


\subsection{M\"obius}

\begin{definition}[M\"obius~contour]\guid{MBYIEQP}
 A M\"obius contour in a hypermap
$(D,e,n,f)$ is an
injective contour path $P=[x_0;\ldots]$ that satisfies
\begin{equation}
\tlabel{eqn:mobius}
x_j = n x_0\quad x_k = n x_i
\end{equation}
for some $0 < i\le j< k$ (Figure~\ref{fig:mobius}).
\indy{Index}{contour!M\"obius}%
\end{definition}


\begin{remark}[Four-Color theorem]\guid{ROIPZSU}
G. Gonthier devised the notion of M\"obius contour as a way to prove
the Four-Color theorem without appeal to topology.  (The Appel-Haken
proof of the Four-Color theorem relies on the Jordan curve theorem.)
This chapter uses a significant amount of material from ~\cite{Gonthier:2005:FourColor}.
\end{remark}

\begin{figure}[htb]
\centering
\szincludegraphics[width=50mm]{\pdfp/mobius.eps}
\caption{A M\"obius contour}
\label{fig:mobius}
\end{figure}

\begin{figure}[htb]
\centering
\szincludegraphics[width=30mm]{\pdfp/3m.eps}
\caption{The face map on this hypermap gives a M\"obius contour with
three darts}
\label{fig:3m}
\end{figure}

\begin{remark}[M\"obius strip]\guid{NGALZAC}
 Heuristically, a M\"obius contour is a 
combinatorial M\"obius strip that
twists to make 
its left-hand side into
its right-hand side.  A planar hypermap has no such contour.  
Figure~\ref{fig:violate-jct}
redraws a violation of the Jordan curve theorem
as a M\"obius contour.   
\end{remark}

\begin{figure}[htb]
\centering
\szincludegraphics[width=80mm]{\pdfp/violate-jct2.eps}
\caption{A path that tunnels from the interior to the exterior
of a simple closed curve
is analogous to a M\"obius contour.}
\label{fig:violate-jct}
\end{figure}

\begin{figure}[htb]
\centering
\szincludegraphics[width=80mm]{\pdfp/mobius_contour.eps}
\caption{Some M\"obius contours}
\label{fig:mobius-contour}
\end{figure}






\begin{lemma}[planar-non-M\"obius]\guid{LIPYTUI}\rating{300}\tlabel{lemma:no-mobius}
A planar hypermap does not have a M\"obius contour.
\end{lemma}
\indy{Index}{hypermap!planar}%

\begin{proof} For a contradiction, assume that there exist planar
hypermaps with M\"obius contours.  An edge walkup carries
planar hypermaps into planar hypermaps. An edge walkup
at a dart that is not on the M\"obius contour carries the
M\"obius contour to a M\"obius contour 
and reduces the number of darts.  
In the M\"obius Condition~\ref{eqn:mobius},
a walkup at a dart that is not at position $0$, $i$, $j$, $k$
along the contour carries the M\"obius contour to a M\"obius contour
and reduces the number of darts. Thus, a counterexample with
the smallest possible number of darts contains no
darts except those on the M\"obius contour, and its only darts
are at positions $0$, $i=j=1$, $k=2$.

This is a three darted hypermap (Figure~\ref{fig:3m}.)  
The M\"obius condition, the
definition of contours, together with $e\ocirc n\ocirc f=I_D$ force
$e=n=f$, all permutations of order three.  This hypermap is not planar:
\begin{displaymath}\# e + \# n + \# f = 3~~\ne~~ 5 = \# D + 2\,
\#c.\end{displaymath}
\end{proof}



%\subsection{interior}
%\indy{Index}{interior}%
%
%\begin{definition}[interior]\guid{NQXQALU}\label{def:interior} 
%A dart $y$ lies in the \newterm{interior} of a contour
%loop $L$ if there is a an injective contour path
%$x_0,x_1,\ldots,x_k=y$ such that $x_1 = f x_0$ (or $k=0$), and
%such that $x_i$ lies on the loop $L$ if and only if $i=0$.
%\indy{Index}{interior!contour loop}%
%\indy{Notation}{L@$L$ (contour loop)}%
%\indy{Notation}{y@$y$ (dart)}%
%\end{definition}
%

\begin{lemma}[step coherence]\guid{ILTXRQD}\rating{100}\tlabel{lemma:contour-path-type}
Suppose that a hypermap has no M\"obius contours. Let $L$ be a
contour loop.  Let $P$ be any injective contour path with at least
$3$ darts, that starts and ends on $L$, but visits no other darts of
$L$.  Then the first and last steps of $P$ are both of the same type
($n^{-1}$ or $f$).
\end{lemma}
\indy{Notation}{P@$P$ (contour path)}%

\begin{figure}[htb]
\centering
\szincludegraphics[width=80mm]{\pdfp/interior_nf.eps}
\caption{A path must enter and depart from a contour loop with the
same type of step.}
\label{fig:interior_nf}
\end{figure}


\begin{proof} The proof shows the contrapositive.  Suppose $P=[n x;f n
x;\ldots;n y;y]$.  The successor of $n x$ on $L$ is $x$.  Starting
at $x$, follow $L$ to $y$, and on to $n x$.  Follow $P$ back to $n
y$.
%\begin{displaymath}
%x\cooln L[x:n x] \opat P[ n x;ny].
%\end{displaymath}  
This is a M\"obius contour $x\ldots y\ldots n x\ldots n y$.

Suppose $P=[n x;x;\ldots;f^{-1} y;y]$.  Starting at $x$, follow $P$ to
$y$, then follow $L$ to $n x$, and on to $n y$.  This is a M\"obius
contour.\xfootnote{The second statement can also be deduced from the first statement
by the duality $(D,e,n,f)\leftrightarrow (D,e^{-1},f^{-1},n^{-1})$ that swaps
$f$-steps with $n^{-1}$-steps in a path.}
\end{proof}

%\begin{lemma}[]\guid{UMYSGDB}\rating{80}\tlabel{lemma:dart-interior}
%  Let $L$ be a contour loop on a plain hypermap without M\"obius
%  contours.  Assume a dart $x$ lies in the interior of the loop $L$.
%  Then every dart in its $f$-orbit lies in the interior of the loop.
%  Moreover, if the dart $x$ does not lie on the same node as any dart
%  in $L$, then every dart in the $n$-orbit of $x$ lies in the
%  interior of $L$.
%\end{lemma}

%\begin{proof} Let $P= [x_0;\ldots;x]$ be an injective path that
%  certifies that $x$ lies in the interior of $L$.  If $f x$ lies
%  along this path already or if it lies on $L$, then it is clearly
%  interior.  Otherwise, $[x_0;\ldots,x;f x]$ is a certifying path for
%  $f x$.  Similarly, use the certifying path $[x_0;\ldots;x;n^{-1}
%  x]$ for $n^{-1} x$.
%\end{proof}
%
%%
%\begin{definition}[interior~face,~node]\guid{JUXKJTU}
% A face or a node is interior
%  to a loop in a hypermap if all of its darts are interior.
%  \indy{Index}{interior!face}%
%  \indy{Index}{interior!node}%
%\end{definition}


\begin{lemma}[loop separation]\guid{ICJHAOQ}\rating{180}\tlabel{lemma:contour-f}
Suppose that a hypermap has no M\"obius contours.  Let $L$ be a
contour loop.  Then there does not exist a contour path
$[x_0;\ldots;x_k]$, for $k\ge 1$ with the following properties:
\begin{enumerate}
\item $x_i$ lies on $L$ if and only if $i=0$.
\item $x_1 = f x_0$.
\item $x_0$ and $x_k$ lie in different nodes.
\item Some dart of $L$ is at the node of $x_k$.
\end{enumerate}
\end{lemma}

%\begin{figure}[htb]
%  \centering
%  \szincludegraphics[width=40mm]{\pdfp/no_node_path.eps}
%  \caption{No path exists from a node of $L$ to the interior.}
%  \label{fig:no-node-path}
%\end{figure}

\begin{proof} Assume for a contradiction that the path $P$ exists.
Some sublist is injective and satisfies the same conditions.  Again,
without loss of generality, shrinking the path if needed, $k$ is the
smallest index for which the last two conditions are met.  Append
$n^{-1}$-steps to $P$ to reach a dart of $L$.  This is contrary to
Lemma~\ref{lemma:contour-path-type}.
\end{proof}

\begin{lemma}[three darts]\guid{EUXPBPO}\rating{ZZ}\label{lemma:3dart}  
Assume that each face of a hypermap  has at least three darts.
Then every contour loop that meets at least two nodes has at least
three darts.
\end{lemma}

\begin{proof} Let $P=\lp{x;y}$ be a contour loop meeting two nodes.  Then
$y = f x$ and $x = f y$, so that the face has size two.
\end{proof}


\section{Quotient}
\indy{Index}{quotient}%

\subsection{definition}

\begin{definition}[isomorphism]\guid{GUDUERI}
 Two hypermaps $(D,e,n,f)$ and
$(D',e',n',f')$ are \newterm{isomorphic} when there is a bijection
$G:D\to D'$ such that
\begin{displaymath}h'\circ G = G\circ h\end{displaymath}
for $(h,h')=(e,e'), (f,f'), (n,n')$.
\indy{Index}{isomorphic hypermaps}%
\indy{Notation}{G@$G$ (morphism of hypermaps)}%
\end{definition}


\begin{definition}[normal family]\guid{RQSVFLE}
Let $(D,e,n,f)$ be a hypermap. 
%Assume that 
%there are no darts fixed by $e$ 
%(so that $f x \ne n^{-1} x$ at each dart). 
Let $\cal L$ be a family of contour
loops.  The family $\cal L$ is  normal if the following
conditions hold of its loops. \begin{enumerate}
\item  No dart is visited by two different loops.
\item  Every loop visits at least two nodes.
\item  If a loop visits a node, then every dart at that node is visited
by some loop.
\end{enumerate}
\indy{Index}{normal family}%
\end{definition}

A normal family determines a new hypermap.  A dart in the new set $D'$
of darts is a maximal sublist $[x;n^{-1} x; n^{-2} x;\ldots;n^{-k}
x]$ of $n^{-1}$ steps appearing in some loop in $\cal L$. The map $f'$
takes the maximal path $[x;n^{-1}x;\ldots;y]$ to the maximal path (in
the same contour loop) starting at $f y$. The map ${n'}^{-1}$ takes
the maximal path $[\ldots;y]$ to the maximal sequence (in some other
contour loop) starting $[n^{-1}y;\ldots]$. Equivalently, $n'$ takes
the maximal path $[x;\ldots]$ to the maximal path ending $[\ldots;n
x]$. The map $e'$ is defined by $e'\ocirc n'\ocirc f' = I_{D'}$.
\indy{Index}{path!maximal} %
\indy{Notation}{D@$D$ (dart)}%

\begin{definition}[quotient]\guid{AJENHSB}
 The hypermap $(D',e',n',f')$
constructed from the normal
family $\cal L$ of $H=(D,e,n,f)$ 
is called the \newterm{quotient} of $H$ by $\cal L$, and is denoted
$H/{\cal L}$.  If $x$ is a dart visited by some loop in $\cal L$, then
the maximal path $[\ldots;x;\ldots]$ is called the \newterm{quotient dart} of $x$.
%The hypermap $H$ is said to be a \newterm{cover} of $H/{\cal L}$.  
\indy{Index}{quotient}%
\end{definition}
\indy{Notation}{L@$\cal L$}%
\indy{Notation}{H@$H/{\cal L}$}%
%\indy{Index}{cover}%
\indy{Index}{hypermap!quotient}%
\indy{Index}{normal family}%

Intuitively, the quotient hypermap is represented as a graph whose
cycles under $f'$ are precisely the contour loops in the normal family
(Figure~\ref{fig:quot}).


\begin{figure}[htb]
\centering
\szincludegraphics[width=70mm]{\pdfp/quot.eps}
\caption{The contour loops in a normal family become faces in the
quotient}
\label{fig:quot}
\end{figure}

\subsection{properties}

This subsection explores some of the properties of a quotient hypermap.
The first two lemmas describe the faces and the nodes of the quotient
in terms of the combinatorics of the normal family.

\begin{lemma}[quotient face,~$\F$]\guid{WIRLCNL}\label{lemma:quotient-bijection}
  Let ${\cal L}$ be a normal family of the hypermap $H$.  Then ${\cal
    L}$ is in natural bijection with the set of faces of the quotient
  $H/{\cal L}$.  If $x'=[x;\ldots;n^{-k}x]$ is a maximal path of
  $n^{-1}$ steps in the contour loop $L\in{\cal L}$, then the
  corresponding face $\F(L)$ of $H/{\cal L}$ is the
  one containing the quotient dart $x'$.
\end{lemma}
\indy{Notation}{F@$\F$ (quotient bijection)}

\begin{proof}  This is left as an exercise to the reader.
\end{proof}


\begin{lemma}[quotient node]\guid{UDJNSHH}\label{lemma:quotient-node}
Let $H$ be a hypermap and let ${\cal L}$ be a normal family of $H$.
Then there is a natural bijection between  the set of nodes of
$H/{\cal L}$ and the set of nodes of $H$ that
are visited by some contour loop in ${\cal L}$.   
The bijective function sends the node in $H/{\cal L}$ of 
the dart $x' = [x;n^{-1} x;\ldots;n^{-k}x]$ to the node of $x$ in $H$.
\end{lemma}

\begin{proof}  The proof is an elementary verification.
Let $H/{\cal L} = (D',e',n',f')$.
\claim{This function is well-defined.}  Indeed, 
 \begin{displaymath}
(n')^{-1} x' = [n^{-(k+1)} x;\ldots]
\end{displaymath}
 is also sent to the node of $x$ in $H$.

\claim{This function is onto.}  Indeed,
If $L\in {\cal L}$ visits $x$,  and $x'$ is the image of $x$ in $D'$.
Then the node of $x'$ clearly maps to the node of $x$.  

\claim{Finally, the function is one-to-one.}  Indeed, if the nodes of
two quotient darts $x'$, $y'$ map to the same node of $H$, then
$x'=[x;\ldots]$ and $y'=[y;\ldots]$, where $n^j x = y$ for some $j$.
It follows by the definition of the node map on the quotient, that
$x'$ and $y'$ belong to the same node.
\end{proof}

The next two lemmas look at properties of the quotient that are
inherited from the orignal hypermap.

\begin{lemma}[plain quotient]\guid{JMKRXLA}\rating{280}\tlabel{lemma:quotient-plain}
Let $H$ be a plain hypermap, and let $\cal L$ be a
normal family.  Then $H/{\cal L}$ is a plain hypermap.
\end{lemma}

\begin{proof}  Write $H=(D,e,n,f)$ and $H/{\cal L} = (D',e',n',f')$.  
Write $[\ldots; x]$ for the node in
the quotient ending in dart $x\in D$ and $[x;\ldots]$ for the node
in the quotient starting with dart $x\in D$.  Plainness gives $e^2 x
= x$, so that for any dart $[\ldots x]$ in the quotient:
\begin{displaymath}\begin{array}{lll}
{e'}^{-2} [\ldots; x] &= n' f' n' f' [\ldots; x] = n' f' n' [f x; \ldots] \\&=
n' f' [\ldots; n f x] = n' [f n f; \ldots] = [\ldots; n f n f x]\\ &=
[\ldots; e^{-2} x] = [\ldots; x].
\end{array}\end{displaymath}
Thus, $e'$ has order $2$ on the quotient.
\end{proof}




\begin{definition}[no double joins]\guid{EDUYIEA}
A hypermap $H$ has no \newterm{double joins}, if for every two nodes
in $H$, there is at most one edge that meets both of them.
\end{definition}


\begin{lemma}[quotient-no-double-joins]\guid{KSRDPTZ}
Let $H$ be a plain hypermap with no double joins and let ${\cal L}$ be a normal
family of $H$.  Then $H/{\cal L}$ has no double joins.
\end{lemma}

\begin{proof} By Lemma~\ref{lemma:quotient-plain}, the quotient
  $H/{\cal L}$ is plain.  Let $\{x',e'x'\}$ and $\{y',e'y'\}$ be edges
  with the property that $x'$ and $y'$ lie at one node of $H/{\cal L}$
  and $e'x'$ and $e'y'$ lie at a second (different) node.  Write $x' =
  [\ldots;x]$ and $y' = [\ldots;y]$.  Then $e'x' = [\ldots;e x]$ and
  $e'y' = [\ldots;e y]$.  According to
  Lemma~\ref{lemma:quotient-node}, there is an injective map from
  nodes of $H/{\cal L}$ to nodes of $H$.  It follows that $x$ and $y$
  belong to the same node and that $e x$ and $e y$ belong to the same
  (different) node.  By the assumption that $H$ has no double joins,
  it follows that $x=y$.  Hence also $x' = y'$, and $H/{\cal L}$ has
  no double joins.
\end{proof}


\begin{lemma}[nodal fixed point]\guid{PYOVATA}\label{lemma:nfp}
Let $H=(D,e,n,f)$ be a hypermap in which the edge map has no fixed points.
Let ${\cal  L}$ be a normal family of $H$, with quotient $H/{\cal L} = (D',e',n',f')$.  
Then the following are equivalent conditions:
\begin{itemize}
\item $n'$ has a fixed point in $D'$.
\item The dart set of some $L\in {\cal L}$ contains a node.
\end{itemize}
\end{lemma}

\begin{proof}
If $x'=[x;n^{-1} x;\ldots;n^{-k} x]$ is a dart in $D'$, then $(n')^{-1}x'$ is
$[n^{-(k+1)} x;\ldots]$.  The dart $x'$ is a fixed point if and only if
$x = n^{-(k+1)} x$.  This holds if and only if the dart set of $x'$ is an entire node.
\end{proof}

\subsection{example}

\begin{example}[maximal normal family]\label{ex:Hall} 
  Assume that $H=(D,e,n,f)$ is a
  hypermap. % with no fixed points under $e$.
  Assume that every face meets at least two nodes. Then the set of all
  faces defines a normal family of contour loops: follow $f$ around
  each face $[x;f x;\ldots]$.  If $e$ acts without fixed points, then
  each dart of the quotient is just a unit path consisting of a single
  dart of $H$, and the quotient is isomorphic to $H$ itself.
\end{example}

\begin{example}[minimal normal family]\label{ex:H2} 
  Assume that $H=(D,e,n,f)$ is a plain hypermap.  Let $F = \{x,f x,\ldots\}$ be a face
  that visits at least three nodes and that meets each node in at most
  one dart.  Let $\cal L$ be the family with two contour loops: $\lp{x;f x;\ldots}$ 
and its complement $L^c = \lp{n^{-1} x;\ldots}$.
%\begin{displaymath}
%  [n^{-1} x;n^{-2} x;\ldots;n x;f n x = y;n^{-1} y; n^{-2} y;\ldots; n y; f ny;\ldots]
%\end{displaymath}
The family $\cal L$ is normal. The quotient hypermap $H/{\cal L}$ has
two faces: $F$ and a back side $F'$ of the same cardinality $k$.
\indy{Notation}{F@$F$ (hypermap face)}%
\end{example}

\begin{example}[dihedral]\label{ex:D2k} 
There is a hypermap $\op{Dih}_{2k}$ whose dart set has cardinality $2k$.
The permutations $f,n,e$ have  orders $k$, $2$, and $2$ respectively, and 
$e n f = I$.
The set of darts is given by
\begin{displaymath}
\{x, f x,f^2 x,\ldots,f^{k-1} x\}\cup \{n x, n f x, n f^2 x,\ldots, n f^{k-1} x\},
\end{displaymath}
for any dart $x$.
If a hypermap is isomorphic to $\op{Dih}_{2k}$ for
some $k$, then it is \newterm{dihedral}.\xfootnote{The three permutations generate the dihedral
group of order $2k$, acting on a set of $2k$ darts under the left action of the group upon itself.}   
In particular,
the hypermap constructed in the previous example is dihedral.
\indy{Index}{hypermap!dihedral}%
\end{example}

\begin{lemma}\guid{CODE}\rating{ZZ}\label{lemma:dih-iso}
A hypermap $H$ is isomorphic to $\op{Dih}_{2k}$ if and only if the hypermap is connected,
the dart set has cardinality $2k$, and the permutations $f,n,e$ have orders $k$, $2$, and $2$ respectively.
\end{lemma}

\begin{proof} Let $y$ be any dart of $H$, and $x$ any dart of $\op{Dih}_{2k}$.
By the connectedness of $H$, and the relations between $f$ and $n$, every dart of $H$ is
equal to one of the following: $f^i y$, $n f^i y$, for $i=0,\ldots,k-1$.  By the cardinality assumption,
these darts are all distinct.  By the relations between $f$ and $n$ on the hypermap $\op{Dih}_{2k}$,
the bijection $f^i x \mapsto f^i y$, $n f^i x \mapsto n f^i y$ is an isomorphism of hypermaps.
\end{proof}

%\begin{lemma}[three darts]\guid{BPAKIYP}
%Let $H$ be a plain hypermap with no double joins in which every
%face has at least three darts.  Let ${\cal L}$ be a node-free normal family of $H$.
%Then every face of $H/{\cal L}$ has at least three darts.
%\end{lemma}
%
%\begin{proof}  By definition, every contour loop in a normal family meets
%at least two nodes.  It follows that every face of $H/{\cal L}$ has at least two
%darts.
%
%Suppose for a contradiction that a face of $H'=H/{\cal L}$ only has two darts $x'$
%and $y'$, exchanged by the face map $f'$ of $H'$:
%\begin{displaymath}
%\begin{array}{lllll}
%y' &= [y;n^{-1} y;\ldots; n^{-k} y] &= f' x' &= [f n^{-\ell} x;\ldots;f x],\\
%x' &= [x;n^{-1}x;\ldots; n^{-\ell} x] &= f' y' &= [f n^{-k} y;\ldots; f y].\\
%\end{array}
%\end{displaymath}
%Using the relation $e^{-1} = n f$, it follows that $\{n^{-\ell} x, n y\}$
%and $\{n^{-k} y, n x\}$ are edges.  Both meet the node of $x$ and the node of $y$.
%In a hypermap with no mutiple joins, this implies that the edges are equal:
%\begin{displaymath}
%n y = n^{-k} y,\qquad n x = n^{-\ell} x.
%\end{displaymath}
%This shows that the dart set of $x'$ is the entire node.  This contradicts
%the assumption that ${\cal L}$ is node-free.
%\end{proof}
%

%\begin{lemma}[]\guid{ZOKKAOI}\tlabel{lemma:quotient-planar}
%Let $H$ be a plain planar hypermap, and let $\cal L$
%be a normal family.  Then $H/{\cal L}$ is a plain planar hypermap.
%\end{lemma}
%
%\begin{proof} Suppose $H/{\cal L}$ is not planar.
%Let $P$ be a M\"obius contour on $H/{\cal L}$.  It lifts uniquely to
%a contour on $H$ with the property that the darts visited on $H$ are
%precisely the darts that belong to a dart in the quotient.  This is
%compatible with the node map $n$.  So the contour path lifts to a
%M\"obius contour on $H$.  Thus, $H$ is not planar.
%\end{proof}

\section{Generation}
\indy{Index}{generation}%


The final section of this chapter presents an algorithm that generates all
simple, plain, planar hypermaps satisfying certain general conditions
(Definition~\ref{def:restricted}).  The algorithm proceeds by adding more and
more
edges and nodes to a dihedral hypermap by a sequence of reverse double walkup
transformations.

\begin{definition}[restricted]\guid{INCRVQC}\label{def:restricted}
A restricted hypermap $H = (D,e,n,f)$ is one with the following
properties.
\begin{enumerate}
\item The hypermap $H$ has no double joins, and is nonempty, planar,
  connected, plain and simple.
\item The edge map $e$ has no fixed points.  % Needed in Lemma:[flag quotient]
\item The node map $n$ has no fixed points.
\item The size of every face is at least $3$.
%%  (All hypoth. Needed?)
\end{enumerate}
\indy{Index}{hypermap!restricted}%
\indy{Notation}{H@$H$ (hypermap)}%
\end{definition}

\begin{remark}[step type]\guid{BMZYWKV}
The assumption that $e x \ne x$ implies that $f x \ne n^{-1} x$ so that $f$-steps of a 
path can be distinguished from $n^{-1}$-steps.
\end{remark}


\subsection{boolean value}
\indy{Index}{flag}%

The algorithm  marks certain faces as `true.'
Roughly, this  means that the the face cannot be modified
at any later stage of the algorithm.   When all of its faces
are true, the hypermap stands in final form.
The function that marks each face as true or false is a
\newterm{flag}.  For the algorithm to work properly, it is necessary
to impose some constraints, as captured in Definition~\ref{def:flag}.
\indy{Index}{hypermap!algorithm}%

%% XX \hat doesn't get used.

Under the bijection $\F$ between a normal family ${\cal L}$ and the set of
faces of a quotient $H/{\cal L}$, any function $\hat\varphi$ on ${\cal L}$
can be identified with a function $\check\varphi$ on the set of faces of $H/{\cal L}$.
(The {\it hat} points up to $H$, and the {\it check} points down to the quotient.)
This identification of functions will be used frequently in this section.

\begin{definition}[canonical function]\guid{CRUDEHU}
 Let $H$ be a hypermap with
  normal family $\cal L$.  The \newterm{canonical function}
  $\hat\varphi_{can}$ is the boolean-valued function on ${\cal L}$
  that is true on $L$ exactly when the dart set of $L$ maps
  bijectively to the face $\F(L)$ of $H/{\cal L}$, under $x\mapsto [x]$.  The
  corresponding function $\check\varphi_{can}$ is also called the
  canonical function.  A contour loop $L$ (or face $\F(L)$)
  is said to be canonically true or false, according to the value of the canonical
  function.  
\indy{Index}{function!canonical boolean}%
\indy{Notation}{zzP@$\check\varphi_{can}$}
\indy{Notation}{zzP@$\hat\varphi_{can}$}
\end{definition}

In other words, the face in the quotient is canonically true, exactly
when the corresponding contour loop $L\in {\cal L}$ has no $n^{-1}$
steps.  The dart set of such a contour loop $L$ is a face of $H$.  
%Based on this observation, we make the following definition.


\begin{definition}[flag]\guid{HFTAHWB}\label{def:flag} 
  Let $S$ be a set of darts in a hypermap $H$.  An $S$-\newterm{flag}
  on $H$ is a boolean-valued function $\check\varphi$ on the set of faces that
  satisfies the following two constraints.
\begin{enumerate}
\item If darts $x,y$ belong to true faces,
then there is a contour path from $x$ to $y$ that remains
in true faces.
\item Each edge of the hypermap meets a true face or $S$.
\end{enumerate}
An $\emptyset$-flag is simply called a flag.
%An isomorphism of flagged hypermaps is an isomorphism of
%hypermaps that respects the flags.
\indy{Index}{flag}%
\indy{Notation}{S@$S$ (set ofdarts)}%
\end{definition}



\begin{example}[dihedral hypermap flag] 
The dihedral hypermap of Example~\ref{ex:D2k}, carries a
flag that marks one face true and the other false.
\end{example}

\begin{example}[maximal quotient flag]\label{ex:Hall-flag} 
Let $H$ be a connected hypermap, and let $\cal L$ be the example of
Example~\ref{ex:Hall}, then the canonical map takes value
$\op{true}$ on every face.  This is a flag. In fact,
Lemma~\ref{lemma:connect-contour} provides the contour paths that
are required in the definition of flag.
\end{example}


%\begin{definition}[canonically true]\guid{PYAOBOS}
%  A contour loop $L$ in a hypermap is \newterm{canonically true} 
%if its
%  dart set is a face of $H$.
%\end{definition}


%\begin{lemma}[quotient isomorphism criterion]\guid{STKBEPH}\rating{100}
%\tlabel{lemma:all-dart}  
%  Let $H$ be a hypermap in which $e$ acts without fixed points, and
%  let ${\cal L}$ be a normal family of $H$. If the canonical boolean
%  function on the set of faces of $H/{\cal L}$ has at least as many
%  true values as there are faces of $H$, then $\cal L$ is the normal
%  family in Example~\ref{ex:Hall}. In particular, $H/{\cal L}$ is
%  isomorphic to $H$.
%\end{lemma}
%
%\begin{proof} If a face of  $H/{\cal L}$ is $\op{true}$,  then
%its darts are unit paths, and the face of $H/{\cal L}$ is in natural
%bijection with a face in $H$.  This is an injective map from the
%set of true faces of $H/{\cal L}$ to the set of faces of $H$.  The
%hypothesis of the lemma implies that this injective map is
%bijective. All of the darts of $H$ are accounted for under this
%bijection. Thus, the quotient has no false faces.  The result
%follows.
%\end{proof}
%


There is a standard way of constructing the sets $S$ of darts that
will be used in $S$-flags.  


\begin{definition}[S]\guid{FDRMSZG}
Let $H$ be a hypermap, $L$ a contour loop of the hypermap,
and $x$ an element of the dart set of $L$.
If  $L$ is  canonically true, then let $S=\emptyset$.
Otherwise,
let $m\ge0$ to be the largest $m$ 
such that 
\begin{displaymath}
[x;f x; f^2 x;\ldots;f^{m+1} x]
\end{displaymath}  
is a sublist of $L$, and
%Set $y = f^{m+1} x$
set $S(H,L,x) = \{f^i x \mid 1 \le i\le m\}$.
\end{definition}
\indy{Notation}{S@$S(H,L,x)$ (flag set)}

\begin{lemma}[flag quotient]\guid{KHGAQRG}\label{lemma:flag-set-quotient}
Let $H$ be a hypermap in which $e$ acts without fixed points, 
$L$ a contour loop, and $x$ and element of the dart set of $L$.
Let ${\cal L}$ be a normal family of $H$ that contains $L$.
Then $S(H,L,x)$ maps bijectively to a set $S'$ of darts in the quotient $H/{\cal L}$.
\end{lemma}

\begin{proof} The darts of the quotient are maximal sublists
  $[y;n^{-1} y;\ldots;n^{-k} y]$ of contour loops $L'\in {\cal L}$
  made entirely of $n^{-1}$ steps.  Each $y\in S(H,L,x)$ is preceded
  by an $f$-step and is followed by an $f$-step in $L$.  Hence the
  maximal sublist of $L$ containing $y$ is a unit path $[y]$.  The
  bijection follows.
\end{proof}


\subsection{markup}\label{sec:face-insert}
\indy{Index}{extension}%


%This section describes the transition from $H/{\cal L}$ to $H/{\cal
%M}$.  The algorithm carries along various auxiliary data satisfying
%various assumptions, enumerated as follows:

%\begin{definition}[node free]\guid{RZNSLOS}
%A  family ${\cal L}$ of contour loops in a hypermap $H$  
%is \newterm{node free}, if no node of $H$ is contained the dart set of any 
%$L\in {\cal L}$.
%\end{definition}



\begin{definition}[marked hypermap]\guid{TUFOKWK}\label{def:marked}
Let $(H,{\cal L},L,x)$ be a tuple, consisting of 
\begin{itemize}
\item a hypermap $H=(D,e,n,f)$ with no M\"obius contours, and in which
  $e$ acts without fixed points.  % Mobius needed for HQY...
\item a normal family ${\cal L}$, 
\item a contour loop $L\in{\cal L}$, and
\item a dart $x$ visited by $L$.
%and 
%\item the canonical boolean-valued function $\varphi'$ on $H/{\cal L}$
%  (identified with a function $\varphi$ on ${\cal L}$ by
%  Lemma~\ref{lemma:quotient-bijection}).
\end{itemize}
Such a tuple is a \newterm{marked hypermap} if
the following conditions hold.
\begin{enumerate}
\item The quotient $H'=H/{\cal L} = (D',e',n',f')$ is simple.  
\item $n'$ has no fixed points on $D'$.
\item $x$ is followed by an $f$-step in the loop: $L = \lp{x;fx;\ldots}$.
%\item $\varphi$ coincides with the natural boolean function on ${\cal L}\setminus \{L\}$.
%\item $\varphi$ is false on the loop $L$.
\item The contour loop $L'\in {\cal L}$ that visits%
\xfootnote{$L$ visits  $f x$.  By the definition of normality, some contour loop in
${\cal L}$ visits the dart $n f x$ at the same node.} 
$n f x$ is canonically true.
\item 
  $\check\varphi_{can}$ is an $S'$-flag on $H'$, where $S'$ is the image of 
  $S(H,L,x)$ in $D'$.  %Lemma~\ref{lemma:flag-set-quotient}.
  %under the identification of $\varphi$ with a boolean-valued
  %function on $H'$ ().
\end{enumerate}
\end{definition}

%\item $H$ is a hypermap.
%\item $\cal L$ is a normal family of $H$.
%\item $L$ is a contour loop in ${\cal L}$.
%\item $x$ belongs to the dart set of $L$.
%\item $\varphi$ is a boolean function on the faces of $H'$.
%\item The set $S(H,L,x)$ maps bijectively to a set $S'$
%of darts in the quotient.






\begin{example}[illustration]\label{ex:graph-gen}  
  This example illustrates the markup (Figure~\ref{fig:graph-gen}).
  In the figure, the hypermap $H$ is represented as a planar graph.
  The contour loops are represented by left-side shadings of the edges
  of the planar graph.  The shaded edges give the edges of a planar
  graph representing the quotient.  The polygons that are fully shaded
  are true.  Two polygons in the quotient are false.  A dart of $H$ in
  a false contour loop $L$ in $H'$ is marked $x$.  By inspection,
  $S(H,L,x)=\{f x,f^2 x,f^3 x\}$.  By inspection,
  $\check\varphi_{can}$ is an $S'$-flag.  (In fact, the darts in true
  faces form a connected region.  Every edge in the quotient except
  $\{f' x', e' f' x'\}$ meets a true face, and this one edge meets
  $S'$.)
%% XX recheck
\end{example}

\begin{figure}[htb]
\centering
\szincludegraphics[width=90mm]{\pdfp/graph_gen.eps}
\caption{An example of the current situation.}
\label{fig:graph-gen}
\end{figure}

\begin{definition}[$m$,$p$,$q$,$y$,$z$]\guid{BVUFRRE}\label{def:yz}
Let $(H,{\cal L},L,x)$ be a marked hypermap.
Several lemmas use the following natural numbers $m,p,q$ and darts $y,z$.
Set  $y = f^{m+1} x$, where $m = \card(S(H,L,x))$.
  Set
$z=f^{p+1} y$, where $p$ is the smallest natural number 
%Then there is a smallest $p\ge0$,
such that some contour loop in ${\cal L}$ visits $f^{p+1} y$.
Let $x'$ and $z'$ be the images of $x$ and $z$ respectively in $H/{\cal L} = (D',e',n',f')$.
Let $q$ be the smallest natural number such that $z' = (f')^{q+1} x'$.  
\end{definition}
\indy{Notation}{m@$m$ (face map exponent)}
\indy{Notation}{p@$p$ (face map exponent)}
\indy{Notation}{y@$y$ (dart)}
\indy{Notation}{z@$z$ (dart)}

The existence of $q$ follows from the following lemma showing that $x'$ and $z'$
lie in the same face $\F(L)$ of $H/{\cal L}$.  (The existence of $q$ is trivial
when $L$ is canonically true.)

\begin{lemma}[loop confinement]\guid{HQYMRTX}\rating{200} \label{lemma:yz}
Let $(H,{\cal L},L,x)$ be a marked hypermap.
Assume that  $L$ is canonically false. % canonically true assumption not needed.
Let the natural number $m,p$ and darts $y,z$ be given by Definition~\ref{def:yz}.
Then, $L$ visits $z$, and $z\ne f^k x$ when $0 < k \le {m+1}$.
\end{lemma}

\begin{proof} 
%We break the proof into two cases, depending on whether $L$ is canonically true.
%
%\claim{[$L$ is canonically true.]}  In this case $S=\emptyset$,  $m=0$, $p=0$, $k=1$,
%$y = f x$, and $z = f y = f^2 x$.  Also, $L = \lp{x;f x;f^2 x;\ldots}$, and it clearly
%visits $z$.  If $z=f^k x = y$, then $x=y=z$ is a fixed point of $f$.
%
%This completes the first case.
%
%\claim{[$L$ is canonically false.]} In this case,
  For a contradiction, suppose $f^{p+1} y = f^k x$, for some $0<k\le
  m+1$.  Then also, $f^p y = f^{k-1} x$.  If $p>0$, then this
  contradicts the minimality of $p$.   (Note that $L$ visits $f^{k-1}x$ by the
definition of $S$ and $m$.)  So $p=0$, and $y=f^{k-1} x =
  f^{m+1} x$.  Also, $0\le k-1 < {m+1}$, which implies that the face of $x$ has size at
most $m+1$.  This  forces $L$ to be canonically true, which is contrary to assumption.  This proves the second
  conclusion of the lemma.  In particular, $z\not\in S$, where $S =
  S(H,L,x)$.

Let $L'$ be the contour loop of ${\cal L}$ that visits $z$.  For a contradiction,
assume that $L'\ne L$.

\claim{$L'$ is false.}  Otherwise, $L'$ is true with respect to the
canonical flag and it therefore a loop consisting entirely of
$f$-steps.  In particular, $L'$ visits $z,x$, and $y$.  This is
contrary, to the assumption that the contour loop containing $x$ is
false.

Let $H' = (D',e',n',f') = H/{\cal L}$, and let $S'$ be the image of $S$ in $D'$.
Let $z' = [\ldots;u]\in D'$ be the
image of $z$ in $D'$.    
As $z\not\in S$, we also have $z'\not\in S'$.
By the definition of
$S'$-flag, the dart $e'z'$ lies in a true face or $e'z'\in S'$.  
This disjunction splits
splits the proof into two cases.
\begin{enumerate}
\item\claim{[$e'z'$ lies in a true face.]}  In this case, 
\begin{displaymath}
e'z' = f^{-1} n^{-1} [\ldots,u] = f^{-1} [n^{-1}u;\ldots],
\end{displaymath}
so that $n^{-1} u$ is visited by a true contour loop.
Consider the following
path in $H$:
\begin{displaymath}
[y;fy;\ldots;z] @ [n^{-1}z;\ldots;u] @ [n^{-1} u;\ldots;n^{-1} x].
\end{displaymath}
The first segment consists of $f$-steps; the second of $n^{-1}$-steps;
and the third segment exists within true contour loops of ${\cal L}$
by the connectedness of true faces (by properties of flags).  This
path satisfies all the assumptions of Lemma~\ref{lemma:contour-f}.  (In particular,
the node of $n f x$ (or of $f x$)  is distinct from the node of $x$ by the
assumed simplicity of the quotient $H/{\cal L}$.)
The lemma asserts that the path does not exist.
\item 
\claim {[$e'z'\in S$.]}  In this case,  
$f^{-1}n^{-1}u \in S$ and $L$ visits $n^{-1} u$ at the node of $z$.
Consider the following path of $f$-steps in $H$:
\begin{displaymath}
[y;f y;\ldots;z].
\end{displaymath}
This path satisfies all the enumerated conditions of
Lemma~\ref{lemma:contour-f}.  (In particular, $y$ and $z$ are at
different nodes by the assumed simplicity of the quotient $H/{\cal
  L}$.)  The lemma asserts that the path does not exist.
\end{enumerate}
\end{proof}

\begin{lemma}[parameters]\guid{QRDYXYJ}\label{lemma:parameters}
Let $(H,{\cal L},L,x)$ be a
marked hypermap, where $H$ is restricted. Assume that $L$ is canonically false.
Let $m,p,q$ be the natural numbers and let $x,y,z$ be the darts given by
Definition~\ref{def:yz}.  Let $r = \op{card}(\F(L))$.  Then
\begin{displaymath}
0\le p,\quad 0\le m < q < r,\quad m+1 < p+q.
\end{displaymath}
Furthermore, the darts $x$, $y$ belong to different nodes; the darts
$y$ and $z$ belong to different nodes.
\end{lemma}

\begin{proof}
Let $H/{\cal L}=(D',e',n',f')$.  Let $y'$ and $z'$ be the images of $y$ and $z$
in $D'$, respectively.  Let $S'$ be the image of $S(L,x)$ in $D'$.
By definition $m = \op{card}(S(L,x))$.  Both $m$ and $p$ are natural numbers,
so $0\le p$ and $0\le m$. 

\claim{The darts $x$, $y$ belong to different nodes; the darts $y$ and
  $z$ belong to different nodes.}  Indeed, $x$, $y$, and $z$ belong to
the same face.  By the simplicity of $H$, if two of these darts belong
to the same node, then they are equal to each other.  However, $x\ne
y$, for otherwise the subpath $P=[x;f x;\ldots;f^{m+1}x]$ gives a canonically true contour loop, which is contrary to the assumption that $L$ is
canonically false.  Also, $y\ne z$, for otherwise the face of $y$ is
equal to $\{y,f y,f^2 y,\ldots,f^p y\}$.  It follows that $x = f^k y$
is visited by $L$, for some $1<k\le p$.  This contradicts the defining
minimality property of $p$.

\claim{$q<r$.}  Indeed,
by definition, $z' = (f')^{q+1} x'$ and no smaller natural
number has this property.  Also, $x'$ and $y'$ belong to the face $\F(L)$.
If $q\ge r$, then $(f')^{q+1} = (f')^{q-r+1}$, which contradicts the minimality of $q$.

\claim{$m<q$.} Indeed, if $0\le k< m$, then
\begin{displaymath}
(f')^{k+1} x' = [f^{k+1} x]\in S', \quad z' \not\in S',
\end{displaymath}
by Lemma~\ref{lemma:flag-set-quotient} and Lemma~\ref{lemma:yz}.
Thus, $0\le q< m$.  Also, $q\ne m$, for otherwise $z' = (f')^{m+1} x'
= y'$.  This implies that $z$ and $y$ lie in the same node, which has
been proved impossible.  This completes the proof that $m<q$.
 
\claim{$m+1 < p+q$.}  Indeed, the inequalities $0\le p$ and $m<q$ imply
that either $m+1 < p+q$ or $p=0~\land~m+1=q$.  The second disjunct
cannot hold, for otherwise $z' = (f')^{q+1} x' = f' y'$.  Write $y' = [y;\ldots;u]$.
This is not a unit path by the definitions of $S(L,x)$ and $m$, so $y\ne u$, but
they lie at the same node.  Also from $p=0$ it follows that $z= f^{p+1} y = f y$.
So $e y$ and $e u$ both lie at the node of $z'$.  The existence of
 two edges, $\{y, e y\}$ and
$\{u, e u\}$, between two nodes contradicts the hypothesis
on $H$ of no double joins.  This proves the claim and the lemma.
\end{proof}

\subsection{transform}

\begin{definition}[transform]\guid{YQANQNF}
  From one marked hypermap $(H,{\cal L},L,x)$ in which $L$ is
  canonically false, we construct a new tuple
\begin{displaymath}
T(H,{\cal L},L,x) = (H,{\cal M},L_1,x),
\end{displaymath}
 called the \newterm{transform} of
  $(H,{\cal{L}},L,x)$.  
As the notation indicates, the hypermap $H$ and the dart $x$ are the same for both
tuples.  The data ${\cal M}$ and $L_1$ are specified in
the following paragraphs.
\end{definition}

Let $m$, $p$, $y$, and $z$ be  given by
Definition~\ref{def:yz}.
Let $L_1$ be 
%With improvised notation, write
%\begin{displaymath}
%L_1 = \lp{L[x:y] \opat P[y:z] \opat L[z:x]},
%\end{displaymath}
the contour loop in $H$ that follows $L$ from $x$ to $y$, then takes
$f$-steps from $y$ to $z$, then continues along $L$ back to $x$.  
Let $L_2$ be
%\begin{displaymath}
%L_2 = \lp{L[n^{-1}y:n z] \opat P^c[n z:n^{-1}y]},
%\end{displaymath}
the contour loop in $H$ that follows $L$ from $n^{-1} y$ to $n z$,
then complements the path of $L_1$ from $y$ to $z$, traveling instead
from $n z$ to $n^{-1} y$. 


Set
\begin{displaymath}{\cal M} = ({\cal L}\setminus \{L\}) \cup
\{L_1,L_2\}.\end{displaymath}

\begin{remark}[canonical compatibility]\guid{HBLIYVM}
There is a canonical boolean function on ${\cal L}$ and one on ${\cal M}$.
The canonical boolean functions agree on the intersection ${\cal L}\cap {\cal M}$.
This means, there is a well defined boolean-valued function on 
${\cal L}\cup {\cal M}$.  There is no ambiguity.  
\end{remark}
\indy{Index}{function!boolean}%

%  Under the bijection between the set of faces of a quotient $H/{\cal
%    L}$ and the family ${\cal L}$ (of
%  Lemma~\ref{lemma:quotient-bijection}), an $S$-flag on a quotient
%  $H/{\cal L}$ can be identified with a boolean function on ${\cal L}$
%  satisfying appropriate properties.  With this in mind, 
%\begin{displaymath}
%\begin{cases}
%\psi(L) = \varphi(L), &\text{if } L\ne L_1,L_2,\\
%\psi(L_2) = \op{true}, &\text{if $L_2$ is canonically true},\\
%\psi(L) = \op{false}, &\text{otherwise}.\\
%\end{cases}
%\end{displaymath}




\begin{figure}[htb]
\centering
\szincludegraphics[width=80mm]{\pdfp/L1L2.eps}
\caption{The loop $L$ is replaced with two loops $L_1, L_2$.}
\label{fig:L1L2}
\end{figure}



%Each of the two
%loops $L_i$ meets at least two nodes (those
%of $y$ and $z$) and has length at least three by
%Lemma~\ref{lemma:3dart}.  
% $S(H,L_1,x)$ contains the proper subset $S(H,L,x)$.


\begin{lemma}[markup transform]\guid{AQIUNPP}\rating{600}\tlabel{lemma:flag} 
Let $H$ be a restricted hypermap.
If $(H,{\cal L},L,x)$ is a marked hypermap such that $L$
is canonically false,  then the transform
$(H,{\cal M},L_1,x)$ 
is also a marked hypermap.
\end{lemma}

\begin{proof} Let 
\begin{displaymath}
H=(D,e,n,f),\   H' =(D',e',n',f') = H/{\cal
    L},\  H'' = (D'',e'',n'',f'') = H/{\cal M}.   
\end{displaymath}
Let $S'$ be the image of $S(H,L,x)$ in $D'$.  Let $y,z$ be the darts
constructed in Definition~\ref{def:yz} from the marked hypermap $(H,{\cal
  L},L,x)$.  The dart $z$ is not at the same node as $y$ (by
the simplicity of $H$).

  The proof can be organized into independent parts, according to the separate
  properties of a marked hypermap.  The first part of the proof
  establishes that ${\cal M}$ is a normal family.


\case{normal-1} \claim{No dart is visited by two different loops.}
Indeed by construction, the sets of darts of $L_1$ and $L_2$
are disjoint from each other and disjont from the sets of darts of $L'\in
{\cal L}\setminus \{L\}$.  The result now follows from the normality of  ${\cal L}$.

\case{normal-2} \claim{Every loop visits at least two nodes.}  Indeed, this
is true for $L_1$ and $L_2$ because they visit the nodes of $y$ and
$z$.  It is true of the other loops because they belong to the
normal family ${\cal L}$. 

\case{normal-3} \claim{If a loop visits a node, then every dart at
that node is visited by some loop.} 
% Indeed, the new nodes visited on the
%sublist of $L_1$ from $y$ to $z$ do not contain any darts in any
%loop in ${\cal L}$.  (By the definition of normal family, if a loop
%visits a node, then every dart at that node is visited by some loop
%of ${\cal L}$.)  These new nodes contain two darts, with one dart
%along $L_1$ and the other along $L_2$.  The paths $L_1$ and $L_2$
%visit every dart visited by $L$.  They visit every dart at the nodes
%of $y,\ldots,z$.
Indeed, the nodes that are visited by some loop in ${\cal M}$ are
precisely those visited by some loop in ${\cal L}$, together with the
``new'' nodes; that is, the nodes of $f y,\ldots,f^p y$.  The set of
darts that are visited by some loop of ${\cal M}$ is the union of the
set visited by loops in ${\cal L}$, together with two darts at the
new nodes.  As each new nodes has only two darts, and as ${\cal L}$ itself
it normal, the result follows. It follows that ${\cal M}$ is normal.  


\case{simple} To prove the simplicity of the quotient, it is enough to show that
none of the contour
loops in ${\cal M}$  ever return to a node after leaving it.
 (More precisely, the dart set of any $L'\in{\cal M}$ intersected with a node
is the dart set of a maximal sublist $[z;n^{-1}z;\ldots;n^{-k}z]$ of $n^{-1}$ steps.)
This is true of $L'\in{\cal L}\setminus\{L\}$ by assumption and true of
$L_1$ and $L_2$ by construction.  Simplicity follows.

\case{fixed-point free} By Lemma~\ref{lemma:nfp}, to prove that $n''$
does not have a fixed-point, it is enought to show that no loop in
${\cal M}$ has a dart set containing a node.  It is sufficient to
consider the loops $L_1$ and $L_2$.  The set of darts of $L_1$ and
$L_2$ at the old nodes (that is, those not meeting $\{f y,\ldots,f^p
y\}$) are subsets of the set of darts of $L$ at those nodes.  As the
dart set of $L$ does not contain an old node, neither do $L_1$ and
$L_2$.  At the new nodes, $L_1$ and $L_2$ both have at least one dart,
so neither contains the entire node.  It follows that $n''$ is
fixed-point free.

\claim{$e'y'\not\in S'$, where $y'$ is the image of $y$ in $H'$. }
Otherwise, write $y' = [y;n^{-1}y;\ldots;u]$, a sublist of $L$, and
pick $k$ such that $e'y' = [f^k x]\in S'$.  Then
\begin{displaymath}
(n')^{-1} y' = f'e'y' = f'[f^k x] = [f^{k+1}x;\ldots].
\end{displaymath}
By the construction of $S(H,L,x)$, we know that $L$ visits $f^{k+1}x$.
Hence $y'$ and $(n')^{-1}y'$ both lie in the same node and in the same
face $\F(L)$.  By the simplicity of $H'$, it follows that $y' =
(n')^{-1} y'$.  That is, $y'$ is a fixed point of $n'$.  This is contrary to
assumption. The claim follows.


\claim{$e'y'$ lies in a true face of $H'$.}  Indeed,
 since $\check\varphi_{can}$ is
an $S'$-flag on $H'$, the edge $\{y',e'y'\}$ meets a true face or
$S'$.  However, $y'\not\in S'\subset \F(L)$, by the simplicity of $H$.  Also,
$e'y'\not\in S'$, by the previous paragraph.
The dart $y'$ lies in the false face $\F(L)$.  The only
remaining possibility is that $e'y'$ lies in a true face.  



\case{flag-1} \claim{The true faces of $H''$ are
  connected.}  Indeed,  $L_1$ is connected to a true face
by the contour path $[x;n^{-1} x]$, because $n^{-1} x$ lies
in the same face as $n f x$, which is a true face by assumption.
If $L'\in{\cal M}\setminus\{L_1,L_2\}$ is a true, then $L'\in  {\cal L}$, and
it connects with the true faces of ${\cal L}$ as before.
If $L_2$ is true, then the proof requires more
argument.   Write $y'=[y;\ldots;u]$ as above.
The dart $u''=[n^{-1}y;\ldots;u]$ of $H''$ lies in the face $\F(L_2)$.
Also, 
\begin{displaymath}
(n'')^{-1} u''= [n^{-1} u;\ldots] = [f e u;\ldots] = f'' [\ldots; e u]
  = f'' e' y'
\end{displaymath}
Thus, we have a contour path from $u''$ to $e'' y'$, which lies in a
true face, by an earlier claim.  (The dart $e'y'$ is naturally
identified with the dart $e'' y'$ on in $H''$, because the faces are
true.)  Hence a path exists from the true face $\F(L_2)$ into another
true face, and from there any true face may be reached.

\case{flag-2} \claim{Each edge of $H''$ meets a true face or $S''$, where $S''$ is
the image of $S(H,L_1,x)$ in $D''$.}
Indeed, the function $\check\varphi_{can}$ is an $S'$-flag on $H'$.  The edges of
$H'$ can be identified with a subset of the edges of $H''$.  For this
subset, the flag condition on edges is immediate.  Consider the
new  edges  (that is, edges of $H''$ that are not in $H'$).
They all meet $\{y,f y,\ldots,f^p y\}$.  This set is either contained
ini $S(H,L_1,x)$ (when $F_1$ is false), or is contained in the true face
$F_1$ (when $F_1$ is true).  Hence each new  edge meets
a true face or $S''$.


The other verifications are routine.
\end{proof}


\subsection{digraph}

The aim is to prove that every restricted hypermap with a given bound
on the size of the dart set is generated by a particular algorithm.
The proof is a long induction argument.  The proof starts by showing
how to go from one partially constructed hypermap to another more
fully constructed hypermap.  The hypermap $H$ represents the fully
constructed hypermap and two quotients $H/{\cal L}$ and $H/{\cal M}$
represent the partially constructed hypermaps.  The algorithm involves
the transition from the hypermap $H/{\cal L}$ to $H/{\cal M}$.  The
transition from one quotient to another is given by the transform of
marked hypermaps.  The transform thus represents one step of the
algorithm.  \indy{Notation}{M@$\cal M$ (normal family)}%


\begin{definition}[digraph,~vertex,~edge,~head,~tail,~sink,~path]\guid{AVXKIRW}
  A \newterm{digraph} (directed graph) is an ordered pair $(V,E)$
  where $V$ is any set, and $E$ is a set of ordered pairs of vertices.
  An element of $V$ is called a \newterm{vertex}.  An element of $E$
  is called a \newterm{directed edge}.  If $(v,w)\in E$, then $v$ is
  the \newterm{head} and $w$ is the \newterm{tail} of the directed
  edge.  A vertex $v$ is a \newterm{sink} if it is not the head of any
  directed edge.  A path $P=[v_0;v_1;\ldots;v_{k-1}]$ in a digraph is
  a list of vertices, such that $(v_i,v_{i+1})\in E$ for all $i<k-1$.
\end{definition}

\begin{definition}[digraph of a hypermap]\guid{QQHIKFL}
Let $H$ be a restricted hypermap.  Form a directed graph as follows.
The vertex set $V$ of the digraph is the set of all marked hypermaps 
${\cal H}=(H,{\cal L},L,x)$ such that $L$ is canonically false.  
Write $T{\cal H} = (H,{\cal M},L_1,x)$, where $T$ is the transform on $V$.  
The set of tails of directed edges with head ${\cal H}$ is as follows:
\begin{itemize}
\item If every contour loop of ${\cal M}$ is canonically true, then ${\cal H}$ is a sink.
\item If $M$ is canonically false, 
then there is a single tail $T{\cal H}$.
\item If $M$ is canonically true, but not every loop in ${\cal M}$ is canonically true,
then the tails are $(H,{\cal M},M,y)$, where $M$ is a 
canonically false loop in ${\cal M}$ and $y$ is a dart visited by $M$ such
that $y$ is followed by an $f$-step:  $M = \lp{y;f y;\ldots}$.
\end{itemize}
\end{definition}

\begin{lemma}[digraph sink]\guid{XCOXWYJ}\label{lemma:digraph-sink}
Let $(V,E)$ be the digraph of a restricted hypermap $H=(D,e,n,f)$.  Then every path
in $(V,E)$ reaches a sink after at most $\#D$ steps.  Moreover, if ${\cal H}$ is
a sink, and if $T{\cal H} =(H,{\cal M},L_1,x)$ is its transform, then
$H$ is naturally isomorphic to $H/{\cal M}$, under the map that sends
a dart $y$ of $H$ to the unit path $[y]$.
\end{lemma}

\begin{proof} Each step in the path makes one transform.  Each
  transform increases the number of darts visited by the normal
  family.  The number of darts visited by a normal family is bounded
  by $\#D$.  This bounds the length of a path.

  The condition on a sink ${\cal H}$ is that every contour loop of
  ${\cal M}$ is canonically true. 

\claim{${\cal M}$ visits every dart.}  Indeed, let $y$ be any dart of
$H$. Since $H$ is assumed connected, there exists some contour path
$P=[x;\ldots;y]$.  Let $u$ be the last dart on the path that is
visited by ${\cal M}$.  Let $M\in{\cal M}$ be the contour loop that
visits $u$.  Then $u$ is not followed by an $f$-step, because every contour
loop is canonically true: $M = \lp{u;f u;\ldots}$.
Nor is $u$ followed by a $n^{-1}$ because in a normal family every dart
at the node of $u$ is visited by a loop in ${\cal M}$.  Thus, $u$ is the final
dart in the path $P$, which means that $y$ is visited by ${\cal M}$.

It follows that is a bijection between the dart set of $H$ and the
dart set of $H/{\cal M}$, sending a dart $y$ to the unit path $[y]$.
This bijection induces an isomorphism of hypermaps.
\end{proof}


%In the next lemma there are a number of choices to be made.  Let
%$\op{ch}$ be any function on the set of normal families ${\cal L}$ of $H$ containing
%at least one canonically false loop, returning a dart $x'$ of $H/{\cal L}$
%in a false face (with respect to the canonical boolean function).
%Given a normal family ${\cal L}$ (with a canonically false loop),  write
% $x' = [\ldots;x]$ and let $L$ be the the canonically false contour loop of ${\cal L}$ that
%visits $x'$.  Write ${\cal L}\rightsquigarrow (L,x)$.   The dart $x$ of $H$ has the property
%that $L = \lp{x;f x;\ldots}$; that is, $x$ is followed by an $f$-step in $L$.
%
%
%\begin{lemma}[sequence]\guid{YCFWVQJ}\label{lemma:sequence}  
%  Let $H$ be a restricted hypermap.  Fix a choice function $\op{ch}$ as
%  above, and let $F$ be any face of $H$. Then associated with $(\op{ch},F)$, there is a sequence of
%  marked hypermaps ${\cal H}_i = (H,{\cal L}_i,L_i,x_i)$ for
%  $i=0,\ldots,k-1$ such that
%\begin{itemize}
%\item $\varphi'_i$ is the canonical boolean function on $H/{\cal L}_i$, and
%$\varphi_i$ is its lift to ${\cal L}_i$.
%\item ${\cal L}_0$ is the normal family associated with the given face $F$,
%described in Example~\ref{ex:H2}, and ${\cal L}_0 \rightsquigarrow (L_0,x_0)$.
%\item If the contour loop $L$ in the transform $T{\cal H}_i=(H,{\cal M},L,\ldots)$ 
%is canonically false, then ${\cal H}_{i+1} = T{\cal H}_i$.
%\item If the contour loop $L$ in the transform $T{\cal H}_i=(H,{\cal
%    M},L,\ldots)$ is canonically true, and there exists some other
%  contour loop of ${\cal M}$ that is canonically false, then ${\cal
%    H}_{i+1}=(H,{\cal M},M,y,\varphi_{i+1})$, where ${\cal M}\rightsquigarrow (M,y)$. 
%\item The transform $(H,{\cal L}_k,\ldots)$ of $\,{\cal H}_{k-1}$ has
%  quotient $H/{\cal L}_k$ that is naturally isomorphic to $H$.
%\end{itemize}
%\end{lemma}
%
%\begin{proof}
%  Let ${\cal H}_0=(H,{\cal L}_0,L_0,x_0,\varphi_0)$ be given as
%  follows.  Pick any face $F$ of $H$, and construct the normal family
%  ${\cal L}$ of Example~\ref{ex:H2}.
%
%  The construction of ${\cal H}_{i+1}$ depends on the structure of the
%  transform $(H,{\cal M},L,x)$ of ${\cal H}_i$.  We consider
%  three cases.
%\begin{nomerate}
%\item 
%\claim{[Every contour loop of ${\cal M}$ is canonically true.]}   In this case,  by
%Lemma~\ref{lemma:all-dart}, the quotient map $H\to H/{\cal M}$ is an
%isomorphism.  Set $k= i+1$. The sequence terminates.
%\item \claim{[The contour loop $L$ is canonically false.]}  In this
%  case, from the inductive hypothesis that $\varphi'_i$ is the
%  canonical function and the definition of $\psi$, it follows that
%  $\psi'$ is the canonical function.  By Lemma~\ref{lemma:flag}, the
%  transform is a marked hypermap.  Let ${\cal H}_{i+1}$ equal the
%  transform of ${\cal H}_i$.
%\item \claim{[The contour loop $L$ is canonically true, but not all contour loops in
%     ${\cal M}$ are canonically true.]}  In this
%  case, the canonical function on $H/{\cal M}$ is a flag (with
%  $S=\emptyset$).   
%${\cal H}_{i+1}$, as defined, is a marked
%  hypermap.
%\end{nomerate}
%
%The sequence must terminate eventually, because each step
%constructs a quotient of $H$ with more darts than the previous one and
%the number of steps is bounded by the size of the dart set of
%$H$.
%\end{proof}

Next, we wish to describe the directed edges in a way that relies to a
lesser degree on the structural details of the marked hypermaps.
(These details will not be available to us in the algorithm of the
next subsection.)  The next lemma uses reverse doble walkup
transformations to construct a new hypermap from a given hypermap $H'$
that does not require us to represent it first as a quotient $H' =
H/{\cal L}$ for some normal family.
We can immediately relate the to the digraph we
have constructed.   

\begin{lemma}[walkup-digraph]\guid{ISMLATS}\label{lemma:RDW}
Let ${\cal H} =(H,{\cal L},L,x)$ be a marked hypermap,
  with $H$ restricted.  Assume that $L$ is canonically false, and let 
$(H,{\cal{M}},M,x)$ be 
  the transform of $\cal H$.   
Let $m$, $p$, $q$, $y$, and $z$ be
  the constants of Definition~\ref{def:yz}.  Let $x'$ be the
  image of the dart $x$ in $H/{\cal L} = (D',e',n',f')$.
%Then $m,p,q$ satisfy the constraints of
%  Definition~\ref{def:R}, and 
Then $H/{\cal M}$ is isomorphic to $RDW(H/{\cal{L}},x',m,p,q)$.
\end{lemma}


\begin{proof}
%The constraints on $(m,p,q)$ of Definition~\ref{def:R} are satisfied 
%by Lemma~\ref{lemma:parameters}.
%
  By construction, the passage from $H/{\cal M}$ to $H/{\cal L}$
  consists of a double walkups to eliminate the nodes (of size two) at
  $f y$, $f^2 y, \ldots, f^p y$, and then a double walkup to eliminate
  the edge that runs from the the node of $y$ to the node of
  $z$.  %The passage in the other direction from $H/{\cal L}$ to
%$H/{\cal M}$ comes by adding an edge from the node of $y$ to
%the node of $z$ and inserting $p$ new nodes (of degree two)
%along it.
If we play these double walkups in reverse,
one may also pass from $H/{\cal L}$ to $H/{\cal M}$.  
%If
%$m,p,q$ are chosen as above, 
Then $RDW(H/{\cal L},x',m,p,q)$ is isomorphic to
$H/{\cal M}$.  
\end{proof}


\begin{figure}[htb]
\centering
\szincludegraphics[width=80mm]{\pdfp/L1L2dart.eps}
\caption{$H/{\cal L}$ is obtained from $H/{\cal M}$ by double walkup
transformations.}
\label{fig:L1L2dart}
\end{figure}



\subsection{algorithm}

This final section puts the algorithm a precise form, based on
the Knaster-Tarski fixed point theorem.  The Knaster-Tarski fixed
point theorem is a common way to give precise mathematical form to
an algorithm.

\begin{lemma}[Knaster-Tarski]\guid{EAOGWLE}\rating{ZZ}   
Let $X$ be a set.  Let $f:\powerset(X)\to \powerset(X)$ be a
function from the powerset of $X$ to itself.  Assume that $f$ is
monotonic in the sense that whenever $Y\subset Z\subset X$, it
follows that $f(Y) \subset f(Z)$.  Then $f$ has a least fixed point.
That is there exists a set $\op{fix}(f,X)\subset X$ such that
$f(\op{fix}(f,X)) = \op{fix}(f,X)$ and such that the following
minimality condition holds: if $Y\subset X$ is any set such that
$f(Y) \subset Y$, then $Y\subset \op{fix}(f,X)$.
\end{lemma}
\indy{Notation}{X@$X$ (set)}%
\indy{Notation}{f@$f$ (function on powerset)}%
\indy{Notation}{fix@$\op{fix}$~(Knaster-Tarski fixed point)}%
\indy{Notation}{P@$\powerset(\cdot)$ (powerset)}%
\indy{Index}{Knaster-Tarski fixed point theorem}%

\begin{proof} Let $\op{fix}(f,X)$ be the intersection of all subsets
$Y$ of $X$ such that $f(Y)\subset Y$.  It is easily verified that
this set has the required properties.
\end{proof}

Various data are needed for the application of the Knaster-Tarski
fixed-point theorem to the construction of restricted hypermaps.  This
data is presented in a series of definitions.  The first definition
gives a domain $\Omega$ that will contain all the darts of all the
hypermaps that will be constructed by the algorithm.  In practice, we
take $\Omega$ to be a finite subset of the set of natural numbers.

\begin{definition}[$\Omega$,~$\op{ch}$,~$d$]\guid{LVXTTSP}
  Let $\Omega$ be any fixed finite set.  Fix a choice function
  $\op{ch}:\powerset(\Omega)\to \Omega$ that picks an element from
  each nonempty subset:
\begin{displaymath}
X\ne\emptyset\quad  \Rightarrow \quad  \op{ch}(X)\in X.
\end{displaymath}
For example, when
$\Omega$ is a well-ordered set, let $\op{ch}$ choose the least element of a subset
of $\Omega$.  An $(\Omega,d)$-hypermap is a hypermap whose dart set is a
subset of $\Omega$ and such that $d$ is the maximum face size.  A
$(\card(\Omega),d)$-hypermap is one isomorphic to an $(\Omega,d)$-hypermap.
(In applications later in the book, $\card(\Omega)\le14$ and $3\le d\le 6$.) 
\end{definition}
\indy{Notation}{d3@$d$ (upper bound)}%
\indy{Notation}{zzZ@$\Omega$ (set of darts)}%
\indy{Notation}{ch@$\op{ch}$~(choice)}%


The following constructions depend on $\Omega$ and $d$, although the
notation does not reflect this.  One may think of $X_1$ in the next
definition as holding the output of the algorithm and $X_2$ as the
workspace that holds partially constructed hypermaps.

\begin{definition}[$X$,~$X_1$,~$X_2$]\guid{KFPEPWO}
Define a set $X$ as the disjoint union of $X_1$ and $X_2$ as follows.
Let $X_1$ be the set of all $(\Omega,d)$-hypermaps.
Let $X_2$ be the set of tuples $(H,m,\check\varphi,x)$, where 
\begin{itemize}
\item $H$ is an $(\Omega,d)$-hypermap,
\item $\check\varphi$ is an $S$-flag on $H$,
\item  $x$ is a dart in a false face of $H$ (with respect to $\check\varphi$),
\item $m\in\{0,\ldots,d-1\}$, 
and
\item $S = \{f^i x\mid 1 \le i \le m\}$.
\end{itemize}
\end{definition}


The set $A$ in the following definition gives the initialization of the algorithm.

\begin{definition}[A]\guid{JBUOJMF}
Let $H$ be a fixed hypermap isomorphic to $\op{Dih}_{2d}$, with darts in $\Omega$.
Let $\check\varphi$ be the flag on $H$ (with one true face and one
false face).  Let $x$ be the value of the choice function on the false
face.  Set
\begin{displaymath}
A = \{(H,m,\check\varphi,x) \mid 0\le m \le d-1\} \subset X_2.
\end{displaymath}
\end{definition}

The following set gives the indexing set for the iteration of the algorithm.

\begin{definition}[C]\guid{IDDKWYX}
Let $C$ be the set of of $4$-tuples $(m,p,q,r)$ that satisfy the following
constraints:
\begin{itemize}
\item $0\le m < q < r$.
\item $0\le p$.
\item $m+1 < p+q$.
\item $m+p+2 \le d$.
\item if $q+1< r$, then $m+p+3\le d$.
\end{itemize}
\end{definition}

\begin{definition}[extension]\guid{ZMVBANY}  
  Let $(H,m,\check\varphi,x)\in X_2$.  Choose $p,q$ such that $(m,p,q,r)\in
  C$, where $r$ is the cardinality of the face of $x$.  Let $F$ be the
  face of $x$ in $H$.  It follows by construction, that every face
  $F'\ne F$ of $H$ is naturally identified with a face of
  $RDW(H,x,m,p,q)$.  Say that a boolean-valued function $\check\psi$ on the
  set of faces of $RDW(H,x,m,p,q)$ is an \newterm{extension} of
  $\check\varphi$ on $H$ if $\check\psi(F') =\check\varphi(F')$, when $F'\ne F$.
  \indy{Index}{extension}%
\end{definition}


The functions $f$ and $g$ give one iteration of the algorithm.  The
powerset-valued function $g$ is the heart of the algorithm.  It takes
one partially constructed hypermap and modifies it in various ways to
construct further partially constructed hypermaps.  When the flag
$\check\varphi$ permits, some hypermaps are also fully constructed.


\begin{definition}[g]\guid{DMAMRYR}
 Let $g:X_2 \mapsto \powerset(X)$ be given as
  follows.  Let $r$ be the cardinality of the face $F$ of $x$.  The
  subset $g(H,m,\check\varphi,x)\subset X$ is presented as a union of two
  sets:
\begin{displaymath}
   Y_i = X_i \cap g(H,m,\check\varphi,x).
\end{displaymath}
 $H'\in Y_1$ if and only 
\begin{itemize}
\item there exists
$p,q$ with $(m,p,q,r)\in C$ such that $H'=RDW(H,x,m,p,q)$,
and 
\item $\check\varphi(F')$ is true for all $F'\ne F$.
\end{itemize}
 $(H',m',\check\psi,x')\in Y_2$ if and only if
\begin{itemize}
\item there exists $p,q$ with $(m,p,q,r)\in C$, such that $H'=RDW(H,x,m,p,q)$,
\item  $\check\psi$ is an extension of
$\check\varphi$, and 
\item Let $F_1$ be the face of $x$ in $H'$.  One of the following two
  conditions hold:
\begin{itemize}
\item $\check\psi(F_1)$ is false;  $x' = x$; and  $p+m+1 \le m' < r$.
\item $\check\psi(F_1)$ is true; there exists a false face in $H'$; $x'$ is
  the value of the choice function on the union of false faces of
  $H'$; and $0 \le m' < r$.
\end{itemize}
\end{itemize}
\end{definition}


\begin{definition}[f]\guid{YSJTEDX}
Given the function 
$g:X_2 \to \powerset(X)$, set 
\begin{displaymath}f(S) = A \cup (\bigcup \{g(s) \mid s\in S\cap
X_2\}).\end{displaymath}
\indy{Notation}{g@$g$(function)}%
\end{definition}

Any function $f :\powerset(X)\to \powerset(X)$ of this form is
monotonic.  Thus, we have a Knaster-Tarski fixed point set
$\op{fix}(f,X)$.  The main result of this chapter is that a fixed
point construction generates all restricted hypermaps:

\begin{theorem}[hypermap algorithm]\guid{BRGEFNH}\rating{2000}  
\label{lemma:algorithm}
Define $f $ and $X$ as above (depending on $d\ge 3$ and $\Omega\ne
\emptyset$) .    Then every restricted hypermap with at
most $\card(\Omega)$ darts and such that the largest face has size  $d$
is isomorphic to a hypermap in $\op{fix}(f,X)\cap X_1$.
\end{theorem}


In informal terms, by starting with the \newterm{seed} hypermaps in $A$
one may find all restricted hypermaps (for given $\Omega$ and $d$) by
applying the function $f$ repeatedly:
\begin{displaymath}
A_0 = A = f(\emptyset),\quad A_1 = f(A_0),\quad A_2 = f(A_1),\ldots
\end{displaymath}
and by looking at the output $A_i \cap X_1$.
\indy{Index}{seed}%

The proof will be presented below.  The proof is a matter of correlating the Knaster-Tarski fixed point set with the digraph of a restricted hypermap $H$.
Write $\op{fix}_i$ for $\op{fix}(f,X)\cap X_i$.


\begin{definition}[correlation]\guid{SFBFNVW}
  A marked hypermap $(H,{\cal L},L,x)$ is said to be
  \newterm{correlated} to an element $(H',m',\check\varphi,x')\in X_2$ if
  the following conditions hold:
\begin{itemize}
\item $H/{\cal L}$ is isomorphic to $H'$ by some isomorphism $G$.
\item The image of $x$ in $H/{\cal L}$ maps to $x'$ under $G$.
\item The pull back of $\check\varphi$ under $G$ is the canonical function
$\check\varphi_{can}$ on $H/{\cal L}$.
\item $m = \card(S(H,L,x))$.
\end{itemize}
\end{definition}

\begin{lemma}[correlated seed]\guid{NRDWGYQ}\label{lemma:correlated-seed}
  Let $H$ be a restricted $(\card(\Omega),d)$-hypermap with digraph
  $(V,E)$.  There exist a marked hypermap in
  $V$ and an element in $\op{fix}_2$ that are
  correlated.
\end{lemma}

\begin{proof}  From the definition of $f$, it follows
that $A\subset \op{fix}_2$.  Thus, it suffices to correlate a marked
hypermap with an element of $A$.  Let $(H',\cdot,\check\varphi,x')\in A$.

Let $F$ be a face of $H$ of cardinality $d$.  Form the normal family ${\cal L}$ of example~\ref{ex:H2}.  The quotient $H/{\cal L}$ is isomorphic to $\op{Dih}_{2d}$, hence also isomorphic to $H'$.  The isomorphism can be chosen so that
$\check\varphi$ pulls back to $\check\varphi_{can}$ on the faces of $H/{\cal L}$, and so that $x'$ is the image of a dart $x$ in the canonically false 
face of $H/{\cal L}$.  

Let $m = \card(S(H,L,x))$.  Then $(H,{\cal L},L,x)$ is a marked hypermap that
is correlated with $(H',m,\check\varphi,x')\in A$.
\end{proof}

\begin{lemma}[correlated edge]\guid{FDQZOSJ}\label{lemma:correlated-edge}
  Let $H$ be a restricted $(\card(\Omega),d)$-hypermap with digraph $(V,E)$. 
  Assume that ${\cal H}\in V$ is not a
  sink and is correlated with an element in $\op{fix}_2$.
  Then there exists a directed edge  $({\cal H},{\cal H}')\in E$, such that
 $\cal H'$ is correlated with an element of  $\op{fix}_2$.
\end{lemma}

\begin{proof}  Assume that the marked hypermap ${\cal H}=(H,{\cal L},L,x)$
is correlated with the tuple $(H',m,\check\varphi,x')\in \op{fix}_2$,
by means of an isomorphism
\begin{displaymath}
G: H/{\cal L} \to H'.
\end{displaymath}
Let $T{\cal H} = (H,{\cal M},L_1,x)$ be the transform.  Let $m,p,q,x,y,z$
be the parameters of Definition~\ref{def:yz}.  Let $r$ be the cardinality
of $\F(L)$, which is equal to the cardinality of the face of $x'$ in $H'$.  
Then $(m,p,q,r)\in C$ by Lemma~\ref{lemma:yz}.

Let $H'' = RDW(H',x',m,p,q)$.  The isomorphism $G$ and Lemma~\ref{lemma:RDW} combine to give an isomorphism $G':H/{\cal M} \mapsto H''$.
Push the canonical function on the faces of $H/{\cal M}$ to a function
$\check\psi$ on the faces of $H''$.

If $L_1$ is false, then $({\cal H},T{\cal H})$ is a directed edge of the digraph, and
$T{\cal H}$ is correlated with $(H'',m',\check\psi,x')\in g(H',m,\check\varphi,x')\subset \op{fix}_2$, where $m'=\card(S(L_1,x))$.  

If $L_1$ is true, then $H''$ has a false face, and the choice function $\op{ch}$ picks a dart $x''$ in the union of the false faces of $H''$.  Transport this by $G$ to a dart $y'\in H/{\cal M}$ in a false face.  Write $y'=[\ldots;y]$
with $y$ a dart visited by a false contour loop $M$ of ${\cal M}$.  Then
$(H,{\cal L},L,x),(H,{\cal M},M,y))$ is a directed edge of the digraph,
whose tail is correlated with $(H'',m',\check\psi,x'')\in g(H',m,\check\varphi,x')\subset \op{fix}_2$, where $m'=\card(S(M,y))$.
\end{proof}

\begin{lemma}[correlated sink]\guid{RIZGJVS}\label{lemma:correlated-sink}
Let $H$ be a restricted $(\card(\Omega),d)$-hypermap with digraph $(V,E)$.
Then some sink in the digraph is correlated with some element of
$\op{fix}_2$.
\end{lemma}

\begin{proof} Start with any  correlated pair $({\cal
    H}_0,{\cal K}_0)$ with ${\cal H}_0\in V$ and ${\cal K}_0\in \op{fix}_2$ 
  (Lemma~\ref{lemma:correlated-seed}).  Use
  Lemma~\ref{lemma:correlated-edge}, to produce a sequence $({\cal
    H}_i,{\cal K}_i)$ of correlated pairs, where $[{\cal H}_0;{\cal
    H}_1;\ldots]$ is a path in the vertex set $V$.  By Lemma~\ref{lemma:digraph-sink},
  the path reaches a sink within $\#D$ steps.  The final marked
  hypermap ${\cal H}_k$ in the path is a sink that is correlated with
  ${\cal K}_k\in X_2$.
\end{proof}

\begin{proof} Turn to the proof of Theorem~\ref{lemma:algorithm}.  Let
  $(H,{\cal L},L,x)$ be a sink that is correlated with some tuple
  $(H',m,\check\varphi',y')\in \op{fix}_2$
  (Lemma~\ref{lemma:correlated-sink}).  By
  Lemma~\ref{lemma:digraph-sink}, $H$ is isomorphic to the quotient
  $H/{\cal M}$, where $(H,{\cal M},\ldots)$ is the transform of the
  sink.  This quotient is isomorphic to $\op{RDW}(H/{\cal
    L},x',m,p,q)$, where $x'$, $m$, $p$, $q$ are given by
  Definition~\ref{def:yz}.  By the correlatedness of the sink, $H$ is
  isomorphic to $H''=\op{RDW}(H',x',m,p,q)$.  By
  Lemma~\ref{lemma:parameters}, $(m,p,q,r)\in C$, where
  $r=\op{card}(\F(L))$.  Under the isomorphism, $r$ is the cardinality
  of the face of $y$ in $H'$.  By the definition of $f$ and $g$,
  $H''\in \op{fix}_1$.
\end{proof}





%
%
%\begin{proof} Let $H$ be a restricted hypermap whose dart set belongs
%  to $D$ and whose greatest face size is $d$.  It has a quotient
%  $H/{\cal L}_0$ isomorphic to $\op{Dih}_{2d}$.  Let ${\cal L}_i = (H,{\cal
%    L}_i,L_i,x_i,\check\varphi_i)$ be the sequence of marked hypermaps
%  constructed in Lemma~\ref{lemma:sequence}.  XX.
%\end{proof}
%
%\begin{proof} Let $H$ be a restricted hypermap whose dart set belongs
%to $D$ and whose greatest face size is $d$.  It has a quotient
%$H/{\cal L}_0$ isomorphic to $\op{Dih}_{2d}$.  By repeating the
%construction of Section~\ref{sec:face-insert}, one obtains a
%sequence of hypermaps $\op{Dih}_i = H/{\cal L}_i$, $i=0,\ldots,N$,
%terminating with a hypermap $H/{\cal L}_N$, which is isomorphic to
%$H$.  The data $m_i,\check\varphi_i,x_i$ is also obtained for each $\op{Dih}_i$.
%
%The tuple $(\op{Dih}_0,m_0,\check\varphi_0,x_0)$ is isomorphic to an element of
%$A$.  By construction, $A\subset \op{fix}(f,X)$.  If the lemma is
%false, there is a smallest $i>0$ for which
%$(H_i,m_i,\check\varphi_i,x_i)\not\in \op{fix}(f,X)$ (or if $i=N$, for which
%$H\not\in \op{fix}(f,X)$).  There are isomorphisms
%$RDW(H_i,x_i,m_i,p_i,q_i) \simeq H_{i+1}$ for appropriate choices of
%$p_i,q_i$.  By construction, when the data belongs to $\op{fix}(f,X)$
%for $i-1$, the data belongs to $\op{fix}(f,X)$ for $i$.  Thus, $H$
%belongs to $\op{fix}(f,X)\cap X_1$.
%\end{proof}
%
%\subsection{old algorithm}
%
%If a hypermap is restricted and $x$ is any dart, then $x$ and $n x$
% lie on different faces.  In particular, a restricted hypermap has at
% least two faces.  To begin the process, take ${\cal L}$ to be the
% normal family of Example~\ref{ex:H2} with two contour loops, whose
% quotient hypermap is a polygon $\op{Dih}_{2d}$.  When the the initial
% contour loop is chosen on a face of maximal size, the natural number
% $d$ is an upper bound on size of a face.
%
% In summary, a process starts with a single polygon and then adds
% edges and nodes of degree two along the inserted edges, to obtain a
% restricted hypermap $H$.
%
% A modification of the process avoids explicit reference to the
% hypermap $H$ and to the normal family ${\cal L}$.  Let $D$ be a
% finite set that contains all the darts for all of the restricted
% hypermaps to be constructed.  For each $d=3,\ldots,\# D$, the
% process generates all restricted hypermaps with greatest face-size
% $d$, with darts in $D$.
%
% The algorithms performs the following initialization.  The initial
% hypermap is the polygonal hypermap $\op{Dih}_{2d}$.  A flag $\varphi$ marks
% one face true and the other false.  A distinguished dart $x$ is
% selected on the false face.  For each $m<d$, set $S=S_m = \{f^i
% x\mid 1\le i\le m\}$.
%
% Each iteration processes a finite list ${\cal H}$ of quadruples
% $(H,m,\varphi,x)$, where $H$ is a simple hypermap, $\varphi$ is an
% $S$-flag, $x$ lies in a false face, and $S = \{f^i x\mid 1\le i\le
% m\}$ for some $m$.  The algorithm terminates when every face of
% every hypermap in ${\cal H}$ is true.  At every step of the
% algorithm, one quadruple with some false face is removed from ${\cal
%   H}$ and finitely many quadruples are returned to ${\cal H}$.
%
% At each iteration the chosen $(H,m,\varphi,x)$ is modified in the
% following ways and each modification is placed back in the list
% ${\cal H}$.  As the natural number $m$ depends on the unknown normal
% family ${\cal L}$, all possible $m < d$ are used.  Similarly, the
% natural number $p < d$ depends on ${\cal L}$, and all possible $p$
% are used.  Any quadruple that is isomorphic to one previously
% considered is discarded as redundant.
%
% In summary, the algorithm constructs of hypermaps.  The process must
% terminate, because the set $D$ is finite, so there are only finitely
% many quadruples (or quadruples up to isomorphism) that construct
% their dart sets from $D$.
%
%\begin{lemma}[]\guid{BRGEFNH}\rating{2000}  Fix $d\ge 3$ and $D\ne \emptyset$.
%  This process constructs in a finite number of steps all restricted
%  hypermaps, up to isomorphism, such that the dart set belongs to the
%  finite set $D$, and whose greatest face size is $d$.
%\end{lemma}




%%%%%%%%%%%%%%%%%


    
    \chapter{Hypermap}
    Use the electrical dual charge proof%
    \footnote{\url{http://www.ics.uci.edu/~eppstein/junkyard/euler/all.html}} of the Euler formula for convex polyhedra (a combinatorial Morse function).  This does not seem to need an explicit invocation of the Jordan curve theorem, since the faces of the polyhedron are convex.
    
    
      
    The local index of a vertex is $2$ - number of min/max edges + number of min/max faces. The local index at the absolute max and min = 2 each.  The local index at any other vertex is $0$ (two of the local face elements are not a max or min).  Sum of local indices $=4$.
    
    The sum of the vertex contributions is $2V$, edge contributions $-2E$, and face contributions $2F$ (by convexity each face has a unique max and min). 
	Thus, the graph is planar: $2 = V-E+F$.

	Without loss of generality, the spherical net of edges formed by quasi-regular tetrahedra is 2-connected.  
    
    Conjecture: We can show that the quasi-regular tetrahedra unit sphere arcs belong to the set of edges of the convex hull of the projection of the points to the unit sphere.  
    
    Removing an edge that does not disconnect the graph does not change the Euler characteristic, so we can work our way from the planar graph formed by the convex hull to the planar graph formed by the quasi-regular tetrahedra.


%%%%%%%%%%%%%%%%%%%%%%%%%%%%%%%%%%%%%%%%%%%%%%%%%%%%
    % file started March 22, 2009
% Marchal objective function

\def\lam{\lambda}
\def\Lam{\Lambda}
\def\bl{{\underline{\lam}}}
\def\bm{{\underline{\mu}}}
\def\norm#1{{\|#1\|}}
\def\angle#1#2#3{{\op{angle}(#1,{#2,#3})}}

\chapter{Formulation}

\section{Statement}



\begin{theorem}[Sphere Packing Problem (Kepler Conjecture)]
\label{theorem:kepler}   No packing of congruent balls in
Euclidean three space has density greater than that of the
face-centered cubic packing.
\end{theorem}

\begin{remark}
This density is $\pi/\sqrt{18}\approx 0.74.$  There are other
packings, such as the hexagonal close packing, that attain this
same density.
\end{remark}

The proof of this result is presented in this book. Here, we
describe the top-level outline of the proof and give references to
the sources of the details of the proof.

By a {\it packing}, we mean an arrangement of congruent balls that
are nonoverlapping in the sense that the interiors of the balls are
pairwise disjoint. Consider a \indy{Index}{packing} packing of congruent
balls in Euclidean three space. There is no harm in assuming that
all the balls have unit radius. 

%The density of a packing does not
%decrease when balls are added to the packing. Thus, to answer a
%question about the greatest possible density we may add
%nonoverlapping balls until there is no room to add further balls.
%Such a packing will be said to be {\it saturated}.
%
% \indy{Index}{saturated}
% \indy{Index}{overlap}

Let $\Lambda$ be the set of centers of the balls in a
%\indy{Index}{saturated} saturated 
packing. Our choice of radius for the
balls implies that any two points in $\Lambda$ have distance at
least $2$ from each other. We call the points of $\Lambda$ {\it
\indy{Index}{vertex} vertices}.
%A bijection
%$\varph:\ring{N}\to\Lambda$ is called an {\it enumerated\/}
%packing.

%\begin{lemma}\guid{DQLSGKR} A saturated packing $\Lambda$ is countably infinite.  Hence,
%there exists a bijection $\varphi:\ring{N}\to\Lambda$.
%\end{lemma}

%\begin{proof}  If $\Lambda$ is finite, there is room to add a ball
%far from the finite cluster $\Lambda$, so $\Lambda$ is not
%saturated.

%The map $\Lambda\to\ring{Z}^3$, $$(x,y,z)\mapsto (\lfloor 2x
%\rfloor, \lfloor 2y \rfloor, \lfloor 2z \rfloor)$$ is a one-one
%map from $\Lambda$ into a countable set.  Hence $\Lambda$ is
%countable.
%\end{proof}


  Let $B(x,r)$ denote the open ball in
Euclidean three space at center $x$ and radius $r$.  The open ball
is measurable, with measure $4\pi r^3/3$.



Let $\delta(\Lambda,x,r)$ be the finite density, defined as the
ratio of the volume of $B(\Lambda,x,r)$ to the volume of $B(x,r)$,
where $B(\Lambda,x,r)$ is defined as the intersection with
$B(x,r)$ of the union of all balls in the packing. Set
$\Lambda(x,r) = \Lambda \cap
B(x,r)$ and $\Lambda^*(x,r) = \Lambda(x,r)\setminus \{x\}$.
\indy{Greek}{ZZdelta@$\delta(\Lambda,x,r)$}
\indy{Greek}{ZZLamda@$\Lambda(x,r)$}
\indy{Greek}{ZZLamda@$\Lambda^*(x,r)$}

\begin{lemma}\guid{KIUMVTC}\label{lemma:Lambda-finite}
Let $\Lambda$ be a packing and let $p\in\ring{R}^3$.
Then the set $\Lambda(p,r)$ is finite.
\end{lemma}

\begin{proof}  Let $p = (p_x,p_y,p_z)$. The map
$$v=(v_x,v_y,v_z)\mapsto (\lfloor 2(v_x-p_x)
\rfloor, \lfloor 2(v_y-p_y) \rfloor, \lfloor 2(v_z-p_z) \rfloor)$$
is a one-one map from $\Lambda$ into the set $\ring{Z}^3\cap B(0,2
r
 + 1)$.  By Lemma~\ref{lemma:Zcount}, this gives a one-to-one map
 into a finite set, hence $\Lambda(p,r)$ is finite.
\end{proof}


\begin{definition}[Voronoi~cell,~$\Omega$]\label{def:voronoi}\indy{Index}{Voronoi cell} Voronoi cell
$\Omega(v)=\Omega(\Lambda,v)$\indy{Greek}{ZZzomega@$\Omega(v)$} around a
vertex $v\in \Lambda$ is the set of points closer to $v$ than to
any other ball center.  Let $\Omega_t(\Lambda,v) = \Omega(\Lambda,v)
\cap B(v,t)$ be the truncated Voronoi cell at radius $t$.
\end{definition}

\begin{lemma}\guid{DRUQUFE} Let $\Lambda$ be a saturated packing and let $t>0$.
Each Voronoi cell $\Omega_t(\Lambda,v)$, for $v\in\Lambda$, is
convex, open, bounded, and measurable.
\end{lemma}

\begin{proof}  This is elementary.
\end{proof}

\begin{definition}[negligible]\label{def:negligible}
Let $A:\Lambda\to\R$ be a function.  We say that $A$\indy{Index}{Az@$A$}
is
  {\it negligible\/}\indy{Index}{negligible}
if there is a constant $C_1$ such that for all $r\ge1$ and all
$x\in\ring{R}^3$,
   $$\sum_{v\in\Lambda(x,r)} A(v) \le C_1 r^2.$$
We say that the function $A:\Lambda\to\R$ is
  {\it fcc-compatible\/}\indy{Index}{fcc-compatible}
if for all $v\in\Lambda$ we have the inequality
$$\sqrt{32}\le \op{vol}(\Omega_2(v)) + A(v).$$
\end{definition}





\begin{remark}
The value $\op{vol}(\Omega_2(v)) + A(v)$ may be interpreted as a
{\it corrected\/} volume\indy{Index}{corrected volume} of the Voronoi
cell. Fcc-compatibility asserts that the corrected volume of the
Voronoi cell is always at least the volume of the Voronoi cells in
the face-centered cubic and hexagonal-close packings.
\end{remark}

\begin{remark} In \cite{Hales:2006:DCG}, the full Voronoi cell $\Omega(v)$
is used, rather than $\Omega_2(v)$.  The truncation at radius $2$
is just a matter of convenience to guarantee the boundedness
and hence the finite volume of the (truncated) Voronoi cell.
In \cite{Hales:2006:DCG}, the same effect was achieved by requiring all packings
to be saturated.  We have dropped the assumption of saturation
on $\Lambda$.
\end{remark}


\begin{lemma}\guid{JGXZYGW}
\label{lemma:deltabound} If there exists a \indy{Index}{negligible}
negligible \indy{Index}{fcc-compatible} fcc-compatible function
$A:\Lambda\to\R$ for a 
%saturated 
packing $\Lambda$, then there
exists a constant $C$ such that for all $r\ge1$ and all
$x\in\ring{R}^3$,
    $$
    \delta(\Lambda,x,r)
    \le \pi/\sqrt{18} + C/r.
    $$
The constant $C$ depends on $\Lambda$ only through the constant
$C_1$ of Definition~\ref{def:negligible}.
\end{lemma}



\begin{proof}
The numerator $\op{vol}\, B(\Lambda,x,r)$ of $\delta(\Lambda,x,r)$
is at most the product of the volume of a ball $4\pi/3$ with the
number $|\Lambda(x,r+1)|$ of balls intersecting $B(x,r)$.  Hence
    \begin{equation}
    \op{vol}\, B(\Lambda,x,r) \le |\Lambda(x,r+1)| 4\pi/3.
    \label{eqn:Abound}
    \end{equation}

%In a %saturated packing 
Each truncated Voronoi cell is contained in a ball of
radius $2$ centered at the {\it center} of the cell.  The volume
of the ball $B(x,r+3)$ is at least the combined volume of 
truncated Voronoi
cells whose center lies in the ball $B(x,r+1)$. This observation,
combined with fcc-compatibility and negligibility, gives
    \begin{equation}
    \begin{split}
    \sqrt{32}|\Lambda(x,r+1)|
    &\le \sum_{v\in\Lambda(x,r+1)} (A(v) +
    \op{vol}(\Omega_2(v))) \\
    &\le C_1 (r+1)^2 + \op{vol}\,B(x,r+3) \\
    &\le C_1 (r+1)^2 + (1+3/r)^3 \op{vol}\,B(x,r)
    \label{eqn:Bbound}
    \end{split}.
    \end{equation}
Recall that $\delta(\Lambda,x,r)=
\op{vol}\,B(\Lambda,x,r)/\op{vol}\,B(x,r)$. Divide Inequality
\ref{eqn:Abound} through by $\op{vol}\,B(x,r)$.  Use
Inequality~\ref{eqn:Bbound} to eliminate $|\Lambda(x,r+1)|$ from the
resulting inequality.  This gives
    $$\delta(\Lambda,x,r)
        \le \frac{\pi}{\sqrt{18}} (1+3/r)^3 + C_1 \frac{(r+1)^2}{r^3\sqrt{32}}.
    $$
The result follows for an appropriately chosen constant $C$.
\end{proof}

\begin{remark} \label{remark:precise}
We take the precise meaning of the packing problem to be the
bound $\delta(\Lambda,x,r) \le \pi/\sqrt{18} + C/r$ of the lemma.
Thus, the solution to the packing problem follows, provided a negligible
fcc-compatible function can be found. The strategy will be to
define a negligible function, and then to solve an optimization
problem in finitely many variables to establish that it is
fcc-compatible.
\end{remark}

If $\Lambda$ is a %saturated 
packing, we have an ordered pair
$(\Lambda,v)$ for each vertex $v\in\Lambda$.  The pairs
$(\Lambda,v)$ are called {\it centered packings}.
Let $\op{CentPack}$ be the set of all centered packings.
\indy{Index}{centered packing}\indy{Index}{CentPack}


\section{Rogers's partition}

\begin{definition}[saturated packing]
A packing $\Lam\subset \ring{R}^3$ is a set such that
$$\forall \lam~\mu\in \Lam.~  \norm{\lam-\mu} < 2 \Rightarrow (\lam=\mu).$$
The packing is saturated if for every $v\in\ring{R}^3$,   there exists $\lam\in\Lam$
such that $\norm{\lam-v}\le 2$.
\end{definition}

We think of $\Lam$ as the set of centers of a  packing of congruent balls of radius $1$.
To be saturated means that there is no room for further balls to be added to the packing.
There is no loss in generality in assuming that the packing is saturated, when searching
for the greatest possible density of a packing.

Given a packing $\Lam$, Rogers gives a partition of Euclidean space into
simplices with vertices at $\Lam$ \cite{Rogers:1958:Packing}.   This section describes his partition.

Let $\Omega(\lam)$  be the closed Voronoi cell centered at $\lam$:
$$
  \Omega(\lam) = \{x \mid  \norm{x-\lam} \le \norm{x-\mu},\ \lam\ne\mu\in\Lam\}.
$$
If $S\subset\Lam$, let $\Omega(S)$ be the intersection of the family:
$$\Omega(S)  = \bigcap \{\Omega(\lam)\mid \lam\in S \}.$$
We call $\Omega(S)$ the $S$-face of $\Omega(\lam)$.

For $X\subset\ring{R}^3$, we let $\dim(X)$ be the dimension of the affine hull
of $X$: $\dim(X) = \dim(\op{aff}(X))$.

If $\bl=(\lam_0,\ldots,\lam_k)$ and $j\le k$ write $\bl[j] = (\lam_0,\ldots,\lam_j)$.
Let $\Lam(k)$ be the set of $k+1$-tuples $\bl=(\lam_0,\ldots,\lam_k)$ such
that 
\begin{equation}\label{eqn:omega-dim}
\dim(\Omega(\bl[j]))+1 = \dim\Omega(\bl[j-1])
\end{equation}
for all $0<j\le k$.
In particular, $\Lam(0)=\Lam$.  Also, $\Lam(1)$ is the
set of all pairs $(\lam,\mu)$ such that the Voronoi cells at $\lam$ and $\mu$ meet along
a face of codimension $1$, and
so forth.


Define $v(\bl)$ recursively on $\Lam(k)$ by
$$v(\lam[0]) = \lam[0]$$
and $v(\lam[k+1])$ is the closest point on $\Omega\{\lam[k+1]\}$ to $v(\lam[k])$.  The point is defined whenever $\Omega\{\lam[k+1]\}$ is nonempty.
The set $\Omega\{\lam[k+1]\}$ is convex and compact.  Thus, the point $v(\lam[k+1])$ exists
uniquely.

For $\bl\in\Lam(k)$, let 
$$R(j,\bl) = \op{conv}\{v(\lam[j]), v(\lam[j+1]),\ldots,v(\lam[k])\}.$$  Set $R(\bl)=R(0,\bl)$.


\begin{lemma} 
For any saturated packing $\Lambda$, we have
$$\ring{R}^3 = \bigcup \{ R(\bl) \mid \bl\in \Lam(3)\}.$$
\end{lemma}

\begin{proof}
We have $$\ring{R}^3 = \bigcup \{\Omega(\lam)\mid \lam\in \Lam(0)\}.$$
So it is enough to show that
$$\Omega(\bl[0]) = \bigcup \{ R(\bl') \mid \bl'\in \Lam(3),~\bl[0]=\bl'[0]\}.$$
We have
$$\Omega(\bl[0]) = \bigcup \op{conv}(c(\bl[0]),\{\Omega(\bl[1])\mid (\bl[1])\in \Lam(1)\}).$$
So it is enough to show that
$$\Omega(\bl[1]) = \bigcup \{ R(1,\bl') \mid \bl'[1]=\bl[1],~\bl'\in \Lam(3)\}.$$
Repeating, it is enough to show that
$$\Omega(\bl[3]) = \bigcup\{R(3,\bl')\mid\bl'[3]=\bl[3],~\bl'\in\Lambda(3)\}.$$
The right-hand side is the singleton $\{v(\bl[3])\}$.  The left-hand side
contains this point and is contained in a properly decreasing chain of affine spaces:
$\ring{R}^3$, the bisector of $\lam_0$ and $\lam_1$, and so forth.  This determines a point,
so the two sides are equal.
\end{proof}

\begin{lemma}  The intersection of any two distinct $R(\bl)$ and $R(\bm)$ is
contained in a plane (hence has measure zero).
\end{lemma}

Combined with the previous lemma, we see that the simplices $R(\bl)$ partition Euclidean
space.

\begin{proof}  Let $\bl = (\lambda_0,\ldots)$ and $\mu = (\mu_0,\ldots)$.
If $\lambda_0\ne\mu_0$, the intersection of the simplices lies in $\Omega\{\lam_0,\mu_0\}$
and we are done.  Let $k+1$ the first index such that
$$\op{conv}(v(\lam_0),\ldots,v(\lam_0,\ldots,\lam_{k+1}))\ne
\op{conv}(v(\mu_0),\ldots,v(\mu_0,\ldots,\mu_{k+1})).
$$
By induction $\Omega\{\lam_0,\ldots,\lam_k\} = \Omega\{\mu_0,\ldots,\mu_k\}$
but
$$\Omega\{\lam_0,\ldots,\lam_{k+1}\} \ne \Omega\{\mu_0,\ldots,\mu_{k+1}\}.$$
The intersection lies in the convex hull $C$ of
$\{v(\lam_0),\ldots,v(\lam_0,\ldots,\lam_k)\}$ and
$V=\Omega\{\lam_0,\ldots,\lam_{k+1},\mu_{k+1}\}$.  The set $V$ lies in an affine space of 
codimension $k+2$
and $C$ has codimension $(k+2) - (k+1) = 1$.
\end{proof}

\begin{lemma}\label{lemma:v2} Let $\bl=(\lam_0,\ldots,\lam_3)\in\Lambda(3)$ and fix $0\le j\le 3$.  Assume that
$$
\norm{v(\bl[j]) - v(\bl[0])} < \sqrt2.
$$
Then $v(\bl[j])$ is the circumcenter of the points $S_j =\{\lam_0,\ldots,\lam_j\}$ and
lies in the convex hull of $S_j$.  The points $v(\bl[j])$, for $j=0,1,2,3$, are distinct.
\end{lemma}

\begin{proof} By definition, an element $\bl\in\Lambda(3)$ gives dimensions: 
$$\Omega\{\lam_0,\ldots,\lam_j\} = 3-j.$$  The case $j=0$ of the lemma is trivially
satisfied.  Assume by induction the result holds for numbers less than $j$.

Now $j>0$.
The circumcenter $p=p_{j}$ of the set $S_{j}$ is the point on the plane $\op{aff}(S_{j})$
closest to $p_{j-1}=v(\bl[j-1])$.  We claim that $p$ lies in $\Omega(S_{j})$.  Otherwise
there is a point $\mu\in\Lambda$ closer to $p$ than any $\lambda\in S_{j}$.  
The angles $\angle{p}{\mu}{\lam_i}$ are obtuse for $i\le j$.  We make a case-by-case
argument for each $j$.

If $j=1$, the points $p,\lam_0,\lam_1$ are collinear and cannot give two obtuse angles.

If $j=2$, $\mu$ lies in the intersection of half spaces with bounding plane $\mu$
and normal $-\lambda_i$.  Let $\mu'$ be the projection of $\mu$ to the plane containing
$p,\lam_0,\lam_1,\lam_2$. The four points $\mu',\lam_0,\lam_1,\lam_2$ can be arranged
cyclically around $p$, each forming an obtuse angle with the next.  A circle around $p$
cannot give four obtuse angles.


If $j=3$, assume that $\lam_0,\ldots,\lam_4$ are labeled in cyclic order around the line
$\op{aff}\{p,\mu\}$.  Consider the dihedral angle 
  $$
  \gamma_i=\gamma=\dih(\{p,\mu\},\{\lam_i,\lam_{i+1}\}).
  $$
By the law of cosines the angle $\gamma$ is given in terms of the edges of the spherical
triangle $a,b,c$ by
$$
  \cos c - \cos a \cos b = \sin a \sin b \cos \gamma.
$$
All the terms on the left-hand side are negative, so $\gamma =\gamma_i > \pi/2$.
This is impossible, as the sum of the four dihedral angles $\gamma_i$ is $2\pi$.
This completes the proof that $v(\bl[j])$ is the circumcenter.

Next we show that this point lies in the convex hull of $S_j$.
If $j=0$, there is nothing to show.  If $j=1$, the point is the midpoint of the convex hull.
If $j=2$, the point is the circumcenter of an acute triangle.  If $j=3$, the point is the
circumcenter of a simplex such that every face has positive orientation (see XX).  Thus,
in every case the point lies in the convex hull.

The distances $\norm{v(\bl[j])-v(\bl[0])}$ are increasing, so it is enough to show
that $v(\bl[j])\ne v(\bl[j+1])$.  If we had equality, then the circumcenter $v(\bl[j])$
would have an equally close packing point $\lambda_{j+1}\in\Lambda$, which is impossible
the the first part of the proof.
\end{proof}

\begin{lemma}   Let $\bl\in\Lambda(k)$.  Assume that $\norm{v(\bl)-v(\bl[0])}<\sqrt2$.
Let $\bm$ be any permutation of the components of $\bl$.  Then $\bm\in\Lambda(k)$ and
 $v(\bl) = v(\bm)$.
\end{lemma}

\begin{proof} 
Since the sets $\Omega(v(\bl[j]))$ satisfy (\ref{eqn:omega-dim}), we have that
$\Omega(v(\bl))$ is a point consisting of the singleton $\{v(\bl)\}$, which is the
circumcenter of the simplex with vertices $\{\lam_0,\ldots,\lam_k\}$.  This describes
the point $v(\bl)$ in a way that is independent of the permutation.

We check that the condition (\ref{eqn:omega-dim}) holds for the permutation $\mu$.
The proof of Lemma~\ref{lemma:v2} shows that the midpoint of $\op{conv}(\mu_0,\mu_1)$
lies in $\Omega(\bm[1])$ and since the inequalities are strict, some neighborhood
of this midpoint in the bisecting plane of $\{mu_0,\mu_1\}$ lies in $\Omega(\bm[1])$.
Thus, $\dim(\Omega(\bm[1]))=2$.  We continue in this fashion, showing that each dimension
$\dim(\Omega(\bm[j]))$ drops by exactly one.
\end{proof}

\begin{lemma}\label{lemma:Rconv} Let $\bl = (\lam_0,\ldots,\lam_k)\in\Lambda(k)$.  Assume that $\norm{v(\bl)-v(\bl[0])}<\sqrt2$.
For any permutation $\pi$, let $\pi(\bl)\in\Lam(k)$ be the permuted entry.  Let
$S(k)$ be the group of all permutations on $k+1$ letters.   Then
$$
\op{conv}\{\lam_0,\ldots,\lam_k\} = \bigcup \{ R(\pi(\bl)) : \pi\in S(k)\}.
$$
\end{lemma}

\begin{proof} Let $L = \{\lam_0,\ldots,\lam_k\}$.  We prove it by induction on $k$.
When $k=0$, the result is trivial.  Now assume $k>0$.

The circumcenter $v(\bl)$ of $L$ lies in the convex hull of these
point (see the proof of Lemma~\ref{lemma:v2}).  Thus, the left-hand side is the union
of the cones over the (k+1)-faces:
$$
\op{conv}(S) = \bigcup\{ \op{conv}(v(\bl),L\setminus \{\lam_i\}\mid i=0,\ldots,k) \}.
$$
The sets $L\setminus \{\lam_i\}$ can be identified with  cosets of $S(k)/S(k-1)$.
By induction $L\setminus \{\lam_i\}$ is the union of $R(\bm)$ as $\bm$ runs
over all permutations $S(k-1)$ of $L\setminus \{\lam_i\}$.
The result follows by induction.
\end{proof}

We remark that the Rogers's simplices $R(\bl)$ are compatible with the Voronoi
decomposition of space (by construction).  Also, by Lemma~\ref{lemma:Rconv}, under
mild restrictions on the circumradius, they can also be reassembled into simplices
with vertices at the centers of the packing (the Delaunay simplices).  

\section{Marchal's partition}

C, Marchal has introduced a new approach to the Keper conjecture that is manifestly superior to the one in \cite{Hales:2006:DCG}.  If his approach can be pushed
to completion, it would cut years off the Flyspeck project.

His articles claim to give
a {\it demonstration} of the Kepler conjecture \cite{marchal:2007}, \cite{marchal:2008}.  However, the
mathematically rigorous part of the article only gives a reduction
of the problem to an optimization problem in a finite number of
variables.  The method of gradient descent is then used to explore
the local minima of the optimization problem in finitely many variables.

Marchal's partition of space is a variant of Rogers's partition into simplices
$R(\bl)$.  The main part of construction is 
the decomposition obtained by truncating the Voronoi cells
by a ball of radius $\sqrt2$.  In a few carefully chosen situations he assembles the simplices
into a larger convex hull along the lines of Lemma~\ref{lemma:Rconv}.

Let $r(\bl[i]) = \norm{v(\bl[i])-\lam_0}$.  The numbers $r(\bl[i])$ are increasing with $i$.

\begin{definition} Let $R=R(\bl)$ be a Rogers simplex with $\bl=(\lam_0,\ldots,\lam_3)$.
\hfill\break\smallskip  
{\bf The $0$-cell} of $R$ is
$$
R(\bl)_0 = \{x\in R(\bl) \mid \norm{v(\lam_0)-x} > \sqrt2.
$$
\bigskip
{\bf The $1$-cell} of $R$ is 
$$
R(\bl)_1 = \op{conv}(v(\lam_0),T_1),\quad T_1 = \{x \in R(\bl) \mid \norm{v(\lam_0)-x}= \sqrt2\}.
$$
\bigskip
{\bf The $2$-cell} of $R$ is
$$
\begin{array}{rll}
R(\bl)_2 &= \op{conv}(v(\bl[0]),v(\bl[1]),T_2),\quad \\
  T_2 &= \{x \in R(\bl)\cap \Omega(\bl[1]) \mid \norm{\lam_0-x}=\norm{\lam_1-x} =\sqrt2\}.
\end{array}
$$
\bigskip
{\bf The $3$-cell} of $R$ is defined to be empty unless 
$$
r(\bl[2]) <\sqrt2 \le r(\bl[3]).
$$
When this inequality holds,
there is a a unique point $p(\bl[2])$ in
$\op{conv}(\bl[2],\bl[3])$ at distance exactly $\sqrt2$ from $\lam_0$.  
Define the $3$ cell to be
$$
R(\bl)_3 = \op{conv}\{p(\bl[2]),\lam_0,\lam_1,\lam_2\}.
$$
\bigskip
{\bf The $4$-cell} of $R$ is defined to be empty unless
$$
r(\bl[3]) <\sqrt2.
$$
When this inequality holds, define the $4$ cell to be
$$
R(\bl)_4 = \op{conv}\{\lam_0,\lam_1,\lam_2,\lam_3\}.
$$
\end{definition}

Note the the $0$, $1$, and $2$-cells are always subsets of $R$.  However, the $3$ and
$4$-cells are finite unions of Rogers simplices.

\begin{lemma}  The $0,\ldots,4$-cells of all Rogers simplices $R(\bl)$, for $\bl\in \Lam(3)$
give a partition of $\ring{R}^3$.  An $i$-cell is never equal to a $j$-cell when $i\ne j$.
Two $3$-cells $R(\bl)_3 = R(\bm)_3$ are equal exactly when
$(\lam_0,\lam_1,\lam_2)$ is a permutation of $(\mu_0,\mu_1,\mu_2)$.  Two $4$-cells are
equal exactly when $\bl$ is a permutation of $\bm$.
\end{lemma}

\begin{proof}  First we will prove that cells are either disjoint or equal.  At the same
time, we determine when $4$-cells or $3$-cells are equal to one another.
By Lemma~XX, the $4$-cell $R(\bl)_4$ is a union of the Rogers simplices
$R(\bm)$, as $\bm$ runs over permutations of $\bl$. 
 Two $4$-cells
that meet are equal. 
Every $4$-cell comes from a simplex
$R(\bl)$ with $r(\bl)<\sqrt2$.  This condition gives $R(\bl)_i=\emptyset$, for $i=0,1,2,3$.

Similarly, the $3$-cell is a union
of the convex hulls $\op{conv}(p(\bm[2]),R(\bm[2]))$ as $\bm[2]$ runs over permutations of $\bl[2]$.  Note that the point $p(\bm[2])=p(\bl[2])$ is independent of the permutation, since
it is determined as the point at distance $\sqrt2$ from $\lam_0$ along the line through
from $v(\bm[2])=v(\bl[2])$ perpendicular to the plane $\op{aff}\{\lam_0,\lam_1,\lam_2\}$
(in the half-space of $\bl[3]$). Two $3$-cells that meet are equal and come from parameters that are permutations of
one another as described.   The intersection $R(\bl)_3$ cannot meet $R(\bl)_i$, for $i<3$, in
a set of full dimension, because the plane $P(\bl[2])=\op{aff}\{\lam_0,\lam_1,p\}$ separates them.

The $2$-cell $R(\bl)_2$ is separated from the $0$ and $1$-cells $R(\bl)_i$
by the cone $P(\bl[1])$ with apex $\lam_0$
passing through $T_2$. The $1$-cell is separated from $0$-cells the the sphere $P(\bl[0])$
of radius
$\sqrt2$, centered at $\lam_0$.

Finally, we show that every point $x$ in $\ring{R}^3$ belongs to a cell.  By the Rogers
partition, the point $x$ belongs to a Rogers simplex $R(\bl)$.  If $r(\bl)<\sqrt2$, then
$x$ belongs to a $4$-cell.  Otherwise $x\in R(\bl)$ belongs to $R(\bl)_i$ according to
which side of the plane $P(\bl[2])$, cone $P(\bl[1])$, and sphere $P(\bl[0])$ the point
lies, as described above.
\end{proof}


\section{Kissing number estimates}

This section shows how Conjecture~\ref{conj:m1} would follow from an explicit estimate
related to the Gregory-Newton problem


We define the following constants and functions.
$$
\begin{array}{lll}
\alpha=\alpha_{cm} &= \arccos(1/3)\\
K=K_{cm} &= (3\alpha-\pi)\sqrt2/(12\pi - 30\alpha)\\
M = M_{cm} &= (18\alpha-7\pi)\sqrt2/(144\pi-360\alpha)\\
f(r) = f_{cm}(r) &=
\begin{cases}
 (\sqrt2-r) (r-1.3254) (9r^2 - 17 r + 3)/(1.627 (\sqrt2-1))& r\le\sqrt2\\
 0 & r >\sqrt2.
\end{cases}
\\
\end{array}
$$
We have 
\begin{equation}\label{eqn:km}K - 12M = \sqrt{1/2}\end{equation}
and
\begin{equation}f(1) = 1,\quad f(\sqrt2) =0\end{equation}


\begin{conjecture}[Marchal-2]\label{conj:m1} For any saturated packing $\Lam$, and
any $\lam_0\in\Lam$, we have
$$
\sum_{\bl=(\lam_0,\lam_1)\in\Lam(1)} f(r(\bl)) \le 12.
$$
\end{conjecture}

\begin{theorem}\label{theorem:mk}
Conjecture~\ref{conj:m1} implies the Kepler conjecture.
\end{theorem}

\begin{proof} 
We show that $A(\lam_0)  = 8 K - \sum 8 M f(r(\bl)) -\op{vol}(\Omega(\lam_0))$ is fcc-compatible and negligeable.  The
result then follows from \cite[Lemma~3.3]{Hales:2006:DCG}.  This is fcc-compatible directly
by equation~(\ref{eqn:km})
and Conjecture~\ref{conj:m1}.  The issue is to prove it negligeable.  Explicitly, we need
to show there exists a constant  $C$ such that or all $r\ge 1$ and all $x\in\ring{R}^3$:
\begin{equation}\label{eqn:neg}
\sum 8K - \sum (\sum 8 M f(r)) - \sum \op{vol}(\Omega(\lam))) \le C r^2,
\end{equation}
where the outer sum runs over $\lam\in \Lam(x,r)$.


Define the {\it total solid angle} and the {\it basic dihedral angles} of $i$-cells as follows for
$i>0$.  The total solid angle is the sum of the solid angles of the $i$-cell, summed
over all the extreme points of the cell that belong to $\Lambda$.  For a $0$-cell
the total solid angle is zero. For $1$ and $2$-cells,
it is just the solid angle of the cell at $\lam_0$.  For a $3$-cell, it is the sum of the
solid angles at $\lam_0,\lam_1,\lam_2$.  And for a $4$-cell, it is the sum of the solid
angles at $\lam_0,\ldots,\lam_3$. Write $\op{tsol}(X)$ for the total solid angle of a cell $X$.


The  basic dihedral angles of a cell $X = R(\bl)_i$ are indexed by the set $D(X)$ of
edges $\op{conv}(\lam_0,v(\bl[1]))$ that lie on the boundary of the cell.  The value
$\dih(X,d)$ of
of the basic dihedral angle indexed by $d$ is just the dihedral angle along that edge
of the boundary.  The indexing set $X$ is empty for  $0$ and $1$-cells.
The (single) basic  dihedral angle of a $2$-cell is
the same as the radian angle subtended by the arc $T_2$ in the construction of $2$-cells.
The basic dihedral angles of a $3$-cell are the dihedral angles along the
six directed edges $(\lam_i,\lam_j)$, $i,j=0,1,2$.  The basic dihedral angles of a $4$-cell are
the dihedral angles along any directed edge.  Each element $d\in D(X)$ also determines
the real number $r(d) = \norm{v(\bl[1])-\lam_0}$.

We have the Marchal's fundamental estimate for any cell $X$:
\begin{equation}\label{eqn:mfe}
\op{vol}(X) \ge \left(\frac{2K}{\pi}\right) \op{tsol}(X) - \left(\frac{4M}{\pi}\right)
\sum_{d\in D(X)} \dih(X,d) f(r(d)).
\end{equation}
We write $\op{vol}_e(X)$ for the volume estimate appearing on the 
right-hand-side, to write this in the form:
\begin{equation}\label{eqn:vole}
\op{vol}(X) \ge \op{vol}_e(X).
\end{equation}
(Note that this is an inequality in at most six variables; the most difficult case to prove
is that of a $4$-cell.)  The volumes and solid angles and so forth are given by the 
explicit formulas.

Sum the inequality~(\ref{eqn:mfe}) over all cells in a large ball $B(x,r)$ to get an
inequality of the form $0\ge T_1 + T_2 + T_3$ for three sums $T_i = T_i(r)$.  We compare this term-by-term
with the three terms $T'_i(r)$ of the desired equation~(\ref{eqn:neg}).  The solid angles around each point sum to
$4\pi$, and the dihedral angles around each directed edge sum to $2\pi$, so we see that
the three terms satisfy
$$T'_i(r) \le T_i(r) + C r^2,$$
for some constant $C$ and error term $C r^2$ coming from the boundary effects of the cells $X$ that meet the boundary of $B(x,r)$.  The result follows.
\end{proof}

\subsection{Old proof of negligeability}

For reference, we insert the old proof of negligeability.\FIXX{edit}

Let $B(x,r)$ be the closed ball of radius $r\in\ring{R}$ centered
at $x$.  Let $\Lambda(x,r)=\Lambda\cap B(x,r)$.

Recall from Definition~\ref{def:negligible} that a function
$A:\Lambda\to\ring{R}$ is said to be {\it negligible} if there is a
constant $C_1$ such that for all $r\ge1$,
   $$\sum_{v\in\Lambda(x,r) } A(v) \le C_1 r^2.$$
%
 \indy{Index}{negligible}


Recall the function $A: \Lambda\to\ring{R}$ given by
Equation~\ref{eqn:A}.  Explicitly, let
   $$A(v) = A_0(\Lambda,v),$$
where $A_0$ in turn depends on functions $A_1$ and $\sigma$, as
determined by Equations~\ref{eqn:A1} and \ref{eqn:a1-sigma}, and
Definition~\ref{def:sigma}.

\begin{claim}\label{claim:negbounds}
    The absolute values of the functions $\sol$, $\op{volan}$, and $\op{vol}$
are bounded on the set of all quarters and
    quasi-regular tetrahedra.\FIXX{Give a proof somewhere. Give a reference.}
\end{claim}


\begin{lemma}\guid{SVGACAQ}\label{lemma:A1bound}  $A_1(v,Q,c)$ is bounded by a constant that is
independent of the choice of packing $\Lambda$, simplex
$Q$, context $c$, and vertex $v\in\Lambda$.
\end{lemma}

\begin{proof} By the definition of $A_1$:
    $$
    A_1(v,Q,c) \le \sol(v,Q)/3 + |\sigma(v,Q,c)|/(4\doct).
    $$
By Claim~\ref{claim:negbounds}, it is enough to bound $|\sigma|$.
For all contexts and choices of simplex, we have
    $$|\sigma(v,Q,c)|\le \sum_{i=1}^4 |\op{svan}(v_i,S)| +
        |\op{sv}_0(v,S)| 
      + |\op{sv}_0(\hat v,S)|.
    $$
It is thus enough to bound $|\op{svan}|$ and $|\op{sv}_0|$.  Referring
to their definitions, it is enough to bound $|\sol(v,S)|$,
$|\op{volan}(v,S)|$, and $\op{vol}(\op{VC}(v,S))$. The first two
are bounded by Claim~\ref{claim:negbounds}.  The set
    $$\op{VC}(v,S) = \op{VCM}(\Lambda,v)\cap \op{conv}(S)$$
is the intersection of two measurable sets.  Hence it is a
measurable subset of $\op{conv}(S)$.  The volume of $\op{VC}(v,S)$ is
dominated by that of $\op{conv}(S)$, which is given by
Claim~\ref{claim:negbounds}.  The result follows.
\end{proof}

\begin{lemma}\guid{FNJPVVT} \label{lemma:quarterfiber}
The number of quarters at a vertex is bounded by a constant that
is independent of the packing and the vertex $v$.
\end{lemma}

\begin{proof}  If $Q=\{v,v_1,v_2,v_3\}$, then $Q\subset
\Lambda(v,2\sqrt2)$.  The set $\Lambda(v,2\sqrt2)$ is finite by
Lemma~\ref{lemma:Lambda-finite} and is bounded by a constant $n$
that is independent of $v$ and $\Lambda$.  So $Q$ is an element of
the powerset of $\Lambda(v,2\sqrt2)$, which has cardinality at
most $2^n$.
\end{proof}

\begin{lemma}\guid{DGLKVWJ}\label{lemma:qr2} There exists a constant $C$ (independent of the
packing and $p\in\ring{R}^3$) such that the number of
quarters in the $Q$-system that have at least one vertex in
$B(p,r)$ and at least one vertex out is at most $C r^2$, for all
$r\ge 1$.
\end{lemma}

\begin{proof}  For each such quarter, pick a vertex $v_Q=v=(v_x,v_y,v_z)\in
\Lambda(p,r)$.  By Lemma~\ref{lemma:quarterfiber}, The map
    $$
    Q\mapsto (\lfloor 2(v_x-p_x)\rfloor, \lfloor
    2(v_y-p_y)\rfloor,\lfloor 2(v_z-p_z)\rfloor)
    $$
is a finite to one map to $\ring{Z}^3 \cap ( B(0,2r+1) \setminus
B(0,2r-k))$ (for some $k\ge 0$) whose fibers have bounded
cardinality independent of the data.  The result now follows from
Lemma~\ref{lemma:Zr2}.
\end{proof}

\begin{lemma}\guid{PZQEATR}\label{lemma:negA1} There exists a constant $C$ that is independent of
the packing and the point $x\in\ring{R}^3$ such that for
all $r\ge 1$ we have
$$\sum_{v\in\Lambda(x,r)} \sum_{Q\in\CalQ(\Lambda,v)}
      A_1 (v,Q,c) \le C r^2$$
\end{lemma}

\begin{proof}
Each quarter $Q=\{v_1,v_2,v_3,v_4\}$ in the $Q$-system occurs in
four sets $\CalQ(\Lambda,v_i)$.  By
Lemma~\ref{lemma:A1-cancel} the sum cancels, except when some
vertex of $Q$ lies inside $\Lambda(x,r)$ and another lies outside.
The number of such quarters is at most quadratic in $r$ by
Lemma~\ref{lemma:qr2}.  The contribution of each such quarter at
each vertex is at most the constant of Lemma~\ref{lemma:A1bound}.
The result follows.
\end{proof}

\begin{theorem}\label{lemma:negligible}
The function $A$ of Equation~\ref{eqn:A} is negligible.
\end{theorem}

\begin{proof}   
Referring to the definition in Equation~\ref{eqn:A1}, we find that
the terms coming from $A_1$ are negligible by
Lemma~\ref{lemma:negA1}.  Thus, it is enough to prove the
negligibility of the function $A'$ defined by
      $$
      A'(\Lambda,v) = -\op{vol}(\Omega_2(\Lambda,v))+
         \op{vol}(\op{VCM}(\Lambda,v)).$$
%The Voronoi cells partition $\ring{R}^3$, as do the $V$-cells. 
We
have $\Omega_2(\Lambda,v)\subset B(v,2)$ and by
construction, 
$\op{VC}(\Lambda,v)\subset B(v,2)$. 

Let 
  $$
  B_2(\Lambda,x,r) = \{y \in B(x,r)\mid \exists v\in \Lambda.\ 
         |v-y| < 2\}.
  $$
Hence the Voronoi cells with
$v\in \Lambda(x,r)$ cover $B_2(\Lambda,x,r-2)$ (except for a null set
consisting of the union of planes equidistant from two vertices in
$\Lambda(x,r+2)$). Moreover, the $V$-cells $\op{VCM}(\Lambda,v)$
with $v\in \Lambda(x,r)$ are contained in $B_2(\Lambda,x,r+d)$.  Hence by
the volume formula for a ball, we have
   $$
   \begin{array}{lll}
   \sum_{v\in\Lambda(x,r)} A'(v) &\le 
     -\op{vol}\,B_2(\Lambda,x,r-2)
      +\op{vol}\,B_2(\Lambda,x,r+d)\\&\le 
       -\op{vol}\,B(x,r-2)
      +\op{vol}\,B(x,r+d)\\& \le 
      C r^2
      \end{array}
   $$
for some constant $C$. This completes the proof.
\end{proof}



\section{More kissing number estimates}

This section shows how to improve on the estimates of the previous section
by combining various Marchal cells into {\it supercells}.

\begin{definition}
Set
$$
\begin{array}{lll}
  h_0  &= 1.26\\
  h_+  &= 1.3254\\
\end{array}
$$
Let $L:[1,\sqrt{2}]\to\ring{R}^2$ be the linear function interpolating the
values:
$$
L(1) = 1\quad L(h_0) = 0.
$$
Let $h_- = 1.23175\ldots$ be the unique root of the quartic polynomial
$f(r)-L(r)$ lying in the interval $[1.2,1.3]$.
Define the function $f_2:[1,\sqrt{2}]\to\ring{R}$ by
$$
f_2(r) = \max(0,\min(f(r),L(r))) = 
\begin{cases}
  f(r) & r \in [1,h_-]\\
  L(r) & r \in [h_-,h_0]\\
  0 & r \in [h_0,\sqrt2]\\
\end{cases}
$$
\end{definition}

The function $f_2$ is continuous and piecewise smooth.  We have
$f_2(r)\le f(r)$ except when $r\in [h_-,h_+]$.  The aim of this section is to prove a variant of Theorem~\ref{theorem:mk} that uses the function $f_2$ rather than $f$.  For this, we need to combine cells into larger groups, called supercells.

\begin{definition}
Let $X$ be a $k$-cell, with $k\ge 2$.  Let $e$ be an edge of $X$ that
has an endpoint at a vertex of $\Lambda$.  The edge is determined by
a pair of point $\lambda_1,\lambda_2\in\Lambda$.  We call
$|\lambda_1-\lambda_2|/2$ the {\it reduced length} of the edge.
(For a $2$-cell this is the actual length of the edge, but for a $3$-cell
or $4$-cell this is half the length of the edge.)
We say the edge is {\it critical} if the reduced length lies in the
range $[h_-,h_+]$.
Let $X$ be a $k$-cell with a critical edge.  The weight of $X$  
is $1/m$, where $m:1..6$ is the
number of critical edges on $X$.
(A $2$-cell with a critical edge has weight $1$; a $3$-cell can have
weights $1$, $1/2$, $1/3$; a $4$-cell has possible weights $1/m$, $1\le m\le 6$.)
\end{definition}

\begin{definition}
Let $e$ be a critical edge of a $k$-cell for some $k\ge 1$.
A supercell is a formal linear combination
$$
\sum_X [X] w(X),
$$
where $w(X)$ is the weight of $X$ and the sum runs over all cells with an edge along $e$.  If $Z = \sum_X [X] w(X)$ is a formal linear combination of cells, define
$$
\begin{array}{lll}
\Gamma(Z) &= \sum_S \gamma(X) w(X),\quad\hbox{ and }\\
\gamma(X) &=  (\op{vol}(X)-\op{vol}_e(X)).
\end{array}
$$

\end{definition}

\begin{theorem}\label{lemma:superineq} 
Let $Z$ be any supercell.  Then $\Gamma(Z)\ge 0$.
\end{theorem}

The proof of this theorem will occupy the rest of the section.  Before giving the proof, we give its relation to the Kepler conjecture.


\begin{conjecture}\label{conj:m2} For any saturated packing $\Lam$, and
any $\lam_0\in\Lam$, we have
\begin{equation}\label{eqn:f12}
\sum_{\bl=(\lam_0,\lam_1)\in\Lam(1)} f_2(r(\bl)) \le 12.
\end{equation}
\end{conjecture}

\begin{theorem}\label{theorem:mk2}
Conjecture~\ref{conj:m1} implies the Kepler conjecture.
\end{theorem}

\begin{proof}  For cells $X$ that do not form part of a supercell,
the proof is just as in the proof of Theorem~\ref{theorem:mk2}.
Each cell $X$ that has a critical edge belongs to $1/w(X)$ different
supercells, each with weight $w(X)$.  In supercells, we replace inequality
(\ref{eqn:vole}) with Theorem~\ref{lemma:superineq}.
\end{proof}

Now we turn to the proof of Theorem~\ref{lemma:superineq}.
We fix one critical edge and consider the supercell of all cells around it.  Consider all the faces along the edge $\{\lambda_0,\lambda_1\}$, consisting of three vertices $\{\lambda_0,\lambda_1,\lambda_2\}$ with circumradius less than $\sqrt2$.  Call such faces {\it blades}.  The proof will be divided into cases, according to the number of blades along the critical edge.

We call a $4$-cell a quarter, when it has exactly one
critical edge and all other edges of the simplex have length
at most $2 h_-$.
By calculations\footnote{XX}, if $X$ is any cell, then
$$
 \gamma(X) \ge 0,
$$ 
except possibly when $X$ is a quarter.  Thus, it
suffices to prove the inequality~(\ref{eqn:superineq})
when there is at least one quarter along the critical
edge.  We note that the weight of any quarter is $1$.

Adjacent to a $4$-cell along a blade is always a $2$-cell or $3$-cell.
Adjacent to a $2$-cell is a $3$-cell.

The $2$-cells appear in pairs along a critical edge,
one stacked on the other along the edge.

\subsection{two blades}

We may assume that the two blades are edges of a quarter
$X_4$. 
The azimuth angle of what remains outside the quarter
is greater than $\pi$.  Thus, there must be a $3$-cell
along each face.  Let $X_3$ be one of these $3$-cells.
It has weight $1$.
Then $$\Gamma(Z)\ge \gamma(X_4)+\gamma(X_3)\ge 0,$$
by a calculation\footnote{XX}.

\subsection{three blades}

Every $4$-cell along the critical edge has azimuth
angle\footnote{XX} at most $2.8$.  Assume a quarter with $\gamma(X) <0$;
this implies\footnote{XX} 
its azimuth angle is at most $1.65$.
As
$$
2 (1.65) + 2.8 < 2 \pi,
$$
the quarter $X$ is unique.  If there is a $3$-cell adjacent
to it, we argue as in the case of two blades.  Thus, without loss of generality, there are $4$-cells adjacent to the quarter $X$.  One of the $4$-cells $Y$ has angle at least
$$
(2\pi - 1.65)/2 > 2.3.
$$
Such a $4$-cell has weight $1$.
We have\footnote{XX},
$$
\Gamma(Z)\ge \gamma(X) + \gamma(Y) \ge -0.0056 + 0.0057 \ge0.
$$

\subsection{six or more blades}

Let $X_1,\ldots,X_r$ be all the $k$-cells between a pair of
blades.  (This will be either a single $4$-cell, or a stacked pair of $2$-cells sandwiched between $3$-cells along
the two blades.)  Calculations\footnote{XX} give
$$\sum_{i=1}^r \gamma(X_i) w(X_i) \ge 0.2147 - 0.20495\sum_{i=1}^4 \op{azim},$$
where the azimuth angles runs between the two blades.
We sum this inequality over the supercell, blade by blade:
$$
\Gamma(Z) \ge 6 0.2147 - 0.20495 (2\pi) \ge 0.
$$


\subsection{five blades}


A quarter satisfies an inequality of the form\footnote{The details here need to be checked more fully.}
$$
\gamma(X) \ge a - b\, \op{azim}(X),
$$
and any other material $X_1,\ldots,X_r$
between a pair of blades
satisfies a similar inequality:
$$
\sum \gamma(X_i)w(X_i) \ge a' - b\,\op{azim}.
$$
If there are $j$-quarters $(j\ge1)$, we get
$$
\Gamma(Z) \ge j a + (5-j) a' - b\, (2\pi) \ge 0.
$$

\subsection{four blades}

We consider subcases, according to the number of quarters along the critical edge.

If there are four quarters, we use\footnote{The details here need to be checked more fully.}
$$
\gamma(X) \ge a - b\, \op{azim}(X),
$$
and
$$
\Gamma(Z) \ge 4 a  - b\, (2\pi) \ge 0.
$$

If there are three quarters. XX FINISH.

If there are two quarters, the quarters might be adjacent or not.  Assume first that they are not adjacent.  Then
on the alternate simplices, all but possibly the edge
opposite the critical edge has length in $[2,2h_-]$.
Calculations\footnote{XX} of the opposite edge shows that
one of these edges has length at most $2h_-$, making it a quarter, contrary to the assumption that there are only two quarters.

Now consider the case of two quarters that are adjacent.  We argue as in the previous case, if all of the edges of the blades, except the diagonal, have lengths in $[2,2h_-]$.  Assume an edge on a blade of length at least $2h_-$.  Then calculations\footnote{XX} give the result.

Finally, if there is a single quarter $X$, 
$$
\gamma(X) \ge a - b\, \op{azim}(X),
$$
and any other material $X_1,\ldots,X_r$
between a pair of blades
satisfies a similar inequality:
$$
\sum \gamma(X_i)w(X_i) \ge a' - b\,\op{azim}.
$$
We get
$$
\Gamma(Z) \ge  a + 3 a' - b\, (2\pi) \ge 0.
$$

\section{bounding the number of spheres}

The purpose of the remainder of this work is to prove the inequality~(\ref{eqn:f12}).   Let $L_2(r) = \max(f_2(r),L(r))$.  We will actually aim for the 
following stronger
inequality, although the weaker inequality would suffice for the proof of the Kepler conjecture.  This stronger inequality has the advantage of being piecewise linear.

\begin{conjecture}
\begin{equation}\label{eqn:L12}
\sum_{\bl=(\lam_0,\lam_1)\in\Lam(1)} L_2(r(\bl)) \le 12.
\end{equation}
\end{conjecture}
Since $L_2(r) = f_2(r) = 0$, for $r\ge h_0$, we may discard all summands for which $r(\bl)\ge h_0$.  Since $L_2(r)\le 1$, it is clear that the inequality holds whenever the number of summands is at most $12$. The following is a variant of a lemma of Marchal.


\begin{lemma}  If the number of nonzero summands is $15$ or greater, then inequality~\ref{eqn:L12} holds.
\end{lemma}

\begin{proof} 
Consider a configuration with $15$ or more nonzero summands, $\lambda_1,\ldots,\lambda_N$. 
We fix the origin $\lambda_0=0$.  Set $r_i = \|\lambda_i\|/2$.  Set
$$
a(r) = \arccos(r/2) - \pi/6.
$$
On the unit sphere,  consider the disks $D_i$ of radii $a(r_i)$, centered at $\lambda_i/\|\lambda_i\|$.  These disks do not overlap; this follows from the easy inequality 
$$
a(r_i) + a(r_j) \le \op{arc}(2r_i,2r_j,2).
$$
Let $P_i$ be the half-space containing the origin, bounded by the plane through the circular boundary of $D_i$.  The intersection of these half-spaces is a convex polytope with faces $F_i$.  The radial projection of $F_i$ to the unit sphere is a spherical polygon $R_i$ containing $D_i$.  Let $n_i$ be the number of sides to the polygon $R_i$.  The area of $R_i$ is at least the area $A(a(r_i),n_i)$ of the smallest $n_i$-gon containing $D_i$.  The smallest spherical $n$-gon containing a disk of radius $a$ is regular.  This leads to the formula
$$
A(a,n) = 2\pi - 2 n (\arcsin(\cos(a)\sin(\pi/n))).
$$
We can verify directly that
$$
A(a(r),n) \ge c_0 + c_1 n + c_2 L_2(r),\quad
n = 3,4,\ldots,\quad 1\le r\le h_0,
$$
where
$$c_0 = 0.6327,\quad c_1 = -0.0333,\quad c_2 = 0.4754.$$
The number of faces, edges, and vertices of the convex
polyhedron satisfy the Euler relation $\sum_i n_i \le (6N-12)$.
Summing over $i=1..N$, for $N\ge 15$, we get the follow
estimate on $\bar L_2 = \sum_i L_2(r_i)$:
$$
\begin{array}{lll}
4\pi &= \sum_i\op{area}(R_i)\\
     &\ge \sum_i A(a(r_i),n_i) \\
     &\ge c_0 N +c_1\sum_i n_i + c_2 \bar L_2\\
     &\ge c_0 N +c_1 (6N-12) + c_2 \bar L_2\\
     &\ge c_0 (15) +c_1 (78) + c_2 \bar L_2\\
11.94 &\ge \bar L_2. 
\end{array}
$$
\end{proof} 


\begin{lemma}\label{lemma:D'}  
If some $\lambda_1$ satisfies $\|\lambda_1\|=2$ and
$\|\lambda_1-\lambda'\|\ge 2.52$ for all $\lambda'\ne0\in\Lambda$,
then  inequality~\ref{eqn:L12} holds.
\end{lemma}

\begin{proof}  We create
a larger disk $D_1'$ centered at $\lambda_1/2$ and argue otherwise
as in the previous proof.  We assume that we have $N\ge 13$.


By the conditions of the lemma, we may take 
$$a'=\arc(2,2,2.52)-a(1.26) \approx 0.797$$
for the arcradius of this disk.  We can verify directly that
$$A(a',n) \ge c_0 + c_1 n + c_2 L_2(1) + c_3$$
where $c_3 = 0.85$.
Then 
$$
\begin{array}{lll}
4\pi &= \sum_i\op{area}(R_i)\\
     &\ge A(a',n_1)+\sum_{i>1} A(a(r_i),n_i) \\
     &\ge  c_0 N +c_1\sum_i n_i + c_2 \bar L_2 + c_3\\
     &\ge c_0 N +c_1 (6N-12) + c_2 \bar L_2 + c_3\\
     &\ge  c_0 (13) +c_1 (66) + c_2 \bar L_2 + c_3\\
11.97 &> \bar L_2. 
\end{array}
$$
\end{proof}


    %% Flyspeck
% Thomas C. Hales
% Starting with Chapter on Tame Hypermaps

% Sep 3, Removed 2 enclosed in quad properties.

\chapter{Tame Hypermap}
\indy{Index}{hypermap!tame}%

\label{sec:tame}
\indy{Index}{tame}%

\begin{summary}
This chapter is the second of the two core chapters that are devoted to the proof of the Kepler Conjecture.  If $V$ is a finite set of vectors in $\ring{R}^3$, let
$$\CalL(V) = \sum_{v\in V} L(\normo{v}/2).$$
Let $\bar B(0,r)$ be the closed ball of radius $r$.
By Corollary~\ref{cor:CE}, it suffices to show that every finite packing $V$
contained in $\bar B(0,2h_0)\setminus B(0,2)$
satisfies
\begin{equation}\label{eqn:CCE}
\CalL(V) \le 12.
\end{equation}

This chapter assumes that there exists a counterexample $V$ to this inequality and reaches a contradiction.  A subset of extremal counterexamples will be selected that are particularly well-suited for further analysis.  Every extremal counterexample gives rise to a fan and a corresponding hypermap.  
A detailed study of these hypermaps leads to a long list of properties that all such hypermaps must
possess.   A {\it tame} hypermap is defined to be precisely the set of hypermaps that have all of these
properties.   Tameness is thus an umbrella term that covers a long list of loosely related properties.

An earlier chapter on hypermaps gives an algorithm that generates all restricted hypermaps (with a given bound on the number of nodes).   Every tame hypermap is restricted and has at most $14$ nodes.  Hence a list of all tame hypermaps can be obtained by generating all restricted hypermaps and filtering out those that are not tame.   This algorithm has been implemented in computer code and run.  The result is that there is an explicit complete list of all tame hypermaps (up to isomorphism).  
This list gives the classification of tame hypermaps up to isomorphism.   This classification solves a
major step of the packing problem.

Each tame hypermap $H$ gives rise to a nonlinear optimization problem.   Maximize $\CalL(V)$ subject to the constraint that the hypermap associated with  $V$ 
is isomorphic to $H$.  This nonlinear optimization problem has a linear relaxation.  This linear program has
a maximum that is at least as large as the maximum of the nonlinear program.  The linear programs
have been solved by computer.  In every case, the maximum is less than $12$.  This means
that inequality~\ref{eqn:CCE} always holds, so that the Kepler Conjecture is confirmed.
\end{summary}
\indy{Index}{hypermap!tame}%
\indy{Index}{tame!hypermap}%



\section{Definition}


\begin{definition}[triangle,~quadrilateral]
Faces of cardinality $3$ in a hypermap are called {\it triangles}, those of
cardinality $4$ are called {\it quadrilaterals}, and so forth.  
 %
 \indy{Index}{triangle}%
 \indy{Index}{quadrilateral}%
 \indy{Notation}{pZ@$p_v$}%
\end{definition}

\begin{definition}[type,~$(p,q,r)$]\label{definition:type}
The {\it type\/} of a node is defined to be a triple of
non-negative integers $(p,q,r)$, where $p$ is the number of
triangles meeting the node, $q$ is the number of quadrilaterals
meeting it, and $r$ is the number of other faces meeting it, so that
$p+q+r$ is the total number of faces meeting the node.
%
 \indy{Index}{node!type}%
\indy{Notation}{pqr@$(p,q,r)$}%
\end{definition}


\subsection{weight assignment}\label{sec:wtassign}
\indy{Index}{weight assignment}%

Call the constant $\op{tgt}=1.541$, which arises repeatedly in
this chapter, the {\it target}. 
%
 \indy{Index}{target}%
\indy{Notation}{tgt@$\op{tgt}=1.541$}%


\begin{definition}[b]
  Define $b:\ring{N}\times \ring{N}\to \ring{R}$ by $b(p,q)=\op{tgt}$,
  except for the values in the following table:
  {
  \def\tx{\op{tgt}}
  $$\begin{matrix}  &q=0&1&2&3&4\\
           p=0&\tx&\tx&\tx&0.618&1.0\\
           1&\tx&\tx&0.66&0.618&\tx\\
           2&\tx&0.8&0.412&1.2851&\tx\\
           3&\tx&0.315&0.83&\tx&\tx\\
           4&0.35&0.374&\tx&\tx&\tx\\
           5&0.04&1.144&\tx&\tx&\tx\\
           6&0.689&\tx&\tx&\tx&\tx\\
           7&1.450&\tx&\tx&\tx&\tx
   \end{matrix}
   $$
   }
\indy{Notation}{b@$b$ (tame)}%
\end{definition}


\begin{definition}[d]
    Define $d:\ring{N}\to \ring{R}$ by
  $$d(k) = \begin{cases}
    0 & k\le3, \\
    0.206 & k=4, \\
    0.4819 & k=5, \\
    0.7578 & k=6, \\
    %1.038 & k=7, \\
    %1.315 & k=8,\\
    \op{tgt}=1.541 & \text{otherwise}.
  \end{cases}
  $$
\indy{Notation}{d@$d$ (tame)}%
\end{definition}




\begin{definition}[weight~assignment]
%
A {\it weight assignment\/} of a hypermap $H$ is a function $\tau$ on
the set of faces of $H$, taking values in the set of non-negative
real numbers. A weight assignment is {\it admissible} if the
following properties hold:
%
 \indy{Index}{weight assignment}%
 \indy{Index}{weight assignment!admissible}%
 \indy{Index}{cardinality}%
 \indy{Index}{node}%
 \indy{Notation}{ZZtau@$\tau$ (weight assignment)}%
\begin{nomerate}
 \item \case{bound~a} Let $v$ be any node of type $(5,0,1)$, and let $A$ be the set of
triangles meeting that node.
        Then
        $$\sum_{F\in A} \tau(F)
            \ge  0.63.$$
%   where $a=0.63$.
%        \label{definition:admissible:excess}
 \item \case{bound~b} If a node $v$ has type $(p,q,0)$, then
        $$\sum_{F:\,v\cap F\ne\emptyset} \tau(F) \ge b(p,q).$$
%        \label{admissible:b}
  \item \case{bound~d} If the face $F$ has cardinality $k$, then
        $\tau(F) \ge d(k)$.
\end{nomerate}
\indy{Notation}{v@$v$ (node)}%
\indy{Notation}{A@$A$ (set of triangles)}%

The sum $\sum_F \tau(F)$ (over all faces) is called the {\it total weight}. % of $\tau$.
\indy{Index}{weight!total}%
%\indy{Notation}{total@$\sum\tau$}%
\end{definition}





\subsection{hypermap property}
\label{sec:graphproperty}

A hypermap is {\it tame\/} if it satisfies the following conditions.
%
 \indy{Index}{tame}%

\begin{itemize}
    \label{definition:tame}
    %1
    \item \case{planar}  The hypermap is plain and planar.
    \item \case{simple} The hypermap is connected and simple.  In particular, every intersection of a face with a node contains at most one dart.
    \item \case{nondegenerate} The edge map $e$ has no fixed points.
    \item \case{no loops} The two darts of each edge lie in different nodes.
    \item \case{no double joins} At most one edge meets any two given nodes.
%    \label{definition:tame:40}
%    \item \case{blank}
\indy{Notation}{edgemapz@$e$ (edge map)}%
%    \item \case{triangles} If $L$ is a contour loop with three face steps, and if there exists a node in
%    the exterior of $L$, then $L$ is a face of the hypermap.
%    \label{definition:tame:3-circuit}
%\item \case{blank}
 %   \item \case{quadrilaterals} If $L$ is a contour loop with four face steps, and there are at least two nodes    in the exterior of $L$, then the interior of $L$ takes one of the forms     illustrated in Figure
 %   \ref{fig:fourcircuit}.
 %   \label{definition:tame:4-circuit}
  %  \begin{figure}[htb]
  %      \centering
  %      \myincludegraphics{\pdfp/fourcircuit.eps}
  %      \caption{Tame $4$-circuits}
  %      \label{fig:fourcircuit}
  %  \end{figure}
  \item \case{face count} The hypermap has at least three faces.
    \item \case{face size} The cardinality of each face is at least $3$ and at most $6$.
    \item \case{node count} There are either $13$ or $14$ nodes.
    \item \case{node size} The cardinality of every node is at least $3$ and at most
    $7$.
  %    \label{definition:tame:degree}
    \item \case{node types} If a node has cardinality $7$, then the type of the
       node is $(p,q,0)$ for some $p,q\ge0$.   If the
        cardinality of the node is exactly $6$, then the type of the node
        is $(5,0,1)$.
%    \label{definition:tame:degreeE}
    \item \case{weights} There exists an admissible weight assignment
        of total weight less than the target, $\op{tgt}=1.541$.
\end{itemize}
%
\indy{Index}{planar}%
\indy{Index}{simple}%
\indy{Index}{nondegenerate}%
\indy{Index}{no loops}%
\indy{Index}{no double joins}%
\indy{Index}{face}%
\indy{Index}{weight}%

\section{Classification}
    \label{sec:proof-classification}

%\section{Statement of the Theorem}
\label{sec:classification}




\begin{definition}[opposite] The opposite of a hypermap $(D,e,n,f)$ is the
hypermap $(D,f n,n^{-1},f^{-1})$.
\indy{Index}{hypermap!opposite}%
\end{definition}

\begin{lemma}\guid{PPHEUFG}
\oldrating{300}
\rating{0}
\formalauthor{Diep Trieu Thi} 
A hypermap is tame if and only if its opposite hypermap is tame as well.
\end{lemma}

A list of hypermaps appears at \cite{website:Hales:1998:Code}.

\begin{theorem}\guid{WTEMDTA}
\label{theorem:classification} Every tame hypermap is isomorphic to
a hypermap in the list~\cite{website:Hales:1998:Code}, or is isomorphic to the opposite of a
hypermap in the list.
\indy{Index}{isomorphic}%
\end{theorem}

\begin{note}%XX
The (java version of the) graph generator has been run with the new parameters described in this version of the book.  It produces finitely many graphs (about 25,000).  This new list will eventually replace the archive from 1998.  The Bauer-Nipkow formalization of the classification has not yet been adapted to this new definition of tameness.
\indy{Index}{Java}%
\end{note}


Computers are used to generate a list of all hypermaps and to check
them against the archive of tame hypermaps.  The computer program is
based on the hypermap generation construction of Section~\ref{sec:face-insert}.  According
to the results of that section, every restricted hypermap can be generated by an
elementary edge-insertion algorithm.  Every tame hypermap is a restricted hypermap.
The list of restricted hypermaps can be filtered to obtain a complete list of tame hypermaps.
\indy{Index}{computer calculation}%

\section{Contravening hypermap}

\indy{Index}{hypermap!contravening}%
%The aim is to prove Conjecture~\ref{conj:L12}.  For a contradiction, assume a counterexample. 
A counterexample to the Kepler conjecture leads to a finite packing $V\subset \bar B(0,2h_0)\setminus B(0,2)$ that violates the inequality~(\ref{eqn:CCE}).
 The purpose of this section is to select a special class of counterexamples to this inequality.
Let
\begin{equation}\label{eqn:fan-edge}
\begin{array}{lll}
 E_{std} &= \{\{v,w\}\in V^2\mid 0 < \norm{v}{w}\le 2\hm \}\\
 E_{ctc} &= \{\{v,w\}\in V^2\mid \norm{v}{w}= 2 \}\subset E_{std}.
\end{array}
\end{equation}

\begin{lemma}\guid{UBHDEUU}\rating{140}
$(V,E)$ is a fan, if $E=E_{std}$ or $E=E_{ctc}$.
\end{lemma}
$(V,E_{std})$ and $(V,E_{ctc})$ are called the standard fan and contact fan, respectively.
Let $\op{hyp}(V,E)$ be the associated hypermap.
\indy{Index}{fan!standard}%
\indy{Index}{fan!contact}%
\indy{Notation}{V@$V$ (packing)}%
\indy{Notation}{V@$( V, E ) $ (fan)}%
\indy{Notation}{hyp@$\op{hyp}$ (hypermap)}%
\indy{Notation}{E@$E_{std}$}%
\indy{Notation}{E@$E_{ctc}$}%

\begin{definition}[isolated,~surrounded]
Let $(V,E)$ be a fan.
Say that $v\in V$ is {\it isolated} in the fan if $E(v)$ is empty.
(That is, the degree of the corresponding node in the hypermap is $0$.) Say that $v\in V$ is {\it surrounded} in the fan if the azimuth angles of all darts at the node $v$ are less than $\pi$.  (In particular, the cardinality of $E(v)$ is at least three.)
\end{definition}
\indy{Index}{isolated}%
\indy{Index}{surrounded}%
\indy{Index}{angle!azimuth}%
\indy{Index}{azimuth}%
\indy{Index}{cardinality}%
\indy{Notation}{E@$E(v)$ (fan)}%

The following lemma appears in Sch\"utte and van der Waerden~\cite{vanderWaerden:1951}.

\begin{lemma}\guid{FATUGPD}\rating{200}
Given any packing $V\subset \bar B(0,2\hm)\setminus B(0,2)$,
there exists a  packing $V'$ 
in bijective correspondence with $V$:
$$
\phi:V\to V'
$$
such that $\normo{v} = \normo{\phi(v)}$ and
such that every vertex $v'$ in the contact fan of $V'$
is either isolated or surrounded.
\end{lemma}
\indy{Index}{packing!finite}%
\indy{Index}{fan!contact}%
\indy{Notation}{V@$V'$ (packing)}%

\begin{proof} Consider all finite packings in 
bijective correspondence with $V$, such that the
bijection preserves distances of points from the origin.
Among these packings, pick one $V'$ with the largest number
of isolated points in the contact graph.  If there is a point $v\in V'$ that
is not isolated and not surrounded, then it has a contact
dart with angle at least $\pi$.   Perturb $v$ away from the contacts, making it isolated, while preserving its distance to $\orz$.  The perturbed packing $V''$ has one more isolated point than $V'$, which is contrary to the supposed maximality of $V'$.  Hence, the conclusion of the 
lemma holds for $V'$.
\end{proof}

The following lemma shows that the corresponding result
holds for the standard fan.  Reader beware!  The set of isolated and surrounded vertices depends on a 
choice of fan.  The next proof makes heavy use of two different fans $(V,E_{std})$ and $(V,E_{ctc})$,
which have different sets of isolated and surrounded vertices, even though the set $V$ is
the same in both cases.



\begin{lemma}\guid{FJLBXS}\rating{1000}\label{lemma:surrounded}  % was rating{500}.
Given any packing $V\subset \bar B(0,2\hm)\setminus B(0,2)$,
there exists a  packing $V'$ 
in bijective correspondence with $V$:
$$
\phi:V\to V'
$$
such that $\normo{v} = \normo{\phi(v)}$ and
such that every vertex $v'$ in the standard fan of $V'$
is isolated or surrounded.
\indy{Index}{packing!finite}%
\end{lemma}

\begin{proof}  For any packing $V'\subset \bar B(0,2\hm)\setminus B(0,2)$, let $V'_{iso},V'_{sur}$ be the
sets of isolated and surrounded vertices of $V'$ respectively, in the {\bf contact}
fan.  
By the previous lemma,  assume
without generality that every vertex $v$ in the contact fan of $V$ is either isolated
or surrounded:
$$
V = V_{iso} \cup V_{sur}.
$$
Consider the set $\CalV$ of pairs $(V',\phi)$ consisting of a packing $V'$
and bijection $\phi:V\to V'$ such that
\begin{itemize}
\item $\phi$ preserves distances from the origin, 
\item $V_{sur}=V'_{sur}$, and 
\item the restriction of $\phi$ to $V_{sur}$ is the identity map $I$.
\end{itemize}
\indy{Notation}{V@$\CalV$ (family of packings)}%%
\indy{Notation}{V@$V_{iso}$ (contact isolated vertices)}%%
\indy{Notation}{V@$V_{sur}$ (contact surrounded vertices)}%%
\indy{Notation}{I@$I$ (identity map)}%

\claim{The set $\CalV$ is a nonempty compact topological space.}  Indeed, by fixing an enumeration
$V = \{v_0,\ldots,v_r\}$,  
the set $\CalV$ embeds into Euclidean space of dimension $3r$ under the map that sends
the pair $(V',\phi)$ to 
the point $(\phi^{-1}v_0,\ldots,\phi^{-1}v_r)$.
Thus, $\CalV$ carries a metric space topology as a subset of Euclidean space.  With a little
argument, it follows from the boundedness of $V\subset B(0,2\hm)$ that $\CalV$ is compact.
Also, $(V,I)\in \CalV$, so $\CalV\ne\emptyset$.


For any pair $(V',\phi)$ in $\CalV$ and $w\in \phi(V_{iso})$, let
$c(V',\phi,w)$, 
be the minimum distance of $w$ to a point
in $V'\setminus \{w\}$.   
For $i=1,\ldots,n=\card(V_{iso})$
define constants $c_i=c_i(V',\phi)$,
by ordering the real numbers $c(V',\phi,w)$, $w\in\phi(V_{iso})$, in increasing order:
$$
c_0 \le c_1 \le c_2 \le \cdots \le c_n.
$$
The constants $c_i:\CalV\to\ring{R}$ are continuous functions on the compact configuration space $\CalV$.
There is a nonempty 
compact subset $\CalV_0$ of $\CalV$ on which
$c_0$ attains its maximum value. Continuing recursively,
there is a nonempty compact subset $\CalV_{i+1}$ of
$\CalV_{i}$ on which $c_i$ attains its maximum value.
\indy{Notation}{c i@$c_i$ (constant)}%
\indy{Notation}{V@$\CalV_0$}%

For the configuration $(V,I)\in \CalV$, it follows that $c_0(V,I) >2$.
Hence, $c_0(V',\phi)>2$ for $(V',\phi)\in\CalV_0$.  It follows
that $V'_{iso}$ is in bijective correspondence with
$V_{iso}$, for all $(V',\phi)\in \CalV$.

\claim{Any $(V',\phi)\in \CalV_n$ has the desired property.} Otherwise,
there exists some vertex $v'\in V'_{iso}$ (in contact isolation)
is neither isolated nor surrounded in the standard fan.  
Any such vertex satisfies $c(V',\phi,\phi(v'))=c_i(V',\phi)$ for some $i$.
Among such vertices, pick the vertex with smallest index
$i$.  Write 
$$
c(V',\phi,\phi(v')) = c_i(V',\phi) = c_{i+1}(V',\phi) =\cdots= c_j(V',\phi) < c_{j+1}(V',\phi),
$$
for some $j\ge i$.  As $v'$ is not isolated in the
standard fan, it follows that $c_j(v') < 2\hm$.  As $v'$ is not surrounded in the standard fan,
in the cyclic order on
$$
\{(v',w')\in V'^2 \mid \norm{v'}{w'} = c_j(v')\},
$$
some azimuth angle is at least $\pi$.
Thus, there is a direction in which $v'$ can be perturbed
that increases $c_j$, while keeping $c_0,\ldots,c_{j-1}$
fixed.  This is contrary to the defining property of
$\CalV_n\subset\CalV_j$.  This establishes the claim.
\indy{Notation}{v@$v$ (vertex)}%
\end{proof}



\begin{lemma}\guid{FCDJDOT}\rating{100}\label{lemma:CE} 
Assume that there exists a counterexample to inequality~\ref{eqn:CCE}.  Then there also exists a counterexample $V$ with the following properties.
\begin{itemize}
\item $V\subset \bar B(0,2h_0)\setminus B(0,2)$
\item $\CalL(V) > 12$, and no finite packing in $\bar B(0,2h_0)$ attains a value larger than $\CalL(V)$.
\item The cardinality of $V$ is $13$ or $14$.
\item Every vertex $v$ is surrounded in the standard fan $(V,E_{std})$.
\item Every vertex $v$ that is not surrounded in the contact
fan $(V,E_{ctc})$ satisfies $\normo{v}=2$.
\end{itemize}
\end{lemma}

\begin{proof}  Assume that a counterexample exists.
The set of counterexamples $V\subset\bar B(0,2h_0)\setminus B(0,2)$ is a compact set.  The function $\CalL$ is a continuous function on this compact set.  Hence, there exists $V$ that maximizes the left-hand side of this inequality.

The set $V$ has cardinality $13$ or $14$ (Lemma~\ref{lemma:13-14}). Lemma~\ref{lemma:surrounded} gives the existence of a counterexample $V$ in which every vertex is surrounded or isolated in the standard fan.  By Lemma~\ref{lemma:D'}, if there are any isolated verices in the standard fan, then it is not a counterexample.  Hence, in fact every vertex is surrounded in the standard fan.  

\claim{A vertex $v$ that is not surrounded in the contact fan satisfies $\norm{v}=2$.}  Otherwise, the counterexample does not maximize the left-hand side of~\eqn{eqn:CCE}. 
\indy{Notation}{V@$V$ (set of vertices)}%
\end{proof}


\begin{definition}[contravening]
A finite packing is $V$ a {\it contravening} packing if it satisfies the properties
of Lemma~\ref{lemma:CE}.  
Its standard hypermap is $\op{hyp}(V,E_{std})$ a {\it contravening} hypermap.
\indy{Index}{packing!centered contravening}%
\indy{Index}{hypermap!contravening}%
\end{definition}



The next section studies further properties of contravening hypermaps $H$.



\section{Contravention is Tame}
\indy{Index}{tame!contravention}%
    \label{sec:contraproof}

Let $V$ be a centered packing with
standard fan $(V,E)$ and hypermap $H=\op{hyp}(V,E)=(D,e,n,f)$
be the hypermap attached to $(V,E)$.
The hypermap $H$ is plain, planar, connected, and simple.
The set of topological components of $Y(V,E)$ is in bijection with
the set of faces of $H$.  
\indy{Notation}{H@$H$ (hypermap)}%
For each face of $H$, the corresponding component $U_F$
is eventually radial with solid
angle
\indy{Notation}{U@$U_F$ (component)}%
  $$
  \sol(U_F) = 2\pi + \sum_{x\in F} (\op{azim}(x) -\pi).
  $$
Recall that
    $$\sum_{F} \sol(U_F) = 4\pi.$$
Recall the a map $v:D\to V$ that maps each dart to its vertex:
$$
v \mapsto v(x); \quad   x = (v(x),\ldots).
$$
Set 
$$h(x) = \normo{v(x)}/2.$$
Define the weight function
\begin{equation}
\begin{array}{lll}
\tau(V,E,F) &=\sum_{x\in F} \op{azim}(x)\left(1 + \dfrac{\sol_0}{\pi}(1- L(h(x)))\right) + \left(\pi+{\sol_0}\right) (2- k(F))\vspace{6pt}\\
  &= \sol(U_F) + (2- k(F))\sol_0 - \dfrac{\sol_0}{\pi}\sum_{x\in F}\op{azim}(x) (L(h(x)) - 1)\vspace{6pt}\\
&= \sol(U_F) \left( 1 + \dfrac{\sol_0}{\pi}\right) - \dfrac{\sol_0}{\pi} \sum_{x\in F} \op{azim}(x)(L(h(x))),\\
\end{array}
\end{equation}
where $\sol_0$ is the solid angle of a spherical equilateral triangle of side $\pi/3$, and $k(F)$ is the cardinality of $F$.
% 
These formulas are equivalent.  The proof of equivalence rests on the Euler formula for planar hypermaps and the solid angle formula for topological components $U_F$.
The first expression for $\tau(V,E,F)$ is particularly convenient, because it expresses $\tau$ as a sum of local contributions from each dart.
\indy{Notation}{L@$L$}%
\indy{Notation}{ZZtau@$\tau$}%
\indy{Notation}{ZZDeltanaught@$\sol_0$}%
The main conjecture may be expressed in the following alternative form:

\begin{lemma}\guid{HRXEFDM}\rating{80}
$$
\sum \tau (F) \ge 4\pi - 20\sol_0
$$
if and only if
$$
\sum L(h(x)) \le 12.
$$
\end{lemma}

\begin{proof}
The solid angles over the sphere sum to $4\pi$ and the azimuth angles at each vertex sum to $2\pi$. 
% The Euler relation for connected plain planar hypermaps gives
%$$
%\sum_F (2- k(F)) = 2\#f - 2\#e = 4 - 2\#n.
%$$
Thus,
\begin{equation}\label{eqn:delta0}
\begin{array}{lll}
\sum \tau (F) 
&= 4\pi (1 + \dfrac{\sol_0}{\pi}) - (\dfrac{\sol_0}{\pi}) 2\pi \sum_{V} L(h)\vspace{6pt}\\

&= (4\pi - 20\sol_0) + 2\sol_0 (12 - \sum_V L(h(v,\orz))).
\end{array}
\end{equation}
The result follows.
\end{proof}

The significance of the constant $\op{tgt}$ is that it is approximately equal to $4\pi - 20\sol_0$.
\indy{Notation}{tgt@$\op{tgt}=1.541$}%

\begin{theorem}\guid{MQMSMAB} \label{theorem:contravene}
Let $V$ be a contravening  packing.  Let $(H,\tau(V,E,\cdot))$ be
the hypermap and function on its faces attached to $V$ as above.
Then $H$ is a tame hypermap with admissible weight function $\tau$.
\end{theorem}
\indy{Notation}{H@$(H, \tau)$}%
\indy{Notation}{H@$H$ (hypermap)}%
\indy{Notation}{ZZtau@$\tau$ (weight)}%



\subsection{general properties}
    \label{sec:startame}

Many of the properties of tameness are trivial or have been established in earlier sections.
The following lemma quickly disposes of many of the properties of tameness.

\begin{lemma}\guid{JGTDEBU}\rating{100} %100=without (vi), %(deprecated: 300=with (vi))
A contravening hypermap $H$ satisfies Properties
\case{planar}, \case{simple}, \case{nondegenerate}, \case{no loops}, \case{no double joins}, \case{face count}, 
\case{node count}, and the first part of \case{node size}
%~(i)--(v),  % was (i)-(vi).
%(viii),  (x), and the first part of (xi).
of tameness.
\end{lemma}

\begin{proof}
The hypermap is plain, planar, connected, and simple by the general results established in the chapter on fan.  That chapter also shows that the hypermap attached to a fan satisfies properties \case{nondegenerate}, \case{no loops}, and \case{no double joins}.

%The property~(vi) is established in \cite[Lemma~3.7]{sp1}.
\claim{Properties~\case{face count} and the first half of property~\case{node size} hold}.  Indeed, every vertex is surrounded, meaning that the azimuth angles of the darts at the vertex are less than $\pi$.  As the angles around the vertex sum to $2\pi$, there are at least three darts in the node. Each of the darts in the node leads into a different face by property~\case{simple}.

%% 13 or 14 nodes.
Finally, property~\case{node count} has already been established in Lemma~\ref{lemma:CE}.
\end{proof}

There remain properties \case{face size}, %(vii), (ix), the second part of (xi), (xii), and (xiii).
\case{node types}, and \case{weights}, and the second part of \case{node size}.


\subsection{properties of nodes}
\indy{Index}{node!properties}%



\begin{lemma}\guid{CDTETAT}\rating{140} \label{lemma:0.852}
Let $H$ be a contravening
hypermap. For every dart $x$ in a triangular face of $H$,
    $$0.852\le \azim(x)\le 1.9.$$
For every dart $x$ in a nontriangular face of $H$, 
    $$1.15\le\azim(x)\le 3.27.$$
\indy{Notation}{H@$H$ (hypermap)}%
\indy{Notation}{x@$x$ (dart)}%
\indy{Notation}{v@$v$ (vertex)}%
Consequently, if a vertex $v$ has type $(p,q,0)$, then $(p,q)$
must be one of the following pairs:
$$
\begin{array}{lll}
&(0,2),~(0,3),~(0,4),~(0,5),~(1,2),~(1,3),~(1,4),~(2,1),~(2,2),~(2,3),\\
&(3,1),~(3,2),~(3,3),~(4,0),~(4,1),~(4,2),~(5,0),~(5,1),~(6,0),~(6,1),~(7,0)
\end{array}
$$
\end{lemma}
 %
 
\begin{proof}
The angle bounds are a calculation.  The sum of the azimuth angles
around a vertex satisfies:
$$
  p (0.852) + q (1.15) \le 2\pi \le p (1.9) + q (3.27),
$$
and the pairs satisfying these constraints are listed.
\end{proof}

\begin{lemma}\guid{SZIPOAS}\oldrating{80}
\rating{0}
\formalauthor{Vu Thanh}
%dcg{Lemma~21.4}{223} 
Formally contravening hypermaps satisfy the second part of property \case{node size}
%\ref{definition:tame:degree} 
of tameness: The cardinality of every
node is at most $7$.
\end{lemma}

\begin{proof}  The azimuth angle bound
$$
 (p+q+r) 0.852 \le 2\pi
$$
implies $p+q+r < 8$.
\end{proof}




\begin{lemma}\guid{KCBLRQC}\rating{300} 
Let $v$ be a node of type $(p,q,0)$.  Let $A$ be the set of faces meeting that node.  Then the property~\case{bound b} of a admissible weight assignment holds:
$$
\sum_{F\in A} \tau(V,E,F) \ge  b(p,q).
$$
\end{lemma}
\indy{Notation}{A@$A$ (faces)}%
\indy{Notation}{v@$v$ (node)}%
\indy{Notation}{pqr@$(p,q,r)$}%

\begin{proof} There is a collection of nonlinear inequalities
for $\tau(V,E,F)$ when $F$ is a triangle or quadrilateral~\cite[FUSDSPJ]{hales:2009:nonlinear}. Each nonlinear inequality has the form
\indy{Notation}{F@$F$ (polygon)}%
$$\tau(V,E,F) \ge a~\op{azim}(x) + b,$$
where $x$ is the uniquely determined dart at the node $v$ in the face $F$.  A linear programming problem comes by relaxing these nonlinear inequalities as follows.  For each nonlinear inequality,  write a linear inequality
\indy{Notation}{ZZtau@$\tau$}%
\indy{Notation}{x@$x$ (dart)}%
\indy{Notation}{v@$v$ (node)}%
$$
t(F) \ge a~z(F) + b,
$$
where $t(F)$ and $z(F)$ are variables indexed by the set $A$.
\indy{Notation}{A@$A$ (index set)}%
\indy{Notation}{t@$t$ (variable)}%
\indy{Notation}{z@$z$ (variable)}%
Then  minimize 
$$\sum_{F\in A} t(F)$$
subject to these linear inequalities and the constraint
$$
2\pi = \sum_{F\in A} z(F).
$$
Run this linear program for each of the types $(p,q,0)$ of Lemma~\ref{lemma:0.852}. The given constants are obtained from the (downward rounded) solutions to these linear programs.
\end{proof}

\begin{lemma}\guid{BDJYFFB}\rating{200}\label{lemma:deg5}
Every contravening hypermap satisfies Property \case{node types}
%\ref{definition:tame:degreeE} 
of tameness: 
If a node has cardinality $7$, then the type of the
       node is $(p,q,0)$ for some $p,q\ge0$.   If the
        cardinality of the node is exactly $6$, then the type of the node
        is $(5,0,1)$.
If the type is $(5,0,1)$, let $A$ be the set of five triangles at the
node $v$.  Then
\indy{Notation}{A@$A$ (set of triangles)}%
\indy{Notation}{v@$v$ (node)}%
$$
\sum_F \tau(V,E,F) > a.
$$
\end{lemma}



\begin{proof} These facts also come from linear programming.
The same set of nonlinear inequalities is used, and the linear
relaxation is constructed in the same way.  The linear programming
bounds are over $\op{tgt}$ in the cases excluded in the conclusion
of the lemma.  The constant $a$ is the downward rounding of the solution to the linear program for $(5,0,1)$.
\end{proof}
\indy{Notation}{tgt@$\op{tgt}=1.541$}%

\indy{Notation}{a@$a$ (constant)}%

\subsection{faces}



\begin{lemma}\guid{CRTTXAT}\rating{140}  
Every face has has cardinality at least $3$ and at most $6$.
\end{lemma}

\begin{proof} The lower bound holds, because the hypermap has no loops or double joins.  For a contradiction, let $F$ be a face of the hypermap of cardinality at least $7$.  Divide the proof into cases depending on whether the
following inequality holds:
$$
\sum _{x\in F} (\normo{v(x)}-2) \ge 4(\hm-1).
$$
If the inequality holds, then since $L$ is the linear interpolation between the points $(1,1)$ and $(\hm,0)$, and there are at most $14$ points in $V$, it follows
$$\sum_{v\in V}L(\normo{v}/2) \le 12 L(1) + 2 L(\hm) =12$$
and the main inequality holds.

Now assume that the inequality is false.
The edge $\{v,w\}$ has arclength at least
$$
\arc(\normo{v},\normo{w},\norm{v}{w}) \ge \arc(\normo{v},\normo{w},2). 
$$

A calculation~\cite[cc:arc]{hales:2009:nonlinear} gives
$$\arc(\normo{v},\normo{w},2)\ge 1 - 0.6076 (\normo{w}/2 - 1) - 0.6076 (\normo{v}/2 - 1).$$ %%CC:arc
The sum over a face of size at least $7$ gives:
$$
\begin{array}{lll}
\sum \arc(\normo{v_i},\normo{v_{i+1}},\norm{v_i}{v_{i+1}})&\ge
7 - 0.6076 \sum (\normo{v_i}-2) \\
   &\ge 7 - 0.6076 (2) (0.52) \\
   &> 2\pi,
\end{array}
$$
which is contrary to the upper bound on the edge length
sum in Lemma~\ref{lemma:convex-hyp}.
\end{proof}



\section{Main Estimate}\label{sec:weight}


The main result of this section is the following:

\begin{theorem}\guid{THJPDQA}\rating{0}  %points for OLNSWLK below.
The weight assignment $\tau$ on a contravening hypermap is admissible of  total weight less than $\op{tgt}=1.541$.
\indy{Index}{weight assignment}%
\end{theorem}
\indy{Notation}{tgt@$\op{tgt}=1.541$}%

%\subsection{admissibility}
\label{sec:admissibility}




\begin{lemma}\guid{OLNSWLK}\rating{3000} %including lemmas that lead up to it.  Interval ineq may be assumed.  
Let $F$ be a face of cardinality $k$ in a contravening hypermap.  Then
        $\tau(V,E,F) \ge d(k)$.
\end{lemma}
\indy{Notation}{F@$F$ (face)}%
\indy{Notation}{k@$k$ (cardinality)}%

In the original 1998 proof, the corresponding result
is called the ``Main Theorem.''  The proof of that 
theorem takes about 30 pages and relies on many
long computer calculations.  The proof given here
is substantially simpler than the proof of the
Main Theorem, but
it is still nontrivial.  Here, we have the advantages
of knowing that the polygons are convex, the hypermap
is simple, and each face has at most six sides.

The proof adopts the following convention for
configurations.  Let $v_0,\ldots,v_{n-1}$ be the vertices
of the polygon, and write
$$
y_i = \normo{v_i},\quad y_{ij} = \norm{v_i}{ v_j}.
$$
\indy{Notation}{v@$v$ (vertex)}%

It is helpful to keep in mind the origin of the constants $d(k)$.
Although the proofs do not produce sharp lower bounds on $\tau(V,E,F)$, the
statement of the lemma is motivated by the following configurations.
Consider a polygon with $k$ sides all of length $y_{i,i+1}=2$, heights
$y_i=2$, and $k-3$ diagonals of length $2\hm$: $y_{0,j}=2\hm$, for
$j=2,\ldots,k-2$.  Evaluating $\tau$ on these rigid configurations gives
$$
\tau(V,E,F) = \begin{cases}
0.20612\ldots & k=4\\
0.48356\ldots & k=5\\
0.760993\ldots &k=6
\end{cases}
$$
These constants $d(k)$ were chosen to be slightly smaller than these values.\footnote{Note $\tau(2.1028,2,2,2,2.52,2.52) = 0.275951\ldots < 0.277433\ldots = \tau(2,2,2,2,2.52,2.52)$.}
\indy{Notation}{d@$d$ (constant)}%
\indy{Notation}{ZZtau@$\tau$}%
\indy{Notation}{k@$k$ (cardinality)}%

\begin{proof}  Consider the cycle in the fan
giving the face $F$.  For the purposes of this
estimate,  discard all the vertices of $V$
that do not belong to the cycle as well as all edges
of $E_{std}$ that lie outside the cycle.
Write $\tau(V)$ for the value of $\tau(V,E,F)$ as a
function of the positions of the vertices in $V$,
keeping the set of edges of the cycle fixed.  The
azimuth angles associated with the face $F$ will be
call the interior angles.  
\indy{Notation}{F@$F$ (fan)}%
\indy{Notation}{E@$E_{std}$}%
\indy{Notation}{ZZtau@$\tau$}%
\indy{Index}{angle!interior}%

The idea of the proof is simple.   Deform $V$ in a way to decrease $\tau(V)$.  Pick the deformations to lie in subspaces of small dimension to allow us to verify directly that the deformations are nonincreasing in $\tau(V)$.  As the deformations continue, the configuration $V$ moves into a subspace of smaller and smaller dimension.  Eventually, the dimension becomes sufficiently small that a direct interval arithmetic calculation may show that $\tau(V)$ satisfies the desired bounds.
\indy{Notation}{V@$V$ (configuration)}%

Break the proof into several parts in the following subsections
\end{proof}


\subsection{halting conditions}
\indy{Index}{halting conditions}%

The deformation is required to maintain the following
constraints:
\begin{itemize}
\item $V\cup\{\orz\}$ is a packing.
\item The interior angles are at most $\pi$.
\item The distances satisfy $\norm{v}{w}\ge 2\hm$, if $\{v,w\}$ is
not an edge.  (Call these the diagonals.)
\end{itemize}
\indy{Notation}{V@$V\cup\{\orz\}$ (packing)}%
\indy{Index}{diagonal}%
If any of the constraints become binding, the
constraint is frozen, and the deformation continues along the remaining degrees of freedom of the configuration.  The deformations are described in detail below.
\indy{Index}{deformation}%

\subsection{recursion}
\indy{Index}{recursion}%

Initially, the diagonals satisfy $\norm{v}{w}>2\hm$.
If after deformation, equality holds;  the deformation stops, and the cycle is cut into two smaller along the diagonal
and continue recursively with deformations for each smaller cycle.  A cut edge has length $\norm{v}{w}=2\hm$, and
 all further deformations are required to keep this distance fixed.  

On a smaller cycle, let $r$ be the number of original edges and $s$ be the number of edges produces by cuts along a diagonal fixed at $2\hm$.  (Call this second kind of edge a {\it cut} edge.  Let $\tau(V(r,s))$ be the functon $\tau$ on a configuration $V$ with parameters $r$ and $s$.  Starting from polygon with at most six sides, the values $(r,s)$ that might be obtained are
\indy{Notation}{r@$r$ (parameter)}%
\indy{Notation}{s@$s$ (parameter)}%
\indy{Notation}{ZZtau@$\tau$}%
$$
(3,0),~(2,1),~(1,2),~(0,3),~
(4,0),~(3,1),~(2,2),~
(5,0),~(4,1),~
(6,0)
$$
That is, $0\le s\le 3$ and $3-s\le r\le 6-2s$.
The recursive bound is
\begin{equation}\label{eqn:drs}
\tau(V(r,s)) \ge d(r,s) = 0.103 (2-s) + 0.2759 (r+2s-4) 
\end{equation}
Note that $d(k,0) = d(k)$. Also, note that cutting
$V(r,s)$ along a new diagonal to produce $V(r_1,s_1)$
and $V(r_2,s_2)$ gives parameters that satisfy $r_1+r_2=r$ and $s_1+s_2 = 2+s$.
Also,
\begin{equation}\label{eqn:drs-add}
\begin{array}{lll}
d(r,s) &= d(r_1+r_2,s_1+s_2-2) \\
  &=0.103 (4-s_1-s_2) + 0.2759 (r_1+r_2+2s_1+2s_2-8) \\
  &=d(r_1,s_1) + d(r_2,s_2).\\
\tau(V(r,s)) &= \tau(r_1+r_2,s_1+s_2-2)\\
  &=\tau(V(r_1,s_1)) +\tau(V(r_2,s_2))\\
\end{array}
\end{equation}
The definition of $d(r,s)$ has been
chosen so that the recursive 
bound~\eqn{eqn:drs} implies the
lemma.
\indy{Notation}{d@$d$ (function)}%

In the proof that follows,  assume for a contradiction that the
inequality~\eqn{eqn:drs} fails, and that the counterexample is minimal in the sense that $r+s$ as small as possible.   Minimality allows the assumption that no diagonals develop as the deformation progresses.

\subsection{deformations}

There are the following $\tau$-nonincreasing deformations:
\begin{itemize}
\item {\bf (Vertex push)} Push one vertex radially toward $\orz$.  By the formula for $\tau$, this deformation decreases $\tau$.
\item {\bf (Lexell)} Fix all the heights $\normo{v}$. Then $\tau$ depends only on the area of the convex spherical polygon.  Consider an {\it ear} of the polygon (a triangle formed by two adjacent edges and a diagonal).  By Lexell's theorem, as one increases the length of one of the edges of the polygon, the area of the ear has a unique local maximum and no local minimum.  Thus,  deform until each edge is as long or as short as possible.
\indy{Index}{deformation}%
\indy{Index}{vertex push}%
\indy{Index}{Lexell's Theorem}%
\end{itemize}


\begin{note}%XX 
These deformations are not exactly Lexell's theorem, but a slight modification that applies to the function $\tau$.  Lexell itself is not used.  This will be corrected in the text.  
\end{note}

\subsection{flat vertices}


If the interior angle has increased to $\pi$, call the vertex {\it flat}. If there are $k$ consecutive flat vertices, there are  $k+1$ corresponding edges that form a linear series, all lying in a common plane through the origin $\orz$.  When the vertex is flat, require the Lexell triangle deformations to preserve the flat vertex.  Thus,  deform along the linear series as a whole, until each is as long or as short as possible.  The next few lemmas use the triangle inequality to constrain configurations with flat vertices.
\indy{Index}{Lexell's triangle}%
\indy{Index}{vertex!flat}%

\begin{lemma}\guid{TESVAFW}
There cannot be three consecutive flat vertices.
\end{lemma}

\begin{proof} This is because of the triangle inequality.  Three flat vertices produces a linear series of length at least
$$
4\arc(2\hm,2\hm,2) > 3,
$$
but the remaining two edges (on a hexagon) have combined length at most
$$
2\arc(2,2,2\hm) < 3.
$$
\end{proof}

\begin{lemma}\guid{SDCCMGA}
If there are two consecutive flat vertices, there is no cut edge among the corresponding set of three edges.
\end{lemma}

\begin{proof}  Argue by contradiction.  From the constraints on $r$ and $s$ given above, if $s>0$, then $r+s\le 5$.  Thus, $r+s=5$, forming a triangle with a linear series of three, and the two other edges.  If $y=\normo{v}$ where $v$ is a vertex of the triangle formed by an end of the linear series, the linear series has length at least
$$
\arc(y,2\hm,2)+\arc(2\hm,2\hm,2) +\arc(2\hm,2\hm,2\hm)
$$
which is greater than the maximum sum of the other two lengths:
$$
\arc(y,2,2\hm)+\arc(2,2,2\hm).
$$
This violates the triangle inequality.
\end{proof}

\begin{lemma}\guid{CFJSRQH}  If there are two consecutive flat vertices, then (following the Lexell triangle deformations) the edges are as short as possible.
\end{lemma}

\begin{proof} The Lexell argument either increases edges as much as possible or compresses them as much as possible.  If the vertices $v_1,v_2,v_3,v_4$ of the linear series are stretched, it has length at least
\begin{equation}\label{eqn:3side}
\arc(y_1,2,2\hm)+\arc(2,2,2\hm)+\arc(2,y_4,2\hm),
\end{equation}
where $\normo{v_i}=y_i$.
However, there are at most two other vertices and three other edges.  By the triangle inequality, the sum of these three lengths is at most \eqn{eqn:3side}.
Thus, equality is obtained in the triangle inequality, and the polygon reduces to a linear segment.  This forces the two other vertices to be precisely equal to $v_2$ and $v_3$, which is contrary to the assumption of a packing with distance separations at least $2$.
\end{proof}

\begin{lemma}\guid{OUCPLRI} If there are two consecutive flat vertices, then without loss of generality,  assume that one of them has minimal or maximal height:
$$\normo{v}\in \{2,2\hm\}.$$
\indy{Index}{minimal height}%
\indy{Index}{maximal height}%
\end{lemma}

\begin{proof}
Assume there are two adjacent flat angles $v_2,v_3$, forming a linear series $v_1,v_2,v_3,v_4$.
Assume without loss of generality (by previous reductions) that
\indy{Index}{angle!flat}%
$\norm{v_i}{v_{i+1}}=2$, for $i=1,2,3$.
Let $y_i = \normo{v_i}$.
Pull one vertex $v_2$ away from $\orz$ and push the other $v_3$ in such a way that fixes $y_2+y_3$, while decreasing the arclength of the linear series.  This follows by concavity.
In fact, the arclength is given by a sum of three terms:
  $$
  \sum_{i=1}^3\arc(y_i,y_{i+1},2).
  $$
This follows from a calculation showing that the second derivative of $\arc(t,s,2)$ with respect to $t$ is negative, for $t,s\in[2,2\hm]$.  Thus, the
sum is also concave, so that arclength is decreased as much as possible when the heights are extremal.
\end{proof}

In summary,  the following possible configurations of flat angles may occur:
\begin{itemize}
\item There are two consecutive flat angles contracted as much as possible.  They form a linear series $v_1,v_2,v_3,v_4$ where
$$
y_2\in\{2,2\hm\},\quad
y_{i,i+1}=2,\quad i=1,2,3.
$$
\item There is a single flat angle contracted as much as possible.  There
is a linear series $v_1,v_2,v_3$ where
$$
y_{i,i+1}=2,\quad i=1,2.
$$
\item There is a single flat angle stretched as much as possible.  There
is a linear series $v_1,v_2,v_3$ where
$$
\normo{v_2}=2,\quad
\norm{v_i}{v_{i+1}}=2\hm,\quad i=1,2.
$$
\indy{Index}{linear!series}%
\end{itemize}


\subsection{Adjusting heights}

Fix three consecutive vertices $u,v,w$ of the cycle.
\indy{Notation}{uvertex@$u$ (vertex)}%
\indy{Notation}{v@$v$ (vertex)}%
\indy{Notation}{wz@$w$ (vertex)}%
Assume $\norm{v}{u}$ and $\norm{v}{w}$ are already at their extremal value (either $2$ for  $2\hm$).  Then the function $\tau(V)$ may be considered as a function
of the edges of the simplex $\orz,u,v,w$ as $v$ moves with the other points of $V$ fixed.  Fix $5$ edges of the simplex as parameters and vary $\normo{v}$, so that $\tau$ becomes a function of a single variable.
\indy{Notation}{ZZtau@$\tau$}%

\begin{lemma}\guid{DFSLRHA} As a function of $\normo{v}$
 $\tau$ has negative second derivative whenever the derivative is zero.  Thus, $\tau$ has no local minimum.
\indy{Index}{derivative}%
\indy{Index}{local minimum}%
\end{lemma}

Consequently,  deform by either increasing or decreasing $\normo{v}$ as much as possible.  

\begin{proof}
The proof is an interval arithmetic calculation over a four-dimensional space~\cite[cc:d2a]{hales:2009:nonlinear}.  %%cc:d2a
\end{proof}


The same can be done when there is one flat vertex forming a linear series of length $2$, with compressed edges:

\begin{lemma}\guid{DCEETTF}
As a function of $\normo{v}$
 $\tau$ has negative second derivative whenever the derivative is zero.  Thus, $\tau$ has no local minimum.
\end{lemma}

\begin{proof}
The proof is an interval arithmetic calculation over a five-dimensional space~\cite[cc:d2b]{hales:2009:nonlinear}. %%cc:d2b
\end{proof}


Thus, if there a sequence of a flat vertex $u$, an acute vertex $v$, and another acute vertex, then  assume that $u$ or $v$ has extremal height $\in\{2,2\hm\}$.
\indy{Notation}{u@$u$ (vertex)}%
\indy{Notation}{v@$v$ (vertex)}%


\subsection{Adjusting quadrilaterals}

Suppose that there are four consecutive vertices $v_1,v_2,v_3,v_4$ that are not flat, and consider the chain of four consecutive vertices $v_1,v_2,v_3,v_4$. By Lexell arguments, each edge $\norm{v_i}{v_{i+1}}$ is $2$ or $2\hm$. Set $y_i = \normo{v_i}$ and $y_{ij} = \norm{v_i}{v_j}$. After adjusting heights $y_i$ is $2$ or $2\hm$, for $i=2,3$. This leaves four degrees of freedom:
$$
y_1,y_4,y_{14},
$$
and a diagonal to the quadrilateral, say $y_{13}$. Assume without loss of generality that  $y_{14}\ge 2\hm$.  (If it is smaller, it is an uncut edge of the polygonal, of length $2$, so that the polygon is actually a convex quadrilateral.  This special case is most easily dealt with separately.)
\indy{Notation}{y@$y$ (diagonal)}%

\begin{lemma}\guid{CMBZAOZ}
In this context, the function $\tau$ as a function of $y_1,y_4,y_{14},y_{13}$ does not have an interior point local minimum.
\end{lemma}
\indy{Notation}{ZZtau@$\tau$}%

\begin{proof} This is an interval arithmetic calculation~\cite[cc:qua]{hales:2009:nonlinear} in four variables.%% cc:quad tau-estimate, deform a quadrilateral.
\end{proof}

As a consequence, the minimum occurs when a new flat vertex forms.   Then eliminate any case with four consecutive vertices that are not flat (after treating the quadrilateral with four edges $y_{ij}=2$ as a special case).

\subsection{Cases}

More than one of these can be combined in a single polygon.   Represent the sizes of the linear series in a polygon with $k$ sides as a partition of $k$.
The number of parts of the partition gives the number of sides of the polygon when each linear series is considered a single side. 
\indy{Index}{linear!series}%

Let us review the combinatorial possibilities.  In each case,  check that the inequality holds, to complete the proof of the main estimate.  Label each 
case by the partition.  If the partition is $(\mu_1,\mu_2,\ldots)$, 
index the vertices in the same order as the partition, with $v_0$
the vertex occurring just before the first flat vertex.  For example,
if the partition is $(3,1,1)$, the polygon is a pentagon, flattened into
an effective triangle, with vertices $v_0,\ldots,v_4$, where $v_1$ and $v_2$
are the flat vertices.
\indy{Notation}{ZZmu@$\mu$ (partition)}%

The following constraints hold on the partition: it is
a partition of $k$, where $3\le k\le 6$; there
are at least three parts; the largest part is at most $3$;  there are
at most two consecutive parts that equal $1$, unless it is a special triangle $(1,1,1)$ or
quadrilateral case $(1,1,1,1)$, discussed below.  In only one case
does the order of the parts matter: $(2,1,2,1)$ versus $(2,2,1,1)$.

\begin{itemize}
\item {\bf (1,1,1)}  This is a triangle.  There are three degrees of
freedom $y_i$.  Since this is arising from a polygon that was not originally a triangle,  assume that for at least one of the edges: $y_{ij}=2\hm$.
\item {\bf (1,1,1,1)}  This is a convex quadrilateral with four sides $y_{ij}=2$ and four extremal heights $y_i$.  There is one degree of freedom given by the diagonal.
\item {\bf (2,1,1)} This is a quadrilateral, flattened into an effective triangle.  There are two degrees of freedom: one of the three variables $y_0,y_1,y_2$ is extremal.  All of the other variables are extremal.
\item {\bf (3,1,1)}  This is a pentagon, flattened into an effective triangle.  There are three degrees of freedom $y_0,y_{03},y_3$ along the flattened side, and no freedom in the rest of the figure.  It follows that $y_{01}=y_{12}=y_{23}=2$.
\item {\bf (2,2,1)} This is a pentagon, flattened into an effective triangle.  The edge lengths $y_{i,i+1}$ are all extremal.  Because of height adjustments, there are three degrees of freedom in the heights $y_i$.
\item {\bf (3,2,1)} This is a hexagon, flattened into an effective triangle.  There are four degrees of freedom.
\item {\bf (2,2,2)}  This is a hexagon, flattened into an effective triangle.  There are six degrees of freedom, given by all heights $y_i$.
\item {\bf (2,1,2,1)} This is a hexagon, flattened into an effective quadrilateral.  There are five degrees of freedom: a diagonal $y_{03}$ to the quadrilateral and four independent heights.
\item {\bf (2,2,1,1)}  This is a hexagon, flattened into an effective quadrilateral.  There are four degrees of freedom: a diagonal $y_{03}$ to the quadrilateral and three independent heights.
\indy{Index}{triangle}%
\indy{Index}{quadrilateral}%
\indy{Index}{pentagon}%
\indy{Index}{hexagon}%
\end{itemize}


Interval arithmetic calculations~\cite[cc:tau]{hales:2009:nonlinear} %% cc:par partition cases for tau[r,s]. 
for each of these cases completes the proof. The proof that contravening hypermaps are tame is complete.

\section{Linear Programs}

One can attach a linear program to each tame hypermap.
For each tame hypermap $H$ there is a configuration space $D(H)$ of all
finite packings $V\subset \bar B(0,2\hm)\setminus B(0,2)$ whose standard fan is
isomorphic to $H$.
\indy{Notation}{H@$H$ (hypermap)}%
\indy{Notation}{D@$D(H)$ (configuration space)}%

A nonlinear optimization problem asks for the maximum of
\begin{equation}\label{eqn:L2}
\sum_{v\in V} L(\normo{v}/2)
\end{equation}
over all $V\in D(H)$.

The linear program comes as a linear relaxation of this nonlinear
optimization problem on $D(H)$. That is, the optimal solution of the
the linear program has value at least as great as the corresponding
nonlinear problem.  By showing that the value of each linear program
is at most $12$, one may conclude that the maximum of \eqn{eqn:L2}
is at most $12$.


\begin{note}%XX
The linear programming part of the proof has not been completed. There are six cases that remain.
\end{note}



%%%%%%%%%%%%%%%%%%%%%%%%%%%%%%%%%%%%%%%%%%%%%%%%%%%%

    %%% June 2006.



\section{Introduction}

Until now, the most poorly redacted part of the proof of the
Kepler conjecture has been the linear programming.  The proof
involves approximately 100,000 linear programs, each involving
some 200 variables and 2000 constraints.  The data for these
linear programs fills about three gigabytes of storage.

The purpose of this document is to give a specification of the
linear programs. This specification has never been produced,
except in highly abbreviated form in the original 1998 proof. The
linear programs are left so vague that it would be extremely
difficult from that description alone to give independent
corroboration of the results.

The original proof appeared as a series of papers named {\it
Sphere Packings I} through {\it Sphere Packings VI}.  This paper
may be viewed as {\it Sphere Packings VII}.

\subsubsection{Flyspeck}

An important part of the motivation for the work comes from the
Flyspeck project, which aims to give a formal proof of the Kepler
conjecture.  The word flyspeck has the form F*P*K, which is an
acronym for the Formal Proof of the Kepler conjecture.

Steven Obua has made important progress toward the formal
verification of linear programs, as needed for the Flyspeck
project.  What he lacks to complete the project is a specification
of the linear programs in a form suitable for formal verification.
This document is the first step toward bridging that gap.

\subsubsection{Geometry}

This document avoids all geometry.  Instead we work at a purely
combinatorial level.    A hypermap (which is just a finite set
with two permutations) encodes most of the combinatorial
structure.  The rest is encoded as discrete collection of flags.
The linear programs can be specified directly from the hypermap
and flags.  There is a database of a few thousand hypermaps that
are used in the construction of the systems of inequalities.  The
flags are abstractions of various geometric conditions, and yet
it is not necessary to understand anything about these geometric
conditions to work with the flags.

\subsubsection{Formal Inequalities}

There is a tight relationship between the linear inequalities that
appear in the linear programs and certain nonlinear inequalities
that appear in a database.  A simple example is illustrated in
Figure~\ref{XX}.  The three angles $\alpha_i$ of the triangle
satisfy
    \begin{equation}
    \label{eqn:2pi}
    \alpha_1 + \alpha_2 + \alpha_3 = 2\pi.
    \end{equation}
This equality can be viewed two ways.  It can be viewed linearly
as the equation of an affine plane in $\ring{R}^3$.  Or, it can be
viewed nonlinearly, where each $\alpha_i$ is a nonlinear function
of the edge lengths of the triangles, and the equation asserts a
nonlinear relation among edge lengths.

For humans it is a simple matter to switch back and forth between
these points of view.  To understand the 1998 proof of the Kepler
conjecture, it is necessary to jump continually back and forth
between these two views of an equation.

For computers, it is not such a simple matter to jump from one
point of view to another.  To assist in the full automation of the
proof, we create a new object called {\it formal inequalities}
whose purpose is to provide the connection between the two types
of inequalities.  A formal inequality can be {\it deformalized}
into an ordinary linear inequality.  It can also admits a
nonlinear {\it interpretion}.  A formal inequality thus appears as
a structure that unifies a linear inequality with a nonlinear
inequality (Figure~\ref{fig:formal}). With this new structure, we
have the framework that would allow us to completely automate the
process of retrieving linear programs from a database of nonlinear
inequalities.

\begin{figure}[htb]
  \centering
  %% XX \myincludegraphics{\ps/formal.eps}
  \caption{Unification of nonlinear and linear inequalities}
  \label{fig:formal}
\end{figure}

\subsubsection{Infeasibility}

To prove the Kepler conjecture, it is enough to show that no
counterexamples exist.  In technical terms, it is enough to show
that no contravening centered packings exist (except for those
attached to the face-centered cubic and hexagonal close packings).

The strategy for proving the nonexistence of counterexamples is
remarkably naive. In simple terms, it goes as follows.  First, put
all the counterexamples in a box. Then measure the width of the
box. If the width of the box is negative, the box contains no
counterexamples.

This basic idea can be embellished in various ways.  The box can
be replaced by more general polyhedra defined by hyperplane
inequalities.  Instead of showing that the box has negative width,
we can show that the system of hyperplane inequalities admits no
feasible solutions.

A further refinement of this idea allows us to cover the set of
counterexamples with a finite collection of polyhedra and then
show that each polyhedron is empty.  We can do this in with an
iterative refinement.  If a given cover of the space of
counterexamples contains a polyhedron that is not empty, we can
replace it with two others that give a finer cover of the space of
counterexamples.  If by a series of  refinements we eventually
arrive at a cover by empty polyhedra, our goal (of eliminating all
counterexamples) is accomplished.

In the original 1998 proof, the series of refinements was guided
by detailed geometric information about the space of centered
packings.  In this revision of the proof, the series of
refinements is purely algebraic and combinatorial.

\subsubsection{Linear Programs}

This idea of covering the space of counterexamples with a finite
collection of empty polyhedron can be expressed somewhat
differently in terms of linear programming.  In fact, it is
possible to describe the problem in such a way that no reference
is made to the space of counterexamples.

The basic aim of the linear program is to show the infeasibility
of a certain system.  However, the primary linear programs are not
infeasible in the traditional sense.  It is necessary to weaken
the definition of infeasibility to encompass systems that
eventually become infeasible after certain branch and bound
operations.  These operations are somewhat subtle, because the
underlying combinatorial structure of the linear program changes
as branching operations progress.   In Section~\ref{sec:weak}, we
will say that a system is weakly infeasible, if a combinatorial
tree of systems can be constructed starting with the given system
as the root, and infeasible systems as the outermost leaves.  The
hypermap and its flags encode the possible branchings of the tree.
The main theorem of this document will then assert that each
primary system is weakly infeasible.

This weak infeasibility result is one of the main steps of the
proof of the Kepler conjecture.  Another document will show that
if there is a counterexample to the Kepler conjecture, then it can
be used to construct a strongly feasible solution to a primary
system.  The existence of a strongly feasible solution is
incompatible with weak infeasibility.  Thus, weak infeasibility
leads to a proof that there are no counterexamples to the Kepler
conjecture.

\subsubsection{Functional Language}

Our aim has been to give a description of the linear programming
part of the Kepler conjecture in a way that is conducive to
implementation in a functional programming language.  We have
allowed ourselves to drift somewhat from conventional mathematical
syntax toward expressions that have a bit more the look and feel
of a functional programming language.  We note that there are
scarcely any theorems in this document.  This is actually a highly
desirable thing, from the point of view of formalization.

\subsubsection{Modular design}

In a long and complex proof, it is desirable to modularize and
encapsulate the various parts of the proof to the greatest extent
possible.  We have been able to make this document almost entirely
self contained: it can be read with almost no reference to other
chapters in this book.  To achieve this, we have repeated a few
results from earlier chapters, such as material on hypermaps.
There are a few outside facts that we make reference to, such as
the existence of a database of nonlinear inequalities, and the
existence of an archive of planar hypermaps.  For the purposes of
this document, is not necessary to understand the origin of this
database or this archive. They can simply be accepted as part of
the givens of our situation.

The main theorem (Theorem~\ref{thm:lpbound}) of this document is
one of the main parts of the proof of the Kepler conjecture.
Although the theorem seems at first glance to be rather simple,
this apparent simplicity is deceptive, because the definitions on
which the theorem rests, are extremely long and cumbersome.  This
was the price of encapsulation.

\subsubsection{Warning and Hope}

There are some minor incompatibilities between how things are set
up here and how things were set up in the original 1998 proof of
the Kepler conjecture.  These incompatibilities were introduced to
simplify matters.  My hope is that they have not broken anything.
However, the only way to be sure of this is to generate all the
linear programs and check that they all hold.  I have not carried
this out, so the reader is encouraged to be a bit skeptical of
everything that is written here, at least until those
verifications have been made.

As part of the Flyspeck project, it is our hope that one day there
will be a formal proof of the main theorem of this document.  S.
Obua can provide details about the current status of this project.

\section{Review of Hypermaps}

The combinatorial structure of the Kepler conjecture is expressed
through various hypermaps.   This chapter takes
the connected sums of hypermaps, as described in Section~\ref{sec:csum}.

%
%\subsection{Basic Definitions}
%
%\begin{definition}  A {\it hypermap} $H=(D,e,n,f)$ is a finite set $D$
%together with three permutations $e,n,f:D\to D$ satisfying the
%identity $e\circ n\circ f = I$.  The elements of $D$ are called
%{\it darts}.  The permutations $e,n,f$ are called the {\it edge},
%{\it node}, and {\it face} permutations, respectively.
%\end{definition}
%

\begin{definition}  A {\it face} of a hypermap $H=(D,e,n,f)$ is an orbit of $D$
under $f$.  We write $\op{face}(\alpha)$ for the face of
$\alpha\in D$.  We write $H/f$ for the set of faces.
\end{definition}

\begin{definition}  An {\it edge} of a hypermap $H=(D,e,n,f)$ is an orbit of $D$
under $e$.  We write $\op{edge}(\alpha)$ for the edge of
$\alpha\in D$.  We write $H/e$ for the set of edges.
\end{definition}

\begin{definition}  A {\it node} of a hypermap $H=(D,e,n,f)$ is an orbit of $D$
under $n$.  We write $\op{node}(\alpha)$ for the node of
$\alpha\in D$.  We write $H/n$ for the set of nodes.
\end{definition}


\section{Linear Programming}


\subsection{Preliminaries}

There is a vast mathematical literature describing algorithms to
solve the constrained optimization problem of finding a vector
$x\in\ring{R}^n$ that minimizes the objective function
    $$x\mapsto c\cdot x$$
where $c\in\ring{R}^n$ subject to a system of linear inequalities,
$A x\le b$ for some real-valued matrix $A$ and vector $b\in
\ring{R}^m$.  These algorithms have been implemented in various
software packages.

The basic linear program admits many variations.  We can maximize
the objective function
$$x \mapsto c\cdot x$$ rather than maximize it.
Any linear equation
    $a\cdot x = b$ can be written as a system of inequalities
        $$a \cdot x \le b \text{ and } -a \cdot x \le -b.$$
This allows us to reduce constrained linear optimization with
equality constraints to a constrained linear optimization with
inequality constraints.

Because of this extensive and well-developed literature on the
subject of linear programming, we will take it as given that we
can solve large-scale linear programs in a highly reliable way.

\subsection{Infeasibility}


The linear programming problems in this document are {\it
infeasibility} problems.  We assume a system of inequalities in
the variables $x$ of the form
    \begin{equation}
    \label{eqn:lpsys}
    \begin{array}{lll}
    A x &\le c\\
    a &\le x \le b,\\
    \end{array}
    \end{equation}
with explicitly given lower and upper bounds $(a\le x\le b)$ on
the variables $x$. (An inequality of vectors $a\le b$ means
component-wise inequality $a_i\le b_i$, for all $i$.) The problem
is to verify that the system is infeasible; that is, there does
not exist an $x$ satisfying the system of inequalities.

\begin{definition}
  A certificate of infeasibility for the system \ref{eqn:lpsys},
  is a pair $(u,v)$ satisfying
    $$
    \begin{array}{lll}
    0&\le u\\
    0&\le v\\
    0&\le u A + v\\
    0& > u(c-Aa) + v(b-a). \\
    \end{array}
    $$
\end{definition}

\begin{example}
If there is an index $i$ for which $b_i < a_i$, then the system is
clearly infeasible.  In this case, we have an obvious certificate
of infeasibility given by $(u,v)=(0,e_i)$, where $e_i$ is the
standard basis vector at index $i$.
\end{example}

\begin{lemma}
  If a certificate of infeasibility $(u,v)$ exists for the system
  \ref{eqn:lpsys}, then the system is infeasible.
\end{lemma}

\begin{proof}
    Suppose for a contradiction that $x$ is a feasible solution.
    Then
    $$
    \begin{array}{lll}
    0 &\le (u A + v)(x-a) \\
      &= [u (c- A a) + v (b- a)] - [u (c - A x)] - [v (b - x)]\\
      &< 0.
    \end{array}
    $$
\end{proof}

To prove a system infeasible, it is not necessary to know how a
certificate of infeasibility is produced.  It can be produced, for
instance, by an untrusted computer algorithm.  All that is needed
to be known is that it is a certificate.

If an approximation $(u',v')$ to a certificate of infeasibility is
produced (say by computer), it can often be adjusted to give a
true certificate.  Given an approximation $(u',v')$, define
$(u,v)$ by
    $$
    \begin{array}{lll}
    u &= \max(0,u')\\
    v &= \max(0,v',-u A)\\
    \end{array}
    $$
If $(u,v)$ satisfies
    $$u( c-A a) + v(b-a) <0,$$
then it is a certificate of infeasibility.

The theory of duality for linear programming insures that if a
system is infeasible, then a certificate of infeasibility exists.
Since the equations defining a certificate of infeasibility are
linear in $u$ and $v$, linear programming algorithms may be used
to produce certificates:
    $$
    \begin{array}{rll}
        \min_{u,v} &u (c-A a) + v(b-a) \text{ such that }\\
        0 &\le u\\
        0 &\le v \\
        0& \le u A + v\\
    \end{array}
    $$
This has a feasible solution $u=v=0$.  If we add the constraint
that the objective function is at least $-1$, then there exists a
bounded feasible solution.

\begin{remark}
\label{rem:bounded}
If explicit upper bounds $b_j$  are not known
for particular variables $x_j:j\in J$, then we can search for a
certificate $(u,v)$ that satisfies $v_j = 0:j\in J$.  This
eliminates the dependence on $b_j$.  If explicit lower bounds
$a_j$ are not known for particular variables $x_j:j\in J$, then we
can search for a certificate that satisfies $v_j = 0:j\in J$ and
$u A e_j = 0: j\in J$. This eliminates the dependence on the
variables $x_j:j\in J$. Again, this augmented system has a
feasible solution $u=v=0$, and we can impose a constraint to make
the objective function bounded.
\end{remark}

\subsection{The Abstract Theory}

We assume a finite indexing set $I$ that will be used  to index
all the variables for a linear program. Let $\ring{R}^I$ be the
vector space of real-valued functions on $I$. We define the
standard dot product $(\cdot)$ on $\ring{R}^I$ by
    $$v\cdot w = \sum_{i\in I} v_i w_i.$$

\begin{definition}
A formal inequality (on the indexing set $I$) is a triple
    $$(f:I\to\ring{R},s\in\{(\le),(=),(\ge)\},c\in\ring{R}).$$
The second coordinate $s$ is a binary relation on $\ring{R}$.
\end{definition}

Let $\op{FOR}_I$ be the set of formal inequalities.  We drop the
subscript $I$ when the indexing set $I$ is fixed.  There is a
function $\op{def}:\op{FOR}\to \ring{R}^I \to \bool$ (short for
deformalize) given by
    $$\op{def}:\ (f,s,c)\ v \mapsto s(f\cdot v,c).$$

\begin{remark}
We use the shorthand of writing $(f,\op{sgn},c)$ as
   $$f\,\,\op{{\bf sgn}},\, c,$$
where $\op{{\bf sign}}$ is the curly version $(\seq,\sle,\sge)$ of
the relation symbol $(=,\le,\ge)$. For example, the formal
inequality
    $(\optt{y}(\alpha),(\ge),0)$ is written simply as $\optt{y}(\alpha)\sge 0$.
The curly relation symbol reminds us that the inequality is
formal.  Without the modified relation symbol, this shorthand
would be dangerously ambiguous, because a formula such as
    $$\optt{y}(\alpha) = 0$$
could be read either as a formal identity
$(\optt{y}(\alpha),(=),0)$, or as an ordinary vector equality
between the vector $\optt{y}(\alpha)$ and the zero vector $0$.

We generalize the notation to allow statements such as
    $$v \seq w$$
for the formal inequality $(v-w,(=),0)$.
\end{remark}

%The image of $\op{def}$ on a formal inequalities is called an
%inequality (or equality if we wish to stress that it comes from an
%equality constraint). Let $\op{INEQ}$ be the set of inequalities.

\begin{definition}
A formal disjunction $A$ is a finite set of finite sets of
$\op{FOR}_I$. (Intuitively, the set is to be thought of as
representing a formula in disjunctive normal form, whose atomic
formulas are formal inequalities.) Let $\op{LFOR}_I$ be the set of
formal disjunctions. Again, we drop the subscript $I$, when the
indexing set $I$ is fixed.  We may consider a formal inequality as
a formal disjunction under the map
    $$a \mapsto \{\{a\}\}.$$
\end{definition}

\begin{remark}
The formal disjunction
    $$A = \{\{a_{11},a_{12},\ldots,\},\ldots,\{a_{k1},\ldots\}\}$$
will be written more suggestively as
    $$A = (a_{11}\wedge a_{12}\wedge\cdots) \vee \cdots \vee (a_{k1}\wedge
    \cdots)$$
as a formula in disjunctive normal form.  For example, the formal
disjunction
    $$A =  \{\{(y,(\le),0)\},\{y,(\ge),0\}\}$$
will be written
    $$A = (y\sle 0) \vee (y \sge 0).$$
If $A$ is a singleton: $A = \{\{\{a_{11},a_{12},\ldots,\}\}$ then
we write it as a conjunction:
    $$
    A = (a_{11}\wedge a_{12}\wedge\cdots)
    $$
\end{remark}

We extend $\op{def}$ to a function $\op{ldef}$ on $\op{LFOR}$ by
    $$\op{ldef}\ A \ v  =
    (a'_{11}\wedge a'_{12}\wedge\cdots) \vee \cdots \vee (a'_{k1}\wedge
    \cdots).$$
where $a'_{ij} = \op{def}(a_{ij}) v$ and $A\in \op{LFOR}$ has the
form
    $$
    A = \{\{a_{11},a_{12},\ldots,\},\ldots,\{a_{k1},\ldots\}\}$$

If $S$ is a set of formal disjunctions, let $\op{single}(S)$ be
the union of singleton sets (that is, conjunctions):
    $$
    \op{single}(S) = \bigcup \{a \mid \{a\} \in S\}
    $$

\begin{definition}
Let $S$ be a set of formal disjunctions (on an indexing set $I$).
We say that $S$ is {\it infeasible} and write
    $$\op{INFEAS}(S)$$
if $\op{single}(S)$ is an infeasible system of linear
inequalities.
\end{definition}


\subsection{Inequalities from a Database}
\label{sec:lookup}

There is a large database of nonlinear inequalities in a small
number of variables that has been proved correct by means of
interval arithmetic.  This subsection indicates how to extract
linear inequalities from that database to be used for linear
programming.

%Suppose that we are given a set of  configurations $C$ and wish to
%generate inequalities that are valid on a subset $S\subset C$.

We assume that the indexing set $I$ contains a subset $I'$ (called
the set of nonlinear indices). Assume a nonlinearization function
    $$
    \op{interp}:\ I' \to \ring{R}^{I} \to \ring{R}.
    $$
Extend it to
    $$
    \begin{array}{rll}
    \op{interp\_vec}:&\ \ring{R}^{I} \to \ring{R}^I \to
    \ring{R}\\
    u\ v \mapsto& \sum_{i\in I'} v_i (\op{interp}_i u) +
        \sum_{j\in I\setminus I'} v_j u_j.
    \end{array}
    $$
Extend it to
    $$
    \begin{array}{rll}
    \op{interp\_for}:&\ \op{FOR}\to\ring{R}^I\to\bool\\
    (f,s,c)\ u\mapsto& s(\op{interp\_vec} f u,c)\\
    \end{array}
    $$
Finally, extend it to
    $$
    \begin{array}{rll}
    \op{interp\_lfor}:\ \op{LFOR}\to&\ring{R}^I\to\bool\\
    ((a_{11}\wedge a_{12}\wedge\cdots) \vee \cdots \vee (a_{k1}\wedge
    \cdots)\ u\mapsto& (a'_{11}\wedge a'_{12}\wedge\cdots) \vee \cdots \vee (a'_{k1}\wedge
    \cdots)\\
    \end{array}
    $$
where $a'_{ij} = \op{interp\_for} (a_{ij}) u$.


Let $\op{nonlin}(C,L)$ be the formula
    $$
    \forall y\in\ring{R}^\ring{N}.\ (\forall c\in C.\
    c(y)) \Rightarrow L(y).
    $$
Suppose that we have a database $B$ of such nonlinear
inequalities.  Given a map $\psi:\ring{N}\to I$, there is a dual
map
    $$
    \pi_\psi: \ring{R}^I \to \ring{N}\to \ring{R},\quad v\ n\mapsto
    v_{\psi n}.
    $$


\begin{definition}
    We say that the formal disjunction $A$ is a {\it lookup} in
    database $B$ with domain constraints in $S$ (a set of formal disjunctions),
    if there exists
        $$
        \op{nonlin}(C,L)\in B
        $$
    and there exist $\psi:\ring{N}\to I$ and
    $r:C\to\op{FOR}_I$
    such that the following two conditions hold
    \begin{itemize}
    \item The system $S$ implies the domain restrictions of
    $\op{nonlin}(C,L)$.  That is, for every $c \in C$,
        $$
        \forall v\in\ring{R}^I.\ \op{interp\_for} (r(c)) v
        \Rightarrow (c(\pi_\psi v))
        $$
    \item The predicate $L$ comes by interpretation of $A$.  That
    is,
        $$
        \forall v\in \ring{R}^I.\ (\forall c\in C.\ \op{interp\_for}(r(c)) v)
        \Rightarrow (L(\pi_\psi v) =
        \op{interp\_lfor}(A) v).
        $$
    \end{itemize}
\end{definition}


In brief, a formal disjunction $A$ is a lookup of
$\op{nonlin}(C,L)$ if the interpretation of $A$ coincides with the
predicate $L$ on a suitably restricted domain, and if those domain
constraints are in the system $S$.

It should be possible to write a completely automated tool that
finds the lookups from a database $B$ whose domain constraints lie
in $S$.

\begin{example}  We give a somewhat contrived example of lookup.
Later, nonlinear lookup will be performed on a massive database of
inequalities.  Suppose that our database contains the nonlinear
inequality:
    $$\forall y_1,\ldots,y_6.\ (2\le y_i \le 2.1,\ i=1,\ldots
    6) \Rightarrow (\Gamma(y_1,\ldots,y_6) \le \pt),
    $$
    for some nonlinear function $\Gamma$ and constant $\pt$.
We can put this in the form $N = \op{nonlin}(C,L)$ with
    $$
    \begin{array}{llll}
    %J &=\{1,\ldots,6\}\\
    C &= \{\\
        &y&\mapsto (y_1\ge 2)\\
        &y&\mapsto (y_1\le 2.1)\\
        &&\cdots\\
        &y&\mapsto (y_6\le 2.1)\}\\
    L &= (y&\mapsto (\Gamma(y_1,\ldots,y_6)\le \pt)).
    \end{array}
    $$
Suppose, further, for the sake of an example that our indexing
sets contain elements $g\in I'$ and $i_1,\ldots,i_6\in I\setminus
I'$ and that
    $$
    \op{interp}(g) = (v \mapsto \Gamma(v_{i_1},\ldots,v_{i_6})).
    $$
Let $\optt{gamma} = e_g \in\ring{R}^{I}$. Then we see that the
formal disjunction (in this case a single formal inequality)
$\optt{gamma} \sle \pt$ is a lookup in the database $B = \{N\}$
with domain constraints in
    $$
    S = \{e_{i_1}\sge 2,\ldots,e_{i_6}\sge 2,
          e_{i_1}\sle 2.1,\ldots, e_{i_6}\sle 2.1\},
    $$
\end{example}


\section{Hypermap Systems}

The previous sections have described the general theory of linear
programming without reference to the specifics of the Kepler
conjecture.  It is now time to develop linear programming in the
specific context of  hypermaps.

\subsection{Statement of the Main LP Bound}

Let $\mathcal H$ be the set of  hypermaps that are connected,
plain, planar, and simple. For each $H\in \mathcal H$, there is a
set of combinatorial flags $\op{fl}$.  This set is defined in
Section~\ref{sec:flag}.  Section~\ref{sec:var-index} will specify
an indexing set $I$, a subset of nonlinear indices $I'\subset I$,
and an interpretation
  $$\op{interp}:\ I'\to \ring{R}^I\to\ring{R}$$
for $I'$.  The data $I$, $I'$, and $\op{interp}$ depend on the
parameters $(H,\op{fl})$: $I = I(H,\op{fl})$, etc.

\begin{definition}
A {\it hypermap system} is a triple $(H,\op{fl},S)$ where
    $H\in {\mathcal H}$,
    $\op{fl}$ is a
    combinatorial flag,
    and $S\subset \op{LFOR}_I$, with $I = I(H,\op{fl})$.
    Let $\op{HS}$ be the set of hypermap systems.
\end{definition}

Section~\ref{sec:PTHS} defines a subset $\op{PTHS}$ of $\op{HS}$
(the set of primary tame hypermap systems).  There is a database
of nonlinear inequalities (B) that is described in
Section~\ref{XX}. Section~\ref{sec:basic} will describe the
systems $S$ for $(H,\op{fl},S)\in\op{PTHS}$.  A predicate
$\op{weak\_infeas}$ (weak infeasibility) on hypermap systems
$\op{HS}$ will be described in Section~\ref{sec:weak}.

\begin{theorem}[Main LP Bound]\label{thm:lpbound} Assume the nonlinear
inequalities (B).  Then for every $\rho\in \op{PTHS}$, we have
$\op{weak\_infeas}(\rho)$.
\end{theorem}

Briefly put, every primary tame hypermap system is weakly
infeasible.  The result is obtained by linear relaxation and
linear programming.\footnote{The warning that appears at the
beginning of the document should be repeated.  There are some mild
incompatibilities with the linear programming verifications done
in 1998, and the new linear programs have not yet been executed.
Thus, it is possible that some adjustments may have to be made
before claiming the Main LP Bound as a theorem.}

\begin{remark}  Each of the nonlinear inequalities in the database
(B) was proved by interval arithmetic as part of the 1998 proof.
Thus, we can state the Main LP Bound with this hypothesis.  The
hypothesis is included simply to keep interval arithmetic out of
the proof.
\end{remark}

\subsection{Flags}
\label{sec:flag}

The second component of a  hypermap system $(H,\op{fl},S)$ is a
combinatorial flag.  This subsection describes the flag and some
auxiliary functions.

\begin{definition}
 A flag is a tuple
    $$
    \op{fl}=(\op{qrn},\op{qre},\op{hex\_loop\_diag},
    \op{typ}(4),\op{sum\_typ}(4),\ldots,
    \op{typ}(8),\op{sum\_typ}(8))
    $$
 with flag components $\op{qrn},\ldots,\op{sum\_typ}(8)$ to be listed below.
 \end{definition}

We give a brief description of each of the flag components.  The
interpretations of the flag components are irrelevant to the proof
of the Main LP Bound.  Interpretations are merely provided as a
guide to the intuition.  All that is relevant is the domain and
range of each function.  For instance, in the following
definition, the comment that $\op{qrn}$ ``detects quasi-regular
nodes'' is irrelevant.  Any function $\op{qrn}:H/n\to\bool$
qualifies as a flag component.

 \begin{definition}
 The flag components $\op{qrn}$ and $\op{qre}$ detect {\it quasi-regular
 nodes and edges}.
        $$\op{qrn}:H \to \bool$$
        $$\op{qre}:H \to \bool$$
 The function $\op{qrn}$ is constant on orbits of $n$, so that it
 descends to a function on $H/n$.  The function $\op{qre}$ is
 constant on the orbits of $e$, so that it descends to a function
 on $H/e$.
\end{definition}

\begin{definition} The function $\op{std}$ detects standard faces
in a hypermap.
    $$\op{std}:H/f \to \bool$$
$$\op{std}(F) = \forall\alpha\in F.\
\op{qrn}(\alpha)\wedge \op{qre}(\alpha)
$$
\end{definition}

\begin{notation}
We write $(H/f)_X$ to denote the subset of $H/f$ described by the
satisfying the given flags in $X$.  If the subscript is a natural
number, we let $(H/f)_n$ be the set of faces in $H$ of cardinality
$n$.  For example, $(H/f)_{4,\op{std}}$ is the set
    $$
    \{F \in H/f \mid \card(F) = 4 \wedge \op{std}(F)\}.
    $$
\end{notation}

%\begin{definition}
%The flag component $\op{pri}:(H/f)_{\op{std}}\to\bool$ tells
%whether a standard face is primary.
%\end{definition}

%\begin{definition}
%The flag component $\op{is\_refined}$ keeps tracks of changes to
%an initial configuration hypermap.
%   $$\op{is\_refined}:H/f\to \bool$$
%\end{definition}

In the following definitions, we give symbolic names to elements
of various sets.  For example, in the following definition
    $$\{\op{patchable},\op{puremix}\}$$
is to be understood as a finite set of cardinality two, with
elements $\op{patchable}$ and $\op{puremix}$.  Also, some of the
values depend on parameters.  For example, in the following
definition, if $F=\{\alpha_1,\ldots,\alpha_4\}$, then
    $$
    \op{sum\_typ}(4)(F)
    $$
can take on five distinct values:
    $$
    \op{upright}, \op{flat}(\alpha_1), \ldots,
    \op{flat}(\alpha_4).
    $$
This is indicated in abbreviated form in the following
definitions.

\begin{definition}
The flag components $\op{typ}(4)$ and $\op{sum\_typ}(4)$
classify the type of a standard quad face.
    $$
    \begin{array}{rll}
    \op{typ}(4):& F\in (H/f)_{4,\op{std}}
    \to
    \{\op{patchable},\op{puremix}\}\\
    \op{sum\_typ}(4):& F\in (H/f)_{4,\op{patchable}}
    \to
    \{\op{upright},\op{flat}(\alpha:F)\}.
    \end{array}
    $$
\end{definition}

\begin{remark}
The data for $\op{flat}(\alpha)$ and $\op{flat}(f^2\alpha)$ are
equivalence in a sense that I won't spell out in detail here.
Roughly, they both correspond to a connected sum obtained by drawing the
diagonal of a quadrilateral. In the 1998 proof, these two cases
were combined into a single proof. This creates fewer cases, but
leads to a proof obligation of proving their equivalence.  (This
is geometrically obvious, but it is not obvious in formal
verification.) Following Nipkow's work on the formalization of the
classification of tame plane graphs, it is my hunch that the
formal proof will go more smoothly if the proof obligation is
dropped, and the two are treated as separate cases.
\end{remark}

\begin{definition}
The flag components $\op{typ}(5)$ and $\op{sum\_typ}(5)$
classify the type of a standard pentagonal face.
    $$
    \begin{array}{rll}
    \op{typ}(5):& F\in (H/f)_{5,\op{std}}
    \to
    \{\op{patchable},\op{trunc},\op{aggA},\op{aggB},\op{crowd3},\op{notpri}\}\\
    \op{sum\_typ}(5):& F\in (H/f)_{5,\op{patchable}}
    \to
    \{\op{loop5},\op{flat}(\alpha:F),\op{flat2}(\alpha:F),\op{loop4}\}.
    \end{array}
    $$
\end{definition}
%% not primary is standard but occurring in a connected sum.

\begin{definition}
The flag components $\op{typ}(6)$ and $\op{sum\_typ}(6)$
classify the type of a standard hexagonal face.
    $$
    \begin{array}{rll}
    \op{typ}(6):& F\in (H/f)_{6,\op{std}}
    \to
    \{\op{patchable},\op{trunc},\op{agg},\op{crowd3},\op{crowd4},\op{loop51}\}\\
    \op{sum\_typ}(6):& F\in (H/f)_{6,\op{patchable}}
    \to
    \{\text{list connected sumes}\ldots\}.
    \end{array}
    $$
\end{definition}

%% XX list connected sumes 6,7,8

\begin{definition}
The flag components $\op{typ}(7)$ and $\op{sum\_typ}(7)$
classify the type of a standard heptagonal face.
    $$
    \begin{array}{rll}
    \op{typ}(7):& F\in (H/f)_{7,\op{std}}
    \to
    \{\op{patchable},\op{trunc},\op{crowd3},\op{crowd4}\}\\
    \op{sum\_typ}(7):& F\in (H/f)_{7,\op{patchable}}
    \to
    \{\text{list connected sumes}\ldots\}.
    \end{array}
    $$
\end{definition}

\begin{definition}
The flag components $\op{typ}(8)$ and $\op{sum\_typ}(8)$
classify the type of a standard octagonal face.
    $$
    \begin{array}{rll}
    \op{typ}(8):& F\in (H/f)_{8,\op{std}}
    \to
    \{\op{patchable},\op{trunc},\op{crowd3},\op{crowd4}\}\\
    \op{sum\_typ}(8):& F\in (H/f)_{8,\op{patchable}}
    \to
    \{\text{list connected sumes}\ldots,\op{agg51}\}.
    \end{array}
    $$
\end{definition}


\begin{definition} The function $\op{upq}$ detects upright
quarters in a hypermap.
    $$\op{upq}:\op{dart}(H) \to \bool$$
    $$\op{upq}(\beta) =  (\card(\op{face}\beta)=3)\wedge
      \forall\alpha\in \op{face}(\beta).\
         \op{qre}(\alpha)\wedge
         (\op{qrn}(\alpha)\Leftrightarrow
         (\beta\ne\alpha))$$
\end{definition}

\begin{definition}
The function $\op{flatq}$ detects flat quarters in a hypermap.
     $$\op{flatq}:H \to \bool$$
    $$\op{flatq}(\beta) = (\card(\op{face}\beta)=3)\wedge
       \forall\alpha\in \op{face}(\beta).\
         \op{qrn}(\alpha)\wedge
         (\op{qre}(\alpha) \Leftrightarrow
         (\beta\ne\alpha)).$$
\end{definition}
(Note: the dart $\beta$ on which $\op{flatq}$ is true is the one
for which $\op{edge}(\beta)$ is not quasi-regular; that is, the
diagonal.)


\begin{definition}
The function
    $$\op{inquad}:(H/n)\to \bool$$
gives whether an upright diagonal lies on upright quarter of type
$(4,0)$ inside a quad cluster.
    $$\op{inquad}(N) =
    (\card(N)=4) \wedge (\forall \alpha\in N.\
        \op{upq}(\alpha))
    $$
\end{definition}

\begin{definition}
The function
    $$\op{quad\_diag}:\op{dart}(H)_{\neg\op{qre}} \to \bool$$
tells whether an edge that is not quasi-regular is the diagonal of
a quadrilateral.  The function $\op{quad\_diag}$ is constant on
orbits of $e$, so it descends to a function on a subset of $H/e$.
$$
    \op{quad\_diag}(\alpha) = \op{flatq}(\alpha)\wedge
    \op{flatq}(e\alpha).
$$
\end{definition}

\begin{definition}
The flag component
    $$
    \op{hex\_loop\_diag}:\op{dart}(H)_{\neg\op{qre}}\to\bool
    $$
tells whether an edge that is not quasi-regular is the diagonal of
a hexagon, along a slice in a loop.  The function
$\op{hex\_loop\_diag}$ is constant on orbits of $e$, so it
descends to a function on a subset of $H/e$.
\end{definition}


\begin{definition}
Let $\op{fqex}$ be the predicate that picks out a dart in flat
quarters that are not part of a quadrilateral pair.
\begin{equation}
    \begin{array}{llll}
    \op{fqex}(\beta) &= \exists F\in(H/f)_{3,\neg\op{std}}.\ \op{flatq}(\beta) &\wedge\\
        \beta\in F &\wedge
        \ \neg\op{quad\_diag}(\beta).
    \end{array}
\end{equation}
\end{definition}




\subsection{Bifurcation}



\begin{definition}
We define a bifurcation relation $\op{bif}(\rho,\rho',\rho'')$ on
triples of  hypermap systems.  Let $\rho=(H,\op{fl},S)$, and
similarly for $\rho'$ and $\rho''$ with appropriately primed
notation.  We say that $\rho$ {\it bifurcates} into $\rho'$ and
$\rho''$, and write $\op{bif}(\rho,\rho',\rho'')$ if there exists
a formal disjunction $A\in S$, and formal disjunctions $A',A''$
such that
    $$
    \begin{array}{rlll}
    H &= H' &= H''\\
    \op{fl} &= \op{fl}' &= \op{fl}''\\
        A &= A'&\cup\ A''\\
    S'\setminus  \{A'\} &\subset S &\cup\ \op{single}(S)\\
    S''\setminus  \{A''\} &\subset S&\cup\ \op{single}(S)\\
    \end{array}
    $$
\end{definition}
The simplest way to bifurcate is to break a disjunction into two
disjuncts, putting one disjunct into each of the two branches of
the bifurcation.



\subsection{Weak infeasibility}
\label{sec:weak}

We define a notion of infeasibility of  hypermap systems that is
broad enough to encompass branch and bound.  We describe the
predicate $\op{weak\_infeas}$ on the set of hypermap systems.

We define a tree structure that keeps track of the
branch-and-bound operations used in the search for infeasible
hypermap systems.  We define a type of an $\op{tree}$ by
    $$
    \begin{array}{lll}
    |\ \op{LEAF}&\\
    |\ \op{DATA} &\text{of } \alpha * \op{tree}\\
    |\ \op{NODE} &\text{of } \op{tree} * \op{tree}\\
    \end{array}
    $$
%
We define a predicate $\op{is\_data}$ on trees which is true
exactly when it has the form $\op{DATA}(\cdot,\cdot)$. We define a
relation $\op{head}:(\op{tree},\alpha)\to \bool$ by
    $$
    \begin{array}{lll}
    |\ \op{DATA}(\rho,s),\rho &\to \true\\
    |\  \op{\_} &\to \false
    \end{array}
    $$
Section~\ref{sec:refine} defines a binary relation
$\op{ref}(\rho,\rho')$ (called refinement) on hypermap systems. We
define a recursive predicate $\op{wft}$ (well-formed trees) on the
trees by
    $$
    \begin{array}{lll}
    |\ \op{LEAF} &\to\true\\
    |\ \op{NODE}(s,t) &\to \op{wft}(s) \wedge \op{wft}(t)
    \wedge \op{is\_data}(s) \wedge \op{is\_data}(t)\\
    |\ \op{DATA}(\rho,s) &\to \op{wft}(s) \wedge\\
        &\quad \text{case}\ s\ \text{of}\\
        &\qquad|\op{LEAF}\to\true\\
        &\qquad|\op{DATA}(\rho',s')\to \ \op{ref}(\rho,\rho')\\
        &\qquad|\op{NODE}(\op{DATA}(\rho',s'),\op{DATA}(\rho'',s''))\to
        \ \op{bif}(\rho,\rho',\rho'')\\
        &\qquad|\op{\_} \to \false\\
    \end{array}
    $$
In brief, in a well-formed tree, a child must be obtained from the
parent by bifurcating or refinement.
%
By abuse of notation, we write $\op{INFEAS}(H,\op{fl},S)$ for
$\op{INFEAS}(S)$.

We define a predicate $\op{infeas}$ (infeasible trees) on trees by
    $$
    \begin{array}{lll}
    |\ \op{LEAF} &\to \false\\
    |\ \op{DATA}(\rho,\op{LEAF}) &\to \op{INFEAS}(\rho)\\
    |\ \op{NODE}(s,s') &\to \op{infeas}(s) \wedge \op{infeas}(s')\\
    |\ \op{DATA}(\rho,s) &\to \op{infeas}(s)\\
    |\ \op{\_} &\to \false\\
    \end{array}
    $$
In brief, the data closest to the leaves must be infeasible
systems of equations.

\begin{definition}
    Let $\op{weak\_infeas}$ be the predicate on hypermap systems given by
        $$
        \op{weak\_infeas}(\rho) = \exists t.\
        \op{head}(t,\rho) \wedge \op{wft}(t) \wedge
        \op{infeas}(t).
        $$
\end{definition}

In the actual linear programming, the trees are not used directly.
Rather, all the needed properties of the function
$\op{weak\_infeas}$ are given by the following three lemmas.

\begin{lemma}
If $\op{bif}(\rho,\rho',\rho'',L)$, $\op{weak\_infeas}(\rho')$,
and
    $\op{weak\_infeas}(\rho'')$ then $\op{weak\_infeas}(\rho)$.
\end{lemma}

\begin{proof} If $t'$ and $t''$ are the trees for $\rho'$ and
$\rho''$, then use the tree $\op{DATA}(\rho,\op{NODE}(t',t''))$
for $\rho$.
\end{proof}

\begin{lemma}
If $\op{INFEAS}(\rho)$, then $\op{weak\_infeas}(\rho)$.
\end{lemma}

\begin{proof} Use the tree $\op{DATA}(\rho,\op{LEAF})$.
\end{proof}

\begin{lemma}
If $\op{ref}(\rho,\rho')$ and $\op{weak\_infeas}(\rho')$, then
  $\op{weak\_infeas}(\rho)$.
\end{lemma}

\begin{proof} If $t'$ is a tree for $\rho'$, then for $\rho$ use
    $$
    \op{DATA}(\rho,t')
    $$
\end{proof}


\subsection{Variable indexing}
\label{sec:var-index}

We describe the finite set $I = I(H,\op{fl})$ of variable indices.
Extend notation by setting $I(\rho) = I(H,\op{fl})$, for $\rho =
(H,\op{fl},S)$.  If $e_i\in \ring{R}^I$ is a standard basis
vector, we let $\iota(e_i)=i\in I$ be the corresponding index.  We
generally give names to the basis vectors, and refer to the
corresponding indices by means of the function $\iota$.





\begin{definition}
Let $D$ be the set of darts of hypermap $H$.  The set $I$ (for
$(H,\op{fl})$ is defined as the disjoint union of the following
indexing sets:
    $$\begin{array}{llll}
        \text{index} &\text{basis vector name}&\text{namesake?}
        \\
        %
        \alpha\in D, &\optt{azim}(\alpha), &\text{yes}
        \\
        %
        \alpha\in D, &\optt{yn} (\alpha)
        \\
        %%
        \alpha\in D,
        &\optt{ye} (\alpha),
        \\
        %%
        F\in H/f,
        &\optt{sol} (F), &\text{yes}
        \\
        %%
        F\in H/f
        &\optt{sc} (F)
        \\
        %%
        F\in H/f
        &{\optt{tau\_sc}} (F)
        &\text{ }\\
        %%
        F\in (H/f)_{3,flq}
        &\optt{sigmahat} (F)
        &  \text{yes}\\
        %%
        F\in (H/f)_{3,flq}
        &\optt{tauhat} (F)
        &\text{yes}\\
        %%
        F\in (H/f)_{3,\neg\op{std}}
        &\optt{eta}(F)
        &\text{yes}\\
        %%
        \alpha\in H
        &\optt{eta}'(\alpha)
        &\text{yes}\\
        %%
        F\in (H/f)_{3,std}
        &\optt{sigma\_qrtet\_x} (F)
        &\text{yes} \\
        %%
        F\in (H/f)_{3,std}
        &\optt{tau\_sigma\_x} (F)
        &\text{yes} \\
        %%
        F\in (H/f)_{\ge4,\neg\op{std}}
        &\optt{vor0}(F)
        &\text{yes}\\
        %%
        F\in (H/f)_{\ge4,\neg\op{std}}
        &\optt{tau0}(F)
        &\text{yes}\\
        %%
        \alpha\in D
        &\optt{Adih}(\alpha)
        &\text{yes}\\
        %%
        \alpha\in D
        &\optt{quo}(\alpha)
        &\text{yes}\\
        %%
        \alpha\in D
        &\optt{quob}(\alpha)
        &\\
        %%
        \alpha \in \op{flatq}
        &\optt{mu}(\alpha)
        &\text{yes}\\
        %%
        \alpha\in \op{upq}
        &\optt{nu}(\alpha)
        &\text{yes}\\
        %%
        \text{[FINISH LIST]} %% XX FINISH
    \end{array}
    $$
\end{definition}

\subsection{Interpretation}

There are a number of interpretations of variables that have the
general form
    $$\op{interp}(\iota\optt{mango})\ v\mapsto
    \op{mango}(u_0\cdot v,u_2\cdot v,\ldots,u_k\cdot v)$$
for some nonlinear function $\op{mango}$, and vectors $u_i\in
\ring{R}^I$.  We typeset the name ($\optt{mango}$) of the basis
vector in $\optt{tt}$ font and typeset name of the nonlinear
interpreting function $\op{mango}$ as an operator. Greek symbols
are expanded as text to make the correspondence with function
names in the file
    {\it definitions\_kepler.ml} explicit. (Thus, the nonlinear
    function $\sigma$ appears
    as $\optt{sigma\_qrtet\_x}$ and so forth.)
When the correspondence between a basis vector and nonlinear
interpreting function occurs on a name-by-name basis in this way,
we will omit details.  We will simply say that it is a {\it
namesake} interpretation.

In namesake interpretations, the vectors $u_i$ are general tightly
related to the index $\iota\optt{func}\in I$.  Generally, if
$\alpha$ is a dart, $F=\op{face}(\alpha)$,  and $\card(F)=3$, and
if $\optt{func}(\alpha)$ has nonlinear namesake
$\op{func}(y_1\cdot v,\ldots,v_{i_6})$, then $u_i$ follow the
standard convention:
    $$
    \begin{array}{llll}
    u_i &= \optt{yn}_i(f^i\alpha), &i &=0,1,2\\
    u_i &= \optt{ye}_i(f\alpha), &i &=3,4,5
    \end{array}
    $$
This indexing convention is followed with the nonlinear functions:
    $$
    \optt{gamma},\ \optt{nu},\ \optt{mu},\ \text{list.}
    %
    $$
(N.B. This list is incomplete.) %% XX Complete and correct.


\subsection{Penalty Constants}
\label{sec:pc}


We define a function $\xi:(H/f)_{\neg\op{std}}\to\ring{R}$ as
follows.

If $\op{qrn}(\alpha)$ for all $\alpha\in F$, and $n=\card(F)\ge
4$, then we define $\xi$ as follows.  Let $k$ be the number of
edges of $F$ such that $\neg\op{qre}(\alpha)$.  Set
    $$
    \begin{array}{lll}
    \xi(F) & (n,k)\\
    6 \xi_\Gamma & (8,0)\\
    6 \xi_\Gamma & (7,1)\\
    4 \xi_\Gamma + 2 \xi_V & (6,2)\\
    2 \xi_\Gamma + 4 \xi_V & (5,3)\\
    0 & (4,4)\\
    6 \xi_\Gamma & (7,0)\\
    5 \xi_\Gamma & (6,1)\\
    3\xi_\Gamma + 2 \xi_V & (5,2)\\
    \xi_\Gamma + 4\xi_V & (4,3)\\
    2(0.008) & (6,0) \\
    0.008 & (5,1) \\
    0 & (4,2) \\
    0.008 & (5,0) \\
    0 & (4,1) \\
    0 & (4,0)
    \end{array}
    $$
%
\begin{itemize}
\item Set
    $\xi(F) = 3\xi_\Gamma$, if $F$ has the properties $\card(F)=5$
    and
    $$\forall\alpha\in F.\
    \op{qre}(\alpha)$$
    %
    $$
    \exists\beta\in F.\forall\alpha\in F.\
    \op{qrn}(\alpha)\Leftrightarrow
(\beta\ne\alpha).$$
%
%
\item Set $\xi(F) = \xi_\Gamma + 2\xi_V$, if $F$ has the
properties
    $\card(F)=4$ and
    $$\exists\beta\in F.\forall\alpha\in F.\
    \op{qre}(\alpha)\Leftrightarrow (\beta\ne\alpha)
    $$
    $$\exists\beta\in F.\ \op{qrn}(\alpha)\Leftrightarrow
    (\beta\ne\alpha).$$
\item Set
    $$
    \begin{array}{lll}
    \xi(F) & \card(F)\\
    3 \xi_\Gamma & 4\\
     \xi_\Gamma + 2\xi_V & 3\\
    \end{array}
    $$
if $F$ has the properties
    $$
    \exists\alpha\in F.\ \op{hex\_loop\_diag}(\alpha)
    $$
    $$
    \forall\alpha\in F.\ \op{qrn}(\alpha).
    $$
\item In all other cases, set $\xi(F)=0$.
\end{itemize}


\section{Inequalities}
%
\label{sec:basic}

This section gives a set $S=\op{basic}(H,\op{fl})$ of formal
disjunctions on the indexing set $I=I(H,\op{fl})$. Elements of
this set are call {\it basic formulas}.

The following are basic formulas.  The first formal inequality is
key.

\begin{equation}
    \sum_{F\in (H/f)} \optt{sc}(F) \sge 8\,\pt.
    \label{eqn:8pt}
\end{equation}

The next family of formal disjunctions allows quite arbitrary
bifurcations. For every $f\in\ring{R}^I$ and $c\in \ring{R}$, we
have basic formulas:
\begin{equation}
    (f \sle c) \vee (f \sge c).
\end{equation}

Some of the variables are constant on orbits:

\begin{equation}
    \begin{array}{lll}
  \optt{yn}(\alpha) &\seq \optt{yn}(n\alpha),\\
  \optt{ye}(\alpha) &\seq \optt{ye}(e\alpha),\\
  \optt{eta}'(\alpha) &\seq \optt{eta}'(e\alpha),\\
  \end{array}
\end{equation}


\begin{equation}
\optt{tau\_sc}(F) = \optt{sol}(F)\,\zeta\,\pt - \optt{sc}(F).
\end{equation}

Many of the basic formulas hold subject to various restrictions on
the indexing set.  These restrictions are indicated in a line
following the basic formula.  For example, the following basic
formula holds under the restriction that $F$ is a non-standard
face of cardinality at least four:
\begin{equation}
    \label{eqn:xi}
    \begin{array}{lll}
    \optt{sc}(F) \sle \optt{vor0} + \xi(F),\\
    &F\in (H/f)_{\ge4,\neg\op{std}}
    \end{array}
\end{equation}

\begin{equation}  % Type potA along a hex loop.
    \begin{array}{lll}
    \optt{sc}(F) &\sle \optt{vor0} + \xi(F),\\
    &&F\in (H/f)_{3,\exists\alpha\in F.\ \op{hex\_loop\_diag}(\alpha)}
    \end{array}
\end{equation}

\begin{equation}  % DCG 25.6.1 page 250, one flat quarter.
    \begin{array}{lll}
    \optt{sc}(F) & \sle (s_{1+\card(F)} - Z(3,1)) + \xi_\Gamma + 2 \xi_V,\\
    &&F\in (H/f)_{\ge6}\\
    &&\exists!\alpha\in F.\ \neg\op{qre}(\alpha)
    \end{array}
\end{equation}

\begin{equation}  % DCG 25.6.1 page 251, one flat quarter.
    \begin{array}{llll}
    \optt{yn}(\beta)&\sge 2.2&\vee\\
    \optt{ye}(\alpha)&\sge 2.7&\vee\\
    \optt{sc}(F) & \sle (s_{1+\card(F)} - Z(3,1)),\\
    &&&\neg\op{qre}(\alpha),\\
    &&&\beta=f^2e\alpha,\\
    &&&\alpha\in F\in (H/f)_{\ge6}\\
    &&&\exists!\alpha\in F.\ \neg\op{qre}(\alpha)
    \end{array}
\end{equation}

%% Two flat quarters, DCG 25.6.2, page 252.
\begin{equation}
    \begin{array}{lllll}
    \optt{yn}(\beta')&\sge 2.2&&\vee\\
    \optt{yn}(\beta')&\sge 2.2&&\vee\\
    \optt{ye}(\alpha')&\sge 2.7&&\vee\\
        \optt{ye}(\alpha')&\sge 2.7&&\vee\\
    (\optt{sc}(F) & \sle (s_{2+\card(F)} - 2 Z(3,1))&\wedge \\
    \optt{tau\_sc}(F) & \sge (t_{2+\card(F)} - 2 D(3,1))&\wedge \\
    \optt{sc}(F) & \sle \optt{vor0} + 3\xi_\Gamma), \\
    &&&&\neg\op{qre}(\alpha),\\
    &&&&\beta=f^2e\alpha,\ \beta'=f^2e\alpha',\\
    &&&&\alpha\ne\alpha'\in F\in (H/f)_{\ge5}\\
    &&&&\forall\gamma\in F.\ \op{qre}(\gamma)\Leftrightarrow (\gamma\ne\alpha,\alpha')\\
    \end{array}
\end{equation}

%% DCG 25.6.2, page 252, loop of type (n,k) = (4,2), or not.
%% Constants from table 25.1.
\begin{equation}
    \begin{array}{lllll}
    \optt{sc}(F) + \optt{sc}(\op{face}(\beta)) +
    \optt{sc}(\op{face}(\beta')) &\sle -0.1999&\wedge \\
    \optt{tau\_sc}(F) + \optt{tau\_sc}(\op{face}(\beta)) +
    \optt{tau\_sc}(\op{face}(\beta')) &\sge 0.5309 &\wedge \\
    \optt{sc}(F) & \sle s_{2+\card(F)} - 2 Z(3,1) + 2(\xi_\Gamma + 2\xi_V), \\
    &&\neg\op{qre}(\alpha),\\
    &&\beta=f^2e\alpha,\ \beta'=f^2e\alpha',\\
    &&\alpha\ne\alpha'\in F\in (H/f)_{\ge5}\\
    &&\forall\gamma\in F.\ \op{qre}(\gamma)\Leftrightarrow (\gamma\ne\alpha,\alpha')\\
    \end{array}
\end{equation}



% Page 239 DCG
Let $(H,\op{fl},S)\in\op{HS}$.   Let $F\in H/f$ and $N\in H/n$. We
have the following formal equalities:
\begin{equation}
    \begin{array}{lll}
    \tau_{sc}(F) &\seq \sol(F)\zeta\pt - \op{sc}(F)\\
    %%
    \sum_{\alpha\in N}\op{azim}(\alpha) &\seq
    2\pi\\
    %%
    \op{sol}(F) &\seq \sum_{\alpha\in F}\op{azim}(\alpha) - (\card(F)
    - 2)\pi\\
    %%
    \end{array}
    \label{eqn:tau-sc}
\end{equation}


\begin{remark}  For every  primary tame hypermap system $(H,\op{fl},S)$, the
hypermap $H$ is connected, plane, and planar.  Thus, the Euler
relation holds in the form:
    $$
    \card(H/f) - \card(H/e) + \card(H/n) = 2.
    $$
From this and the previous formal inequality, it follows that
    \begin{equation}
    \sum_{F\in H/f} \optt{sol}(F) \seq 4\pi.
    \end{equation}
As a linear consequence of other equations program, it is
redundant.  Thus, it does not need to be inserted explicitly into
the linear programs.
\end{remark}



%%%%%%%%%%%%%%%%%%%%%%%%%%%%%%%%%%%%%%%%%%%%%%%%%%%%%%%%%%%%%%%%%
\subsection{Variable Bounds}
%% Page 239 DCG.
\begin{equation}
    \begin{array}{lll}
        \optt{yn}(\alpha) &\sge0\\
        \optt{yn}(\alpha) &\sle 2\sqrt2\\
        \optt{yn}(\alpha) &\sle 2t_0 &\text{if }
        \op{qrn}(\alpha)\\
        \optt{ye}(\alpha)&\sge0\\
        \optt{ye}(\alpha)&\sle 2\sqrt2\\
        \optt{ye}(\alpha)&\sle 2t_0&\text{if }
        \op{qre}(\alpha)\\
        \optt{azim}(F) &\sge 0\\
        \optt{azim}(F) &\sle 2\pi\\
        \optt{sol}(F) &\sge 0\\
        \optt{sol}(F) &\sle 4\pi\\
    \end{array}
\end{equation}





\begin{remark}
The way the infeasibility certificates are set up in
Inequality~\ref{eqn:lpsys}, it is necessary to have lower and
upper bounds on all the variables, not just those shown here.  The
need for bounds on the other variables can be removed by the
comments of Remark~\ref{rem:bounded}.
\end{remark}




%%%%%%%%%%%%%%%%%%%%%%%%%%%%%%%%%%%%%%%%%%%%%%%%%%%%%%%%%%%%%%%%%
\subsection{Triangles}

%
\begin{equation}
    \optt{sc}(F) = \optt{sigma\_qrtet\_x}(F),\quad\text{if } F\in
    (H/f)_{3,\op{std}}.
\end{equation}
%
\begin{equation}
        \optt{tau\_sc}(F) =
        \optt{tau\_sigma\_x}(F),\quad\text{if } F\in
        (H/f)_{3,\op{std}}.
\end{equation}
%
\begin{equation}
    \begin{array}{lll}
    \optt{sc}(\op{face}(\alpha)) &\sle \optt{sigmahat}(\alpha),\\
    \optt{tau\_sc}(\op{face}(\alpha)) &\sge \optt{tauhat}(\alpha),\\
    &&\alpha\in \op{fqex}
    \end{array}
\end{equation}
%
\begin{equation}
    \begin{array}{lll}
    \optt{sc}(\op{face}(\alpha)) &\sle \optt{mu}(\alpha),\\
    \optt{tau\_sc}(\op{face}(\alpha)) &\sge \optt{tau\_mu}(\alpha),\\
    &&\alpha\in \op{quad\_diag}
    \end{array}
\end{equation}



%%25.7 page 252, Branching on Upright Quarters (triangles with a potC).
%%XX Put in as a condition that the long flat edge is not a hex-loop-diag.



\subsection{Quadrilaterals}

\begin{equation}
 \optt{sc}(F) = \optt{sigma\_quad}(F),\quad\text{if } F\in
    (H/f)_{4,\op{std}}.
\end{equation}
%
\begin{equation}
 \optt{tau\_sc}(F) = \optt{tau\_quad}(F),\quad\text{if } F\in
    (H/f)_{4,\op{std}}.
\end{equation}
%

There is a list of nonlinear inequalities that appears in the
database for quadrilateral regions.  The first of these reads



\begin{equation}
\label{eqn:exdih}
\begin{array}{lll} \op{let}\ \op{J\_310151857} &= (`\forall\ v_0\ v_1\ v_2\ v_3\ v_4.\
  (\op{is\_quad\_cluster\_v}\ v_0\ v_1\ v_2\ v_3\ v_4\ \ \Rightarrow\\
        &\ \ ((\op{sigma\_quad\_approx1}\ v_0\ v_1\ v_2\ v_3\ v_4) &<
        -5.7906\\
                + &&\ 4.56766*(\op{dih\_or\_v}\ v_0\ v_1\ v_2\ v_4)))`);;
\end{array}
\end{equation}

Let $C$ be a constant less than the right hand sides of all the
inequalities.  The function $\op{dih\_or\_v}$ takes non-negative
values.  So for instance, from Inequality~\ref{eqn:exdih}, we have
the constraint
    $$C < -5.7906.$$
Then let $\op{squad}(\op{sol},\op{azim}_1,\op{azim}_2)$ be the
maximum of $\op{sigma\_quad\_approx1}$ subject to the constraints
that the simplex formed by $v_0,v_1,v_2,v_3,v_4$ has solid angle
$\op{sol}$, and azimuth angles $\op{azim}_i$.  Or let
$\op{squad}=C$, if no such simplices exist. More precisely, set
    $$\begin{array}{rll}
    \op{squad}(s,d_1,d_2) =
        \sup \{\sigma&\mid \exists v_0\cdots v_4.\ \\
           \sigma &= \op{sig\_quad\_approx1}\ v_0\cdots v_4 \\
           &\quad \op{is\_quad\_cluster}\ v_0\cdots v_4\\
           s &= \op{solid\_v}\ v_0\cdots v_4\\
           d_1 &= \op{azim}\\
           d_2 &= \op{azim}
           \}\cup \{C\}
    \end{array}
    $$




With this definition, we can rewrite the quadrilateral
inequalities and then write them as formal inequalities that take
the following shape:
    \begin{equation}
    \label{eqn:quad}
    \op{QU}(\optt{squad},\optt{sol},\optt{azim}_1,\optt{azim}_2,i) :
    (\optt{squad} \sle a_i \optt{azim}_1 + b_i \optt{azim}_2 + c_i \optt{sol}
    + d_i),
    \end{equation}
for constants $a_i,b_i,c_i,d_i$, $i=1,\ldots,N$.  The formal
inequalities $\op{QU}$ (\ref{eqn:quad}) hold for all indices $i$
and the following parameter values:


\begin{equation}
  \begin{array}{rll}
    \optt{squad} &= \optt{sc}(F),\\
    \op{sol} &= \optt{sol}(F),\\
    \op{azim}_1 &= \optt{azim}(\alpha),\\
    \op{azim}_2 &= \optt{azim}(f\alpha),\\
    &&\alpha\in F\in (H/f)_{4,\op{std}}
   \end{array}
\end{equation}

\begin{equation}
\begin{array}{rll}
    \optt{squad} &=\sum_{i=0}^3\optt{sc}(\op{face}(n^i\alpha)),\\
    \op{sol} &=
          \sum_{i=0}^3\optt{sol}(\op{face}(n^i\alpha)),\\
    \op{azim}_1 &=\optt{azim}(f\alpha)+\optt{azim}(n f\alpha),\\
     \op{azim}_2 &= \optt{azim}(f^{-1}\alpha)+\optt{azim}(n^{-1}f^{-1}\alpha),
         \\
    &&\alpha\in N\in (H/n)_{\op{inquad}},\\
\end{array}
\end{equation}

\begin{equation}
\begin{array}{rll}
    \optt{squad} &=\optt{sc}(\op{face}(\alpha))+\optt{sc}(\op{face}(\beta)),\\
    \op{sol} &= \optt{sol}(\op{face}(\alpha))+\optt{sol}(\op{face}(\beta)),\\
    \op{azim}_1 &=\optt{azim}(\alpha)+\optt{azim}(f\beta),\\
     \op{azim}_2 &= \optt{azim}(f^{-1}\beta),
         \\
    &&\alpha,\beta \in (H)_{\op{quad\_diag}},\ \beta=e\alpha.\\
\end{array}
\end{equation}

\begin{equation}
\begin{array}{rll}
    \optt{squad} &=\optt{sc}(\op{face}(\alpha))+\optt{sc}(\op{face}(\beta)),\\
    \op{sol} &= \optt{sol}(\op{face}(\alpha))+\optt{sol}(\op{face}(\beta)),\\
    \op{azim}_1 &=\optt{azim}(f^{-1}\beta),\\
     \op{azim}_2 &= \optt{azim}(\beta)+\optt{azim}(f\alpha),
         \\
    &\alpha,\beta \in (H)_{\op{quad\_diag}},\ \beta=e\alpha\\
\end{array}
\end{equation}

\begin{remark}  There is a small conflict between nonlinear and
linear variables here that is easy to resolve.  The variables
$\optt{azim}$ are nonlinear and are interpreted as the azimuth
function.  However, in the quadrilateral inequalities give above,
they are independent variables to the function $\optt{squad}$, and
should be treated as linear variables.  To overcome these
conflicting constraints, we introduce two copies of each variable,
say $\optt{azim}$ and $\optt{azim}'$, related by the equation
    $$
    \optt{azim}(\alpha) = \optt{azim}(\alpha)'.
    $$
We then use the nonlinear copy everywhere but in the quadrilateral
inequalities, we switch to the linear copy.
\end{remark}

\subsubsection{Specialized Inequalities and Constraints}



\begin{equation}
    \begin{array}{lll}
        \optt{sc}(F) &\sle s_8 + \pt\\
        \optt{tau\_sc}(F) &\sge t_8\\
        &&F \in (H/f)_{5,\op{aggA}}.
    \end{array}
\end{equation}

\begin{equation}
    \begin{array}{lll}
        \optt{sc}(F) &\sle s_7\\
        \optt{tau\_sc}(F) &\sge t_7\\
        &&F \in (H/f)_{5,\op{aggB}}.
    \end{array}
\end{equation}

\begin{equation}
    \begin{array}{lllll}
        \optt{yn}(\alpha) &\sge 2.168,&\text{some }\alpha\in F
        &\vee\\
        \optt{azim}(\alpha) &\sge 2.89,&\text{some }\alpha\in
        F&\vee\\
        \optt{ye}(\alpha)+\optt{ye}(f\alpha)&\sge
        4.804,&\text{some }\alpha\in F&\vee\\
        \sum_{\alpha\in F} \optt{ye}(\alpha) &\sge
        11.4707&&\vee\\
        \optt{sc}(F) &\sle -0.2345&&\vee\\
        \optt{tau\_sc}(F) &\sge 0.6079,\\
        &&&&F\in (H/f)_{5,\op{aggB}}.
    \end{array}
\end{equation}

\begin{equation}
    \begin{array}{llll}
        \optt{sc}(F) &\sle -0.4339&\vee\\
        \optt{sc}(F) &\sge 0.5606\\
        &&&F\in (H/f)_{n,\op{crowd3}},\quad n\ge 5.
    \end{array}
\end{equation}

\begin{equation}
    \begin{array}{llll}
        \optt{sc}(F) &\sle -0.25\\
        \optt{tau\_sc}(F) &\sge 0.4\\
        &&F\in (H/f)_{n,\op{crowd4}},\quad n\ge 6.
    \end{array}
\end{equation}

\begin{equation}
    \begin{array}{lllll}
        \optt{sc}(F) &\sle -0.31547&&\vee\\
        \optt{yn}(\alpha) + \optt{yn}(f^4\alpha) &\sge 4.6 &\text{some
        }\alpha\in F,\\
        &&&F\in (H/f)_{n,\op{crowd4}},\quad n\ge 6.
    \end{array}
\end{equation}

\begin{equation}
    \begin{array}{llll}
        \optt{sc}(F) &\sle s_8\\
        \optt{tau\_sc}(F) &\sge t_8\\
        &&F\in (H/f)_{6,\op{agg}}.
    \end{array}
\end{equation}

\begin{equation}
    \begin{array}{llll}
        \optt{sc}(F) &\sle -0.37595\\
        \optt{tau\_sc}(F) &\sge 0.65995\\
        &&F\in (H/f)_{6,\op{loop51}}.
    \end{array}
\end{equation}





\begin{equation}
    \begin{array}{lllll}
        \optt{yn}(\alpha) &\sge 2.14 &\text{some }\alpha\in F&\vee\\
        \optt{ye}(\alpha) &\sge 2.77 &\text{some }\alpha\in
        F&\vee\\
        (\optt{sc}(F) &\sle \optt{gamma}(F) &\wedge\\
         \optt{yn}(\alpha) &\sle 2.14 \text{ all }\alpha\in
         F&\wedge\\
        \optt{eta}(F) &\sle \sqrt2) &&\vee\\
        (\optt{ye}(\alpha) &\sle 2.7 &\wedge\\
        \optt{yn}(\alpha) &\sle 2.14 \text{ all }\alpha\in
         F&\wedge\\
        \optt{eta}(F) &\sge \sqrt2 &\wedge\\
        \optt{sc}(F) &\sle \optt{svor0}(F)) &&\vee\\
        (\optt{ye}(\alpha) &\sge 2.7 &\wedge\\
        \optt{yn}(\alpha) &\sle 2.14 \text{ all }\alpha\in
         F&\wedge\\
        \optt{sc}(F) &\sle \optt{svor0}(F)),\\
        &&&&\alpha\in F\in (H/f)_{3,\op{fqex}}.
    \end{array}
\end{equation}

 Let $\op{potA}$ be the predicate on $(H/f)_{3,\neg\op{std}}$
that picks out potentially type A.  That is, $\op{qrn}(\alpha)$
holds for all three darts, $\op{qre}(\alpha)$ holds for exactly
one of the three darts, and no dart is in a hex diagonal.

\begin{equation}
    \begin{array}{lll}
    \op{potA}(F) &= \exists\beta\in F.\ \forall\alpha\in F. \\
        &\op{qrn}(\alpha)\\
        &\op{qre}(\alpha) \Leftrightarrow (\alpha=\beta)\\
        &\neg\op{hex\_loop\_diag}(\alpha).
    \end{array}
\end{equation}

\begin{equation}
 \begin{array}{lllll}
    \optt{yn}(\alpha) &\sge 2.77 &&\vee\\
      \optt{yn}(\beta) &\sge 2.77 &&\vee\\
    (\optt{yn}(\alpha) &\sle 2.77 &\wedge\\
      \optt{yn}(\beta) &\sle 2.77 &\wedge\\
      \optt{sc}(F) &\sle \optt{svoran}) &&\vee\\
   (\optt{yn}(\alpha) &\sle 2.77 &\wedge\\
      \optt{yn}(\beta) &\sle 2.77 &\wedge\\
      \optt{eta}(F) &\sge \sqrt2) &&\vee\\
     &&&&\alpha\ne\beta,\gamma\in F\in (H/f)_{3,potA}\\
     &&&&\neg\op{qre}(\alpha),\neg\op{qre}(\beta),\op{qre}(\gamma).
 \end{array}
\end{equation}

Let $\op{potC}$ be the set of triangles that are potentially of
type $C$. That is, for $F\in (H/f)_{3,\neg\op{std}}$:
\begin{equation}
    \begin{array}{lll}
    \op{potC}(\beta,F) &= \beta\in F \wedge\ \forall\alpha\in F. \\
        &\op{qrn}(\alpha)\Leftrightarrow (\alpha\ne\beta)&\wedge \\
        &\op{qre}(\alpha) \Leftrightarrow (\alpha\ne f\beta).
    \end{array}
\end{equation}













\section{Final Specs}


\subsection{Primary Tame Hypermap System}
\label{sec:PTHS}

This subsection gives a specification of primary tame hypermap
systems ($\op{PTHS}$).

Define $\op{type}(N):H/n\to\ring{N}^3$ by $\op{type}(N) = (p,q,r)$
if there are $p$ triangles, $q$ quadrilaterals, and $r$
exceptional faces that meet the node $N$.  (An exceptional face is
defined to be one with at least $5$ darts.)


\begin{definition}
Let $H$ be a hypermap.  A flag $\op{fl} = (\op{qrn},\ldots)$ is a
{\it primary flag} for $H$ if it satisfies the following
conditions:
    \begin{enumerate}
        \item $\op{qrn}(\alpha)$, for every dart $\alpha$.
        \item $\op{qre}(\alpha)$, for every dart $\alpha$.
        \item $\op{typ}(5)\ne\op{notpri}$.  (This is a double negation: the flag is
        {\it not not\/} primary.)
        \item $\exists\alpha\in F.\ \op{type}(\op{node}(\alpha)) \ne
    (3,0,1),\quad F\in (F/f)_{5,\op{aggA}}$.
        \item Except for the preceding two restrictions on $\op{typ}(n)$, the flags
        $\op{typ}(n)$ and $\op{sum\_typ}(n)$ are arbitrary, for $n=4,\ldots,8$.
    \end{enumerate}
\end{definition}

\begin{remark} By the first two properties, every face is
standard.  Furthermore, the domains of $\op{quad\_diag}$ and
$\op{hex\_loop\_diag}$ are empty, and so they leave nothing to
specify.
\end{remark}

\begin{definition}
A triple $\rho=(H,\op{fl},S)$ is a primary tame hypermap system if
    \begin{enumerate}
    \item $H$ is a tame hypermap in the archive.
    \item $\op{fl}$ is a primary flag for $H$.
    \item $S=\op{basic}(H,\op{fl})$ is the set of basic formulas for $I=I(H,\op{fl})$.
    \end{enumerate}
\end{definition}


\subsection{Refinement -- Lookup}
%
\label{sec:refine}

The refinement predicate $\op{ref}(\rho,\rho')$ is a disjunction
    \begin{equation}
    \op{ref}(\rho,\rho') =
    \op{lookup\_database}(\rho,\rho')\vee
    \op{sum\_hyper}(\rho,\rho').
    \end{equation}
This subsection describes the predicate $\op{lookup\_database}$,
and the following subsections  will describe the more complicated
predicate $\op{path\_hyper}$.

As the name suggests, $\op{lookup\_database}(\rho,\rho')$ holds
when $\rho'$ is obtained from $\rho=(H,\op{fl},S)$ by adding
lookups from a database.

\begin{definition}
We say that $\rho'$ is derived from $\rho=(H,\op{fl},S)$ by lookup
in the database $B$ and write
    $$\op{lookup\_database}(\rho,\rho')$$
if $\rho' = (H,\op{fl},S \cup S')$ where $S'$ is a set of lookups
in $B$ with domain constraints in $\op{single}(S)$. (The database
$B$ is fixed once and for all and hence is omitted from notation.)
\end{definition}

\subsection{Connected Sum -- Hypermaps}

The second mode of refinement is the connected sum of a hypermap.  This
subsection gives the details of this construction.

Let $\rho=(H,\op{fl},S)$ and $\rho'=(H',\op{fl}',S')$.   The rough
idea is that $\op{sum\_hyper}(\rho,\rho')$ holds if we can
obtain $H'$ from $H$ by a connected sum. That is,
    $$H' = H\,\#_\phi\, B.
    $$
The actual
construction of the predicate $\op{sum\_hyper}$ is somewhat more
involved because it involves a change in hypermap and
corresponding flags. This means that the indexing sets for the
system of formal disjunctions also change.  Part of the connected sum
operation involves relating the formal disjunctions $S$ on
indexing set $I(H,\op{fl})$ to formal disjunctions $S'$ on the
indexing set $I(H',\op{fl'})$.

\subsubsection{Sum Kit}.


A connected sum of hypermaps $H\,\#\,B$  requires an
orientation reversing bijection $\phi:F\to F_B$ between a face of
$H$ and a face of $B$.  Some of the flags have been designed to
give such bijections.  In fact,  if $n=\card(F)$, every flag
$\op{sum\_typ}(n)(F)$ of a patchable face $F\in (H/f)_{n}$
uniquely determines a hypermap $B$ and an orientation reversing
bijection $\phi:F\to F_B$, for some face $F_B$ of $F$.

The data of $B,\phi$ and some auxiliary information about flags is
what we call a {\it connected sum kit}.  The connected sum kits are described by
various diagrams.  We list the diagrams (Figure~\ref{XX}) and then
say how to interpret them.

(N.B. This section is incomplete.  It is necessary to give all the
figures, labeled according to the interpretive key.)  % XX DO THIS.

The following comments provide an interpretive key to the
diagrams.
\begin{itemize}
    \item $B$ is the connected, plain, planar
    hypermap obtained by inserting darts at each angle of the
    planar graph, in the usual way, as illustrated in
    Figure~\ref{XX}. %% XX Provide figure.
    (N.B. Figure is to be inserted later.) %% XX DO THIS
    \item The face $F_B$ is the outermost face; that is, the one
    bounding the unbounded region.
    \item The face $F_B$ has a distinguished dart $x$ marked on it,
    indicated as a small black triangle.
    \item In cases when the hypermap $B$ is not isomorphic to the
    mirror-image of $B$, there is a second hypermap that is
    relevant. In these cases, the diagram for mirror 
    connected sum is obtained as the
    mirror image of the diagram.  Thus, we draw a single diagram
    and give it a second mirror name that is used to refer to the
    mirror data.
    \item The orientation-reversing bijection $\phi:F\to F_B$ is
    determined by choosing $\alpha\in F$ and specifying that
    $\phi(\alpha) =x$.  The map $\phi$ extends by the rule
        $$f_B^i\phi(f^i\alpha ) = x.$$
    \item Each diagram has a name that we use as the name of the connected sum kit.
\end{itemize}

\begin{remark} As just mentioned, the orientation-reversing bijection $\phi:F\to F_B$
is determined by $\alpha\in F$.  The $\alpha$ and $f^i\alpha$ give
isomorphic connected sums if there is an automorphism of the hypermap
carrying $\alpha$ to $f^i\alpha$.  In this case, the infeasibility
for the parameter $f^i\alpha$ can be deduced from the
infeasibility for the parameter $\alpha$.  Thus, the two cases can
be combined into one.  From a mathematical point of view this
reduces the amount of verification to be done.  However, Nipkow's
experience with the formal verification of tame graphs suggest
that case mergings generate additional proof obligations, and
sometimes actually increase the amount of work required for a
formal proof.
\end{remark}








\subsubsection{Compatibility with Flags}
\label{sec:com}

Let $H'= H \,\#_\phi\, B$.   There is a natural inclusion of
darts.
    $$
    \iota: H\setminus F \to H'.
    $$
There is a natural inclusion of faces, which we will express by
the same symbol:
    $$
    \iota: H/f \setminus \{F\}\to H'/f.
    $$

\begin{definition}
We say that a flag $\op{fl}'$ is compatible with
$(H,\op{fl},B,\phi)$  if the following conditions hold:
    \begin{enumerate}
    \item $\op{fl}'$ is a flag on $H'= H\,\#_\phi\, B$.
    \item The domain $F$ of $\phi$ is a patchable face.  That is,
    $\op{type}(\card(F))(F) = \op{patchable}$.
    \item The connected sum kit $(B,F_B,\phi,\ldots)$ of $F$ is given by
    $\op{sum\_typ}(\card(F))(F)$.
    \item We have $\op{qrn}(\alpha) = \op{qrn}'(\iota\alpha)$.
    (We recall also $\op{qrn}'$ is required to be constant on each
    node.)
    \item The node $N$ of $H'/n'$ (there is at most one) that does not lie in the image of
    $\iota$ has
        $\op{qrn}'(\beta)$ for $\beta\in N$ exactly when that node
        is {\it not} marked $\otimes$ in the Figure~\ref{XX}.  (That is, in
        every case except Figure~\ref{XX}, darts at a new node satisfy
        $\neg\op{qrn}'(\beta)$.)
    \item We have $\op{qre}(\alpha)= \op{qre}'(\iota\alpha)$.
    (We recall that $\op{qre}'$ is required to be constant on each
    edge.)
    \item The edges $E$ of $H'/e'$ that do not lie in the image of
    $\iota$ satisfy $\op{qre}'(\beta)$ for $\beta\in E$ exactly
    when that edge is marked as a single line in Figure~\ref{XX}.
    (That is, double lines give $\neg\op{qre}'(\beta)$.)
    \item We have $\op{hex\_loop\_diag}(\alpha) =
    \op{hex\_loop\_diag}(\iota\alpha)$.
    \item The edges $E$ of $H'/e'$ that do not lie in the image of
    $\iota$ satisfy $\op{hex\_loop\_diag}(\alpha)$ for $\alpha\in E$
    exactly when the edge is explicitly marked as such in the
    figures.  (The figures also shed some light on the terminology: such
    edges always run between opposite corners of a hexagonal face
    and run alongside a loop of four triangles.)
    \item The flags $\op{typ}(i)$ and
    $\op{sum\_typ}(i)$ are preserved on faces:
        $$
        \begin{array}{lll}
        \op{typ}(\card(F))(F) &= \op{typ}'(\card(F))(\iota F),\\
        \op{sum\_typ}(\card(F))(F) &= \op{sum\_typ}'(\card(F))(\iota F),\\
        \end{array}
        $$
    \item Usually, a face that is not in the image of $\iota$ is
    not standard, and hence lies outside the domain of the flags
    $\op{typ}(i)$ and $\op{sum\_typ}(i)$.  In this case there is
    no compatibility condition.
    \item If $\card(F)=8$, $\op{typ}(8)(F)=\op{patchable}$,
    $\op{sum\_typ}(8)(F) = \op{agg51}$, then there is a
    pentagonal face $F'$ in $B$ that carries over to $H'$.  In this case, we
    give the compatibility condition
        $$\op{typ}'(5)(F') = \op{notpri}.$$
    (This is, in fact, the only way to obtain the flag value
    $\op{notpri}$.)
    \end{enumerate}
\end{definition}


\begin{remark}
We note that if $H'= H\,\#_\phi\, B$ and $\op{fl}$ are given,
then there is at most one flag $\op{fl}'$ that gives
compatibility.
\end{remark}

\begin{remark}
We also note that if $H' = H\,\#_\phi\, B$, then by the
compatibility conditions, the faces coming from $B$ never have
type $\op{patchable}$. This means that we can never apply one
connected sum upon another.
\end{remark}

\subsubsection{Indexing Sets}

We assume that $H' = H\,\#_\phi\, B$.  As in the previous
subsection (\ref{sec:com}) we write $\iota$ for the natural
inclusion  on darts and faces.

Under the natural inclusion
    $$\iota:H\setminus F \to H'$$
there is a correspondence between a subset of $I=I(H,\op{fl})$ and
$I'=I(H',\op{fl}')$.  We make this correspondence more explicit
here.



All the indexing sets for $I=I(H,\op{fl})$ run over a subset of
darts (for $\op{azim}$, $\op{yn}$, $\op{ye}$, and so forth) or
over a subset of faces (for $\op{sc}$, $\op{sol}$, and so forth).
Let $I^F\subset I$ be the subset of the indexing set that does not
involve any of the darts of $F$, and does not involve the face $F$
itself.

We then have a one-to-one map $I^F \to I'$.  This gives a
one-to-one map between the subset of the set of formal
inequalities on the indexing set $I$ (those involving only the
indices $I'$) and the set of formal inequalities on the indexing
set $I'$. Extending to formal disjunctions, we get a one-to-one
map $\op{LFOR}_{I^F}\to \op{LFOR}_{I'}$.

\begin{definition}  We say that a formal disjunction on the
indexing set $I'$ is {\it inherited} from one of $I$, if it lies
in the image of this one-to-one map with domain
    $$\op{LFOR}_{I^F}\subset \op{LFOR}_I.$$
\end{definition}



\begin{definition}
We say that $\rho'=(H',\op{fl}',S')$ is obtained by a hypermap
connected sum $\rho =(H,\op{fl},S)$  and write
$\op{sum\_hyper}(\rho,\rho')$ if there exists a patchable face
$F\in H/f$ such that
    \begin{enumerate}
    \item $(B,\phi)$ is the data determined by $\op{sum\_type}(\card(F))(F)$.
    \item $H' = H\,\#_\phi\, B$.
    \item The flag $\op{fl}'$ is compatible with $(H,\op{fl},B,\phi)$.
    \item Every formal disjunction $A'$ in $S'$ is either a basic
    formula for $(H,\op{fl})$ or is inherited from $S$.
    \end{enumerate}
\end{definition}

\subsection{Constants}

This document has used a number of constants: $\xi_\Gamma$,
$\xi_V$, $\pt$, $s_i$, $t_i$. $D(n,k)$, $Z(n,k)$, $t_0$, $\phi_0$,
$\delta_\op{oct}$.
%
These constants are all the same as in \cite{DCG}.  They are
listed in its index.


\section{Implementing the Proof of Weak Infeasibility}

By the definition of weak infeasibility, there exists a tree that
describes the branches followed in the search for infeasibility.
Ultimately the proof of the main LP bound is constructive, in the
sense that the proof gives all of the data needed to construct the
search tree.

Of course, many different infeasibility trees can be used as
certificate of weak infeasibility.  Making less than optimal
decisions about how to construct the tree may increase the amount
of computation needed, but if a reasonable strategy is followed,
the exact choices are not too important.  What we give in this
section are some search strategies that seem to be helpful, and
that were in fact helpful in the original 1998 proof of the Kepler
conjecture.  We wish to emphasize that it may very well be
possible to devise much more efficient search strategies than
those described here.

At every stage of the search there are many possible options about
how to proceed.  Some of the options are to
    \begin{itemize}
    \item Obtain a new hypermap by a connected sum of the current one.
    \item Bifurcate on a disjunctive formula.
    \item Enhancing the system of formal disjunctions by adding
    lookups.
    \end{itemize}
The simple act of bifurcating on a formula has infinitely many
possibilities.

To construct the infeasibility trees, it is necessary to be guided
by various heuristics.  We are not looking for a precise
algorithm, merely heuristics that can guide a computer program to
construct the appropriate trees. What actions should have high
priority and what actions should be used as a last resort?



Of course, it may be the the original system is infeasible.  The
first thing to try is the infeasibility of the system obtained
without including any inequalities that depend on the value of the
flags.  (It is often best to ignore the inequalities depending on
flags in the early stages.)  If this system is infeasible, it
implies in one fell swoop the weak infeasibility of all primary
tame hypermap systems $(H,\op{fl},S)$ with a given hypermap $H$.
In fact, thousands of hypermaps can be quickly eliminated by this
one simple strategy.  We now confine our remarks to primary tame
hypermap systems for which this simple strategy fails.

Sometimes the heuristic can be guided by geometry.  The nodes of
the hypermap represent the centers of a finite cluster of spheres
surrounding the origin.
 The variables $\optt{yn}(\alpha)$ represent the distances of
 these sphere centers from the origin.  The idea is move the
 system in the direction of a cluster of spheres in which all of
 the distances $\optt{yn}(\alpha)$ are constrained to be about
 2.0:
    $$2 \sle \optt{yn}(\alpha) \sle 2+\epsilon.$$
 Once these variables are heavily constrained, the focus moves to the
 variables $\optt{ye}(\alpha)$.  The strategy is then to constrain
 these variables as tightly as possible between an upper and lower
 bound.  Once strong constraints are obtained on these variables,
 the weak infeasibility often follows quite easily.

Another rough guide is to keep the number of branches under
control.  For example, we could easily get exponential growth in
the tree if we repeatedly bifurcate all active branches of the
tree.  Large exponential growths should be avoided.  Thus, it is
often good to bifurcate in a way that one of two possibilities is
infeasible.  For example, we can choose a constant $c$ (obtained
by linear programming) for which
    $$\optt{yn}(\alpha)\sge c$$ is infeasible.  Then a bifurcation
    on
$$(\optt{yn}(\alpha) \sge c) \vee (\optt{yn}(\alpha) \sle c)$$
has a single continuation.  This allows us to add
    $$\optt{yn}(\alpha)\sle c$$
to the system.  Adding such an inequality has further
consequences, because there are many nonlinear lookups whose
domain constraints take precisely this shape.

One rough guide is that the most useful inequalities come from
faces of the smallest size.  In particular, the inequalities for
triangles are by far the most useful.  Thus, one strategy is to
make triangles the focus of the bifurcation for as far into the
search tree as possible, bringing in large faces only as necessary
late in the game.

Since all triangles in a primary tame hypermap system are
standard, the opening moves should concentrate on standard
triangles.  A good opening strategy might be to narrow the
feasible range on the variables $\optt{yn}(\alpha)$.

\section{Conclusions}

This document is still sketchy in many places.  However, the
general framework of the linear programming has been placed on
solid foundation.  There are other references that may be
consulted for further details.  These references include the
published proof of the Kepler Conjecture in \cite{DCG} (Section
23, pages 236ff). There is the database of nonlinear inequalities
at Hales's web page. There is the archive of tame planar graphs.

\subsection{Mathematica}

There are Mathematica files that were used to generate the linear
programs.  The most important of these is {\it MathToCplex.m}
which appears with the linear programming materials in the
original 1998 proof.  This file is important because it generates
all of the basic linear programs.  There are also supplementary
files {\it MathToCplexExcept.m} and {\it MathToCplexPent.m} that
handle the more sophisticated branch and bound operations.  All
the linear programs were generated by these files.

Interactive Mathematica sessions were used to run and analyze the
linear programs.  The Mathematica programs generate text files in
cplex format and stored as external files. The Cplex linear
programming package was run on these files, by issuing a command
from Mathematica to the operating system.  The results from the
linear programs were stored as text files, then parsed back into
the Mathematica session for analysis.  The human controller
(Hales) made interactive decisions about how to branch the linear
programs, based on heuristics and intuition about how to establish
infeasibily in the quickest way possible.  Various Mathematica
procedures were written to automate the branch and bound process.

These interactive sessions were summarized in the form of log
files, to make it possible to reconstruct the moves that led to
infeasibility. These log files consist of notes written by Hales.
Mixed in with the notes are snippets of the Mathematica code that
was used to direct the branch and bound operations.  Thus, in
principle, the log files contain all that is necessary to
reconstruct the sequence of linear programs.  These are all
available from Hales's web page.

The files terminating in .m are Mathematica files.  Those in .exec
are files that control the execution of Cplex.  Those in .lp are
cplex format linear programs.  Those in .log are the output from
Cplex sessions that get parsed back into a Mathematica session.
Those in .html contain my notes and snippets of Mathematical code
that describe the branch and bound for the linear programs.

It would be clearly very valuable to rework this, so that the
process is automated from beginning to end.  The Flyspeck project
hopes to carry this even further, to give a formal proof of the
main theorem of this paper.

\subsection{An example of a log file}

The following is the file from Hales's webpage
(/SHORT/OCT/oct.zip/OCT/index.html) that describes how the
hypermaps containing an octahedral face were handled.  It is a
mixture of plain text and Mathematica code.  A few comments have
been added.  Otherwise, it is identical to the 1998 log.  The
hypermaps with octahedral faces would be a good place to start in
a second generation treatment of the linear programs for hypermaps
with an exceptional face.

\bigskip

{\obeylines\tt

(* How to do the octagonal cases.  This first block of code
generates all the linear programs for the octahedra.  The files
are stored in the SHORT/OCT/LP directory, and the file names all
start with cplexE.lp.  *) \

\

\

f[i\_] := (
        exceptStem = "SHORT/OCT/LP/cplexE.lp";
        Initialize[8,i];
        Array[(Print[{8,i, \#1}];
        WRITEOUTexcept[8,i, Length[GBLregions], \#1]) \& , Length[AllOct]]
        );
Map[f,locthi];

\

\

(* The following block of code generates code for an interactive
cplex session.  It tells cplex to read each of the files and run
the linear programming optimization in the file. *)

\

\

optText[n\_,i\_]:= Module[{},
Initialize[n,i];
        Array[{"read ",
            fileText["cplexE.lp",n,i,Length[GBLregions],\#], " lp$\backslash$n",
            "opt$\backslash$n"
            }\&, Length[AllOct]
            ]//StringJoin
        ];

\

\

(* This stores the cplex instructions to an external file called
cplexE.exec *)

\

\

Module[{vv,stream},
    vv = Map[optText[8,\#]\&,locthi]//StringJoin;
    stream=OpenWrite["SHORT/OCT/cplexE.exec"];
    WriteString[stream,vv];
    Close[stream]
    ];

\

\


(* After running the cplex program, the results are parsed back
into Mathematica as an array called values.  The following command
finds that there is just one case that didn't pass on this first
attempt.  The command is everything up to the symbol ==.
Everything following the symbol "==" is the output from the
command.  The constant 0.4429 is approximately 8 pt.  *)

\

\

hh= Select[values,\#[[1]]>0.4429\&]== {{0.459297, 8, 14, 17, 21}};

\

\

(* The rest gives a few hints about how to eliminate this final
case.  It is based on the inequalities from DCG page 252, section
25.6.2. We are able to get a smaller value of the penalty from the
final lines of DCG section 25.6.2.  What follows is a mixture of
plain text and Mathematica.  *)

\

\

LPmax[bf,"-sigma{17}"] == 0.299186; if there is a Z(4,2):   s8 +
ZLP(4,2)-Z(4,2) = -0.31398. So no (4,2)s.  Since the two flats are
opposite, the only penalties
    are from (4,1)s that mask a flat, so pen17 < 2(xiG+2xiV) = 0.045304.
    Edit the file so pen17 = 0.045304.
bf = "SHORT/OCT/LP/cplexE.lp8.14.F17.C21pen"

\

\

(* With this smaller penalty in place, we branch and bound on
eight of the triangular faces, with the usual 6.25 cutoff, as
described in DCG, Sec 24.1, page 241. We find that in every case
the bound comes out less than 8 pt. Again, everything before the
"==" is the Mathematica command, and everything following it is
the output from the interactive session.  Since this constant is
less than 8 pt, we are done. *)

\

\

ww=( Initialize[8,14]; Array[
LPdisplay["SHORT/OCT/LP/cplexE.lp8.14.F17.C21pen",
    Branch[{6,7,9,10,11,12,13,14},\#] ]\&,2\^8])//Max ==  0.441675

\

\

OCTs are done.

}




\subsubsection{Second Generation LP}

A couple of years after the proof was completed, an independent
check was made of the linear programs for tame plane graphs that
contain only triangles and quadrilaterals (2000, unpublished, but
available on Hales's web page). This is an entirely new and
independent implementation of the linear programming, but it was
not done in sufficient generality to handle hypermaps with
exceptional faces. This code is written in Java.  The results of
the linear programs are checked with a Java-based implementation
of interval arithmetic. This code has the advantage that the
branch and bound operations are entirely automated.  This code is
available from Hales's website.

The reason that this code is unable to treat hypermaps with
exceptional faces is that no implementation has been given of the
connected sum operation (called faced refinements in the original 1998
proof).


\subsection {What is still missing from this document}

This document is starting to get quite long, and it still is far
from complete.  Based on what is here, it now looks as if a
detailed specification of the linear programming part of Kepler
may run to some 100 pages. My hope at this point is that it
contains enough about the structure of the task, to make the other
sources of information comprehensible.

This final subsection has been written as a work in progress,
giving lists of topics that must eventually be carefully
specified, but that have not yet been.

\subsubsection{Inequalities}

This document does not contain a complete list of the inequalities
that are used in the linear programming infeasibility.  Consult
the Mathematica files and \cite{DCG} for the list.  There are
various important inequalities that need to be added.

\begin{itemize}
    \item There is an important identity for $\optt{vor0}$ in terms of the
    variables $\optt{quo}$, $\optt{Adih}$, etc.
    \begin{equation}
    \optt{vor0}(F) = \phi_0 \optt{sol}(F) +\sum_{\alpha\in F}
    (\optt{Adih}(\alpha) - 4\delta_\op{oct}(\optt{quo}(\alpha)
    +\optt{quob}(\alpha))).
    \end{equation}
    (Compare \cite{DCG}, Sec 25.12, page 257.)
%
    \item There are various inequalities for $\optt{quo}$ obtained
    by nonlinear lookup.
%
    \item We have
        \begin{equation}
        \optt{quob}(\alpha) = \optt{quo}(e\alpha),\quad
        \forall\alpha.
        \end{equation}
    The nonlinear interpretation of $\optt{quo}(\alpha)$ is
    the function $\op{quo}(v,w)$ of \cite{DCG}, and that of
    $\optt{quob}(\alpha)$ is $\op{quo}(w,v)$.
%
    \item There are various inequalities for $\optt{Adih}(\alpha)$
    obtained by nonlinear lookup.  These are rather
    unconventional, because they require the input of lower and
    upper bounds on the variable $\optt{azim}(\alpha)$.  There is
    a note about this at the end of Sec 25.12 of \cite{DCG}, page 257.
%
    \item When $\optt{eta}'(\alpha)\sge t_0$, then $\optt{quo}(\alpha)
    = \optt{quob}(\alpha)=0$.
    \item There are lower and upper bounds on $\optt{azim}(\alpha)$ for all
    darts $\alpha$, depending on the structure of the face to
    which $\alpha$ belongs.  These are given in a table in
    \cite[VI,p.54]{Hal98d}.  The same bounds appear in the database as
    inequality 853728973.
    \item There are the hexagonal inequalities of \cite[25.13]{DCG}.
  There are some problems -- still unresolved by this document --
  about the interpretation of the variables $\sigma_R$,
  $\sigma_R^+$, $\tau_R$, $\tau_R^+$ of that paper in terms of the variables
  of this document.
  \item There is the ``$3.0$'' quadrilateral inequality of
  \cite[Sec.25.10]{DCG}.
  \item There are the $t(n)$ and $s(n)$ inequalities when there
  are no aggregates, and the weaker version of the same
  inequalities that hold over aggregates.
\end{itemize}

There are several important inequalities that hold over groups of
faces.

\begin{itemize}
  \item There are the ``$0.55$'' and ``$0.48$'' bounds found in
  \cite[Lemma 10.6,p.101]{DCG}.
  \item There are the ``$1.4$'' and ``$1.5$'' bounds found in
  \cite[Lemma~22.12,p.233]{DCG}.  There is a remark at the end of
  Section~4.10 of \cite[VI,p.37]{Hal98d} that shows how these
  inequalities can sometimes be improved by removing certain faces
  from the sums of those inequalities.
\end{itemize}


There are special inequalities on simplices that are potentially
of type C or C'.  See the note below in
Subsection~\ref{sec:wrong}.





\subsubsection{Connected Sums}

This document does not yet draw all of the connected sum possibilities.
Again, this is found in \cite[Sec.25.5--25.6,pp.249--251]{DCG}.
When drawing these connected sum kits, it is necessary to add a mirror
case to each one that does not have reflectional symmetry.  When
drawing these connected sum kits, it is necessary to add markings as
described in the interpretive key in this document.

The penalties from \cite{DCG} have been changed in this document
so that they adhere to a specific face, rather than to the entire
connected sum kit.  Doing so, cleans up the scoring rule of
Equation~\ref{eqn:xi}.




\subsubsection{Lookup}

Detailed information about which inequalities were pulled from the
nonlinear database by lookup can be determined by consulting the
Mathematica files.





\subsubsection{Variables}

Complete the list of variable names.  Check the names of variables
on such things as $\optt{sigmaflat}$ $\optt{tauflat}$,
$\optt{svoran}$, $\optt{nu}$, $\optt{vorx}$.

\subsubsection{Extensions}

There are other linear programs, besides the Main LP Bound that
arise in the proof of the Kepler conjecture.  These are the {\it
small} linear programs, so-called because they involve a small
number of variables and constraints.  An example of small linear
programs are those appearing in the proof of Lemma~10.5 in
\cite{DCG}.  There are similar linear programs scattered
throughout the paper.

The linear programs for the proof of the dodecahedral conjecture
should be developed along similar lines as this.

This document separates the linear programming infeasibility
results from the geometry.  To build a bridge between the Main LP
Bound and the proof of the Kepler conjecture, it is necessary to
prove that any counterexample to the Kepler conjecture leads to a
strong solution to a tame hypermap system.  The basic ingredients
to prove this are already in \cite{DCG}, but it will be necessary
to write this up, to show exactly how it follows from results in
those papers.




\subsubsection{Things that may be wrong or incomplete}
\label{sec:wrong}

The treatment of linear programming variables related to the score
is very muddy in \cite{DCG}.   I am not aware of any explicit
errors, but it is extremely difficult to piece together precise
linear programming inequalities, based on what is written there.
Section 25.12 of \cite{DCG} is particularly confusing, and should
be avoided whenever possible.  The description of penalties in
Section~\ref{sec:pc} and Equation~\ref{eqn:xi} are intended to be
clarified versions of \cite{DCG}.

There is one important case of the scoring function that has not
been entirely specified in this document.  This concerns the
inequalities for $\optt{sc}(F)$ for triangular faces $F$ with one
edge that is not quasi-regular and one node that is not
quasi-regular.  These include the cases of simplices of type C and
type C' from \cite[Sec.9.4,pp.94--99]{DCG}, which have special
scoring rules.  On such simplices there will be an inequality that
takes the general form
    $$
    \optt{sc}(F)\sle \optt{vorx}(F).
    $$
There are a number of lookup inequalities for $\optt{vorx}$.  What
I still need to confirm is that nothing funny is going on with the
switches to type C, C', etc.

All the linear programs use weak inequalities $(\le),(\ge)$.  When
we branch over a disjunction
    \begin{equation}
    \label{eqn:dis}
    (\optt{A}\sle C)\vee (\optt{A}\sge C),
    \end{equation}
there is overlap between the two cases when equality occurs.
Normally, this is not a problem; but in a couple of special
situations it creates a problem.  (In the 1998 proof this was
handled carefully at a conceptual level, but it does not find its
way into the code.  For the fully automated proof, it has to find
its way into the code.)

Even though this issue may seem minor, I think that it may be
necessary to add some components to the flags to resolve the
problems this creates.  What I am hoping to avoid is a
modification of the flag that makes it so that in a connected sum, there
can be more than one compatible flag $\op{fl}'$.  I prefer to
maintain uniqueness for $\op{fl}'$.

One place where it creates a problem is when there is explicit
branches between scoring systems.   For example, sometimes it is
necessary to consider two cases of scoring (compressed and 
decompressed) on quarters.  Since the scoring type cannot be
determined by weak inequalities (in the ambiguous case when a face
circumradius is exactly $\sqrt2$), weak inequalities are
inadequate in such cases.  The branching should really run over a
flag giving the scoring type, rather than over weak inequalities
for the circumradius of a face.


There is a trick that can often be used to avoid the case of
equality in the Formula~\ref{eqn:dis}.   One helpful strategy is
to transform disjunctions $A\vee (\optt{x} < B)$ into implications
    \begin{equation}
    (\optt{x} \ge B) \Rightarrow A.
    \end{equation}
(where there is no weakening), rather than into weak inequalities
    \begin{equation}
    A \vee (\optt{x}\le B)
    \end{equation}
Implications can be handled with nonlinear lookups,
$\op{nonlin}(C,L)$, where the antecedent is placed in $C$ and the
consequent is placed in $L$.

\begin{example}  For example, consider the branching over edge
lengths at the cutoff $2.45$ as discussed in \cite[Sec.25.7]{DCG}.
The inequalities $\optt{ye}(\alpha)\sle 2.45$ should transformed
as antecedents in the nonlinear lookup.  %% XX Check this. Does it really work?
\end{example}
 
 
%%%%%%%%%%%%%%%%%%%%%%%%%%%%%%%%%%%%%%%%%%%%%%%%%%%%
    %
%% XX issues.

We mark warnings with the letters WW in 
the TeX file.

We mark serious issues with the letters XX.

\section{Notation}

We use the notation $]x,y[$ in trig.tex, then move away from it.

\section{Definitions}

Some definitions are repeated ($\Delta$, $\dtet$, etc.)
Mark these in some way.

Is the notation $(p_x,p_y,p_z)$ for coordinates 
of $p\in R^3$ consistent throughout?

Use the word simplex combinatorially for four vertices.
Use the word tetrahedron geometrically for the convex hull.

barrier should be closed. Intersecting segment should be open ended.
(Open ends prevent meeting at a vertex.)
Everything will be done up to a null set
anyway for measure related things.

\section{Define before use}

Make sure circle, great circle are defined before use.
A good place would be in vector geoemtry chapter.

Define basic notions of $\op{aff}$, $\op{conv}$, $\op{cone}$
before first use, as in vector geometry chapter.


\section{Policy Things}

Here we list some things that have been changed.  Make sure that
all occurences have been eliminated.

There should be no mention of enumerated packings.

Indexing: Things before the ``AT'' symbol are used for ordering entries.
Things after used for display in the actual index.
Subentry is!  After vertical slash stuff gets passed to page
number.  \\index{policy|bold}  puts the page in bold.
\\index{policy|(} starts a subrange \\index{policy|)} ends a page range.

Make all Deltas take same args in same order (to simplify formal proofs).
$\chi$, $\Delta$, $\ups$ should appear only in tarski.tex. 

\section{Incomplete Things in Tarski}

Eventually, I want to separate out the Tarski stuff, quoting
only the results that are needed.  Quote by footnote. Index every use.

Need an introduction to Tarski. What it is all about.

Replace $y_i$ by $y_{jk}$ globally.







    % backmatter

%\renewcommand\partname{Back Matter}
%\renewcommand\chaptername{Appendix}

\backmatter


\bibliographystyle{abbrv}
\begin{bibliography}{CHMS94}

% Hales:2006:DCG
\bibitem[Hal06]{DCG} Thomas C. Hales, A Proof of the
Kepler Conjecture (unabridged version), 
Discrete and Computational Geometry, 35:1, 2006.

% marchal:2007
\bibitem[Mar07]{Marc07} C. Marchal, D\'emonstration de la conjecture de Kepler,
http://hal.archives-ouvertes.fr/hal-00160350/en/, 2007.

% marchal:2008
\bibitem[Mar08]{Mar08} C. Marchal, Demonstration of Kepler's conjecture. The problem of the closest packing, http://publications.onera.fr/exl-php/notices/236481-demonstration-kepler-conjecture-the-problem-the-cl.htm, 2008.

% Rogers:1958:Packing
\bibitem[Rog58]{Rog58} C. A. Rogers, The packing of equal spheres, {\it Proc. London Math.
    Soc.} (3) 8 (1958), 609--620.

%vanderWaerden:1951


\end{thebibliography}




    %%


\end{document}
